% =====================================================================
%  Chapitre 6 : Applications et algorithmes quantiques
%  Fichier importé par main.tex avec % =====================================================================
%  Chapitre 6 : Applications et algorithmes quantiques
%  Fichier importé par main.tex avec % =====================================================================
%  Chapitre 6 : Applications et algorithmes quantiques
%  Fichier importé par main.tex avec % =====================================================================
%  Chapitre 6 : Applications et algorithmes quantiques
%  Fichier importé par main.tex avec \include{applications}, entre chapitre5 et
%  chapitre6 (exercices). Il ne contient ni préambule ni \begin{document} : les
%  environnements (definition, resultat, remarque, aretenir, savoirfaire), les
%  macros (\ii, \ee, \dd, \ket, \bra, \braket), les symboles de circuits (ctl,
%  cible, croix, mesure, porte) et les couleurs (insarouge, insagris) sont définis
%  dans main.tex. Un chapitre = une \section.
%  Le numéro d'un chapitre est donné par l'ordre des \include, pas par le nom du
%  fichier : ce chapitre est le chapitre 6, les exercices (chapitre6.tex) le 7.
% =====================================================================

% Macros de Dirac (déjà définies dans main.tex ; ce bloc évite l'erreur
% « Undefined control sequence \ket » si main.tex est une ancienne version).
\providecommand{\ket}[1]{|#1\rangle}
\providecommand{\bra}[1]{\langle #1|}
\providecommand{\braket}[2]{\langle #1|#2\rangle}

% =====================================================================
\section{Applications et algorithmes quantiques}\label{chap:applications}
% =====================================================================

Les chapitres précédents ont construit le formalisme : qubits, portes, intrication, mesure. Ce
chapitre montre à quoi il sert. On y voit comment un circuit exploite l'interférence pour résoudre
certains problèmes plus vite qu'un ordinateur classique (sections~\ref{sec:interference}
à~\ref{sec:shor}), pourquoi cela menace une partie de la cryptographie actuelle, RSA et TLS en tête
(section~\ref{sec:crypto}), ce qu'il faut pour construire une telle machine, du cryostat aux qubits
de chat d'Alice \& Bob (section~\ref{sec:machine}), et quelles autres applications et quels autres
algorithmes existent (section~\ref{sec:autres}). Les circuits sont dessinés comme aux chapitres~4 et~5 :
le fil du haut est le dernier chiffre du ket.

\begin{remarque}[Chiffres d'actualité]
Les estimations de ressources, les calendriers de migration et les records expérimentaux cités ici
datent d'octobre 2026 et évoluent vite. Chaque chiffre est accompagné de sa source (bibliographie) ;
les annonces d'entreprises sont signalées comme telles et les prépublications aussi.
\end{remarque}

% ---------------------------------------------------------------------
\subsection{Calculer par interférence}\label{sec:interference}
% ---------------------------------------------------------------------

Un registre de $n$ qubits est décrit par $2^n$ amplitudes (chapitre~\ref{chap:operateurs}). Une porte
agit sur toutes ces amplitudes à la fois, mais la mesure n'en révèle qu'un seul résultat, tiré au
hasard avec les probabilités de Born. Dire que «~le circuit calcule toutes les valeurs en parallèle~»
est donc trompeur. Un algorithme quantique est une \emph{chorégraphie d'interférences} : les
amplitudes des mauvaises réponses se compensent (signes ou phases opposés), celles de la bonne
réponse s'additionnent, et la mesure finale donne la bonne réponse avec une probabilité proche de $1$.

\begin{definition}{Oracle}
Pour $f:\{0,1\}^n\to\{0,1\}$, l'\emph{oracle} de $f$ est la porte à $n+1$ qubits
\[
  U_f\ket{x}\ket{y}=\ket{x}\ket{y\oplus f(x)}.
\]
Elle est unitaire et $U_f^2=\mathrm{Id}$ (car $y\oplus f(x)\oplus f(x)=y$). Elle se construit à partir
d'un circuit réversible de $f$ (portes CNOT et Toffoli, section~\ref{sec:controlees}). Le coût d'un
algorithme se mesure en \emph{nombre d'appels} à $U_f$.
\end{definition}

\begin{resultat}{Rebond de phase de l'oracle}
Avec l'ancilla dans l'état $\ket{-}=\frac{1}{\sqrt2}(\ket0-\ket1)$ :
\[
  U_f\ket{x}\ket{-}=(-1)^{f(x)}\ket{x}\ket{-}.
\]
\end{resultat}

En effet $U_f\ket{x}\ket{-}=\frac{1}{\sqrt2}\big(\ket{x}\ket{f(x)}-\ket{x}\ket{1\oplus f(x)}\big)$.
Si $f(x)=0$ on retrouve $\ket{x}\ket{-}$ ; si $f(x)=1$, les deux termes sont échangés et le signe
change. La valeur $f(x)$ est écrite dans la \emph{phase} de $\ket{x}$, sans toucher à l'ancilla : c'est
le même mécanisme que le rebond de phase du CNOT (section~\ref{sec:cnot}), et la brique de tous les
algorithmes de ce chapitre.

\subsubsection{Algorithme de Deutsch}

Soit $f:\{0,1\}\to\{0,1\}$. Elle est \emph{constante} si $f(0)=f(1)$ et \emph{équilibrée} si
$f(0)\neq f(1)$. On veut savoir laquelle des deux est vraie. Classiquement, il faut évaluer $f$ deux
fois. Le circuit de la figure~\ref{fig:deutsch} répond avec \emph{un seul} appel à $U_f$.

\begin{figure}[!ht]
\centering
\begin{tikzpicture}[>=Stealth, font=\small]
  % fils : ancilla (second chiffre) en haut, requête (premier chiffre) en bas
  \draw (0,1.5) -- (8.4,1.5);
  \draw (0,0) -- (7.4,0);
  \node[anchor=east] at (0,1.5) {$\ket{1}$};
  \node[anchor=east] at (0,0) {$\ket{0}$};
  \node[porte] at (1.3,1.5) {$H$};
  \node[porte] at (1.3,0) {$H$};
  \node[draw=insarouge, fill=insarouge!8, thick, minimum width=13mm, minimum height=22mm]
    at (3.9,0.75) {$U_f$};
  \node[porte] at (6.0,0) {$H$};
  \pic at (7.2,0) {mesure};
  \node[anchor=west] at (7.7,0) {$f(0)\oplus f(1)$};
  \node[anchor=west, text=insagris, font=\footnotesize] at (8.5,1.5) {ancilla};
  \foreach \x/\lab in {0.55/0, 2.2/1, 5.15/2, 6.75/3}
    {\draw[dashed, insagris!70] (\x,2.1) -- (\x,-0.7);
     \node[anchor=north, font=\footnotesize] at (\x,-0.7) {$\Pi_\lab$};}
\end{tikzpicture}
\caption{Circuit de Deutsch. La requête $x$ est le premier chiffre du ket (fil du bas), l'ancilla
$y$ le second (fil du haut). La mesure du seul qubit de requête donne $f(0)\oplus f(1)$ : $0$ si $f$ est
constante, $1$ si elle est équilibrée.}
\label{fig:deutsch}
\end{figure}

\begin{align*}
  \ket{\Pi_0}&=\ket{0}\ket{1},\qquad
  \ket{\Pi_1}=(H\otimes H)\ket{\Pi_0}=\ket{+}\ket{-},\\[1ex]
  \ket{\Pi_2}&=U_f\ket{\Pi_1}
    =\frac{(-1)^{f(0)}\ket{0}+(-1)^{f(1)}\ket{1}}{\sqrt2}\ket{-}
    =(-1)^{f(0)}\,\frac{\ket{0}+(-1)^{f(0)\oplus f(1)}\ket{1}}{\sqrt2}\,\ket{-},\\[1ex]
  \ket{\Pi_3}&=(H\otimes I)\ket{\Pi_2}=(-1)^{f(0)}\,\ket{f(0)\oplus f(1)}\ket{-}.
\end{align*}
La deuxième ligne utilise le rebond de phase ; la troisième, $H\ket{+}=\ket0$ et $H\ket{-}=\ket1$.
Les deux amplitudes de la requête interfèrent : elles s'additionnent si $f$ est constante (état
$\ket{+}$, devenu $\ket{0}$) et se retranchent si $f$ est équilibrée (état $\ket{-}$, devenu $\ket{1}$).
Le calcul ne dit pas les valeurs $f(0)$ et $f(1)$ : il extrait seulement leur parité, et c'est ce qui
permet d'économiser un appel.

\begin{remarque}[Des oracles à $n$ qubits : Deutsch--Jozsa, Bernstein--Vazirani, Simon]
Le même schéma ($H^{\otimes n}$, oracle, $H^{\otimes n}$) généralise l'algorithme de \textcite{deutsch1985}. Pour $f$ sur $n$ bits promise
constante ou équilibrée, \textcite{deutschjozsa1992} décident avec \emph{un} appel là où un algorithme classique
déterministe en demande $2^{n-1}+1$. \textcite{bernstein1997} retrouvent en un appel un mot caché
$s$ défini par $f(x)=s\cdot x$ (classiquement : $n$ appels). \textcite{simon1997} trouve un masque caché avec $O(n)$
appels, contre $2^{\Omega(n)}$ classiquement : c'est la première séparation exponentielle, et l'idée
de l'algorithme de Shor. Ces problèmes sont artificiels (ils portent sur une boîte noire), mais
ils prouvent que l'interférence offre un gain réel.
\end{remarque}

% ---------------------------------------------------------------------
\subsection{Recherche dans une liste : l'algorithme de Grover}\label{sec:grover}
% ---------------------------------------------------------------------

Un oracle $f$ vaut $1$ en un unique élément $w$ parmi $N=2^n$, et $0$ ailleurs ; on cherche $w$.
Classiquement, il faut en moyenne $N/2$ évaluations de $f$, et $N$ dans le pire cas. \textcite{grover1996}
le trouve avec environ $\frac\pi4\sqrt N$ appels.

\begin{resultat}{Algorithme de Grover}
On part de $\ket{s}=H^{\otimes n}\ket{0\ldots0}=\frac{1}{\sqrt N}\sum_x\ket{x}$, puis on répète $k$
fois l'itération $G=D\,O_w$ avec
\[
  O_w=I-2\ket{w}\bra{w}\ \ (\text{oracle, par rebond de phase}),\qquad
  D=2\ket{s}\bra{s}-I=H^{\otimes n}\,(2\ket{0}\bra{0}-I)\,H^{\otimes n}.
\]
La mesure donne $w$ avec la probabilité $\sin^2\!\big((2k+1)\theta\big)$, où $\sin\theta=\frac{1}{\sqrt N}$.
Elle est proche de $1$ pour $k\simeq\frac\pi4\sqrt N$. Aucun algorithme quantique ne fait mieux que
$\Omega(\sqrt N)$ appels : le gain est quadratique, et optimal.
\end{resultat}

\paragraph{Pourquoi cela marche.} Posons $\ket{s'}=\frac{1}{\sqrt{N-1}}\sum_{x\neq w}\ket{x}$. Alors
$\ket{s}=\sin\theta\,\ket{w}+\cos\theta\,\ket{s'}$, et toute l'évolution reste dans le plan
engendré par $\ket{w}$ et $\ket{s'}$. L'oracle change le signe de la composante sur $\ket{w}$ : c'est
la symétrie par rapport à l'axe $\ket{s'}$. La diffusion $D=2\ket{s}\bra{s}-I$ est la symétrie par
rapport à $\ket{s}$. Deux symétries d'axes séparés de l'angle $\theta$ composent une rotation de
$2\theta$ (figure~\ref{fig:grover-geo}). Après $k$ itérations,
\[
  G^k\ket{s}=\sin\!\big((2k+1)\theta\big)\ket{w}+\cos\!\big((2k+1)\theta\big)\ket{s'}.
\]
On s'arrête quand $(2k+1)\theta\simeq\frac\pi2$. Contrairement à un algorithme classique, itérer plus
longtemps n'améliore rien : on continue de tourner et on dépasse $\ket{w}$.

\begin{figure}[!ht]
\centering
\begin{tikzpicture}[>=Stealth, font=\small]
  \draw[->, insagris] (-0.4,0) -- (4.5,0) node[right] {$\ket{s'}$};
  \draw[->, insagris] (0,-0.5) -- (0,4.5);
  \draw[insagris!50, dashed] (3.6,0) arc (0:90:3.6);
  \draw[dashed, insagris!60] (210:1.8) -- (30:3.6);
  \draw[->, very thick, insarouge] (0,0) -- (30:3.6) node[right] {$\ket{s}$};
  \draw[->, very thick, insagris] (0,0) -- (-30:3.6) node[right] {$O_w\ket{s}$};
  \draw[->, very thick, insarouge] (0,0) -- (90:3.6) node[above right] {$G\ket{s}=\ket{w}$};
  \draw[insagris] (1.1,0) arc (0:30:1.1);   \node[text=insagris] at (15:1.45) {$\theta$};
  \draw[insagris] (1.1,0) arc (0:-30:1.1);  \node[text=insagris] at (-15:1.45) {$\theta$};
  \draw[->, insarouge, thick] (30:2.2) arc (30:90:2.2);  \node[text=insarouge] at (60:2.65) {$2\theta$};
  \node[text=insagris, font=\footnotesize, anchor=north west, align=left] at (-1.7,-1.2)
    {symétrie $O_w$ : $\ket{s}\mapsto O_w\ket{s}$\\ symétrie $D$ d'axe $\ket{s}$ (pointillés)};
\end{tikzpicture}
\caption{Une itération de Grover pour $N=4$ ($\theta=\frac\pi6$). L'oracle renvoie $\ket{s}$ en
$O_w\ket{s}$ (angle $-\theta$), la diffusion le renvoie en $G\ket{s}$ (angle $3\theta=\frac\pi2$) :
une seule itération amène exactement sur $\ket{w}$.}
\label{fig:grover-geo}
\end{figure}

\subsubsection{Exemple complet : $N=4$ (deux qubits)}

Prenons $w=11$. L'oracle est la porte $Z$ contrôlée $\mathrm{CZ}=\operatorname{diag}(1,1,1,-1)$
(section~\ref{sec:controlees}). La diffusion s'écrit avec les portes du cours, car
$2\ket{00}\bra{00}-I=\operatorname{diag}(1,-1,-1,-1)=\mathrm{CZ}\,(Z\otimes Z)$ :
\[
  D=(H\otimes H)\;\mathrm{CZ}\,(Z\otimes Z)\;(H\otimes H).
\]
Le circuit de la figure~\ref{fig:grover-circ} s'écrit donc $H\otimes H$, $\mathrm{CZ}$,
$H\otimes H$, $Z\otimes Z$, $\mathrm{CZ}$, $H\otimes H$. Calcul pas à pas :
\begin{align*}
  \ket{s}&=\tfrac12\big(\ket{00}+\ket{01}+\ket{10}+\ket{11}\big),\\
  O_w\ket{s}&=\tfrac12\big(\ket{00}+\ket{01}+\ket{10}-\ket{11}\big),\\
  \braket{s}{O_w s}&=\tfrac14(1+1+1-1)=\tfrac12,\\
  D\,O_w\ket{s}&=2\braket{s}{O_w s}\ket{s}-O_w\ket{s}=\ket{s}-O_w\ket{s}
    =\tfrac12\big(0,0,0,2\big)^{\mathsf T}=\ket{11}.
\end{align*}
La mesure donne $11$ \emph{avec certitude} après un seul appel à l'oracle, là où un algorithme
classique en demande en moyenne $2{,}25$. C'est cohérent avec la formule : $\sin\theta=\frac12$,
$\theta=\frac\pi6$ et $\sin(3\theta)=1$. Pour un autre élément marqué, on change l'oracle (un
CZ entouré de portes $X$) et le résultat est le même.

\begin{figure}[!ht]
\centering
\adjustbox{max width=\linewidth}{%
\begin{tikzpicture}[>=Stealth, font=\small]
  \draw (0,1.4) -- (10.3,1.4);
  \draw (0,0) -- (10.3,0);
  \node[anchor=east] at (0,1.4) {$\ket{0}$};
  \node[anchor=east] at (0,0) {$\ket{0}$};
  \foreach \y in {0,1.4} {
    \node[porte] at (1.1,\y) {$H$};
    \node[porte] at (4.3,\y) {$H$};
    \node[porte] at (5.7,\y) {$Z$};
    \node[porte] at (8.5,\y) {$H$};
    \pic at (9.7,\y) {mesure};
  }
  \foreach \x in {2.7,7.1} {
    \draw[insarouge, thick] (\x,1.4) -- (\x,0);
    \pic at (\x,1.4) {ctl};
    \pic at (\x,0) {ctl};
  }
  \foreach \x/\lab in {1.9/{\ket{s}}, 3.5/{O_w\ket{s}}, 9.1/{\ket{11}}} {
    \draw[dashed, insagris!70] (\x,2.0) -- (\x,-0.6);
    \node[anchor=south, font=\footnotesize] at (\x,2.0) {$\lab$};
  }
  \draw[decorate, decoration={brace, amplitude=4pt, mirror}, insagris] (2.3,-0.8) -- (3.1,-0.8)
    node[midway, below=5pt, font=\footnotesize] {oracle $O_w$};
  \draw[decorate, decoration={brace, amplitude=4pt, mirror}, insagris] (3.9,-0.8) -- (9.0,-0.8)
    node[midway, below=5pt, font=\footnotesize] {diffusion $D=(H\otimes H)\,\mathrm{CZ}\,(Z\otimes Z)\,(H\otimes H)$};
\end{tikzpicture}}
\caption{Algorithme de Grover pour $N=4$ et $w=11$ : une itération suffit. Les deux points
reliés par un trait vertical forment la porte CZ.}
\label{fig:grover-circ}
\end{figure}

\paragraph{Ordres de grandeur.} Pour $N=2^{20}\approx10^6$, Grover demande
$\frac\pi4\cdot2^{10}\approx804$ itérations, contre $5\times10^5$ évaluations en moyenne. Pour une
recherche de clé AES-128 ($N=2^{128}$), il faut $\frac\pi4\,2^{64}\approx1{,}4\times10^{19}$
itérations \emph{séquentielles}, dont chacune contient un circuit AES complet ; et partager le travail
entre $m$ machines ne fait gagner qu'un facteur $\sqrt m$. L'effet est tout de même de diviser par deux le
nombre de bits de sécurité d'un chiffrement symétrique : on se protège en doublant la taille des
clés (AES-256).

% ---------------------------------------------------------------------
\subsection{Transformée de Fourier quantique et algorithme de Shor}\label{sec:shor}
% ---------------------------------------------------------------------

Le chiffrement RSA repose sur la difficulté de factoriser un entier $N$ (section~\ref{sec:crypto}).
L'algorithme de Shor, conçu en 1994 et publié en 1997 \parencite{shor1997}, factorise $N$ en un temps polynomial en $\log N$.
L'algorithme a deux étages : une réduction \emph{classique} de la factorisation à la recherche d'une
période, puis une recherche \emph{quantique} de cette période par transformée de Fourier.

\subsubsection{Réduction : de la factorisation à la période}

Soit $N$ impair et non puissance d'un nombre premier. On tire $a$ avec $1<a<N$. Si
$\operatorname{pgcd}(a,N)>1$, on tient déjà un facteur. Sinon, la suite $x\mapsto a^x\bmod N$ est
périodique : sa période $r$ est le plus petit entier $r\geq1$ tel que $a^r\equiv1\pmod N$ (l'\emph{ordre}
de $a$ modulo $N$).

\begin{resultat}{Factorisation à partir de l'ordre}
Si $r$ est pair et si $a^{r/2}\not\equiv-1\pmod N$, alors $\operatorname{pgcd}(a^{r/2}-1,N)$ et
$\operatorname{pgcd}(a^{r/2}+1,N)$ sont deux facteurs non triviaux de $N$. Pour $a$ tiré au hasard,
ces deux conditions sont réunies avec une probabilité au moins égale à $\frac12$ \parencite[chap.~5]{nielsen2010}.
\end{resultat}

En effet, posons $u=a^{r/2}$. Alors $(u-1)(u+1)=a^r-1\equiv0\pmod N$ : $N$ divise le produit. Il ne
divise pas $u-1$, sinon l'ordre de $a$ serait $r/2$, ni $u+1$ par hypothèse. Si $\operatorname{pgcd}(u-1,N)$
valait $1$, $N$ diviserait $u+1$ (lemme de Gauss) ; s'il valait $N$, $N$ diviserait $u-1$. Les deux sont exclus, donc
le pgcd est un diviseur strict de $N$ ; de même pour $u+1$. Tout se ramène donc à \emph{trouver $r$}. Classiquement, c'est
aussi difficile que factoriser. C'est la seule étape où l'ordinateur quantique intervient.

\subsubsection{La transformée de Fourier quantique}

\begin{definition}{Transformée de Fourier quantique (QFT)}
Sur l'espace de dimension $M$ de base $\ket{0},\ldots,\ket{M-1}$, avec $\omega=\ee^{2\ii\pi/M}$ :
\[
  \mathrm{QFT}_M\ket{x}=\frac{1}{\sqrt M}\sum_{y=0}^{M-1}\omega^{xy}\ket{y}.
\]
C'est une porte unitaire : ses colonnes sont orthonormées (somme géométrique de raison $\omega^{x-x'}$).
Pour $M=2^m$, elle s'implémente avec $O(m^2)$ portes élémentaires, alors que la transformée de Fourier
rapide classique demande $O(m\,2^m)$ opérations sur les $2^m$ amplitudes.
\end{definition}

Pour $M=2$ on a $\omega=-1$ et $\mathrm{QFT}_2=H$. Pour $M=4$, $\omega=\ii$ ; en écrivant $x=2x_1+x_2$ et
$y=2y_1+y_2$, on a $\ii^{xy}=(-1)^{x_1y_2+x_2y_1}\,\ii^{x_2y_2}$ (car $\ii^4=1$), d'où
\[
  \mathrm{QFT}_4\ket{x_1x_2}=\underbrace{\frac{\ket0+(-1)^{x_2}\ket1}{\sqrt2}}_{=\,H\ket{x_2}\ \text{(chiffre }y_1)}
  \otimes\ \frac{\ket0+(-1)^{x_1}\ii^{x_2}\ket1}{\sqrt2}\ \ \text{(chiffre }y_2).
\]
Le circuit $\mathrm{QFT}_4=\mathrm{SWAP}\,(I\otimes H)\,\mathrm{CS}\,(H\otimes I)$ de la
figure~\ref{fig:qft4} produit exactement cet état : $H$ sur le premier chiffre donne
$\frac{1}{\sqrt2}(\ket0+(-1)^{x_1}\ket1)$, la porte $\mathrm{CS}=\operatorname{diag}(1,1,1,\ii)$ ajoute le
facteur $\ii^{x_2}$ devant $\ket1$, $I\otimes H$ transforme $\ket{x_2}$ en $H\ket{x_2}$, et le SWAP remet les
chiffres dans l'ordre.

\begin{figure}[!ht]
\centering
\begin{tikzpicture}[>=Stealth, font=\small]
  \draw (0,1.2) -- (7.0,1.2);
  \draw (0,0) -- (7.0,0);
  \node[anchor=east] at (0,1.2) {$\ket{x_2}$};
  \node[anchor=east] at (0,0) {$\ket{x_1}$};
  \node[porte] at (1.2,0) {$H$};
  \draw[insarouge, thick] (2.7,1.2) -- (2.7,0);
  \pic at (2.7,1.2) {ctl};
  \node[porte] at (2.7,0) {$S$};
  \node[porte] at (4.2,1.2) {$H$};
  \draw[insarouge, thick] (5.7,1.2) -- (5.7,0);
  \pic at (5.7,1.2) {croix};
  \pic at (5.7,0) {croix};
  \node[anchor=west] at (7.0,1.2) {$\ket{y_2}$};
  \node[anchor=west] at (7.0,0) {$\ket{y_1}$};
\end{tikzpicture}
\caption{$\mathrm{QFT}_4$ : $H$, $S$ contrôlée par l'autre qubit, $H$, puis échange des deux fils.
Le fil du haut est le dernier chiffre.}
\label{fig:qft4}
\end{figure}

\begin{resultat}{La QFT transforme une période en pics}
Si $r$ divise $M$ et si $\ket{\psi}=\sqrt{r/M}\sum_{k=0}^{M/r-1}\ket{x_0+kr}$ (état périodique, de période
$r$), alors
\[
  \mathrm{QFT}_M\ket{\psi}=\frac{1}{\sqrt r}\sum_{j=0}^{r-1}\omega^{x_0\,jM/r}\,\ket{jM/r}.
\]
La mesure donne $y=jM/r$ avec la probabilité $\frac1r$, pour $j=0,\ldots,r-1$, quel que soit le décalage $x_0$.
\end{resultat}

Preuve : l'amplitude de $\ket{y}$ est $\frac{1}{\sqrt M}\sqrt{\frac rM}\,\omega^{x_0y}\sum_{k}\omega^{kry}$.
La somme géométrique de raison $\omega^{ry}=\ee^{2\ii\pi ry/M}$ vaut $M/r$ si $ry/M$ est entier, et $0$ sinon
(la raison est alors différente de $1$ et sa puissance $M/r$ vaut $1$). L'amplitude vaut donc
$\sqrt{r}/M\cdot M/r=1/\sqrt r$ aux multiples de $M/r$, et $0$ ailleurs. Le décalage $x_0$ n'apparaît que dans une
\emph{phase}, qui disparaît à la mesure : c'est ce qui permet de lire la période sans connaître le point de départ.

\subsubsection{L'algorithme}

\begin{resultat}{Algorithme de Shor : recherche de l'ordre de $a$ modulo $N$}
On choisit $M=2^t\geq N^2$ (soit $t\approx2n$ avec $n=\lceil\log_2N\rceil$), un registre de contrôle de $t$ qubits et un
registre de travail de $n$ qubits.
\begin{enumerate}[nosep]
  \item Préparer $\ket{0}^{\otimes t}\ket{1}$ et appliquer $H^{\otimes t}$ au registre de contrôle :
        $\ket{\varphi_1}=\frac{1}{\sqrt M}\sum_{x=0}^{M-1}\ket{x}\ket{1}$.
  \item Appliquer l'exponentiation modulaire $U_a\ket{x}\ket{z}=\ket{x}\ket{a^xz\bmod N}$ :
        $\ket{\varphi_2}=\frac{1}{\sqrt M}\sum_{x}\ket{x}\ket{a^x\bmod N}$.
  \item Appliquer $\mathrm{QFT}_M$ au registre de contrôle et le mesurer : on obtient $y$.
  \item (Classique) Développer $y/M$ en fraction continue : on y lit une fraction $j/r'$ proche de
        $y/M$ avec $r'<N$. Si $\operatorname{pgcd}(j,r)=1$, alors $r'=r$ et l'on vérifie $a^{r'}\equiv1\pmod N$.
        Sinon, $r'$ est un diviseur strict de $r$ et $a^{r'}\not\equiv1$ : on recommence.
\end{enumerate}
\end{resultat}

L'oracle $U_a$ est un produit de $t$ multiplications modulaires \emph{contrôlées} par les qubits $x_j$ de $x$ :
comme $a^x=\prod_j(a^{2^j})^{x_j}$, la $j$-ème porte multiplie le registre de travail par $a^{2^j}\bmod N$, un
nombre calculé à l'avance par un ordinateur classique. Le circuit est celui de la
figure~\ref{fig:shor-circ}.

\begin{figure}[!ht]
\centering
\begin{tikzpicture}[>=Stealth, font=\small]
  % registre de contrôle (haut) et registre de travail (bas)
  \draw (0,1.7) -- (8.3,1.7);
  \draw (0,0) -- (4.6,0);
  \draw[dashed] (4.6,0) -- (8.3,0);
  \node[anchor=east] at (0,1.7) {$\ket{0}^{\otimes t}$};
  \node[anchor=east] at (0,0) {$\ket{1}$};
  \draw (0.5,1.6) -- (0.7,1.8); \node[font=\scriptsize, anchor=south] at (0.6,1.78) {$t$};
  \draw (0.5,-0.1) -- (0.7,0.1); \node[font=\scriptsize, anchor=south] at (0.6,0.08) {$n$};
  \node[porte] at (1.5,1.7) {$H^{\otimes t}$};
  \pic at (3.3,1.7) {ctl};
  \draw[insarouge, thick] (3.3,1.7) -- (3.3,0.4);
  \node[draw=insarouge, fill=insarouge!8, thick, minimum width=17mm, minimum height=8mm] at (3.3,0) {$U_a$};
  \node[porte] at (5.4,1.7) {$\mathrm{QFT}_M$};
  \pic at (7.0,1.7) {mesure};
  \node[anchor=west] at (7.5,1.7) {$y$};
  \node[anchor=north west, text=insagris, font=\footnotesize] at (4.7,-0.3) {registre de travail : ignoré};
  \foreach \x/\lab in {2.3/1, 4.3/2, 6.3/3}
    {\draw[dashed, insagris!70] (\x,2.35) -- (\x,-0.3);
     \node[anchor=south, font=\footnotesize] at (\x,2.35) {$\ket{\varphi_\lab}$};}
\end{tikzpicture}
\caption{Circuit de Shor. Les petits traits obliques indiquent un registre de plusieurs qubits ($t$ pour le
contrôle, $n$ pour le travail). $U_a\ket{x}\ket{z}=\ket{x}\ket{a^xz\bmod N}$ est commandé par le registre de
contrôle.}
\label{fig:shor-circ}
\end{figure}

\paragraph{Pourquoi cela marche.} Après l'étape~2, mesurons (en pensée) le registre de travail : il
donne une valeur $f$ de la suite $a^x\bmod N$, et le registre de contrôle se projette sur les $x$ tels que
$a^x\equiv f$, c'est-à-dire $x\in\{x_0,x_0+r,x_0+2r,\ldots\}$ : un état périodique de période $r$. Le résultat sur la QFT
ci-dessus s'applique : la QFT concentre toute la probabilité sur les $y$ voisins de $jM/r$. La mesure du
registre de travail n'est pas nécessaire (les probabilités de $y$ sont les mêmes si on l'ignore), et le
décalage $x_0$ dépend de $f$ mais ne change pas la loi de $y$. Ainsi $y/M\approx j/r$ pour un entier $j$ \emph{aléatoire} :
on ne choisit pas $j$, mais il est premier avec $r$ dans une fraction $\Omega(1/\log\log r)$ des cas, et quelques
répétitions suffisent. Le choix $M\geq N^2$ garantit qu'une seule fraction de dénominateur inférieur à $N$ est
à distance de $\frac{1}{2M}$ de $y/M$ : le développement en fraction continue l'identifie.

\subsubsection{Exemple complet : $N=15$, $a=7$}

\paragraph{Période.} La suite $7^x\bmod15$ vaut $1,7,4,13,1,7,4,13,\ldots$ : l'ordre est $r=4$. Les
multiplicateurs de $U_a$ sont $a^{2^j}\bmod 15=7,\ 4,\ 1,\ 1$ pour $j=0,1,2,3$ (les deux derniers sont triviaux).
On prend ici $t=4$, soit $M=16$. Ce choix est trop court en général ($M\geq N^2=225$ serait exigé) mais suffit
car $r=4$ divise $M$, ce qui rend le calcul exact.

\paragraph{État avant la QFT.} $\ket{\varphi_2}=\frac14\sum_{x=0}^{15}\ket{x}\ket{7^x\bmod15}$. Si le
registre de travail est mesuré et donne $1$, le registre de contrôle est dans $\frac12(\ket0+\ket4+\ket8+\ket{12})$ ;
s'il donne $7$, dans $\frac12(\ket1+\ket5+\ket9+\ket{13})$, etc.

\paragraph{QFT.} On a $\omega=\ee^{2\ii\pi/16}$ et $\omega^4=\ii$. Pour le décalage $x_0=0$ :
\[
  \mathrm{QFT}_{16}\,\tfrac12\big(\ket0+\ket4+\ket8+\ket{12}\big)
  =\frac18\sum_{y=0}^{15}\Big(\sum_{k=0}^{3}\ii^{\,ky}\Big)\ket{y}
  =\frac12\big(\ket0+\ket4+\ket8+\ket{12}\big),
\]
car $\sum_{k=0}^{3}\ii^{\,ky}$ vaut $4$ si $y$ est multiple de $4$ et $0$ sinon. Pour les autres valeurs de
$f$, seules les phases changent. Dans tous les cas la mesure de $y$ donne $0,4,8$ ou $12$, chacun avec la probabilité $\frac14$
(figure~\ref{fig:shor-n15}).

\paragraph{Lecture classique.} On calcule $y/16$ :
\begin{center}
\renewcommand{\arraystretch}{1.2}
\begin{tabular}{@{}cccp{7.3cm}@{}}
\toprule
$y$ & $y/16$ & fraction irréductible & conclusion\\
\midrule
$0$ & $0$ & $0/1$ & inutile ($j=0$) : on recommence\\
$4$ & $\frac14$ & $1/4$ & $r=4$\ ; $7^4=2401\equiv1$ : correct\\
$8$ & $\frac12$ & $1/2$ & dénominateur $2$ : $7^2=49\equiv4\not\equiv1$, rejeté, on recommence\\
$12$ & $\frac34$ & $3/4$ & $r=4$ : correct\\
\bottomrule
\end{tabular}
\end{center}
Un essai réussit donc avec la probabilité $\frac12$ ($y=4$ ou $12$). Avec $r=4$ pair et $7^{2}=49\equiv4\not\equiv-1\pmod{15}$ :
\[
  \operatorname{pgcd}(7^2-1,15)=\operatorname{pgcd}(48,15)=3,\qquad
  \operatorname{pgcd}(7^2+1,15)=\operatorname{pgcd}(50,15)=5,\qquad 15=3\times5.
\]

\begin{figure}[!ht]
\centering
\adjustbox{max width=\linewidth}{%
\begin{tikzpicture}[>=Stealth, font=\small]
  % (a) suite 7^x mod 15, x = 0..15
  \draw[->, insagris] (-0.2,0) -- (6.7,0) node[right] {$x$};
  \draw[->, insagris] (0,-0.1) -- (0,2.6);
  \foreach \x/\v in {0/1,1/7,2/4,3/13,4/1,5/7,6/4,7/13,8/1,9/7,10/4,11/13,12/1,13/7,14/4,15/13}
    \fill[insarouge] ({\x*0.42},{\v*0.17}) circle (1.7pt);
  \foreach \v in {1,4,7,13}
    {\draw[insagris] (-0.05,{\v*0.17}) -- (0.05,{\v*0.17});
     \node[anchor=east, font=\scriptsize] at (-0.08,{\v*0.17}) {$\v$};}
  \foreach \x in {0,4,8,12}
    {\draw[insagris] ({\x*0.42},-0.05) -- ({\x*0.42},0.05);
     \node[anchor=north, font=\scriptsize] at ({\x*0.42},-0.07) {$\x$};}
  \draw[<->, insarouge] (0,2.45) -- (1.68,2.45) node[midway, above, font=\footnotesize] {période $r=4$};
  \node[font=\footnotesize, anchor=south] at (3.2,3.15) {(a) $f(x)=7^x\bmod15$};
  % flèche QFT
  \draw[->, very thick, insarouge] (7.0,0.8) -- (8.1,0.8) node[midway, above, font=\footnotesize] {$\mathrm{QFT}_{16}$};
  % (b) loi de y : P(y)=1/4 pour y = 0, 4, 8, 12
  \draw[->, insagris] (8.5,0) -- (14.9,0) node[right] {$y$};
  \draw[->, insagris] (8.6,-0.1) -- (8.6,2.6);
  \node[font=\footnotesize, anchor=south] at (11.7,3.15) {(b) loi de la mesure $y$ : $P(y)$};
  \foreach \y in {0,4,8,12}
    {\fill[insarouge] ({8.6+\y*0.4-0.1},0) rectangle ({8.6+\y*0.4+0.1},1.5);
     \node[anchor=north, font=\scriptsize] at ({8.6+\y*0.4},-0.07) {$\y$};}
  \draw[dashed, insagris] (8.55,1.5) -- (14.2,1.5);
  \node[anchor=east, font=\scriptsize] at (8.5,1.5) {$\frac14$};
  \node[text=insagris, font=\footnotesize, align=center] at (11.7,2.35) {$y/16=\frac{j}{4}$, $j=0,1,2,3$};
\end{tikzpicture}}
\caption{Shor pour $N=15$, $a=7$, $t=4$. À gauche, la suite $7^x\bmod15$ a la période $4$. À droite,
la mesure du registre de contrôle après la QFT ne donne que $y\in\{0,4,8,12\}$, avec la probabilité
$\frac14$ chacun : la période $r=4$ se lit dans l'écart $M/r=4$ entre les pics.}
\label{fig:shor-n15}
\end{figure}

\subsubsection{Ce que coûte l'algorithme}

Notons $n=\lceil\log_2N\rceil$ le nombre de bits de $N$ ($n=2048$ pour un module RSA-2048). Le coût s'exprime en
fonction de $n$, pas de $N$.
\begin{align*}
  \text{meilleur algorithme classique connu (crible algébrique, GNFS) :}&\quad
    \exp\!\Big(O\big(n^{1/3}(\log n)^{2/3}\big)\Big),\\
  \text{algorithme de Shor :}&\quad O(n^3)\ \text{portes élémentaires.}
\end{align*}
Le premier temps est sous-exponentiel en $n$ \parencite{lenstra1993}, mais supérieur à tout polynôme. Le second est un
polynôme en $n$ avec la multiplication scolaire ($O(n^2\log n\log\log n)$ avec une multiplication rapide).

\begin{remarque}[Où est le « log » ?]
On résume parfois « de $O(N^p)$ à $O(\log N)$ ». La comparaison correcte est la suivante. La division
d'essai coûte $O(\sqrt N)=O(2^{n/2})$ : exponentielle en $n$. Le crible algébrique est sous-exponentiel en $n$. Shor est
polynomial en $\log N=n\ln2$, soit $O((\log N)^3)$ et non $O(\log N)$ : rien que pour lire $N$ et multiplier
modulo $N$, il faut au moins $n$ opérations. Le « $\log$ » désigne donc la \emph{taille} du nombre.
\end{remarque}

Doubler la taille de la clé multiplie le coût de Shor par $8$, alors qu'il augmente de façon
spectaculaire celui de l'attaque classique. Le tableau compare le coût $n^3$ de Shor à la sécurité
classique que le NIST attribue aux modules RSA \parencite{nist80057}.

\begin{center}
\renewcommand{\arraystretch}{1.2}
\begin{tabular}{@{}rccc@{}}
\toprule
Module RSA $n$ (bits) & sécurité classique (NIST) & $n^3$ (portes, ordre de grandeur) & $\log_2(n^3)$\\
\midrule
$1024$ & $2^{80}$ & $1{,}1\times10^{9}$ & $30{,}0$\\
$2048$ & $2^{112}$ & $8{,}6\times10^{9}$ & $33{,}0$\\
$3072$ & $2^{128}$ & $2{,}9\times10^{10}$ & $34{,}8$\\
$7680$ & $2^{192}$ & $4{,}5\times10^{11}$ & $38{,}7$\\
$15360$ & $2^{256}$ & $3{,}6\times10^{12}$ & $41{,}7$\\
\bottomrule
\end{tabular}
\end{center}

Passer de $2048$ à $3072$ bits fait gagner un facteur $2^{16}=65\,536$ sur l'attaque classique, mais le coût de
Shor n'est multiplié que par $(3/2)^3\approx3{,}4$. Face à un ordinateur quantique, \emph{grossir
les clés RSA ne protège pas} : il faut changer d'algorithme (section~\ref{sec:crypto}). Cela dit, le nombre de portes
n'est pas le vrai coût : il faut des qubits \emph{corrigés d'erreurs}, et c'est là que se trouve le verrou technique
(section~\ref{sec:machine}). Les estimations les plus citées sont les suivantes.
\begin{itemize}[nosep]
  \item \textbf{RSA-2048} : environ $3n+0{,}002\,n\log_2n\approx6\,189$ qubits logiques et
        $0{,}3\,n^3+0{,}0005\,n^3\log_2n\approx2{,}6\times10^9$ portes de Toffoli ; avec des qubits physiques
        d'erreur $10^{-3}$ et un cycle de code de surface de $1\,\mu$s, $20$ millions de qubits physiques pendant
        $8$ heures \parencite{gidney2021}. Une optimisation de 2025 (arithmétique approchée, codes « yoked »,
        « magic state cultivation ») descend à \emph{moins d'un million} de qubits physiques, pour moins d'une
        semaine \parencite{gidney2025}.
  \item \textbf{Courbes elliptiques (ECC-256)}, utilisées dans TLS : au plus $2\,330$ qubits logiques
        \parencite{roetteler2017}. Avec des qubits de chat (section~\ref{sec:machine}), $126\,133$ qubits suffiraient
        pour $9$ heures de calcul \parencite{gouzien2023}. Une \emph{prépublication} de 2026, fondée sur des
        hypothèses optimistes pour les atomes neutres, évoque $10\,000$ à $26\,000$ atomes et quelques
        jours \parencite{cain2026}.
\end{itemize}
Ce sont des estimations, pas des mesures : les démonstrations expérimentales de $15=3\times5$ ou
$21=3\times7$ reposent sur des circuits simplifiés grâce à la connaissance préalable de la réponse
\parencite{smolin2013}. À notre connaissance, aucune n'a factorisé par l'algorithme de Shor un nombre hors de portée d'un
calcul à la main.

% ---------------------------------------------------------------------
\subsection{Cryptographie : RSA, TLS et « stocker maintenant, déchiffrer plus tard »}\label{sec:crypto}
% ---------------------------------------------------------------------

La sécurité des échanges sur Internet repose sur deux familles d'algorithmes. La cryptographie
\emph{symétrique} (AES, ChaCha20) chiffre les données avec une clé secrète partagée. La cryptographie
\emph{asymétrique} (RSA, Diffie--Hellman, courbes elliptiques) sert à \emph{établir} cette clé entre deux
personnes qui ne se sont jamais rencontrées, et à \emph{authentifier} l'interlocuteur par une signature. Les
algorithmes asymétriques s'appuient sur deux problèmes réputés difficiles : la factorisation d'un entier
(RSA) et le logarithme discret (Diffie--Hellman, courbes elliptiques). Ce sont deux problèmes de
\emph{période} dans un groupe : l'algorithme de Shor les résout tous les deux.

\subsubsection{RSA et ce que Shor en fait}

\begin{definition}{Chiffrement RSA}
Alice choisit deux nombres premiers $p\neq q$, pose $n=pq$ et $\varphi(n)=(p-1)(q-1)$, choisit $e$ premier avec
$\varphi(n)$ et calcule $d$ tel que $ed\equiv1\pmod{\varphi(n)}$. La \emph{clé publique} est $(n,e)$, la \emph{clé
privée} $d$. Pour un message $m<n$ :
\[
  c=m^e\bmod n\quad(\text{chiffrement}),\qquad m=c^d\bmod n\quad(\text{déchiffrement}).
\]
Le déchiffrement fonctionne parce que $m^{ed}\equiv m\pmod n$ (théorème d'Euler) \parencite{rivest1978}.
\end{definition}

\paragraph{Exemple jouet.} $p=3$, $q=11$, $n=33$, $\varphi(n)=20$, $e=3$ (premier avec $20$) et $d=7$ (car
$3\times7=21\equiv1\pmod{20}$). Pour $m=4$ : $c=4^3\bmod33=64-33=31$. Le déchiffrement donne
$31^7\bmod33$ ; comme $31\equiv-2$, on a $(-2)^7=-128\equiv-128+132=4=m$ : on retrouve le message.

\paragraph{Attaque.} L'attaquant connaît $(n,e)$. Il lui manque $\varphi(n)$ : le calculer équivaut à
factoriser $n$. Avec $n=33=3\times11$, il obtient $\varphi=20$ puis $d=e^{-1}\bmod20=7$ par l'algorithme d'Euclide
étendu (polynomial). Tout le secret de RSA est donc la difficulté de factoriser. Pour un module de $2048$ bits
(617 chiffres décimaux), les meilleures factorisations classiques publiques (RSA-240, $795$ bits, en 2019
\parencite{boudot2020}) sont loin du compte : on a vu à la section~\ref{sec:shor} que
l'algorithme de Shor l'attaque en $O(n^3)$ portes. Le même algorithme, appliqué au groupe d'une courbe elliptique,
résout le logarithme discret et casse Diffie--Hellman, ECDH et ECDSA.

\subsubsection{Ce que cassent Shor et Grover}

\begin{center}
\renewcommand{\arraystretch}{1.25}
\begin{tabular}{@{}>{\raggedright\arraybackslash}p{4.1cm}>{\raggedright\arraybackslash}p{2.9cm}>{\raggedright\arraybackslash}p{2.5cm}>{\raggedright\arraybackslash}p{5.2cm}@{}}
\toprule
\textbf{Primitive} & \textbf{Rôle} & \textbf{Attaque quantique} & \textbf{Conséquence}\\
\midrule
RSA, Diffie--Hellman, ECDH (X25519) & échange de clé, chiffrement & Shor & \textbf{cassés} (temps polynomial)\\
RSA, ECDSA, EdDSA & signature, certificats & Shor & \textbf{cassés}\\
AES, ChaCha20 & chiffrement symétrique & Grover & clé de $k$ bits : $2^{k/2}$ itérations ; AES-128 affaibli, \textbf{AES-256 sûr}\\
SHA-2, SHA-3 & hachage & Grover (pré-image) & SHA-256 : $2^{128}$, \textbf{sûr}\\
ML-KEM, ML-DSA, SLH-DSA & échange de clé, signature & aucune attaque efficace connue & normes NIST de 2024 (voir plus bas)\\
\bottomrule
\end{tabular}
\end{center}

Les trois premières lignes ne sont pas à égalité. Un algorithme asymétrique est \emph{détruit} par Shor : il faut
le remplacer. Un algorithme symétrique n'est que \emph{ralenti} par Grover (racine carrée, section~\ref{sec:grover}) :
il suffit de doubler la taille de la clé. « Aucune attaque connue » n'est pas une preuve de sécurité, d'où la prudence
des agences décrite plus bas.

\subsubsection{TLS 1.3 : où sont les points faibles ?}

Le protocole TLS 1.3 (RFC~8446 \parencite{rfc8446}) protège HTTPS. Il combine les trois mécanismes : un
\emph{échange de clés} éphémère (ECDHE, le plus souvent X25519 : chaque côté tire un secret $a$ ou $b$ et
publie $g^a$ ou $g^b$, d'où le secret partagé $g^{ab}$), une \emph{authentification} du serveur par signature
(certificat RSA ou ECDSA, vérifié par une chaîne d'autorités de certification, elle aussi faite de signatures),
et un \emph{chiffrement symétrique} des données (AES-GCM ou ChaCha20-Poly1305) avec des clés dérivées de $g^{ab}$
(figure~\ref{fig:tls}).

\begin{figure}[!ht]
\centering
\adjustbox{max width=\linewidth}{%
\begin{tikzpicture}[>=Stealth, font=\small,
    tag/.style={rounded corners=2pt, font=\footnotesize, inner sep=3pt, anchor=west, text width=4.6cm, align=left}]
  \node[draw=insarouge, fill=insarouge!8, thick, minimum width=20mm] (C) at (0,0) {Client};
  \node[draw=insarouge, fill=insarouge!8, thick, minimum width=20mm] (S) at (6.2,0) {Serveur};
  \draw[insagris, thick] (0,-0.4) -- (0,-6.2);
  \draw[insagris, thick] (6.2,-0.4) -- (6.2,-6.2);
  \draw[->, thick] (0,-1.0) -- (6.2,-1.0) node[midway, above, font=\footnotesize] {ClientHello + clé publique $g^a$};
  \draw[<-, thick] (0,-1.9) -- (6.2,-1.9) node[midway, above, font=\footnotesize] {ServerHello + clé publique $g^b$};
  \node[font=\footnotesize, text=insagris, anchor=east] at (-0.2,-2.35) {calcule $g^{ab}$, dérive les clés};
  \node[font=\footnotesize, text=insagris, anchor=west] at (6.4,-2.35) {calcule $g^{ab}$};
  \draw[<-, thick] (0,-3.2) -- (6.2,-3.2) node[midway, above, font=\footnotesize, align=center]
       {\{Certificat, signature (CertificateVerify),\\ Finished\}};
  \draw[->, thick] (0,-4.1) -- (6.2,-4.1) node[midway, above, font=\footnotesize] {\{Finished\}};
  \draw[<->, very thick, insarouge] (0,-5.2) -- (6.2,-5.2) node[midway, above, font=\footnotesize] {données chiffrées (AES-GCM, ChaCha20)};
  \node[tag, fill=insarouge!18] at (7.8,-1.45) {\textbf{Échange de clés} : cassé par Shor (logarithme discret)};
  \node[tag, fill=insarouge!18] at (7.8,-3.2) {\textbf{Signature et certificats} : cassés par Shor (RSA, ECDSA)};
  \node[tag, fill=insagris!20] at (7.8,-5.2) {\textbf{Chiffrement symétrique} : \mbox{Grover} seulement, AES-256 suffit};
\end{tikzpicture}}
\caption{Poignée de main TLS 1.3 simplifiée. Les accolades \{\,\} indiquent des messages déjà chiffrés avec les clés
issues de $g^{ab}$. Deux des trois briques tombent face à Shor ; la troisième résiste.}
\label{fig:tls}
\end{figure}

\begin{remarque}[La confidentialité persistante ne protège pas contre l'ordinateur quantique]
Les secrets $a$ et $b$ sont détruits après la session : si la clé privée du serveur est volée plus tard, les sessions
passées ne peuvent pas être déchiffrées (propriété dite de \emph{forward secrecy}). Mais les valeurs publiques
$g^a$ et $g^b$ ont circulé en clair. Un adversaire qui les a enregistrées peut, avec l'algorithme de Shor (logarithme
discret sur la courbe), retrouver $a$ puis $g^{ab}$, donc toutes les clés de la session. La forward secrecy protège
contre le vol d'une \emph{clé}, pas contre l'effondrement de la difficulté du \emph{problème mathématique}.
\end{remarque}

\subsubsection{« Stocker maintenant, déchiffrer plus tard »}

La menace n'attend pas l'existence de la machine. Un adversaire qui capte le trafic chiffré aujourd'hui (il lui
suffit d'écouter une liaison, ce qui est discret et peu coûteux) peut le \emph{conserver} et attendre qu'un
ordinateur quantique assez puissant existe pour le déchiffrer. C'est l'attaque \emph{Harvest Now, Decrypt Later}
(HNDL, ou « stocker maintenant, déchiffrer plus tard »). Elle vise la \emph{confidentialité de long terme} :
dossiers médicaux et génétiques, secrets industriels, secrets d'État, archives juridiques. Elle ne vise pas un
jeton de session valable dix minutes. On ne sait pas qui enregistre quoi ; il faut supposer que cela se fait.

\begin{resultat}{Inégalité de Mosca}
Soit $X$ la durée pendant laquelle les données doivent rester secrètes, $Y$ le temps nécessaire pour migrer vers une
cryptographie résistante, et $Z$ le délai avant qu'un ordinateur quantique capable de casser la cryptographie actuelle existe
(tous en années \parencite{mosca2018}). Les données sont \emph{déjà exposées} si
\[
  X+Y>Z.
\]
\end{resultat}

Les données chiffrées juste avant la fin de la migration (à la date $Y$) doivent encore être secrètes à la date $X+Y$ ;
si $Z$ tombe avant, l'adversaire les lit. Exemple (figure~\ref{fig:mosca}) : des données médicales à protéger
$X=20$ ans, une migration de $Y=7$ ans, et $Z=15$ ans : $X+Y=27>15$, avec douze années d'exposition. On ne connaît pas
$Z$ (ce chapitre ne cherche pas à le prédire), mais on connaît $X$ et $Y$ : plus ils sont grands, plus il faut
commencer tôt, et le coût d'une migration trop précoce est faible devant celui d'une divulgation irréversible.

\begin{figure}[!ht]
\centering
\begin{tikzpicture}[>=Stealth, font=\small, x=0.3cm]
  \draw[->, insagris] (0,0) -- (31,0) node[right] {années};
  \foreach \t in {0,5,...,30}
    {\draw[insagris] (\t,0.08) -- (\t,-0.08);
     \node[anchor=north, font=\scriptsize] at (\t,-0.1) {$\t$};}
  \fill[insagris!35] (0,0.5) rectangle (7,1.1);
  \fill[insarouge!18] (7,0.5) rectangle (27,1.1);
  \fill[insarouge!60] (15,0.5) rectangle (27,1.1);
  \node[font=\footnotesize] at (3.5,0.8) {$Y=7$};
  \node[font=\footnotesize] at (11,0.8) {$X=20$};
  \draw[insarouge, very thick] (15,0) -- (15,2.0);
  \node[anchor=south, font=\footnotesize, text=insarouge] at (15,2.0) {$Z=15$ : ordinateur quantique opérationnel};
  \node[anchor=west, font=\footnotesize, text=white] at (15.4,0.8) {exposé 12 ans};
  \node[anchor=north west, font=\footnotesize, text=insagris] at (0,-0.65)
    {$Y$ : migration (la cryptographie vulnérable est encore utilisée)};
  \node[anchor=north west, font=\footnotesize, text=insagris] at (0,-1.15)
    {$X$ : durée de confidentialité exigée pour les données chiffrées à la fin de la migration};
\end{tikzpicture}
\caption{Inégalité de Mosca pour $X=20$, $Y=7$, $Z=15$ ans. La zone rouge foncé est la période où des données
interceptées sont lisibles par l'adversaire alors qu'elles devraient rester secrètes.}
\label{fig:mosca}
\end{figure}

\begin{remarque}[Signatures : un risque différent]
L'attaque HNDL ne concerne que la confidentialité. Une signature ne sert qu'au moment où on la vérifie : falsifier en 2035 un
certificat TLS expiré ne change rien à une session de 2026. Les signatures à longue durée de vie (micrologiciels,
racines de confiance, documents juridiques) devront néanmoins migrer, mais plus tard : c'est l'échange de clés qui est
urgent.
\end{remarque}

\subsubsection{Première réponse : la cryptographie post-quantique}

La \emph{cryptographie post-quantique} (PQC) regroupe des algorithmes qui s'exécutent sur des ordinateurs
\emph{classiques} mais reposent sur des problèmes contre lesquels on ne connaît aucune attaque quantique efficace :
réseaux euclidiens (problème \emph{Learning With Errors}), fonctions de hachage, codes correcteurs d'erreurs. Le NIST a
publié le 13 août 2024 trois normes \parencite{fips203,fips204,fips205}, puis sélectionné un quatrième algorithme en 2025 \parencite{nist2025hqc}.

\begin{center}
\renewcommand{\arraystretch}{1.2}
\begin{tabular}{@{}llll@{}}
\toprule
\textbf{Norme} & \textbf{Nom} & \textbf{Usage} & \textbf{Problème sous-jacent}\\
\midrule
FIPS 203 & ML-KEM (ex-Kyber) & échange de clé & réseaux (module-LWE)\\
FIPS 204 & ML-DSA (ex-Dilithium) & signature & réseaux\\
FIPS 205 & SLH-DSA (ex-SPHINCS+) & signature & fonctions de hachage\\
(à venir) & HQC & échange de clé (second choix) & codes correcteurs\\
\bottomrule
\end{tabular}
\end{center}

\paragraph{Hybridation.} Dans TLS 1.3, la solution déployée est \emph{hybride} : le client envoie à la fois une clé
publique X25519 et une clé publique ML-KEM-768, et les deux secrets sont concaténés dans la dérivation des clés. L'adversaire
doit casser \emph{les deux} : X25519 (que Shor casse, mais que l'on connaît bien) et ML-KEM (que Shor ne casse pas, mais
que l'on connaît depuis moins longtemps). Le groupe \texttt{X25519MLKEM768} est normalisé par la RFC~10024 d'août 2026
\parencite{rfc10024}, et il est activé par défaut dans les navigateurs depuis fin 2024 (Chrome 131, Firefox 132). Le
prix est la taille : $32$ octets pour une clé X25519, mais $1\,184$ octets de clé publique et $1\,088$ octets de
chiffré pour ML-KEM-768.

\paragraph{Pourquoi pas ML-KEM seul ?} Les algorithmes post-quantiques sont plus jeunes et moins étudiés. En 2022, le
candidat SIKE, arrivé au quatrième tour du NIST, a été cassé par une \emph{attaque classique} : une dizaine de minutes sur un seul
cœur d'ordinateur suffisent pour son plus petit jeu de paramètres \parencite{castryck2023}. L'ANSSI recommande donc de \emph{combiner} cryptographie classique et
post-quantique pendant la transition \parencite{anssi2024}.

\paragraph{Calendrier.} Le projet de document NIST~IR~8547 (version initiale publique de novembre 2024, qui peut
avoir évolué) prévoit de déprécier RSA-2048 et les algorithmes de sécurité équivalente après 2030, et d'interdire tous les
algorithmes vulnérables après 2035 \parencite{nistir8547}. La migration est longue : il faut inventorier où la cryptographie est
utilisée (logiciels, matériel, certificats, protocoles), puis la remplacer. L'« agilité cryptographique », c'est-à-dire
la capacité de changer d'algorithme sans refondre un système, est la qualité qui coûtera le moins cher.

\subsubsection{Seconde réponse : la distribution quantique de clé}

La \emph{distribution quantique de clé} (QKD) ne cherche pas un problème difficile : elle utilise la physique. Dans le
protocole BB84 \parencite{bennett1984}, Alice envoie à Bob des qubits dont l'état est choisi au hasard parmi quatre :
$\ket0,\ket1$ (base $Z$) ou $\ket+,\ket-$ (base $X$). Bob mesure chaque qubit dans une base $Z$ ou $X$ choisie au hasard.
Ensuite, Alice et Bob annoncent \emph{publiquement leurs bases} (pas les bits) et gardent les positions où les bases
coïncident : si la base est la bonne, Bob lit exactement le bit d'Alice ; sinon son résultat est aléatoire. Exemple, sans espion :

\begin{center}
\renewcommand{\arraystretch}{1.15}
\begin{tabular}{@{}l*{8}{c}@{}}
\toprule
position & 1 & 2 & 3 & 4 & 5 & 6 & 7 & 8\\
\midrule
bit d'Alice & 0 & 1 & 1 & 0 & 1 & 0 & 0 & 1\\
base d'Alice & $Z$ & $X$ & $X$ & $Z$ & $X$ & $Z$ & $X$ & $Z$\\
base de Bob & $Z$ & $Z$ & $X$ & $X$ & $X$ & $Z$ & $Z$ & $Z$\\
\midrule
bases identiques ? & oui & non & oui & non & oui & oui & non & oui\\
clé tamisée & 0 & & 1 & & 1 & 0 & & 1\\
\bottomrule
\end{tabular}
\end{center}

La clé tamisée est $01101$ : en moyenne la moitié des qubits est conservée. Pour détecter un espion Ève, Alice et Bob
comparent publiquement un échantillon de la clé tamisée. Une attaque « intercepter puis renvoyer » est détectée parce que
\emph{mesurer perturbe} (postulat de la mesure) et que l'on ne peut pas copier un état inconnu pour mesurer la copie
(théorème de non-clonage, section~\ref{sec:teleportation}). Ève choisit la mauvaise base une fois sur deux ; dans ce cas
le qubit renvoyé est dans la mauvaise base, et Bob, s'il a la bonne base, se trompe une fois sur deux. Le taux d'erreur de la clé tamisée
(\emph{QBER}) vaut donc
\[
  \mathrm{QBER}=\underbrace{\tfrac12}_{\text{Ève se trompe de base}}\times\underbrace{\tfrac12}_{\text{Bob se trompe}}=\tfrac14=25\,\%.
\]
En pratique, quelques pour cent d'erreurs viennent du bruit ; au-delà d'un seuil de l'ordre de $11\,\%$, on abandonne la clé. La
variante E91 \parencite{ekert1991} remplace l'envoi de qubits par des paires intriquées et teste l'absence d'espion par la violation
de l'inégalité CHSH (section~\ref{sec:chsh}).

\begin{remarque}[Pourquoi les agences recommandent la PQC plutôt que la QKD]
En janvier 2024, l'ANSSI, le BSI allemand et leurs homologues néerlandais et suédois ont publié une position commune
\parencite{anssi2024qkd} : la QKD n'est pas assez mûre pour la sécurité, ne convient qu'à des cas d'usage restreints, et la priorité est donnée à
une technologie « beaucoup plus compatible avec les systèmes cryptographiques existants ». Elle demande des liaisons optiques
dédiées de portée limitée, des nœuds de confiance, une authentification classique des messages publics (qui nécessite
elle-même de la cryptographie), et ses implémentations ont des failles de matériel. La QKD protège donc un lien, la PQC un
protocole entier.
\end{remarque}

% ---------------------------------------------------------------------
\subsection{Construire un ordinateur quantique}\label{sec:machine}
% ---------------------------------------------------------------------

Les sections précédentes supposent des qubits parfaits. Les qubits réels sont bruités, et c'est le bruit, bien plus que le
nombre de portes, qui sépare les machines actuelles des estimations de la section~\ref{sec:shor}. Cette section
explique ce qu'il faut pour fabriquer un qubit, pourquoi certains tournent dans un réfrigérateur à
$10\ \mathrm{mK}$, comment on corrige les erreurs, et ce que proposent les qubits de chat d'Alice \& Bob.

\subsubsection{Que faut-il ? Les critères de DiVincenzo}

\textcite{divincenzo2000} a résumé les conditions nécessaires au calcul quantique :
\begin{enumerate}[nosep]
  \item des qubits bien définis, en nombre extensible ;
  \item une initialisation fiable dans l'état $\ket{0\ldots0}$ ;
  \item une cohérence longue devant la durée d'une porte ;
  \item un jeu universel de portes : portes à un qubit et au moins une porte à deux qubits comme le CNOT (chapitre~\ref{chap:portes}) ;
  \item une mesure individuelle des qubits.
\end{enumerate}
S'y ajoutent, pour les réseaux, deux critères de communication. Le critère~3 (cohérence) est le plus difficile à tenir quand le nombre de qubits augmente.

\subsubsection{La décohérence}

Un qubit n'est jamais seul. Supposons un qubit $S$ dans l'état $\frac{1}{\sqrt2}(\ket0+\ket1)$ qui interagit
avec son environnement $E$ de façon que celui-ci « enregistre » l'état de $S$ : $\ket0\ket{E_0}$ et $\ket1\ket{E_1}$.
L'état global est $\frac{1}{\sqrt2}(\ket0\ket{E_0}+\ket1\ket{E_1})$, intriqué. L'environnement n'est pas observé : l'état de $S$ seul
est obtenu par trace partielle (section~\ref{sec:trace}),
\[
  \rho_S=\frac12\begin{pmatrix}1&\braket{E_1}{E_0}\\\braket{E_0}{E_1}&1\end{pmatrix}.
\]
Si les deux états de l'environnement sont orthogonaux, $\rho_S=\frac12 I$ : l'état est maximalement mixte
(section~\ref{sec:pur-mixte}) et la superposition est perdue. En pratique le recouvrement décroît
comme $\braket{E_1}{E_0}=\ee^{-t/T_2}$, d'où, pour la pureté,
\[
  \operatorname{Tr}\rho_S^2=\frac{1+\ee^{-2t/T_2}}{2}\ \xrightarrow[t\to\infty]{}\ \frac12 .
\]
(à $t=T_2$, elle vaut $0{,}57$). Les termes diagonaux, eux, évoluent avec le temps de relaxation $T_1$ (perte d'énergie : $\ket1\to\ket0$). Le rapport
entre $T_2$ et la durée $t_{\mathrm{porte}}$ d'une porte fixe le nombre d'opérations avant que l'information soit perdue. Avec
$T_2=100\ \mu\mathrm{s}$ et des portes de $50\ \mathrm{ns}$, il vaut $2\,000$ : l'erreur par porte est de l'ordre de
$t_{\mathrm{porte}}/T_2=5\times10^{-4}$. Calculer sur des milliards de portes sans corriger les erreurs est
donc impossible.

\subsubsection{Les plateformes}

Plusieurs systèmes physiques servent de qubits. Chacun a ses forces et ses faiblesses ; les valeurs du tableau sont des ordres de
grandeur tirés de la littérature récente et évoluent vite.

\begin{center}
\renewcommand{\arraystretch}{1.2}
\small
\begin{tabular}{@{}>{\raggedright\arraybackslash}p{2.7cm}>{\raggedright\arraybackslash}p{2.2cm}>{\raggedright\arraybackslash}p{1.9cm}>{\raggedright\arraybackslash}p{1.9cm}>{\raggedright\arraybackslash}p{5.5cm}@{}}
\toprule
\textbf{Plateforme} & \textbf{Température} & \textbf{Porte à 2 qubits} & \textbf{Cohérence} & \textbf{Atouts / difficultés}\\
\midrule
Circuits supra\-conducteurs \parencite{kjaergaard2020} & $10$--$20\ \mathrm{mK}$ & $20$--$100\ \mathrm{ns}$ & $\sim0{,}1\ \mathrm{ms}$ & portes rapides, fabrication par lithographie ; cryogénie, câblage, voisinage limité\\
Ions piégés \parencite{bruzewicz2019} & ambiante (vide poussé) & $10\ \mu\mathrm{s}$--$1\ \mathrm{ms}$ & secondes à minutes & meilleures fidélités ($\geq99{,}9\,\%$), connexion totale ; portes lentes, passage à grande échelle\\
Atomes neutres \parencite{saffman2016} & ambiante (vide poussé) & $0{,}1$--$1\ \mu\mathrm{s}$ & $\sim1\ \mathrm{s}$ & réseaux de milliers d'atomes reconfigurables ; erreurs de porte plus élevées\\
Photons \parencite{kok2007} & ambiante (détecteurs $\sim2\ \mathrm{K}$) & $\mathrm{ns}$ & pas de décohérence en vol & communication ; pertes, portes à deux photons probabilistes\\
Spins dans le silicium \parencite{burkard2023} & $0{,}1$--$1\ \mathrm{K}$ & $10$--$100\ \mathrm{ns}$ & $\sim\mathrm{ms}$ & compatible avec la microélectronique ; peu de qubits à ce jour\\
\bottomrule
\end{tabular}
\end{center}

Aucune plateforme ne domine. Les machines actuelles comptent de quelques dizaines à quelques milliers de qubits physiques
\emph{bruités}, sans correction d'erreurs : \textcite{preskill2018} les a appelées « NISQ » (\emph{noisy intermediate-scale quantum}). Elles ne suffisent pas pour Shor.

\subsubsection{Le cryostat : pourquoi $10\ \mathrm{mK}$ ?}

Un transmon (qubit supraconducteur) a une fréquence de transition $f\approx5\ \mathrm{GHz}$ : une micro-onde. À l'équilibre thermique
à la température $T$, la loi de Boltzmann donne $p_1/p_0=\ee^{-hf/k_BT}$ avec $p_0+p_1=1$, d'où
\[
  p_1=\frac{1}{1+\ee^{\,hf/k_BT}},\qquad
  \frac{hf}{k_B}=\frac{6{,}626\times10^{-34}\times5\times10^{9}}{1{,}381\times10^{-23}}\approx0{,}24\ \mathrm{K}.
\]
Le qubit n'est « dans l'état $\ket0$ » que si $T\ll0{,}24\ \mathrm{K}$. Sinon il est dans un mélange statistique de
$\ket0$ et $\ket1$, de matrice densité $\operatorname{diag}(1-p_1,p_1)$ (section~\ref{sec:rho}) :

\begin{center}
\renewcommand{\arraystretch}{1.15}
\begin{tabular}{@{}rccc@{}}
\toprule
Température $T$ & $hf/k_BT$ & population $p_1$ de $\ket1$ & pureté $\operatorname{Tr}\rho^2$\\
\midrule
$300\ \mathrm{K}$ (ambiante) & $0{,}0008$ & $0{,}4998$ & $0{,}500$\\
$4\ \mathrm{K}$ (hélium liquide) & $0{,}060$ & $0{,}485$ & $0{,}500$\\
$1\ \mathrm{K}$ & $0{,}24$ & $0{,}44$ & $0{,}507$\\
$100\ \mathrm{mK}$ & $2{,}4$ & $0{,}083$ & $0{,}85$\\
$20\ \mathrm{mK}$ & $12$ & $6{,}2\times10^{-6}$ & $0{,}99999$\\
$10\ \mathrm{mK}$ & $24$ & $3{,}8\times10^{-11}$ & $1{,}000$\\
\bottomrule
\end{tabular}
\end{center}

À température ambiante, $p_1\approx\frac12$ : l'état est maximalement mixte (pureté minimale d'un qubit,
$\operatorname{Tr}\rho^2=\frac12$, section~\ref{sec:pur-mixte}) et le qubit ne porte aucune information. Il faut descendre vers $20\ \mathrm{mK}$ pour que
l'état initial soit pur à mieux que $10^{-5}$ près. S'y ajoute une seconde raison : les jonctions Josephson sont en aluminium, supraconducteur
seulement sous $\approx1{,}2\ \mathrm{K}$, et les photons thermiques des lignes de commande doivent être éteints. Les autres
plateformes évitent ce problème autrement : ions et atomes sont préparés dans $\ket0$ par pompage optique (un laser les force dans l'état voulu)
et refroidis par laser ; les photons ont des fréquences optiques ($\sim10^{14}\ \mathrm{Hz}$, soit $hf/k_B\approx5\times10^3\ \mathrm{K}$),
sans aucune excitation thermique.

\paragraph{Le réfrigérateur à dilution.} On atteint $10\ \mathrm{mK}$ avec un mélange d'isotopes d'hélium, $^3$He et $^4$He. Sous
$0{,}87\ \mathrm{K}$, le mélange se sépare en deux phases, l'une riche en $^3$He, l'autre diluée. Faire passer des atomes de $^3$He de
la phase riche à la phase diluée \emph{absorbe de la chaleur}, comme une évaporation : c'est le principe de la dilution. Un
refroidisseur à tube pulsé fournit le froid initial (étages à $50\ \mathrm{K}$ puis $4\ \mathrm{K}$), et le circuit de dilution
descend en étages jusqu'à la \emph{chambre de mélange} (figure~\ref{fig:cryostat}). Le prix est la \emph{puissance
froide} : de l'ordre de $10\ \mu\mathrm{W}$ à $20\ \mathrm{mK}$ pour un appareil commercial, ce qui limite la chaleur que les lignes et
l'électronique peuvent apporter à cet étage.

\begin{figure}[!ht]
\centering
\adjustbox{max width=\linewidth}{%
\begin{tikzpicture}[>=Stealth, font=\small,
    plaque/.style={fill=insagris!35, draw=insagris, thick, minimum height=3.5mm, inner sep=0pt},
    legende/.style={font=\footnotesize, anchor=west, text width=7.2cm, align=left}]
  % plaques : (hauteur, largeur, température)
  \node[plaque, minimum width=6.4cm] (p0) at (0,0)    {};
  \node[plaque, minimum width=5.6cm] (p1) at (0,-1.3) {};
  \node[plaque, minimum width=4.8cm] (p2) at (0,-2.6) {};
  \node[plaque, minimum width=4.0cm] (p3) at (0,-3.9) {};
  \node[plaque, minimum width=3.2cm] (p4) at (0,-5.2) {};
  \node[plaque, minimum width=2.4cm, fill=insarouge!30, draw=insarouge] (p5) at (0,-6.5) {};
  \foreach \a/\b in {0/-1.3, -1.3/-2.6, -2.6/-3.9, -3.9/-5.2, -5.2/-6.5}
    {\draw[insagris, thick] (-0.9,\a-0.175) -- (-0.9,\b+0.175);
     \draw[insagris, thick] (0.9,\a-0.175) -- (0.9,\b+0.175);}
  % puce
  \node[draw=insarouge, fill=insarouge!15, thick, minimum width=14mm, minimum height=5mm] (puce) at (0,-7.15) {\footnotesize puce};
  \draw[insarouge, thick] (0,-6.675) -- (0,-6.9);
  % ligne de commande (coaxial) avec atténuateurs
  \draw[insarouge, thick] (0.3,0.35) -- (0.3,-6.9);
  \foreach \y/\att in {-2.6/{$-20$ dB}, -5.2/{$-20$ dB}, -6.5/{$-20$ dB}}
    {\node[draw=insarouge, fill=white, inner sep=1.8pt] at (0.3,\y) {};
     \node[font=\scriptsize, text=insarouge, anchor=west, fill=white, inner sep=1pt] at (1.3,\y) {\att};}
  % températures (gauche)
  \foreach \y/\T in {0/{$300\ \mathrm{K}$}, -1.3/{$50\ \mathrm{K}$}, -2.6/{$4\ \mathrm{K}$}, -3.9/{$\approx1\ \mathrm{K}$}, -5.2/{$\approx100\ \mathrm{mK}$}, -6.5/{$\approx10\ \mathrm{mK}$}}
    \node[anchor=east, font=\footnotesize\bfseries] at (-3.4,\y) {\T};
  % légendes (droite)
  \node[legende] at (3.7,0)    {électronique de commande et de lecture, ambiante};
  \node[legende] at (3.7,-1.3) {1\up{er} étage du tube pulsé, ancrage thermique des câbles};
  \node[legende] at (3.7,-2.6) {2\up{e} étage du tube pulsé ; amplificateur HEMT (lecture)};
  \node[legende] at (3.7,-3.9) {« still » du circuit de dilution};
  \node[legende] at (3.7,-5.2) {plaque froide : filtres et atténuateurs};
  \node[legende] at (3.7,-6.5) {chambre de mélange : qubits, amplificateur paramétrique};
\end{tikzpicture}}
\caption{Réfrigérateur à dilution en coupe schématique (« lustre »). Les plaques sont de plus en plus froides vers le
bas ; les atténuateurs placés sur la ligne de commande (valeurs typiques) ramènent le bruit thermique de $300\ \mathrm{K}$ à
celui de l'étage où ils sont fixés.}
\label{fig:cryostat}
\end{figure}

\paragraph{Le goulot d'étranglement : le câblage.} Chaque qubit demande plusieurs lignes coaxiales (commande, lecture),
qui apportent de la chaleur depuis les étages chauds et du bruit thermique. Des atténuateurs ancrés à chaque étage
ramènent ce bruit : $T_{\mathrm{eff}}\approx T_{\text{étage}}+T_{\text{chaud}}\,10^{-A/10}$, donc $60\ \mathrm{dB}$ au total limitent à
$300\ \mathrm{K}\times10^{-6}=0{,}3\ \mathrm{mK}$ la contribution du bruit de l'étage ambiant \parencite{krinner2019}. Mais des machines de millions de
qubits physiques ne tiennent pas dans un cryostat avec un coaxial par qubit : les pistes de recherche sont le multiplexage des lignes,
l'électronique de commande placée à $4\ \mathrm{K}$ et le raccordement de plusieurs cryostats entre eux.

\subsubsection{Corriger les erreurs : du code à répétition au code de surface}

Classiquement, on se protège du bruit en copiant chaque bit trois fois et en prenant le vote majoritaire. Le théorème de
non-clonage interdit de copier $\alpha\ket0+\beta\ket1$, et la mesure détruirait la superposition. On \emph{intrique} plutôt
le qubit avec deux autres.

\begin{definition}{Code à répétition à trois qubits}
Le qubit logique $\alpha\ket0+\beta\ket1$ est encodé par deux CNOT dont le contrôle est $q_1$ (et les cibles $q_2$, $q_3$) :
\[
  (\alpha\ket0+\beta\ket1)\ket{00}\ \longmapsto\ \alpha\ket{000}+\beta\ket{111}.
\]
Ce n'est pas $\ket\psi^{\otimes3}$ : l'état est intriqué. Dans cette section, les kets s'écrivent $\ket{q_1q_2q_3}$ et les
fils sont étiquetés : la convention « fil du haut = dernier chiffre » ne s'applique pas. Une erreur $X_k$ (inversion de bit sur $q_k$) est détectée en mesurant les
\emph{parités} $Z_1Z_2$ et $Z_2Z_3$, de valeurs propres $\pm1$.
\end{definition}

Sur $\alpha\ket{000}+\beta\ket{111}$ les deux parités valent $+1$. Une inversion fait changer le signe des parités qui contiennent le qubit touché :

\begin{center}
\renewcommand{\arraystretch}{1.15}
\begin{tabular}{@{}lccc@{}}
\toprule
erreur & état obtenu & $(Z_1Z_2,\ Z_2Z_3)$ & correction\\
\midrule
aucune & $\alpha\ket{000}+\beta\ket{111}$ & $(+1,+1)$ & aucune\\
$X_1$ & $\alpha\ket{100}+\beta\ket{011}$ & $(-1,+1)$ & $X_1$\\
$X_2$ & $\alpha\ket{010}+\beta\ket{101}$ & $(-1,-1)$ & $X_2$\\
$X_3$ & $\alpha\ket{001}+\beta\ket{110}$ & $(+1,-1)$ & $X_3$\\
\bottomrule
\end{tabular}
\end{center}

Les deux parités sont identiques pour $\ket{000}$ et $\ket{111}$ : leur mesure localise l'erreur \emph{sans rien apprendre sur
$\alpha$ et $\beta$}, donc sans détruire la superposition. On mesure les parités avec deux qubits
auxiliaires (\emph{ancillas}) $a_1$ et $a_2$ : le CNOT inscrit dans l'ancilla le ou-exclusif des deux qubits qu'il contrôle
(figure~\ref{fig:code3}).

\begin{figure}[!ht]
\centering
\adjustbox{max width=\linewidth}{%
\begin{tikzpicture}[>=Stealth, font=\small]
  \foreach \y in {3.6, 2.4, 1.2} \draw (0,\y) -- (11.3,\y);
  \foreach \y in {0, -1.2} \draw (0,\y) -- (8.7,\y);
  \node[anchor=east] at (0,3.6) {$q_1:\ \alpha\ket0+\beta\ket1$};
  \node[anchor=east] at (0,2.4) {$q_2:\ \ket0$};
  \node[anchor=east] at (0,1.2) {$q_3:\ \ket0$};
  \node[anchor=east] at (0,0) {$a_1:\ \ket0$};
  \node[anchor=east] at (0,-1.2) {$a_2:\ \ket0$};
  % encodage : CNOT q1 -> q2, q1 -> q3
  \draw[insarouge, thick] (1.1,3.6) -- (1.1,2.4);
  \pic at (1.1,3.6) {ctl}; \pic at (1.1,2.4) {cible};
  \draw[insarouge, thick] (2.0,3.6) -- (2.0,1.2);
  \pic at (2.0,3.6) {ctl}; \pic at (2.0,1.2) {cible};
  % erreur éventuelle
  \node[draw=insagris, dashed, minimum width=9mm, minimum height=9mm] at (3.3,2.4) {$X$};
  \node[font=\footnotesize, text=insagris, anchor=south] at (3.3,3.0) {erreur ?};
  % syndrome : a1 <- q1 xor q2 ; a2 <- q2 xor q3
  \draw[insarouge, thick] (4.5,3.6) -- (4.5,0);  \pic at (4.5,3.6) {ctl}; \pic at (4.5,0) {cible};
  \draw[insarouge, thick] (5.4,2.4) -- (5.4,0);  \pic at (5.4,2.4) {ctl}; \pic at (5.4,0) {cible};
  \draw[insarouge, thick] (6.3,2.4) -- (6.3,-1.2); \pic at (6.3,2.4) {ctl}; \pic at (6.3,-1.2) {cible};
  \draw[insarouge, thick] (7.2,1.2) -- (7.2,-1.2); \pic at (7.2,1.2) {ctl}; \pic at (7.2,-1.2) {cible};
  \pic at (8.4,0) {mesure};   \pic at (8.4,-1.2) {mesure};
  \node[anchor=west, font=\footnotesize] at (9.0,-0.6) {syndrome $(s_1,s_2)$};
  \node[porte] at (10.0,3.6) {$X^{?}$}; \node[porte] at (10.0,2.4) {$X^{?}$}; \node[porte] at (10.0,1.2) {$X^{?}$};
  \draw[dashed, insagris] (8.4,0.25) -- (8.4,0.5) -- (10.0,0.5) -- (10.0,0.85);
  \foreach \x/\t in {2.65/{encodage}, 6.0/{extraction du syndrome}}
    \node[font=\footnotesize, text=insagris] at (\x,-2.1) {\t};
\end{tikzpicture}}
\caption{Code à répétition à trois qubits : encodage, erreur éventuelle (ici sur $q_2$), extraction du syndrome
$(s_1,s_2)=(q_1\oplus q_2,\,q_2\oplus q_3)$ dans les ancillas, puis correction par une porte $X$ sur le qubit désigné
par le syndrome. Les kets s'écrivent $\ket{q_1q_2q_3}$.}
\label{fig:code3}
\end{figure}

\paragraph{Gain.} La correction échoue si au moins deux qubits sont inversés. Avec une probabilité $p$ d'inversion par
qubit, la probabilité d'échec est $3p^2(1-p)+p^3=3p^2-2p^3$, inférieure à $p$ si et seulement si $(2p-1)(p-1)>0$,
c'est-à-dire $p<\frac12$. Le code est utile \emph{si le bruit est assez faible}. C'est le principe d'un
\emph{seuil}.

\paragraph{Erreurs de phase.} Une erreur $Z_1$ sur $\alpha\ket{000}+\beta\ket{111}$ donne $\alpha\ket{000}-\beta\ket{111}$ : un état
logique différent, que ce code ne voit pas. On se protège en travaillant dans la base $X$ : $H^{\otimes3}(\alpha\ket{000}+\beta\ket{111})
=\alpha\ket{+++}+\beta\ket{---}$, où une erreur $Z$ (qui échange $\ket+$ et $\ket-$) joue le rôle d'une inversion de
bit. \textcite{shor1995} a montré que l'on peut \emph{concaténer} les deux codes (neuf qubits) pour corriger n'importe quelle erreur sur un
qubit : toute erreur est une combinaison de $I,X,Y,Z$ et la mesure du syndrome l'y projette.

\paragraph{Code de surface.} Les qubits sont disposés sur un réseau plan de $d\times d$ qubits de donnée, entourés de qubits
de mesure ; les syndromes de type $X$ et de type $Z$ sont mesurés à chaque cycle. La \emph{distance} $d$ est le nombre minimal
d'erreurs qui changent l'état logique : le code corrige jusqu'à $\lfloor(d-1)/2\rfloor$ erreurs, et il faut
$2d^2-1$ qubits physiques par qubit logique. Si le taux d'erreur physique est sous le seuil ($\approx1\,\%$), l'erreur
logique \emph{décroît exponentiellement avec $d$} : à chaque passage de $d$ à $d+2$, elle est divisée par un
facteur $\Lambda>1$. Google a observé ce comportement sur des codes de distance 3, 5 puis 7, avec
$\Lambda=2{,}14\pm0{,}02$ et une erreur logique de $0{,}143\,\%$ par cycle à $d=7$ (101 qubits physiques)
\parencite{acharya2025}. En prolongeant avec ce même $\Lambda$, atteindre une erreur logique de $10^{-12}$ par cycle demande
$\Lambda^{k}\geq1{,}43\times10^{-3}/10^{-12}$, soit $k\geq\ln(1{,}43\times10^{9})/\ln(2{,}14)\approx27{,}7$, donc
$k=28$ et $d=7+2\times28=63$ : près de $8\,000$ qubits physiques par qubit logique. C'est une extrapolation (l'erreur n'est pas
exactement exponentielle), mais elle explique les millions de qubits de la section~\ref{sec:shor} : le surcoût de la
correction d'erreurs domine tout, et un meilleur $\Lambda$ (des qubits physiques moins bruités) le réduit fortement.

\subsubsection{Les qubits de chat d'Alice \& Bob}

Le code de surface doit corriger \emph{deux} types d'erreurs ($X$ et $Z$) sur un réseau plan, d'où un surcoût en $d^2$. L'idée des
\emph{qubits de chat} \parencite{mirrahimi2014} est de concevoir un qubit dont une des deux erreurs est déjà supprimée par la physique. Il ne reste
alors qu'un type d'erreur à corriger : un code \emph{unidimensionnel} (la répétition de la section précédente), de $2d-1$ qubits
au lieu de $2d^2-1$, suffit.

\paragraph{Un qubit dans un oscillateur.} On utilise un mode d'un résonateur micro-onde, oscillateur harmonique d'états de
Fock $\ket n$ ($n$ photons). Un \emph{état cohérent} $\ket\alpha=\ee^{-|\alpha|^2/2}\sum_n\frac{\alpha^n}{\sqrt{n!}}\ket n$ est vecteur propre de l'opérateur
d'annihilation : $a\ket\alpha=\alpha\ket\alpha$, avec $|\alpha|^2$ photons en moyenne. Pour $\alpha$ réel, on pose
\[
  \ket{0_L}\approx\ket{\alpha},\qquad\ket{1_L}\approx\ket{-\alpha},\qquad
  \braket{\alpha}{-\alpha}=\ee^{-|\alpha|^2-|\alpha|^2}=\ee^{-2|\alpha|^2}.
\]
Le recouvrement est $1{,}8\times10^{-2}$ pour $|\alpha|^2=2$, $3{,}4\times10^{-4}$ pour $|\alpha|^2=4$ et $1{,}1\times10^{-7}$ pour
$|\alpha|^2=8$ : les deux états sont presque orthogonaux dès que $|\alpha|^2\gtrsim4$. Les états de la base $X$ sont les
\emph{états de chat} de Schrödinger
\[
  \ket{C_\pm}=N_\pm\big(\ket\alpha\pm\ket{-\alpha}\big),\qquad N_\pm=\frac{1}{\sqrt{2\,(1\pm\ee^{-2|\alpha|^2})}},\qquad
  \ket{\pm_L}\approx\ket{C_\pm}.
\]
$\ket{C_+}$ ne contient que des nombres \emph{pairs} de photons et $\ket{C_-}$ que des nombres \emph{impairs} (les termes en
$\alpha^n$ et $(-\alpha)^n$ s'annulent ou s'ajoutent selon la parité de $n$).

\paragraph{Une perte de photon est une erreur de phase.} Le bruit principal d'un résonateur est la perte d'un photon, décrite par $a$.
Comme $a\ket{\pm\alpha}=\pm\alpha\ket{\pm\alpha}$,
\[
  a\ket{C_\pm}=\alpha N_\pm\big(\ket\alpha\mp\ket{-\alpha}\big)=\alpha\,\frac{N_\pm}{N_\mp}\,\ket{C_\mp}.
\]
Après normalisation, $\ket{C_+}\leftrightarrow\ket{C_-}$ : c'est $\ket{+_L}\leftrightarrow\ket{-_L}$, une \emph{inversion de phase} $Z$. Sur
$\ket{\pm\alpha}$, $a$ ne change rien (un simple facteur), donc la perte de photon ne provoque aucune inversion de bit. Elle survient au taux
$\kappa_1\langle n\rangle=\kappa_1|\alpha|^2$ : le taux d'inversion de phase \emph{croît linéairement} avec $|\alpha|^2$.

\paragraph{Supprimer les inversions de bit.} On ajoute une dissipation \emph{à deux photons} : le circuit éjecte des paires de photons
(taux $\kappa_2$) et en réinjecte par paires (pompe d'amplitude $\epsilon_2$). Cela équivaut à un saut d'opérateur
$\sqrt{\kappa_2}\,(a^2-\alpha^2)$ avec $|\alpha|^2$ fixé par le rapport pompe/perte. Or $(a^2-\alpha^2)\ket{\pm\alpha}=0$ car
$a^2\ket{\pm\alpha}=\alpha^2\ket{\pm\alpha}$ : l'espace $\mathrm{Vect}(\ket\alpha,\ket{-\alpha})$ est exactement le noyau de la dissipation. Elle ramène
tout état de l'oscillateur vers cet espace et le \emph{stabilise}. Retirer les photons deux par deux conserve la parité du nombre de photons, donc la
cohérence entre $\ket{C_+}$ et $\ket{C_-}$. Pour qu'un bruit transforme $\ket\alpha$ en $\ket{-\alpha}$, il doit franchir la distance $2\alpha$ dans
le plan de phase contre la dissipation, ce qui est exponentiellement improbable : le taux d'inversion de bit décroît
comme $\ee^{-2|\alpha|^2}$.

\begin{figure}[!ht]
\centering
\begin{tikzpicture}[>=Stealth, font=\small]
  \draw[->, insagris] (-4.6,0) -- (4.8,0) node[right] {$\mathrm{Re}\,\alpha$};
  \foreach \sg in {1,-1}
    \shade[inner color=insarouge!80, outer color=insarouge!5] (\sg*2.6,0) circle (0.75);
  \node[anchor=north, font=\footnotesize] at (2.6,-0.85) {$\ket{\alpha}=\ket{0_L}$};
  \node[anchor=north, font=\footnotesize] at (-2.6,-0.85) {$\ket{-\alpha}=\ket{1_L}$};
  \draw[<->, dashed, insagris, thick] (-1.8,0) -- (1.8,0);
  \node[fill=white, inner sep=1pt, text=insarouge] at (0,0) {$\times$};
  \node[anchor=north, font=\footnotesize, text=insagris, align=center] at (0,-0.2)
    {inversion de bit : $\sim\ee^{-2|\alpha|^2}$\\ (supprimée)};
  \draw[->, insarouge, thick] (-1.5,0.75) to[bend left=40] (1.5,0.75);
  \node[anchor=south, font=\footnotesize, text=insarouge, align=center] at (0,1.3)
    {photon perdu : $\ket{C_+}\leftrightarrow\ket{C_-}$\\ (inversion de phase, $\propto\kappa_1|\alpha|^2$)};
\end{tikzpicture}
\caption{Qubit de chat dans le plan de phase de l'oscillateur (axe réel) : deux états cohérents $\ket{\pm\alpha}$ représentent $\ket{0_L}$
et $\ket{1_L}$. La dissipation à deux photons les stabilise ; passer de l'un à l'autre (inversion de bit) est exponentiellement
rare, alors que la perte d'un photon échange les états de chat $\ket{C_\pm}\approx\ket{\pm_L}$, de parités opposées (inversion de phase).}
\label{fig:chat}
\end{figure}

\paragraph{Résultat : un biais de bruit énorme.} On augmente $|\alpha|^2$ pour supprimer les inversions de bit, au prix d'inversions de phase
plus fréquentes. Mais l'asymétrie est telle que le compromis est favorable. Les expériences sont de plus en plus convaincantes :
des équipes parisiennes (ENS, Inria, Alice \& Bob) ont observé la suppression exponentielle de l'inversion de bit
\parencite{lescanne2020}, puis un temps d'inversion de bit supérieur à $10\ \mathrm{s}$, pour un temps d'inversion de phase supérieur à
$490\ \mathrm{ns}$ \parencite{reglade2024}. Le rapport est de l'ordre de $2\times10^{7}$ ($10/(0{,}49\times10^{-6})$). Une \emph{annonce de l'entreprise} Alice \& Bob (septembre 2025, résultat préliminaire, non évalué par des pairs à
notre connaissance) fait état de plus d'une heure avec $\langle n\rangle=11$ \parencite{aliceandbob2025}.

\paragraph{De là à un ordinateur.} Il reste à corriger les inversions de phase, par un code à répétition de $d$ qubits de chat
\parencite{guillaud2019}, dont les portes à deux qubits (CNOT) doivent \emph{préserver le biais} : c'est le point que les expériences
doivent encore démontrer. Une estimation d'architecture évalue le calcul d'un logarithme discret sur une courbe de 256 bits à
$126\,133$ qubits de chat et $9$ heures \parencite{gouzien2023}. Alice \& Bob (start-up française) et Amazon \parencite{putterman2025} suivent cette voie ;
Alice \& Bob vise, dans sa \emph{feuille de route} d'entreprise, une machine de 100 qubits logiques en 2030. C'est un objectif commercial, pas un résultat.

% ---------------------------------------------------------------------
\subsection{Autres applications et algorithmes}\label{sec:autres}
% ---------------------------------------------------------------------

Shor et Grover sont les algorithmes les plus connus, pas les seuls. Il faut distinguer les gains \emph{prouvés} (on sait démontrer
qu'aucun algorithme classique ne fait aussi bien, ou on connaît un bien meilleur algorithme quantique), les gains
\emph{conditionnels} (valables sous hypothèses) et les méthodes \emph{heuristiques} (qui marchent peut-être, sans preuve).

\subsubsection{Panorama}

\begin{center}
\renewcommand{\arraystretch}{1.2}
\small
\begin{tabular}{@{}>{\raggedright\arraybackslash}p{3.1cm}>{\raggedright\arraybackslash}p{3.8cm}>{\raggedright\arraybackslash}p{3.0cm}>{\raggedright\arraybackslash}p{4.6cm}@{}}
\toprule
\textbf{Algorithme} & \textbf{Problème} & \textbf{Gain} & \textbf{Statut et usage}\\
\midrule
Deutsch--Jozsa, Bernstein--Vazirani, Simon (1985--1997) & fonctions « boîtes noires » & exponentiel (Simon) & prouvé ; peu d'usage direct, base de Shor\\
Shor (1994) \parencite{shor1997} & factorisation, logarithme discret & superpolynomial & prouvé ; casse RSA et courbes elliptiques\\
Grover (1996) & recherche non structurée & quadratique : $\sqrt N$ & prouvé, optimal ; clés symétriques, SAT\\
Estimation de phase \parencite{kitaev1995} & valeur propre d'une porte $U$ & précision $2^{-m}$ avec $m$ qubits & prouvé ; brique de Shor, de la chimie, de HHL\\
Estimation d'amplitude \parencite{brassard2002} & estimer une probabilité & $\varepsilon^{-1}$ au lieu de $\varepsilon^{-2}$ & prouvé ; Monte-Carlo, finance\\
Simulation hamiltonienne \parencite{lloyd1996} & évolution $\ee^{-\ii Ht}\ket\psi$ & exponentiel (systèmes quantiques) & prouvé ; chimie, matériaux : application la plus naturelle\\
HHL \parencite{harrow2009} & système linéaire $Ax=b$ & exponentiel \emph{sous conditions} & conditionnel ; en partie « déquantisé » \parencite{tang2019}\\
VQE \parencite{peruzzo2014} & énergie de l'état fondamental & aucun gain prouvé & heuristique ; machines bruitées\\
QAOA \parencite{farhi2014} & optimisation combinatoire & aucun gain prouvé & heuristique ; machines bruitées\\
\bottomrule
\end{tabular}
\end{center}

\subsubsection{Simuler la matière : l'idée de Feynman}

Pour décrire $n$ qubits, un ordinateur classique stocke $2^n$ amplitudes complexes. Avec $16$ octets par amplitude, $n=50$
demande $2^{50}\times16=2^{54}$ octets, soit $16$ pébioctets ($1{,}8\times10^{16}$ octets) ; $n=300$ demande $2^{300}\approx2\times10^{90}$
amplitudes, plus que d'atomes dans l'univers observable ($\sim10^{80}$). \textcite{feynman1982} en a tiré l'idée
inverse : un système quantique simule naturellement un autre système quantique.

La \emph{simulation hamiltonienne} calcule $\ee^{-\ii Ht}\ket{\psi}$ pour $H=A+B$, où $A$ et $B$ sont des sommes de termes simples dont on sait implémenter les exponentielles.
Si $A$ et $B$ ne commutent pas, $\ee^{-\ii(A+B)t}\neq\ee^{-\ii At}\ee^{-\ii Bt}$. La formule de Baker--Campbell--Hausdorff donne, pour un pas $\tau$,
\[
  \ee^{-\ii A\tau}\ee^{-\ii B\tau}=\exp\!\Big(-\ii(A+B)\tau-\tfrac{\tau^2}{2}[A,B]+O(\tau^3)\Big).
\]
On coupe donc $t$ en $N$ pas $\tau=t/N$ (\emph{formule de Trotter}) :
\[
  \Big\|\ee^{-\ii(A+B)t}-\big(\ee^{-\ii At/N}\ee^{-\ii Bt/N}\big)^N\Big\|\ \leq\ \frac{t^2}{2N}\,\big\|[A,B]\big\|\quad(\text{à l'ordre dominant}),
\]
l'erreur de chaque pas ($\sim\tau^2\|[A,B]\|/2$) s'additionnant sur $N$ pas. Pour $A=X$ et $B=Z$ (un qubit), $[X,Z]=-2\ii Y$ et
$\|[X,Z]\|=2$ : l'erreur vaut au plus $t^2/N$ ; pour $t=1$, $N=100$ pas donnent une erreur inférieure à $10^{-2}$ (le calcul numérique donne $7\times10^{-3}$). Les applications visées sont la chimie
(énergie d'une molécule comme le cofacteur FeMoco de la nitrogénase, qui intervient dans la fabrication des engrais \parencite{reiher2017}),
les matériaux (batteries, supraconducteurs) et la physique des particules. Les méthodes classiques (théorie de la fonctionnelle de
la densité, \emph{coupled cluster}) traitent bien les systèmes faiblement corrélés ; c'est pour les systèmes \emph{fortement
corrélés} qu'un avantage est espéré, avec des machines tolérantes aux fautes.

\subsubsection{Méthodes variationnelles et optimisation}

Sur les machines bruitées actuelles, on ne peut pas faire tourner de longs circuits ; on utilise donc des circuits courts à paramètres
$\theta$ dans une boucle avec un ordinateur classique : le circuit prépare $\ket{\psi(\theta)}$, on mesure l'énergie
$E(\theta)=\bra{\psi(\theta)}H\ket{\psi(\theta)}$, et un optimiseur classique corrige $\theta$. Cela repose sur le
\emph{principe variationnel} : en décomposant $\ket\psi=\sum_kc_k\ket{k}$ sur les états propres de $H$ (section~\ref{sec:spectral}),
$E=\sum_k|c_k|^2E_k\geq E_0\sum_k|c_k|^2=E_0$. Minimiser $E(\theta)$ approche donc l'énergie fondamentale $E_0$ par le haut. Exemple pour un qubit,
$H=X+Z$ (valeurs propres $\pm\sqrt2$) avec $\ket{\psi(\theta)}=\cos\frac\theta2\ket0+\sin\frac\theta2\ket1$ :
\[
  E(\theta)=\langle Z\rangle+\langle X\rangle=\cos\theta+\sin\theta=\sqrt2\cos\!\big(\theta-\tfrac\pi4\big),
\]
minimale pour $\theta=\frac{5\pi}{4}$, de valeur $-\sqrt2=E_0$. Le VQE \parencite{peruzzo2014} généralise cette idée à des molécules, et le QAOA
\parencite{farhi2014} l'applique à des problèmes d'optimisation en alternant deux types de portes. Aucun avantage n'est démontré pour ces méthodes :
elles sont l'objet de la recherche actuelle.

\subsubsection{Algèbre linéaire quantique : une promesse sous conditions}

L'algorithme HHL \parencite{harrow2009} prépare un état $\ket{x}\propto\sum_ix_i\ket{i}$ où $Ax=b$, en un temps polynomial en $\log N$
là où la résolution classique est polynomiale en $N$. Le gain exponentiel apparent est soumis à plusieurs conditions. Il faut
\emph{charger} $b$ dans l'état $\ket{b}$ efficacement ; $A$ doit être creuse et bien conditionnée ; et le résultat est un
\emph{état} : en lire les $N$ composantes demanderait déjà $N$ répétitions, de sorte que l'on n'obtient que des
grandeurs globales de $x$. De plus, des algorithmes classiques « inspirés du quantique » (\emph{déquantisation}) obtiennent le même
résultat sous des hypothèses d'accès analogues \parencite{tang2019}, ce qui a supprimé l'avantage exponentiel annoncé pour
certaines applications d'apprentissage. Le bilan honnête : \emph{pas d'avantage démontré en apprentissage automatique à ce jour}.

\subsubsection{Capteurs, nombres aléatoires, réseaux}

\paragraph{Métrologie quantique.} Les premières technologies quantiques commercialisées sont des capteurs (horloges atomiques,
magnétomètres, gravimètres). L'intrication y améliore la précision. Soit $N$ qubits qui subissent chacun la phase
$\ket1\mapsto\ee^{\ii\varphi}\ket1$ à mesurer. Dans l'état de GHZ $\frac{1}{\sqrt2}(\ket{0\ldots0}+\ket{1\ldots1})$ (généralisation à $N$ qubits de l'état de Bell
$\frac{1}{\sqrt2}(\ket{00}+\ket{11})$, section~\ref{sec:bell}), seuls $\ket{0\ldots0}$ et $\ket{1\ldots1}$ interviennent, et l'état devient
\[
  \frac{1}{\sqrt2}\big(\ket{0\ldots0}+\ee^{\ii N\varphi}\ket{1\ldots1}\big),\qquad
  P\big(X^{\otimes N}=+1\big)=\frac{1+\cos N\varphi}{2}=\cos^2\!\frac{N\varphi}{2}.
\]
La phase est amplifiée par $N$ : la précision atteint $\Delta\varphi\sim1/N$ (limite de Heisenberg) au lieu de $1/\sqrt N$ pour
$N$ qubits indépendants (limite quantique standard) \parencite{giovannetti2004}. En pratique, la décohérence réduit ce gain, qui
se limite souvent à de petits $N$.

\paragraph{Nombres aléatoires.} Mesurer $\ket+$ dans la base $Z$ donne $0$ ou $1$ avec la probabilité $\frac12$ : le postulat de Born
garantit un hasard \emph{intrinsèque}, contrairement à un générateur pseudo-aléatoire déterministe. Les générateurs quantiques
de nombres aléatoires (QRNG), à base de photons, sont déjà commercialisés ; ils alimentent par exemple la génération de clés cryptographiques.

\paragraph{Réseaux quantiques.} Relier des ordinateurs quantiques demande de transmettre des qubits (téléportation, section~\ref{sec:teleportation}).
Une fibre optique standard atténue de $0{,}2\ \mathrm{dB/km}$, soit une transmission $T=10^{-0{,}02L}$ ($L$ en km) :
$T=10^{-2}$ à $100\ \mathrm{km}$ et $10^{-10}$ à $500\ \mathrm{km}$. Un amplificateur est impossible (non-clonage) : on utilise des
\emph{répéteurs quantiques}, qui étendent l'intrication par \emph{échange d'intrication} (une mesure de Bell entre deux paires intriquées, comme dans la
téléportation de la section~\ref{sec:teleportation}), ou des satellites : la distribution
d'intrication sur $1\,200\ \mathrm{km}$ a été démontrée par le satellite chinois Micius \parencite{yin2017}. L'objectif est un « internet
quantique » \parencite{wehner2018} reliant des processeurs, des capteurs et des liens de QKD.

\subsubsection{Jalons et niveaux de preuve}

Le tableau place quelques jalons dans le temps et indique la \emph{nature de la source}, car une annonce d'entreprise ou une
prépublication n'a pas la valeur d'un article évalué par des pairs ou d'une norme.

\begin{center}
\renewcommand{\arraystretch}{1.2}
\small
\begin{tabular}{@{}>{\raggedright\arraybackslash}p{1.9cm}>{\raggedright\arraybackslash}p{8.8cm}>{\raggedright\arraybackslash}p{3.3cm}@{}}
\toprule
\textbf{Date} & \textbf{Jalon} & \textbf{Nature de la source}\\
\midrule
1994--1996 & algorithmes de Shor et de Grover \parencite{shor1997} & articles évalués\\
2019 & Google, $53$ qubits : échantillonnage de circuits aléatoires en $200\ \mathrm{s}$, annoncé hors de portée des machines classiques \parencite{arute2019} (contesté ensuite) & article évalué\\
août 2024 & normes NIST FIPS 203, 204, 205 \parencite{fips203} & norme\\
2024 & qubit de chat : inversion de bit $>10\ \mathrm{s}$ \parencite{reglade2024} & article évalué\\
déc. 2024 & code de surface sous le seuil, $\Lambda=2{,}14$ \parencite{acharya2025} & article évalué\\
mars 2025 & HQC choisi par le NIST \parencite{nist2025hqc} & rapport officiel\\
mai 2025 & RSA-2048 : moins d'un million de qubits \parencite{gidney2025} & prépublication\\
sept. 2025 & qubit de chat : plus d'une heure d'inversion de bit \parencite{aliceandbob2025} & annonce d'entreprise\\
mars 2026 & Shor avec $10\,000$ à $26\,000$ atomes \parencite{cain2026} & prépublication\\
août 2026 & RFC~10024 : TLS hybride X25519 + ML-KEM \parencite{rfc10024} & norme IETF\\
\bottomrule
\end{tabular}
\end{center}

Ce qui est \emph{acquis} : les gains théoriques de Shor et de Grover, la correction d'erreurs sous le seuil sur de petits codes,
et les standards post-quantiques. Ce qui \emph{ne l'est pas} : un ordinateur à plusieurs milliers de qubits \emph{logiques},
un avantage utile sur un problème pratique, et surtout la date $Z$ de l'inégalité de Mosca. C'est pourquoi la migration cryptographique se
décide sur $X+Y$, qui se connaît, et non sur une prédiction de $Z$.

\clearpage
% ---------------------------------------------------------------------
\subsection{Récapitulatif du chapitre}\label{sec:recap-appli}
% ---------------------------------------------------------------------

\begin{center}
\adjustbox{max width=\linewidth}{%
\begin{tikzpicture}[>=Stealth, font=\footnotesize,
    boite/.style={draw=insarouge, fill=insarouge!6, rounded corners=2pt, thick,
                  minimum width=4.4cm, minimum height=1.45cm, align=center, inner sep=3pt},
    lien/.style={->, very thick, insagris},
    etiq/.style={font=\scriptsize\itshape, text=insagris, inner sep=2pt}]
  \node[boite] (int) at (0,0)
    {\textbf{Interférence}\\ oracle, Deutsch, Grover : $\sqrt N$\\ {\scriptsize\S\,\ref{sec:interference}, \ref{sec:grover}}};
  \node[boite] (sho) at (6.3,0)
    {\textbf{Shor}\\ QFT, période $r$, $O(n^3)$\\ {\scriptsize\S\,\ref{sec:shor}}};
  \node[boite] (cry) at (12.6,0)
    {\textbf{Cryptographie}\\ RSA et ECC cassés ; HNDL ; PQC\\ {\scriptsize\S\,\ref{sec:crypto}}};
  \node[boite] (cryo) at (12.6,-2.3)
    {\textbf{Qubits physiques}\\ décohérence, cryostat $10\ \mathrm{mK}$\\ {\scriptsize\S\,\ref{sec:machine}}};
  \node[boite] (qec) at (6.3,-2.3)
    {\textbf{Correction d'erreurs}\\ seuil, code de surface, qubits de chat\\ {\scriptsize\S\,\ref{sec:machine}}};
  \node[boite] (aut) at (0,-2.3)
    {\textbf{Autres usages}\\ simulation, VQE, capteurs, réseaux\\ {\scriptsize\S\,\ref{sec:autres}}};
  \draw[lien] (int) -- (sho) node[etiq, midway, above=1pt] {période};
  \draw[lien] (sho) -- (cry) node[etiq, midway, above=1pt] {casse RSA};
  \draw[lien] (cry) -- (cryo) node[etiq, midway, right=1pt] {quelle machine ?};
  \draw[lien] (cryo) -- (qec) node[etiq, midway, above=1pt] {qubits fiables};
  \draw[lien] (qec) -- (aut) node[etiq, midway, above=1pt] {usages};
\end{tikzpicture}}
\end{center}

\begin{aretenir}
\begin{itemize}
  \item Un algorithme quantique organise des \emph{interférences} : l'oracle écrit $f(x)$ dans une phase ; Deutsch décide avec un appel au lieu de deux.
        Grover : $\approx\frac\pi4\sqrt N$ itérations (rotations de $2\theta$), gain quadratique et optimal ; exact en une itération pour $N=4$.
  \item Shor : la QFT transforme la période $r$ de $a^x\bmod N$ en pics en $y=jM/r$ ; $\operatorname{pgcd}(a^{r/2}\pm1,N)$ donne les facteurs ;
        coût $O(n^3)$ ($n=\log_2N$), contre un temps sous-exponentiel classique (exemple : $N=15$, $a=7$, $r=4$).
  \item RSA, Diffie--Hellman, ECDH et ECDSA tombent face à Shor ; AES-256 et SHA-256 résistent à Grover. Doubler les clés RSA ne protège pas.
  \item « Stocker maintenant, déchiffrer plus tard » : exposé dès que $X+Y>Z$ (Mosca) ; la forward secrecy n'aide pas. On migre d'abord l'échange de clés
        (hybride X25519 + ML-KEM), puis les signatures (FIPS 203--205).
  \item Qubit physique : la décohérence ($T_2$) mélange l'état ; le cryostat à $10\ \mathrm{mK}$ rend $p_1$ négligeable ($hf/k_B=0{,}24\ \mathrm{K}$ à $5\ \mathrm{GHz}$). Correction
        d'erreurs : seuil, $2d^2-1$ qubits physiques par qubit logique (code de surface).
  \item Qubit de chat : $\ket{\pm\alpha}$, perte de photon = inversion de phase, inversion de bit en $\ee^{-2|\alpha|^2}$ ; code à répétition. Une annonce d'entreprise
        n'est pas un résultat évalué.
  \item Autres usages : simulation (Trotter), VQE et QAOA (heuristiques), HHL (conditionnel), capteurs (limite de Heisenberg), réseaux (répéteurs).
\end{itemize}
\end{aretenir}

\begin{center}
\renewcommand{\arraystretch}{1.12}
\small
\begin{tabular}{@{}l l r@{}}
\toprule
\textbf{Notion} & \textbf{Formule} & \textbf{§} \\
\midrule
Rebond de phase & $U_f\ket{x}\ket{-}=(-1)^{f(x)}\ket{x}\ket{-}$ & \ref{sec:interference} \\
Grover & $P=\sin^2\!\big((2k+1)\theta\big)$, \quad $\sin\theta=\frac{1}{\sqrt N}$, \quad $k\simeq\frac\pi4\sqrt N$ & \ref{sec:grover} \\
QFT et période & $\mathrm{QFT}_M\ket x=\frac{1}{\sqrt M}\sum_y\ee^{2\ii\pi xy/M}\ket y$ ; pics en $y=jM/r$, $P=\frac1r$ & \ref{sec:shor} \\
Facteurs, coût & $\operatorname{pgcd}(a^{r/2}\pm1,N)$, \quad $O(n^3)$ & \ref{sec:shor} \\
Mosca & $X+Y>Z$ : données déjà exposées & \ref{sec:crypto} \\
Population thermique & $p_1=\frac{1}{1+\ee^{hf/k_BT}}$, \quad $\operatorname{Tr}\rho^2=\frac{1+\ee^{-2t/T_2}}{2}$ (décohérence) & \ref{sec:machine} \\
Qubit de chat & $\braket{\alpha}{-\alpha}=\ee^{-2|\alpha|^2}$, \quad $a\ket{C_\pm}=\alpha\frac{N_\pm}{N_\mp}\ket{C_\mp}$ & \ref{sec:machine} \\
\bottomrule
\end{tabular}
\end{center}

\begin{savoirfaire}
\begin{itemize}[nosep]
  \item Suivre un circuit de Deutsch ou de Grover ($N=4$) jusqu'à la mesure (sections~\ref{sec:interference}, \ref{sec:grover}).
  \item Trouver l'ordre de $a$ modulo $N$, appliquer le critère du pgcd et prévoir la loi de $y$ (section~\ref{sec:shor}, exemple $N=15$).
  \item Dire ce que casse Shor, ce que Grover affaiblit, et appliquer l'inégalité de Mosca (section~\ref{sec:crypto}).
  \item Calculer $p_1$, la pureté après décohérence et le recouvrement $\ee^{-2|\alpha|^2}$ ; lire un syndrome (section~\ref{sec:machine}).
  \item Distinguer résultat prouvé, conditionnel et annonce (section~\ref{sec:autres}).
\end{itemize}
\end{savoirfaire}
, entre chapitre5 et
%  chapitre6 (exercices). Il ne contient ni préambule ni \begin{document} : les
%  environnements (definition, resultat, remarque, aretenir, savoirfaire), les
%  macros (\ii, \ee, \dd, \ket, \bra, \braket), les symboles de circuits (ctl,
%  cible, croix, mesure, porte) et les couleurs (insarouge, insagris) sont définis
%  dans main.tex. Un chapitre = une \section.
%  Le numéro d'un chapitre est donné par l'ordre des \include, pas par le nom du
%  fichier : ce chapitre est le chapitre 6, les exercices (chapitre6.tex) le 7.
% =====================================================================

% Macros de Dirac (déjà définies dans main.tex ; ce bloc évite l'erreur
% « Undefined control sequence \ket » si main.tex est une ancienne version).
\providecommand{\ket}[1]{|#1\rangle}
\providecommand{\bra}[1]{\langle #1|}
\providecommand{\braket}[2]{\langle #1|#2\rangle}

% =====================================================================
\section{Applications et algorithmes quantiques}\label{chap:applications}
% =====================================================================

Les chapitres précédents ont construit le formalisme : qubits, portes, intrication, mesure. Ce
chapitre montre à quoi il sert. On y voit comment un circuit exploite l'interférence pour résoudre
certains problèmes plus vite qu'un ordinateur classique (sections~\ref{sec:interference}
à~\ref{sec:shor}), pourquoi cela menace une partie de la cryptographie actuelle, RSA et TLS en tête
(section~\ref{sec:crypto}), ce qu'il faut pour construire une telle machine, du cryostat aux qubits
de chat d'Alice \& Bob (section~\ref{sec:machine}), et quelles autres applications et quels autres
algorithmes existent (section~\ref{sec:autres}). Les circuits sont dessinés comme aux chapitres~4 et~5 :
le fil du haut est le dernier chiffre du ket.

\begin{remarque}[Chiffres d'actualité]
Les estimations de ressources, les calendriers de migration et les records expérimentaux cités ici
datent d'octobre 2026 et évoluent vite. Chaque chiffre est accompagné de sa source (bibliographie) ;
les annonces d'entreprises sont signalées comme telles et les prépublications aussi.
\end{remarque}

% ---------------------------------------------------------------------
\subsection{Calculer par interférence}\label{sec:interference}
% ---------------------------------------------------------------------

Un registre de $n$ qubits est décrit par $2^n$ amplitudes (chapitre~\ref{chap:operateurs}). Une porte
agit sur toutes ces amplitudes à la fois, mais la mesure n'en révèle qu'un seul résultat, tiré au
hasard avec les probabilités de Born. Dire que «~le circuit calcule toutes les valeurs en parallèle~»
est donc trompeur. Un algorithme quantique est une \emph{chorégraphie d'interférences} : les
amplitudes des mauvaises réponses se compensent (signes ou phases opposés), celles de la bonne
réponse s'additionnent, et la mesure finale donne la bonne réponse avec une probabilité proche de $1$.

\begin{definition}{Oracle}
Pour $f:\{0,1\}^n\to\{0,1\}$, l'\emph{oracle} de $f$ est la porte à $n+1$ qubits
\[
  U_f\ket{x}\ket{y}=\ket{x}\ket{y\oplus f(x)}.
\]
Elle est unitaire et $U_f^2=\mathrm{Id}$ (car $y\oplus f(x)\oplus f(x)=y$). Elle se construit à partir
d'un circuit réversible de $f$ (portes CNOT et Toffoli, section~\ref{sec:controlees}). Le coût d'un
algorithme se mesure en \emph{nombre d'appels} à $U_f$.
\end{definition}

\begin{resultat}{Rebond de phase de l'oracle}
Avec l'ancilla dans l'état $\ket{-}=\frac{1}{\sqrt2}(\ket0-\ket1)$ :
\[
  U_f\ket{x}\ket{-}=(-1)^{f(x)}\ket{x}\ket{-}.
\]
\end{resultat}

En effet $U_f\ket{x}\ket{-}=\frac{1}{\sqrt2}\big(\ket{x}\ket{f(x)}-\ket{x}\ket{1\oplus f(x)}\big)$.
Si $f(x)=0$ on retrouve $\ket{x}\ket{-}$ ; si $f(x)=1$, les deux termes sont échangés et le signe
change. La valeur $f(x)$ est écrite dans la \emph{phase} de $\ket{x}$, sans toucher à l'ancilla : c'est
le même mécanisme que le rebond de phase du CNOT (section~\ref{sec:cnot}), et la brique de tous les
algorithmes de ce chapitre.

\subsubsection{Algorithme de Deutsch}

Soit $f:\{0,1\}\to\{0,1\}$. Elle est \emph{constante} si $f(0)=f(1)$ et \emph{équilibrée} si
$f(0)\neq f(1)$. On veut savoir laquelle des deux est vraie. Classiquement, il faut évaluer $f$ deux
fois. Le circuit de la figure~\ref{fig:deutsch} répond avec \emph{un seul} appel à $U_f$.

\begin{figure}[!ht]
\centering
\begin{tikzpicture}[>=Stealth, font=\small]
  % fils : ancilla (second chiffre) en haut, requête (premier chiffre) en bas
  \draw (0,1.5) -- (8.4,1.5);
  \draw (0,0) -- (7.4,0);
  \node[anchor=east] at (0,1.5) {$\ket{1}$};
  \node[anchor=east] at (0,0) {$\ket{0}$};
  \node[porte] at (1.3,1.5) {$H$};
  \node[porte] at (1.3,0) {$H$};
  \node[draw=insarouge, fill=insarouge!8, thick, minimum width=13mm, minimum height=22mm]
    at (3.9,0.75) {$U_f$};
  \node[porte] at (6.0,0) {$H$};
  \pic at (7.2,0) {mesure};
  \node[anchor=west] at (7.7,0) {$f(0)\oplus f(1)$};
  \node[anchor=west, text=insagris, font=\footnotesize] at (8.5,1.5) {ancilla};
  \foreach \x/\lab in {0.55/0, 2.2/1, 5.15/2, 6.75/3}
    {\draw[dashed, insagris!70] (\x,2.1) -- (\x,-0.7);
     \node[anchor=north, font=\footnotesize] at (\x,-0.7) {$\Pi_\lab$};}
\end{tikzpicture}
\caption{Circuit de Deutsch. La requête $x$ est le premier chiffre du ket (fil du bas), l'ancilla
$y$ le second (fil du haut). La mesure du seul qubit de requête donne $f(0)\oplus f(1)$ : $0$ si $f$ est
constante, $1$ si elle est équilibrée.}
\label{fig:deutsch}
\end{figure}

\begin{align*}
  \ket{\Pi_0}&=\ket{0}\ket{1},\qquad
  \ket{\Pi_1}=(H\otimes H)\ket{\Pi_0}=\ket{+}\ket{-},\\[1ex]
  \ket{\Pi_2}&=U_f\ket{\Pi_1}
    =\frac{(-1)^{f(0)}\ket{0}+(-1)^{f(1)}\ket{1}}{\sqrt2}\ket{-}
    =(-1)^{f(0)}\,\frac{\ket{0}+(-1)^{f(0)\oplus f(1)}\ket{1}}{\sqrt2}\,\ket{-},\\[1ex]
  \ket{\Pi_3}&=(H\otimes I)\ket{\Pi_2}=(-1)^{f(0)}\,\ket{f(0)\oplus f(1)}\ket{-}.
\end{align*}
La deuxième ligne utilise le rebond de phase ; la troisième, $H\ket{+}=\ket0$ et $H\ket{-}=\ket1$.
Les deux amplitudes de la requête interfèrent : elles s'additionnent si $f$ est constante (état
$\ket{+}$, devenu $\ket{0}$) et se retranchent si $f$ est équilibrée (état $\ket{-}$, devenu $\ket{1}$).
Le calcul ne dit pas les valeurs $f(0)$ et $f(1)$ : il extrait seulement leur parité, et c'est ce qui
permet d'économiser un appel.

\begin{remarque}[Des oracles à $n$ qubits : Deutsch--Jozsa, Bernstein--Vazirani, Simon]
Le même schéma ($H^{\otimes n}$, oracle, $H^{\otimes n}$) généralise l'algorithme de \textcite{deutsch1985}. Pour $f$ sur $n$ bits promise
constante ou équilibrée, \textcite{deutschjozsa1992} décident avec \emph{un} appel là où un algorithme classique
déterministe en demande $2^{n-1}+1$. \textcite{bernstein1997} retrouvent en un appel un mot caché
$s$ défini par $f(x)=s\cdot x$ (classiquement : $n$ appels). \textcite{simon1997} trouve un masque caché avec $O(n)$
appels, contre $2^{\Omega(n)}$ classiquement : c'est la première séparation exponentielle, et l'idée
de l'algorithme de Shor. Ces problèmes sont artificiels (ils portent sur une boîte noire), mais
ils prouvent que l'interférence offre un gain réel.
\end{remarque}

% ---------------------------------------------------------------------
\subsection{Recherche dans une liste : l'algorithme de Grover}\label{sec:grover}
% ---------------------------------------------------------------------

Un oracle $f$ vaut $1$ en un unique élément $w$ parmi $N=2^n$, et $0$ ailleurs ; on cherche $w$.
Classiquement, il faut en moyenne $N/2$ évaluations de $f$, et $N$ dans le pire cas. \textcite{grover1996}
le trouve avec environ $\frac\pi4\sqrt N$ appels.

\begin{resultat}{Algorithme de Grover}
On part de $\ket{s}=H^{\otimes n}\ket{0\ldots0}=\frac{1}{\sqrt N}\sum_x\ket{x}$, puis on répète $k$
fois l'itération $G=D\,O_w$ avec
\[
  O_w=I-2\ket{w}\bra{w}\ \ (\text{oracle, par rebond de phase}),\qquad
  D=2\ket{s}\bra{s}-I=H^{\otimes n}\,(2\ket{0}\bra{0}-I)\,H^{\otimes n}.
\]
La mesure donne $w$ avec la probabilité $\sin^2\!\big((2k+1)\theta\big)$, où $\sin\theta=\frac{1}{\sqrt N}$.
Elle est proche de $1$ pour $k\simeq\frac\pi4\sqrt N$. Aucun algorithme quantique ne fait mieux que
$\Omega(\sqrt N)$ appels : le gain est quadratique, et optimal.
\end{resultat}

\paragraph{Pourquoi cela marche.} Posons $\ket{s'}=\frac{1}{\sqrt{N-1}}\sum_{x\neq w}\ket{x}$. Alors
$\ket{s}=\sin\theta\,\ket{w}+\cos\theta\,\ket{s'}$, et toute l'évolution reste dans le plan
engendré par $\ket{w}$ et $\ket{s'}$. L'oracle change le signe de la composante sur $\ket{w}$ : c'est
la symétrie par rapport à l'axe $\ket{s'}$. La diffusion $D=2\ket{s}\bra{s}-I$ est la symétrie par
rapport à $\ket{s}$. Deux symétries d'axes séparés de l'angle $\theta$ composent une rotation de
$2\theta$ (figure~\ref{fig:grover-geo}). Après $k$ itérations,
\[
  G^k\ket{s}=\sin\!\big((2k+1)\theta\big)\ket{w}+\cos\!\big((2k+1)\theta\big)\ket{s'}.
\]
On s'arrête quand $(2k+1)\theta\simeq\frac\pi2$. Contrairement à un algorithme classique, itérer plus
longtemps n'améliore rien : on continue de tourner et on dépasse $\ket{w}$.

\begin{figure}[!ht]
\centering
\begin{tikzpicture}[>=Stealth, font=\small]
  \draw[->, insagris] (-0.4,0) -- (4.5,0) node[right] {$\ket{s'}$};
  \draw[->, insagris] (0,-0.5) -- (0,4.5);
  \draw[insagris!50, dashed] (3.6,0) arc (0:90:3.6);
  \draw[dashed, insagris!60] (210:1.8) -- (30:3.6);
  \draw[->, very thick, insarouge] (0,0) -- (30:3.6) node[right] {$\ket{s}$};
  \draw[->, very thick, insagris] (0,0) -- (-30:3.6) node[right] {$O_w\ket{s}$};
  \draw[->, very thick, insarouge] (0,0) -- (90:3.6) node[above right] {$G\ket{s}=\ket{w}$};
  \draw[insagris] (1.1,0) arc (0:30:1.1);   \node[text=insagris] at (15:1.45) {$\theta$};
  \draw[insagris] (1.1,0) arc (0:-30:1.1);  \node[text=insagris] at (-15:1.45) {$\theta$};
  \draw[->, insarouge, thick] (30:2.2) arc (30:90:2.2);  \node[text=insarouge] at (60:2.65) {$2\theta$};
  \node[text=insagris, font=\footnotesize, anchor=north west, align=left] at (-1.7,-1.2)
    {symétrie $O_w$ : $\ket{s}\mapsto O_w\ket{s}$\\ symétrie $D$ d'axe $\ket{s}$ (pointillés)};
\end{tikzpicture}
\caption{Une itération de Grover pour $N=4$ ($\theta=\frac\pi6$). L'oracle renvoie $\ket{s}$ en
$O_w\ket{s}$ (angle $-\theta$), la diffusion le renvoie en $G\ket{s}$ (angle $3\theta=\frac\pi2$) :
une seule itération amène exactement sur $\ket{w}$.}
\label{fig:grover-geo}
\end{figure}

\subsubsection{Exemple complet : $N=4$ (deux qubits)}

Prenons $w=11$. L'oracle est la porte $Z$ contrôlée $\mathrm{CZ}=\operatorname{diag}(1,1,1,-1)$
(section~\ref{sec:controlees}). La diffusion s'écrit avec les portes du cours, car
$2\ket{00}\bra{00}-I=\operatorname{diag}(1,-1,-1,-1)=\mathrm{CZ}\,(Z\otimes Z)$ :
\[
  D=(H\otimes H)\;\mathrm{CZ}\,(Z\otimes Z)\;(H\otimes H).
\]
Le circuit de la figure~\ref{fig:grover-circ} s'écrit donc $H\otimes H$, $\mathrm{CZ}$,
$H\otimes H$, $Z\otimes Z$, $\mathrm{CZ}$, $H\otimes H$. Calcul pas à pas :
\begin{align*}
  \ket{s}&=\tfrac12\big(\ket{00}+\ket{01}+\ket{10}+\ket{11}\big),\\
  O_w\ket{s}&=\tfrac12\big(\ket{00}+\ket{01}+\ket{10}-\ket{11}\big),\\
  \braket{s}{O_w s}&=\tfrac14(1+1+1-1)=\tfrac12,\\
  D\,O_w\ket{s}&=2\braket{s}{O_w s}\ket{s}-O_w\ket{s}=\ket{s}-O_w\ket{s}
    =\tfrac12\big(0,0,0,2\big)^{\mathsf T}=\ket{11}.
\end{align*}
La mesure donne $11$ \emph{avec certitude} après un seul appel à l'oracle, là où un algorithme
classique en demande en moyenne $2{,}25$. C'est cohérent avec la formule : $\sin\theta=\frac12$,
$\theta=\frac\pi6$ et $\sin(3\theta)=1$. Pour un autre élément marqué, on change l'oracle (un
CZ entouré de portes $X$) et le résultat est le même.

\begin{figure}[!ht]
\centering
\adjustbox{max width=\linewidth}{%
\begin{tikzpicture}[>=Stealth, font=\small]
  \draw (0,1.4) -- (10.3,1.4);
  \draw (0,0) -- (10.3,0);
  \node[anchor=east] at (0,1.4) {$\ket{0}$};
  \node[anchor=east] at (0,0) {$\ket{0}$};
  \foreach \y in {0,1.4} {
    \node[porte] at (1.1,\y) {$H$};
    \node[porte] at (4.3,\y) {$H$};
    \node[porte] at (5.7,\y) {$Z$};
    \node[porte] at (8.5,\y) {$H$};
    \pic at (9.7,\y) {mesure};
  }
  \foreach \x in {2.7,7.1} {
    \draw[insarouge, thick] (\x,1.4) -- (\x,0);
    \pic at (\x,1.4) {ctl};
    \pic at (\x,0) {ctl};
  }
  \foreach \x/\lab in {1.9/{\ket{s}}, 3.5/{O_w\ket{s}}, 9.1/{\ket{11}}} {
    \draw[dashed, insagris!70] (\x,2.0) -- (\x,-0.6);
    \node[anchor=south, font=\footnotesize] at (\x,2.0) {$\lab$};
  }
  \draw[decorate, decoration={brace, amplitude=4pt, mirror}, insagris] (2.3,-0.8) -- (3.1,-0.8)
    node[midway, below=5pt, font=\footnotesize] {oracle $O_w$};
  \draw[decorate, decoration={brace, amplitude=4pt, mirror}, insagris] (3.9,-0.8) -- (9.0,-0.8)
    node[midway, below=5pt, font=\footnotesize] {diffusion $D=(H\otimes H)\,\mathrm{CZ}\,(Z\otimes Z)\,(H\otimes H)$};
\end{tikzpicture}}
\caption{Algorithme de Grover pour $N=4$ et $w=11$ : une itération suffit. Les deux points
reliés par un trait vertical forment la porte CZ.}
\label{fig:grover-circ}
\end{figure}

\paragraph{Ordres de grandeur.} Pour $N=2^{20}\approx10^6$, Grover demande
$\frac\pi4\cdot2^{10}\approx804$ itérations, contre $5\times10^5$ évaluations en moyenne. Pour une
recherche de clé AES-128 ($N=2^{128}$), il faut $\frac\pi4\,2^{64}\approx1{,}4\times10^{19}$
itérations \emph{séquentielles}, dont chacune contient un circuit AES complet ; et partager le travail
entre $m$ machines ne fait gagner qu'un facteur $\sqrt m$. L'effet est tout de même de diviser par deux le
nombre de bits de sécurité d'un chiffrement symétrique : on se protège en doublant la taille des
clés (AES-256).

% ---------------------------------------------------------------------
\subsection{Transformée de Fourier quantique et algorithme de Shor}\label{sec:shor}
% ---------------------------------------------------------------------

Le chiffrement RSA repose sur la difficulté de factoriser un entier $N$ (section~\ref{sec:crypto}).
L'algorithme de Shor, conçu en 1994 et publié en 1997 \parencite{shor1997}, factorise $N$ en un temps polynomial en $\log N$.
L'algorithme a deux étages : une réduction \emph{classique} de la factorisation à la recherche d'une
période, puis une recherche \emph{quantique} de cette période par transformée de Fourier.

\subsubsection{Réduction : de la factorisation à la période}

Soit $N$ impair et non puissance d'un nombre premier. On tire $a$ avec $1<a<N$. Si
$\operatorname{pgcd}(a,N)>1$, on tient déjà un facteur. Sinon, la suite $x\mapsto a^x\bmod N$ est
périodique : sa période $r$ est le plus petit entier $r\geq1$ tel que $a^r\equiv1\pmod N$ (l'\emph{ordre}
de $a$ modulo $N$).

\begin{resultat}{Factorisation à partir de l'ordre}
Si $r$ est pair et si $a^{r/2}\not\equiv-1\pmod N$, alors $\operatorname{pgcd}(a^{r/2}-1,N)$ et
$\operatorname{pgcd}(a^{r/2}+1,N)$ sont deux facteurs non triviaux de $N$. Pour $a$ tiré au hasard,
ces deux conditions sont réunies avec une probabilité au moins égale à $\frac12$ \parencite[chap.~5]{nielsen2010}.
\end{resultat}

En effet, posons $u=a^{r/2}$. Alors $(u-1)(u+1)=a^r-1\equiv0\pmod N$ : $N$ divise le produit. Il ne
divise pas $u-1$, sinon l'ordre de $a$ serait $r/2$, ni $u+1$ par hypothèse. Si $\operatorname{pgcd}(u-1,N)$
valait $1$, $N$ diviserait $u+1$ (lemme de Gauss) ; s'il valait $N$, $N$ diviserait $u-1$. Les deux sont exclus, donc
le pgcd est un diviseur strict de $N$ ; de même pour $u+1$. Tout se ramène donc à \emph{trouver $r$}. Classiquement, c'est
aussi difficile que factoriser. C'est la seule étape où l'ordinateur quantique intervient.

\subsubsection{La transformée de Fourier quantique}

\begin{definition}{Transformée de Fourier quantique (QFT)}
Sur l'espace de dimension $M$ de base $\ket{0},\ldots,\ket{M-1}$, avec $\omega=\ee^{2\ii\pi/M}$ :
\[
  \mathrm{QFT}_M\ket{x}=\frac{1}{\sqrt M}\sum_{y=0}^{M-1}\omega^{xy}\ket{y}.
\]
C'est une porte unitaire : ses colonnes sont orthonormées (somme géométrique de raison $\omega^{x-x'}$).
Pour $M=2^m$, elle s'implémente avec $O(m^2)$ portes élémentaires, alors que la transformée de Fourier
rapide classique demande $O(m\,2^m)$ opérations sur les $2^m$ amplitudes.
\end{definition}

Pour $M=2$ on a $\omega=-1$ et $\mathrm{QFT}_2=H$. Pour $M=4$, $\omega=\ii$ ; en écrivant $x=2x_1+x_2$ et
$y=2y_1+y_2$, on a $\ii^{xy}=(-1)^{x_1y_2+x_2y_1}\,\ii^{x_2y_2}$ (car $\ii^4=1$), d'où
\[
  \mathrm{QFT}_4\ket{x_1x_2}=\underbrace{\frac{\ket0+(-1)^{x_2}\ket1}{\sqrt2}}_{=\,H\ket{x_2}\ \text{(chiffre }y_1)}
  \otimes\ \frac{\ket0+(-1)^{x_1}\ii^{x_2}\ket1}{\sqrt2}\ \ \text{(chiffre }y_2).
\]
Le circuit $\mathrm{QFT}_4=\mathrm{SWAP}\,(I\otimes H)\,\mathrm{CS}\,(H\otimes I)$ de la
figure~\ref{fig:qft4} produit exactement cet état : $H$ sur le premier chiffre donne
$\frac{1}{\sqrt2}(\ket0+(-1)^{x_1}\ket1)$, la porte $\mathrm{CS}=\operatorname{diag}(1,1,1,\ii)$ ajoute le
facteur $\ii^{x_2}$ devant $\ket1$, $I\otimes H$ transforme $\ket{x_2}$ en $H\ket{x_2}$, et le SWAP remet les
chiffres dans l'ordre.

\begin{figure}[!ht]
\centering
\begin{tikzpicture}[>=Stealth, font=\small]
  \draw (0,1.2) -- (7.0,1.2);
  \draw (0,0) -- (7.0,0);
  \node[anchor=east] at (0,1.2) {$\ket{x_2}$};
  \node[anchor=east] at (0,0) {$\ket{x_1}$};
  \node[porte] at (1.2,0) {$H$};
  \draw[insarouge, thick] (2.7,1.2) -- (2.7,0);
  \pic at (2.7,1.2) {ctl};
  \node[porte] at (2.7,0) {$S$};
  \node[porte] at (4.2,1.2) {$H$};
  \draw[insarouge, thick] (5.7,1.2) -- (5.7,0);
  \pic at (5.7,1.2) {croix};
  \pic at (5.7,0) {croix};
  \node[anchor=west] at (7.0,1.2) {$\ket{y_2}$};
  \node[anchor=west] at (7.0,0) {$\ket{y_1}$};
\end{tikzpicture}
\caption{$\mathrm{QFT}_4$ : $H$, $S$ contrôlée par l'autre qubit, $H$, puis échange des deux fils.
Le fil du haut est le dernier chiffre.}
\label{fig:qft4}
\end{figure}

\begin{resultat}{La QFT transforme une période en pics}
Si $r$ divise $M$ et si $\ket{\psi}=\sqrt{r/M}\sum_{k=0}^{M/r-1}\ket{x_0+kr}$ (état périodique, de période
$r$), alors
\[
  \mathrm{QFT}_M\ket{\psi}=\frac{1}{\sqrt r}\sum_{j=0}^{r-1}\omega^{x_0\,jM/r}\,\ket{jM/r}.
\]
La mesure donne $y=jM/r$ avec la probabilité $\frac1r$, pour $j=0,\ldots,r-1$, quel que soit le décalage $x_0$.
\end{resultat}

Preuve : l'amplitude de $\ket{y}$ est $\frac{1}{\sqrt M}\sqrt{\frac rM}\,\omega^{x_0y}\sum_{k}\omega^{kry}$.
La somme géométrique de raison $\omega^{ry}=\ee^{2\ii\pi ry/M}$ vaut $M/r$ si $ry/M$ est entier, et $0$ sinon
(la raison est alors différente de $1$ et sa puissance $M/r$ vaut $1$). L'amplitude vaut donc
$\sqrt{r}/M\cdot M/r=1/\sqrt r$ aux multiples de $M/r$, et $0$ ailleurs. Le décalage $x_0$ n'apparaît que dans une
\emph{phase}, qui disparaît à la mesure : c'est ce qui permet de lire la période sans connaître le point de départ.

\subsubsection{L'algorithme}

\begin{resultat}{Algorithme de Shor : recherche de l'ordre de $a$ modulo $N$}
On choisit $M=2^t\geq N^2$ (soit $t\approx2n$ avec $n=\lceil\log_2N\rceil$), un registre de contrôle de $t$ qubits et un
registre de travail de $n$ qubits.
\begin{enumerate}[nosep]
  \item Préparer $\ket{0}^{\otimes t}\ket{1}$ et appliquer $H^{\otimes t}$ au registre de contrôle :
        $\ket{\varphi_1}=\frac{1}{\sqrt M}\sum_{x=0}^{M-1}\ket{x}\ket{1}$.
  \item Appliquer l'exponentiation modulaire $U_a\ket{x}\ket{z}=\ket{x}\ket{a^xz\bmod N}$ :
        $\ket{\varphi_2}=\frac{1}{\sqrt M}\sum_{x}\ket{x}\ket{a^x\bmod N}$.
  \item Appliquer $\mathrm{QFT}_M$ au registre de contrôle et le mesurer : on obtient $y$.
  \item (Classique) Développer $y/M$ en fraction continue : on y lit une fraction $j/r'$ proche de
        $y/M$ avec $r'<N$. Si $\operatorname{pgcd}(j,r)=1$, alors $r'=r$ et l'on vérifie $a^{r'}\equiv1\pmod N$.
        Sinon, $r'$ est un diviseur strict de $r$ et $a^{r'}\not\equiv1$ : on recommence.
\end{enumerate}
\end{resultat}

L'oracle $U_a$ est un produit de $t$ multiplications modulaires \emph{contrôlées} par les qubits $x_j$ de $x$ :
comme $a^x=\prod_j(a^{2^j})^{x_j}$, la $j$-ème porte multiplie le registre de travail par $a^{2^j}\bmod N$, un
nombre calculé à l'avance par un ordinateur classique. Le circuit est celui de la
figure~\ref{fig:shor-circ}.

\begin{figure}[!ht]
\centering
\begin{tikzpicture}[>=Stealth, font=\small]
  % registre de contrôle (haut) et registre de travail (bas)
  \draw (0,1.7) -- (8.3,1.7);
  \draw (0,0) -- (4.6,0);
  \draw[dashed] (4.6,0) -- (8.3,0);
  \node[anchor=east] at (0,1.7) {$\ket{0}^{\otimes t}$};
  \node[anchor=east] at (0,0) {$\ket{1}$};
  \draw (0.5,1.6) -- (0.7,1.8); \node[font=\scriptsize, anchor=south] at (0.6,1.78) {$t$};
  \draw (0.5,-0.1) -- (0.7,0.1); \node[font=\scriptsize, anchor=south] at (0.6,0.08) {$n$};
  \node[porte] at (1.5,1.7) {$H^{\otimes t}$};
  \pic at (3.3,1.7) {ctl};
  \draw[insarouge, thick] (3.3,1.7) -- (3.3,0.4);
  \node[draw=insarouge, fill=insarouge!8, thick, minimum width=17mm, minimum height=8mm] at (3.3,0) {$U_a$};
  \node[porte] at (5.4,1.7) {$\mathrm{QFT}_M$};
  \pic at (7.0,1.7) {mesure};
  \node[anchor=west] at (7.5,1.7) {$y$};
  \node[anchor=north west, text=insagris, font=\footnotesize] at (4.7,-0.3) {registre de travail : ignoré};
  \foreach \x/\lab in {2.3/1, 4.3/2, 6.3/3}
    {\draw[dashed, insagris!70] (\x,2.35) -- (\x,-0.3);
     \node[anchor=south, font=\footnotesize] at (\x,2.35) {$\ket{\varphi_\lab}$};}
\end{tikzpicture}
\caption{Circuit de Shor. Les petits traits obliques indiquent un registre de plusieurs qubits ($t$ pour le
contrôle, $n$ pour le travail). $U_a\ket{x}\ket{z}=\ket{x}\ket{a^xz\bmod N}$ est commandé par le registre de
contrôle.}
\label{fig:shor-circ}
\end{figure}

\paragraph{Pourquoi cela marche.} Après l'étape~2, mesurons (en pensée) le registre de travail : il
donne une valeur $f$ de la suite $a^x\bmod N$, et le registre de contrôle se projette sur les $x$ tels que
$a^x\equiv f$, c'est-à-dire $x\in\{x_0,x_0+r,x_0+2r,\ldots\}$ : un état périodique de période $r$. Le résultat sur la QFT
ci-dessus s'applique : la QFT concentre toute la probabilité sur les $y$ voisins de $jM/r$. La mesure du
registre de travail n'est pas nécessaire (les probabilités de $y$ sont les mêmes si on l'ignore), et le
décalage $x_0$ dépend de $f$ mais ne change pas la loi de $y$. Ainsi $y/M\approx j/r$ pour un entier $j$ \emph{aléatoire} :
on ne choisit pas $j$, mais il est premier avec $r$ dans une fraction $\Omega(1/\log\log r)$ des cas, et quelques
répétitions suffisent. Le choix $M\geq N^2$ garantit qu'une seule fraction de dénominateur inférieur à $N$ est
à distance de $\frac{1}{2M}$ de $y/M$ : le développement en fraction continue l'identifie.

\subsubsection{Exemple complet : $N=15$, $a=7$}

\paragraph{Période.} La suite $7^x\bmod15$ vaut $1,7,4,13,1,7,4,13,\ldots$ : l'ordre est $r=4$. Les
multiplicateurs de $U_a$ sont $a^{2^j}\bmod 15=7,\ 4,\ 1,\ 1$ pour $j=0,1,2,3$ (les deux derniers sont triviaux).
On prend ici $t=4$, soit $M=16$. Ce choix est trop court en général ($M\geq N^2=225$ serait exigé) mais suffit
car $r=4$ divise $M$, ce qui rend le calcul exact.

\paragraph{État avant la QFT.} $\ket{\varphi_2}=\frac14\sum_{x=0}^{15}\ket{x}\ket{7^x\bmod15}$. Si le
registre de travail est mesuré et donne $1$, le registre de contrôle est dans $\frac12(\ket0+\ket4+\ket8+\ket{12})$ ;
s'il donne $7$, dans $\frac12(\ket1+\ket5+\ket9+\ket{13})$, etc.

\paragraph{QFT.} On a $\omega=\ee^{2\ii\pi/16}$ et $\omega^4=\ii$. Pour le décalage $x_0=0$ :
\[
  \mathrm{QFT}_{16}\,\tfrac12\big(\ket0+\ket4+\ket8+\ket{12}\big)
  =\frac18\sum_{y=0}^{15}\Big(\sum_{k=0}^{3}\ii^{\,ky}\Big)\ket{y}
  =\frac12\big(\ket0+\ket4+\ket8+\ket{12}\big),
\]
car $\sum_{k=0}^{3}\ii^{\,ky}$ vaut $4$ si $y$ est multiple de $4$ et $0$ sinon. Pour les autres valeurs de
$f$, seules les phases changent. Dans tous les cas la mesure de $y$ donne $0,4,8$ ou $12$, chacun avec la probabilité $\frac14$
(figure~\ref{fig:shor-n15}).

\paragraph{Lecture classique.} On calcule $y/16$ :
\begin{center}
\renewcommand{\arraystretch}{1.2}
\begin{tabular}{@{}cccp{7.3cm}@{}}
\toprule
$y$ & $y/16$ & fraction irréductible & conclusion\\
\midrule
$0$ & $0$ & $0/1$ & inutile ($j=0$) : on recommence\\
$4$ & $\frac14$ & $1/4$ & $r=4$\ ; $7^4=2401\equiv1$ : correct\\
$8$ & $\frac12$ & $1/2$ & dénominateur $2$ : $7^2=49\equiv4\not\equiv1$, rejeté, on recommence\\
$12$ & $\frac34$ & $3/4$ & $r=4$ : correct\\
\bottomrule
\end{tabular}
\end{center}
Un essai réussit donc avec la probabilité $\frac12$ ($y=4$ ou $12$). Avec $r=4$ pair et $7^{2}=49\equiv4\not\equiv-1\pmod{15}$ :
\[
  \operatorname{pgcd}(7^2-1,15)=\operatorname{pgcd}(48,15)=3,\qquad
  \operatorname{pgcd}(7^2+1,15)=\operatorname{pgcd}(50,15)=5,\qquad 15=3\times5.
\]

\begin{figure}[!ht]
\centering
\adjustbox{max width=\linewidth}{%
\begin{tikzpicture}[>=Stealth, font=\small]
  % (a) suite 7^x mod 15, x = 0..15
  \draw[->, insagris] (-0.2,0) -- (6.7,0) node[right] {$x$};
  \draw[->, insagris] (0,-0.1) -- (0,2.6);
  \foreach \x/\v in {0/1,1/7,2/4,3/13,4/1,5/7,6/4,7/13,8/1,9/7,10/4,11/13,12/1,13/7,14/4,15/13}
    \fill[insarouge] ({\x*0.42},{\v*0.17}) circle (1.7pt);
  \foreach \v in {1,4,7,13}
    {\draw[insagris] (-0.05,{\v*0.17}) -- (0.05,{\v*0.17});
     \node[anchor=east, font=\scriptsize] at (-0.08,{\v*0.17}) {$\v$};}
  \foreach \x in {0,4,8,12}
    {\draw[insagris] ({\x*0.42},-0.05) -- ({\x*0.42},0.05);
     \node[anchor=north, font=\scriptsize] at ({\x*0.42},-0.07) {$\x$};}
  \draw[<->, insarouge] (0,2.45) -- (1.68,2.45) node[midway, above, font=\footnotesize] {période $r=4$};
  \node[font=\footnotesize, anchor=south] at (3.2,3.15) {(a) $f(x)=7^x\bmod15$};
  % flèche QFT
  \draw[->, very thick, insarouge] (7.0,0.8) -- (8.1,0.8) node[midway, above, font=\footnotesize] {$\mathrm{QFT}_{16}$};
  % (b) loi de y : P(y)=1/4 pour y = 0, 4, 8, 12
  \draw[->, insagris] (8.5,0) -- (14.9,0) node[right] {$y$};
  \draw[->, insagris] (8.6,-0.1) -- (8.6,2.6);
  \node[font=\footnotesize, anchor=south] at (11.7,3.15) {(b) loi de la mesure $y$ : $P(y)$};
  \foreach \y in {0,4,8,12}
    {\fill[insarouge] ({8.6+\y*0.4-0.1},0) rectangle ({8.6+\y*0.4+0.1},1.5);
     \node[anchor=north, font=\scriptsize] at ({8.6+\y*0.4},-0.07) {$\y$};}
  \draw[dashed, insagris] (8.55,1.5) -- (14.2,1.5);
  \node[anchor=east, font=\scriptsize] at (8.5,1.5) {$\frac14$};
  \node[text=insagris, font=\footnotesize, align=center] at (11.7,2.35) {$y/16=\frac{j}{4}$, $j=0,1,2,3$};
\end{tikzpicture}}
\caption{Shor pour $N=15$, $a=7$, $t=4$. À gauche, la suite $7^x\bmod15$ a la période $4$. À droite,
la mesure du registre de contrôle après la QFT ne donne que $y\in\{0,4,8,12\}$, avec la probabilité
$\frac14$ chacun : la période $r=4$ se lit dans l'écart $M/r=4$ entre les pics.}
\label{fig:shor-n15}
\end{figure}

\subsubsection{Ce que coûte l'algorithme}

Notons $n=\lceil\log_2N\rceil$ le nombre de bits de $N$ ($n=2048$ pour un module RSA-2048). Le coût s'exprime en
fonction de $n$, pas de $N$.
\begin{align*}
  \text{meilleur algorithme classique connu (crible algébrique, GNFS) :}&\quad
    \exp\!\Big(O\big(n^{1/3}(\log n)^{2/3}\big)\Big),\\
  \text{algorithme de Shor :}&\quad O(n^3)\ \text{portes élémentaires.}
\end{align*}
Le premier temps est sous-exponentiel en $n$ \parencite{lenstra1993}, mais supérieur à tout polynôme. Le second est un
polynôme en $n$ avec la multiplication scolaire ($O(n^2\log n\log\log n)$ avec une multiplication rapide).

\begin{remarque}[Où est le « log » ?]
On résume parfois « de $O(N^p)$ à $O(\log N)$ ». La comparaison correcte est la suivante. La division
d'essai coûte $O(\sqrt N)=O(2^{n/2})$ : exponentielle en $n$. Le crible algébrique est sous-exponentiel en $n$. Shor est
polynomial en $\log N=n\ln2$, soit $O((\log N)^3)$ et non $O(\log N)$ : rien que pour lire $N$ et multiplier
modulo $N$, il faut au moins $n$ opérations. Le « $\log$ » désigne donc la \emph{taille} du nombre.
\end{remarque}

Doubler la taille de la clé multiplie le coût de Shor par $8$, alors qu'il augmente de façon
spectaculaire celui de l'attaque classique. Le tableau compare le coût $n^3$ de Shor à la sécurité
classique que le NIST attribue aux modules RSA \parencite{nist80057}.

\begin{center}
\renewcommand{\arraystretch}{1.2}
\begin{tabular}{@{}rccc@{}}
\toprule
Module RSA $n$ (bits) & sécurité classique (NIST) & $n^3$ (portes, ordre de grandeur) & $\log_2(n^3)$\\
\midrule
$1024$ & $2^{80}$ & $1{,}1\times10^{9}$ & $30{,}0$\\
$2048$ & $2^{112}$ & $8{,}6\times10^{9}$ & $33{,}0$\\
$3072$ & $2^{128}$ & $2{,}9\times10^{10}$ & $34{,}8$\\
$7680$ & $2^{192}$ & $4{,}5\times10^{11}$ & $38{,}7$\\
$15360$ & $2^{256}$ & $3{,}6\times10^{12}$ & $41{,}7$\\
\bottomrule
\end{tabular}
\end{center}

Passer de $2048$ à $3072$ bits fait gagner un facteur $2^{16}=65\,536$ sur l'attaque classique, mais le coût de
Shor n'est multiplié que par $(3/2)^3\approx3{,}4$. Face à un ordinateur quantique, \emph{grossir
les clés RSA ne protège pas} : il faut changer d'algorithme (section~\ref{sec:crypto}). Cela dit, le nombre de portes
n'est pas le vrai coût : il faut des qubits \emph{corrigés d'erreurs}, et c'est là que se trouve le verrou technique
(section~\ref{sec:machine}). Les estimations les plus citées sont les suivantes.
\begin{itemize}[nosep]
  \item \textbf{RSA-2048} : environ $3n+0{,}002\,n\log_2n\approx6\,189$ qubits logiques et
        $0{,}3\,n^3+0{,}0005\,n^3\log_2n\approx2{,}6\times10^9$ portes de Toffoli ; avec des qubits physiques
        d'erreur $10^{-3}$ et un cycle de code de surface de $1\,\mu$s, $20$ millions de qubits physiques pendant
        $8$ heures \parencite{gidney2021}. Une optimisation de 2025 (arithmétique approchée, codes « yoked »,
        « magic state cultivation ») descend à \emph{moins d'un million} de qubits physiques, pour moins d'une
        semaine \parencite{gidney2025}.
  \item \textbf{Courbes elliptiques (ECC-256)}, utilisées dans TLS : au plus $2\,330$ qubits logiques
        \parencite{roetteler2017}. Avec des qubits de chat (section~\ref{sec:machine}), $126\,133$ qubits suffiraient
        pour $9$ heures de calcul \parencite{gouzien2023}. Une \emph{prépublication} de 2026, fondée sur des
        hypothèses optimistes pour les atomes neutres, évoque $10\,000$ à $26\,000$ atomes et quelques
        jours \parencite{cain2026}.
\end{itemize}
Ce sont des estimations, pas des mesures : les démonstrations expérimentales de $15=3\times5$ ou
$21=3\times7$ reposent sur des circuits simplifiés grâce à la connaissance préalable de la réponse
\parencite{smolin2013}. À notre connaissance, aucune n'a factorisé par l'algorithme de Shor un nombre hors de portée d'un
calcul à la main.

% ---------------------------------------------------------------------
\subsection{Cryptographie : RSA, TLS et « stocker maintenant, déchiffrer plus tard »}\label{sec:crypto}
% ---------------------------------------------------------------------

La sécurité des échanges sur Internet repose sur deux familles d'algorithmes. La cryptographie
\emph{symétrique} (AES, ChaCha20) chiffre les données avec une clé secrète partagée. La cryptographie
\emph{asymétrique} (RSA, Diffie--Hellman, courbes elliptiques) sert à \emph{établir} cette clé entre deux
personnes qui ne se sont jamais rencontrées, et à \emph{authentifier} l'interlocuteur par une signature. Les
algorithmes asymétriques s'appuient sur deux problèmes réputés difficiles : la factorisation d'un entier
(RSA) et le logarithme discret (Diffie--Hellman, courbes elliptiques). Ce sont deux problèmes de
\emph{période} dans un groupe : l'algorithme de Shor les résout tous les deux.

\subsubsection{RSA et ce que Shor en fait}

\begin{definition}{Chiffrement RSA}
Alice choisit deux nombres premiers $p\neq q$, pose $n=pq$ et $\varphi(n)=(p-1)(q-1)$, choisit $e$ premier avec
$\varphi(n)$ et calcule $d$ tel que $ed\equiv1\pmod{\varphi(n)}$. La \emph{clé publique} est $(n,e)$, la \emph{clé
privée} $d$. Pour un message $m<n$ :
\[
  c=m^e\bmod n\quad(\text{chiffrement}),\qquad m=c^d\bmod n\quad(\text{déchiffrement}).
\]
Le déchiffrement fonctionne parce que $m^{ed}\equiv m\pmod n$ (théorème d'Euler) \parencite{rivest1978}.
\end{definition}

\paragraph{Exemple jouet.} $p=3$, $q=11$, $n=33$, $\varphi(n)=20$, $e=3$ (premier avec $20$) et $d=7$ (car
$3\times7=21\equiv1\pmod{20}$). Pour $m=4$ : $c=4^3\bmod33=64-33=31$. Le déchiffrement donne
$31^7\bmod33$ ; comme $31\equiv-2$, on a $(-2)^7=-128\equiv-128+132=4=m$ : on retrouve le message.

\paragraph{Attaque.} L'attaquant connaît $(n,e)$. Il lui manque $\varphi(n)$ : le calculer équivaut à
factoriser $n$. Avec $n=33=3\times11$, il obtient $\varphi=20$ puis $d=e^{-1}\bmod20=7$ par l'algorithme d'Euclide
étendu (polynomial). Tout le secret de RSA est donc la difficulté de factoriser. Pour un module de $2048$ bits
(617 chiffres décimaux), les meilleures factorisations classiques publiques (RSA-240, $795$ bits, en 2019
\parencite{boudot2020}) sont loin du compte : on a vu à la section~\ref{sec:shor} que
l'algorithme de Shor l'attaque en $O(n^3)$ portes. Le même algorithme, appliqué au groupe d'une courbe elliptique,
résout le logarithme discret et casse Diffie--Hellman, ECDH et ECDSA.

\subsubsection{Ce que cassent Shor et Grover}

\begin{center}
\renewcommand{\arraystretch}{1.25}
\begin{tabular}{@{}>{\raggedright\arraybackslash}p{4.1cm}>{\raggedright\arraybackslash}p{2.9cm}>{\raggedright\arraybackslash}p{2.5cm}>{\raggedright\arraybackslash}p{5.2cm}@{}}
\toprule
\textbf{Primitive} & \textbf{Rôle} & \textbf{Attaque quantique} & \textbf{Conséquence}\\
\midrule
RSA, Diffie--Hellman, ECDH (X25519) & échange de clé, chiffrement & Shor & \textbf{cassés} (temps polynomial)\\
RSA, ECDSA, EdDSA & signature, certificats & Shor & \textbf{cassés}\\
AES, ChaCha20 & chiffrement symétrique & Grover & clé de $k$ bits : $2^{k/2}$ itérations ; AES-128 affaibli, \textbf{AES-256 sûr}\\
SHA-2, SHA-3 & hachage & Grover (pré-image) & SHA-256 : $2^{128}$, \textbf{sûr}\\
ML-KEM, ML-DSA, SLH-DSA & échange de clé, signature & aucune attaque efficace connue & normes NIST de 2024 (voir plus bas)\\
\bottomrule
\end{tabular}
\end{center}

Les trois premières lignes ne sont pas à égalité. Un algorithme asymétrique est \emph{détruit} par Shor : il faut
le remplacer. Un algorithme symétrique n'est que \emph{ralenti} par Grover (racine carrée, section~\ref{sec:grover}) :
il suffit de doubler la taille de la clé. « Aucune attaque connue » n'est pas une preuve de sécurité, d'où la prudence
des agences décrite plus bas.

\subsubsection{TLS 1.3 : où sont les points faibles ?}

Le protocole TLS 1.3 (RFC~8446 \parencite{rfc8446}) protège HTTPS. Il combine les trois mécanismes : un
\emph{échange de clés} éphémère (ECDHE, le plus souvent X25519 : chaque côté tire un secret $a$ ou $b$ et
publie $g^a$ ou $g^b$, d'où le secret partagé $g^{ab}$), une \emph{authentification} du serveur par signature
(certificat RSA ou ECDSA, vérifié par une chaîne d'autorités de certification, elle aussi faite de signatures),
et un \emph{chiffrement symétrique} des données (AES-GCM ou ChaCha20-Poly1305) avec des clés dérivées de $g^{ab}$
(figure~\ref{fig:tls}).

\begin{figure}[!ht]
\centering
\adjustbox{max width=\linewidth}{%
\begin{tikzpicture}[>=Stealth, font=\small,
    tag/.style={rounded corners=2pt, font=\footnotesize, inner sep=3pt, anchor=west, text width=4.6cm, align=left}]
  \node[draw=insarouge, fill=insarouge!8, thick, minimum width=20mm] (C) at (0,0) {Client};
  \node[draw=insarouge, fill=insarouge!8, thick, minimum width=20mm] (S) at (6.2,0) {Serveur};
  \draw[insagris, thick] (0,-0.4) -- (0,-6.2);
  \draw[insagris, thick] (6.2,-0.4) -- (6.2,-6.2);
  \draw[->, thick] (0,-1.0) -- (6.2,-1.0) node[midway, above, font=\footnotesize] {ClientHello + clé publique $g^a$};
  \draw[<-, thick] (0,-1.9) -- (6.2,-1.9) node[midway, above, font=\footnotesize] {ServerHello + clé publique $g^b$};
  \node[font=\footnotesize, text=insagris, anchor=east] at (-0.2,-2.35) {calcule $g^{ab}$, dérive les clés};
  \node[font=\footnotesize, text=insagris, anchor=west] at (6.4,-2.35) {calcule $g^{ab}$};
  \draw[<-, thick] (0,-3.2) -- (6.2,-3.2) node[midway, above, font=\footnotesize, align=center]
       {\{Certificat, signature (CertificateVerify),\\ Finished\}};
  \draw[->, thick] (0,-4.1) -- (6.2,-4.1) node[midway, above, font=\footnotesize] {\{Finished\}};
  \draw[<->, very thick, insarouge] (0,-5.2) -- (6.2,-5.2) node[midway, above, font=\footnotesize] {données chiffrées (AES-GCM, ChaCha20)};
  \node[tag, fill=insarouge!18] at (7.8,-1.45) {\textbf{Échange de clés} : cassé par Shor (logarithme discret)};
  \node[tag, fill=insarouge!18] at (7.8,-3.2) {\textbf{Signature et certificats} : cassés par Shor (RSA, ECDSA)};
  \node[tag, fill=insagris!20] at (7.8,-5.2) {\textbf{Chiffrement symétrique} : \mbox{Grover} seulement, AES-256 suffit};
\end{tikzpicture}}
\caption{Poignée de main TLS 1.3 simplifiée. Les accolades \{\,\} indiquent des messages déjà chiffrés avec les clés
issues de $g^{ab}$. Deux des trois briques tombent face à Shor ; la troisième résiste.}
\label{fig:tls}
\end{figure}

\begin{remarque}[La confidentialité persistante ne protège pas contre l'ordinateur quantique]
Les secrets $a$ et $b$ sont détruits après la session : si la clé privée du serveur est volée plus tard, les sessions
passées ne peuvent pas être déchiffrées (propriété dite de \emph{forward secrecy}). Mais les valeurs publiques
$g^a$ et $g^b$ ont circulé en clair. Un adversaire qui les a enregistrées peut, avec l'algorithme de Shor (logarithme
discret sur la courbe), retrouver $a$ puis $g^{ab}$, donc toutes les clés de la session. La forward secrecy protège
contre le vol d'une \emph{clé}, pas contre l'effondrement de la difficulté du \emph{problème mathématique}.
\end{remarque}

\subsubsection{« Stocker maintenant, déchiffrer plus tard »}

La menace n'attend pas l'existence de la machine. Un adversaire qui capte le trafic chiffré aujourd'hui (il lui
suffit d'écouter une liaison, ce qui est discret et peu coûteux) peut le \emph{conserver} et attendre qu'un
ordinateur quantique assez puissant existe pour le déchiffrer. C'est l'attaque \emph{Harvest Now, Decrypt Later}
(HNDL, ou « stocker maintenant, déchiffrer plus tard »). Elle vise la \emph{confidentialité de long terme} :
dossiers médicaux et génétiques, secrets industriels, secrets d'État, archives juridiques. Elle ne vise pas un
jeton de session valable dix minutes. On ne sait pas qui enregistre quoi ; il faut supposer que cela se fait.

\begin{resultat}{Inégalité de Mosca}
Soit $X$ la durée pendant laquelle les données doivent rester secrètes, $Y$ le temps nécessaire pour migrer vers une
cryptographie résistante, et $Z$ le délai avant qu'un ordinateur quantique capable de casser la cryptographie actuelle existe
(tous en années \parencite{mosca2018}). Les données sont \emph{déjà exposées} si
\[
  X+Y>Z.
\]
\end{resultat}

Les données chiffrées juste avant la fin de la migration (à la date $Y$) doivent encore être secrètes à la date $X+Y$ ;
si $Z$ tombe avant, l'adversaire les lit. Exemple (figure~\ref{fig:mosca}) : des données médicales à protéger
$X=20$ ans, une migration de $Y=7$ ans, et $Z=15$ ans : $X+Y=27>15$, avec douze années d'exposition. On ne connaît pas
$Z$ (ce chapitre ne cherche pas à le prédire), mais on connaît $X$ et $Y$ : plus ils sont grands, plus il faut
commencer tôt, et le coût d'une migration trop précoce est faible devant celui d'une divulgation irréversible.

\begin{figure}[!ht]
\centering
\begin{tikzpicture}[>=Stealth, font=\small, x=0.3cm]
  \draw[->, insagris] (0,0) -- (31,0) node[right] {années};
  \foreach \t in {0,5,...,30}
    {\draw[insagris] (\t,0.08) -- (\t,-0.08);
     \node[anchor=north, font=\scriptsize] at (\t,-0.1) {$\t$};}
  \fill[insagris!35] (0,0.5) rectangle (7,1.1);
  \fill[insarouge!18] (7,0.5) rectangle (27,1.1);
  \fill[insarouge!60] (15,0.5) rectangle (27,1.1);
  \node[font=\footnotesize] at (3.5,0.8) {$Y=7$};
  \node[font=\footnotesize] at (11,0.8) {$X=20$};
  \draw[insarouge, very thick] (15,0) -- (15,2.0);
  \node[anchor=south, font=\footnotesize, text=insarouge] at (15,2.0) {$Z=15$ : ordinateur quantique opérationnel};
  \node[anchor=west, font=\footnotesize, text=white] at (15.4,0.8) {exposé 12 ans};
  \node[anchor=north west, font=\footnotesize, text=insagris] at (0,-0.65)
    {$Y$ : migration (la cryptographie vulnérable est encore utilisée)};
  \node[anchor=north west, font=\footnotesize, text=insagris] at (0,-1.15)
    {$X$ : durée de confidentialité exigée pour les données chiffrées à la fin de la migration};
\end{tikzpicture}
\caption{Inégalité de Mosca pour $X=20$, $Y=7$, $Z=15$ ans. La zone rouge foncé est la période où des données
interceptées sont lisibles par l'adversaire alors qu'elles devraient rester secrètes.}
\label{fig:mosca}
\end{figure}

\begin{remarque}[Signatures : un risque différent]
L'attaque HNDL ne concerne que la confidentialité. Une signature ne sert qu'au moment où on la vérifie : falsifier en 2035 un
certificat TLS expiré ne change rien à une session de 2026. Les signatures à longue durée de vie (micrologiciels,
racines de confiance, documents juridiques) devront néanmoins migrer, mais plus tard : c'est l'échange de clés qui est
urgent.
\end{remarque}

\subsubsection{Première réponse : la cryptographie post-quantique}

La \emph{cryptographie post-quantique} (PQC) regroupe des algorithmes qui s'exécutent sur des ordinateurs
\emph{classiques} mais reposent sur des problèmes contre lesquels on ne connaît aucune attaque quantique efficace :
réseaux euclidiens (problème \emph{Learning With Errors}), fonctions de hachage, codes correcteurs d'erreurs. Le NIST a
publié le 13 août 2024 trois normes \parencite{fips203,fips204,fips205}, puis sélectionné un quatrième algorithme en 2025 \parencite{nist2025hqc}.

\begin{center}
\renewcommand{\arraystretch}{1.2}
\begin{tabular}{@{}llll@{}}
\toprule
\textbf{Norme} & \textbf{Nom} & \textbf{Usage} & \textbf{Problème sous-jacent}\\
\midrule
FIPS 203 & ML-KEM (ex-Kyber) & échange de clé & réseaux (module-LWE)\\
FIPS 204 & ML-DSA (ex-Dilithium) & signature & réseaux\\
FIPS 205 & SLH-DSA (ex-SPHINCS+) & signature & fonctions de hachage\\
(à venir) & HQC & échange de clé (second choix) & codes correcteurs\\
\bottomrule
\end{tabular}
\end{center}

\paragraph{Hybridation.} Dans TLS 1.3, la solution déployée est \emph{hybride} : le client envoie à la fois une clé
publique X25519 et une clé publique ML-KEM-768, et les deux secrets sont concaténés dans la dérivation des clés. L'adversaire
doit casser \emph{les deux} : X25519 (que Shor casse, mais que l'on connaît bien) et ML-KEM (que Shor ne casse pas, mais
que l'on connaît depuis moins longtemps). Le groupe \texttt{X25519MLKEM768} est normalisé par la RFC~10024 d'août 2026
\parencite{rfc10024}, et il est activé par défaut dans les navigateurs depuis fin 2024 (Chrome 131, Firefox 132). Le
prix est la taille : $32$ octets pour une clé X25519, mais $1\,184$ octets de clé publique et $1\,088$ octets de
chiffré pour ML-KEM-768.

\paragraph{Pourquoi pas ML-KEM seul ?} Les algorithmes post-quantiques sont plus jeunes et moins étudiés. En 2022, le
candidat SIKE, arrivé au quatrième tour du NIST, a été cassé par une \emph{attaque classique} : une dizaine de minutes sur un seul
cœur d'ordinateur suffisent pour son plus petit jeu de paramètres \parencite{castryck2023}. L'ANSSI recommande donc de \emph{combiner} cryptographie classique et
post-quantique pendant la transition \parencite{anssi2024}.

\paragraph{Calendrier.} Le projet de document NIST~IR~8547 (version initiale publique de novembre 2024, qui peut
avoir évolué) prévoit de déprécier RSA-2048 et les algorithmes de sécurité équivalente après 2030, et d'interdire tous les
algorithmes vulnérables après 2035 \parencite{nistir8547}. La migration est longue : il faut inventorier où la cryptographie est
utilisée (logiciels, matériel, certificats, protocoles), puis la remplacer. L'« agilité cryptographique », c'est-à-dire
la capacité de changer d'algorithme sans refondre un système, est la qualité qui coûtera le moins cher.

\subsubsection{Seconde réponse : la distribution quantique de clé}

La \emph{distribution quantique de clé} (QKD) ne cherche pas un problème difficile : elle utilise la physique. Dans le
protocole BB84 \parencite{bennett1984}, Alice envoie à Bob des qubits dont l'état est choisi au hasard parmi quatre :
$\ket0,\ket1$ (base $Z$) ou $\ket+,\ket-$ (base $X$). Bob mesure chaque qubit dans une base $Z$ ou $X$ choisie au hasard.
Ensuite, Alice et Bob annoncent \emph{publiquement leurs bases} (pas les bits) et gardent les positions où les bases
coïncident : si la base est la bonne, Bob lit exactement le bit d'Alice ; sinon son résultat est aléatoire. Exemple, sans espion :

\begin{center}
\renewcommand{\arraystretch}{1.15}
\begin{tabular}{@{}l*{8}{c}@{}}
\toprule
position & 1 & 2 & 3 & 4 & 5 & 6 & 7 & 8\\
\midrule
bit d'Alice & 0 & 1 & 1 & 0 & 1 & 0 & 0 & 1\\
base d'Alice & $Z$ & $X$ & $X$ & $Z$ & $X$ & $Z$ & $X$ & $Z$\\
base de Bob & $Z$ & $Z$ & $X$ & $X$ & $X$ & $Z$ & $Z$ & $Z$\\
\midrule
bases identiques ? & oui & non & oui & non & oui & oui & non & oui\\
clé tamisée & 0 & & 1 & & 1 & 0 & & 1\\
\bottomrule
\end{tabular}
\end{center}

La clé tamisée est $01101$ : en moyenne la moitié des qubits est conservée. Pour détecter un espion Ève, Alice et Bob
comparent publiquement un échantillon de la clé tamisée. Une attaque « intercepter puis renvoyer » est détectée parce que
\emph{mesurer perturbe} (postulat de la mesure) et que l'on ne peut pas copier un état inconnu pour mesurer la copie
(théorème de non-clonage, section~\ref{sec:teleportation}). Ève choisit la mauvaise base une fois sur deux ; dans ce cas
le qubit renvoyé est dans la mauvaise base, et Bob, s'il a la bonne base, se trompe une fois sur deux. Le taux d'erreur de la clé tamisée
(\emph{QBER}) vaut donc
\[
  \mathrm{QBER}=\underbrace{\tfrac12}_{\text{Ève se trompe de base}}\times\underbrace{\tfrac12}_{\text{Bob se trompe}}=\tfrac14=25\,\%.
\]
En pratique, quelques pour cent d'erreurs viennent du bruit ; au-delà d'un seuil de l'ordre de $11\,\%$, on abandonne la clé. La
variante E91 \parencite{ekert1991} remplace l'envoi de qubits par des paires intriquées et teste l'absence d'espion par la violation
de l'inégalité CHSH (section~\ref{sec:chsh}).

\begin{remarque}[Pourquoi les agences recommandent la PQC plutôt que la QKD]
En janvier 2024, l'ANSSI, le BSI allemand et leurs homologues néerlandais et suédois ont publié une position commune
\parencite{anssi2024qkd} : la QKD n'est pas assez mûre pour la sécurité, ne convient qu'à des cas d'usage restreints, et la priorité est donnée à
une technologie « beaucoup plus compatible avec les systèmes cryptographiques existants ». Elle demande des liaisons optiques
dédiées de portée limitée, des nœuds de confiance, une authentification classique des messages publics (qui nécessite
elle-même de la cryptographie), et ses implémentations ont des failles de matériel. La QKD protège donc un lien, la PQC un
protocole entier.
\end{remarque}

% ---------------------------------------------------------------------
\subsection{Construire un ordinateur quantique}\label{sec:machine}
% ---------------------------------------------------------------------

Les sections précédentes supposent des qubits parfaits. Les qubits réels sont bruités, et c'est le bruit, bien plus que le
nombre de portes, qui sépare les machines actuelles des estimations de la section~\ref{sec:shor}. Cette section
explique ce qu'il faut pour fabriquer un qubit, pourquoi certains tournent dans un réfrigérateur à
$10\ \mathrm{mK}$, comment on corrige les erreurs, et ce que proposent les qubits de chat d'Alice \& Bob.

\subsubsection{Que faut-il ? Les critères de DiVincenzo}

\textcite{divincenzo2000} a résumé les conditions nécessaires au calcul quantique :
\begin{enumerate}[nosep]
  \item des qubits bien définis, en nombre extensible ;
  \item une initialisation fiable dans l'état $\ket{0\ldots0}$ ;
  \item une cohérence longue devant la durée d'une porte ;
  \item un jeu universel de portes : portes à un qubit et au moins une porte à deux qubits comme le CNOT (chapitre~\ref{chap:portes}) ;
  \item une mesure individuelle des qubits.
\end{enumerate}
S'y ajoutent, pour les réseaux, deux critères de communication. Le critère~3 (cohérence) est le plus difficile à tenir quand le nombre de qubits augmente.

\subsubsection{La décohérence}

Un qubit n'est jamais seul. Supposons un qubit $S$ dans l'état $\frac{1}{\sqrt2}(\ket0+\ket1)$ qui interagit
avec son environnement $E$ de façon que celui-ci « enregistre » l'état de $S$ : $\ket0\ket{E_0}$ et $\ket1\ket{E_1}$.
L'état global est $\frac{1}{\sqrt2}(\ket0\ket{E_0}+\ket1\ket{E_1})$, intriqué. L'environnement n'est pas observé : l'état de $S$ seul
est obtenu par trace partielle (section~\ref{sec:trace}),
\[
  \rho_S=\frac12\begin{pmatrix}1&\braket{E_1}{E_0}\\\braket{E_0}{E_1}&1\end{pmatrix}.
\]
Si les deux états de l'environnement sont orthogonaux, $\rho_S=\frac12 I$ : l'état est maximalement mixte
(section~\ref{sec:pur-mixte}) et la superposition est perdue. En pratique le recouvrement décroît
comme $\braket{E_1}{E_0}=\ee^{-t/T_2}$, d'où, pour la pureté,
\[
  \operatorname{Tr}\rho_S^2=\frac{1+\ee^{-2t/T_2}}{2}\ \xrightarrow[t\to\infty]{}\ \frac12 .
\]
(à $t=T_2$, elle vaut $0{,}57$). Les termes diagonaux, eux, évoluent avec le temps de relaxation $T_1$ (perte d'énergie : $\ket1\to\ket0$). Le rapport
entre $T_2$ et la durée $t_{\mathrm{porte}}$ d'une porte fixe le nombre d'opérations avant que l'information soit perdue. Avec
$T_2=100\ \mu\mathrm{s}$ et des portes de $50\ \mathrm{ns}$, il vaut $2\,000$ : l'erreur par porte est de l'ordre de
$t_{\mathrm{porte}}/T_2=5\times10^{-4}$. Calculer sur des milliards de portes sans corriger les erreurs est
donc impossible.

\subsubsection{Les plateformes}

Plusieurs systèmes physiques servent de qubits. Chacun a ses forces et ses faiblesses ; les valeurs du tableau sont des ordres de
grandeur tirés de la littérature récente et évoluent vite.

\begin{center}
\renewcommand{\arraystretch}{1.2}
\small
\begin{tabular}{@{}>{\raggedright\arraybackslash}p{2.7cm}>{\raggedright\arraybackslash}p{2.2cm}>{\raggedright\arraybackslash}p{1.9cm}>{\raggedright\arraybackslash}p{1.9cm}>{\raggedright\arraybackslash}p{5.5cm}@{}}
\toprule
\textbf{Plateforme} & \textbf{Température} & \textbf{Porte à 2 qubits} & \textbf{Cohérence} & \textbf{Atouts / difficultés}\\
\midrule
Circuits supra\-conducteurs \parencite{kjaergaard2020} & $10$--$20\ \mathrm{mK}$ & $20$--$100\ \mathrm{ns}$ & $\sim0{,}1\ \mathrm{ms}$ & portes rapides, fabrication par lithographie ; cryogénie, câblage, voisinage limité\\
Ions piégés \parencite{bruzewicz2019} & ambiante (vide poussé) & $10\ \mu\mathrm{s}$--$1\ \mathrm{ms}$ & secondes à minutes & meilleures fidélités ($\geq99{,}9\,\%$), connexion totale ; portes lentes, passage à grande échelle\\
Atomes neutres \parencite{saffman2016} & ambiante (vide poussé) & $0{,}1$--$1\ \mu\mathrm{s}$ & $\sim1\ \mathrm{s}$ & réseaux de milliers d'atomes reconfigurables ; erreurs de porte plus élevées\\
Photons \parencite{kok2007} & ambiante (détecteurs $\sim2\ \mathrm{K}$) & $\mathrm{ns}$ & pas de décohérence en vol & communication ; pertes, portes à deux photons probabilistes\\
Spins dans le silicium \parencite{burkard2023} & $0{,}1$--$1\ \mathrm{K}$ & $10$--$100\ \mathrm{ns}$ & $\sim\mathrm{ms}$ & compatible avec la microélectronique ; peu de qubits à ce jour\\
\bottomrule
\end{tabular}
\end{center}

Aucune plateforme ne domine. Les machines actuelles comptent de quelques dizaines à quelques milliers de qubits physiques
\emph{bruités}, sans correction d'erreurs : \textcite{preskill2018} les a appelées « NISQ » (\emph{noisy intermediate-scale quantum}). Elles ne suffisent pas pour Shor.

\subsubsection{Le cryostat : pourquoi $10\ \mathrm{mK}$ ?}

Un transmon (qubit supraconducteur) a une fréquence de transition $f\approx5\ \mathrm{GHz}$ : une micro-onde. À l'équilibre thermique
à la température $T$, la loi de Boltzmann donne $p_1/p_0=\ee^{-hf/k_BT}$ avec $p_0+p_1=1$, d'où
\[
  p_1=\frac{1}{1+\ee^{\,hf/k_BT}},\qquad
  \frac{hf}{k_B}=\frac{6{,}626\times10^{-34}\times5\times10^{9}}{1{,}381\times10^{-23}}\approx0{,}24\ \mathrm{K}.
\]
Le qubit n'est « dans l'état $\ket0$ » que si $T\ll0{,}24\ \mathrm{K}$. Sinon il est dans un mélange statistique de
$\ket0$ et $\ket1$, de matrice densité $\operatorname{diag}(1-p_1,p_1)$ (section~\ref{sec:rho}) :

\begin{center}
\renewcommand{\arraystretch}{1.15}
\begin{tabular}{@{}rccc@{}}
\toprule
Température $T$ & $hf/k_BT$ & population $p_1$ de $\ket1$ & pureté $\operatorname{Tr}\rho^2$\\
\midrule
$300\ \mathrm{K}$ (ambiante) & $0{,}0008$ & $0{,}4998$ & $0{,}500$\\
$4\ \mathrm{K}$ (hélium liquide) & $0{,}060$ & $0{,}485$ & $0{,}500$\\
$1\ \mathrm{K}$ & $0{,}24$ & $0{,}44$ & $0{,}507$\\
$100\ \mathrm{mK}$ & $2{,}4$ & $0{,}083$ & $0{,}85$\\
$20\ \mathrm{mK}$ & $12$ & $6{,}2\times10^{-6}$ & $0{,}99999$\\
$10\ \mathrm{mK}$ & $24$ & $3{,}8\times10^{-11}$ & $1{,}000$\\
\bottomrule
\end{tabular}
\end{center}

À température ambiante, $p_1\approx\frac12$ : l'état est maximalement mixte (pureté minimale d'un qubit,
$\operatorname{Tr}\rho^2=\frac12$, section~\ref{sec:pur-mixte}) et le qubit ne porte aucune information. Il faut descendre vers $20\ \mathrm{mK}$ pour que
l'état initial soit pur à mieux que $10^{-5}$ près. S'y ajoute une seconde raison : les jonctions Josephson sont en aluminium, supraconducteur
seulement sous $\approx1{,}2\ \mathrm{K}$, et les photons thermiques des lignes de commande doivent être éteints. Les autres
plateformes évitent ce problème autrement : ions et atomes sont préparés dans $\ket0$ par pompage optique (un laser les force dans l'état voulu)
et refroidis par laser ; les photons ont des fréquences optiques ($\sim10^{14}\ \mathrm{Hz}$, soit $hf/k_B\approx5\times10^3\ \mathrm{K}$),
sans aucune excitation thermique.

\paragraph{Le réfrigérateur à dilution.} On atteint $10\ \mathrm{mK}$ avec un mélange d'isotopes d'hélium, $^3$He et $^4$He. Sous
$0{,}87\ \mathrm{K}$, le mélange se sépare en deux phases, l'une riche en $^3$He, l'autre diluée. Faire passer des atomes de $^3$He de
la phase riche à la phase diluée \emph{absorbe de la chaleur}, comme une évaporation : c'est le principe de la dilution. Un
refroidisseur à tube pulsé fournit le froid initial (étages à $50\ \mathrm{K}$ puis $4\ \mathrm{K}$), et le circuit de dilution
descend en étages jusqu'à la \emph{chambre de mélange} (figure~\ref{fig:cryostat}). Le prix est la \emph{puissance
froide} : de l'ordre de $10\ \mu\mathrm{W}$ à $20\ \mathrm{mK}$ pour un appareil commercial, ce qui limite la chaleur que les lignes et
l'électronique peuvent apporter à cet étage.

\begin{figure}[!ht]
\centering
\adjustbox{max width=\linewidth}{%
\begin{tikzpicture}[>=Stealth, font=\small,
    plaque/.style={fill=insagris!35, draw=insagris, thick, minimum height=3.5mm, inner sep=0pt},
    legende/.style={font=\footnotesize, anchor=west, text width=7.2cm, align=left}]
  % plaques : (hauteur, largeur, température)
  \node[plaque, minimum width=6.4cm] (p0) at (0,0)    {};
  \node[plaque, minimum width=5.6cm] (p1) at (0,-1.3) {};
  \node[plaque, minimum width=4.8cm] (p2) at (0,-2.6) {};
  \node[plaque, minimum width=4.0cm] (p3) at (0,-3.9) {};
  \node[plaque, minimum width=3.2cm] (p4) at (0,-5.2) {};
  \node[plaque, minimum width=2.4cm, fill=insarouge!30, draw=insarouge] (p5) at (0,-6.5) {};
  \foreach \a/\b in {0/-1.3, -1.3/-2.6, -2.6/-3.9, -3.9/-5.2, -5.2/-6.5}
    {\draw[insagris, thick] (-0.9,\a-0.175) -- (-0.9,\b+0.175);
     \draw[insagris, thick] (0.9,\a-0.175) -- (0.9,\b+0.175);}
  % puce
  \node[draw=insarouge, fill=insarouge!15, thick, minimum width=14mm, minimum height=5mm] (puce) at (0,-7.15) {\footnotesize puce};
  \draw[insarouge, thick] (0,-6.675) -- (0,-6.9);
  % ligne de commande (coaxial) avec atténuateurs
  \draw[insarouge, thick] (0.3,0.35) -- (0.3,-6.9);
  \foreach \y/\att in {-2.6/{$-20$ dB}, -5.2/{$-20$ dB}, -6.5/{$-20$ dB}}
    {\node[draw=insarouge, fill=white, inner sep=1.8pt] at (0.3,\y) {};
     \node[font=\scriptsize, text=insarouge, anchor=west, fill=white, inner sep=1pt] at (1.3,\y) {\att};}
  % températures (gauche)
  \foreach \y/\T in {0/{$300\ \mathrm{K}$}, -1.3/{$50\ \mathrm{K}$}, -2.6/{$4\ \mathrm{K}$}, -3.9/{$\approx1\ \mathrm{K}$}, -5.2/{$\approx100\ \mathrm{mK}$}, -6.5/{$\approx10\ \mathrm{mK}$}}
    \node[anchor=east, font=\footnotesize\bfseries] at (-3.4,\y) {\T};
  % légendes (droite)
  \node[legende] at (3.7,0)    {électronique de commande et de lecture, ambiante};
  \node[legende] at (3.7,-1.3) {1\up{er} étage du tube pulsé, ancrage thermique des câbles};
  \node[legende] at (3.7,-2.6) {2\up{e} étage du tube pulsé ; amplificateur HEMT (lecture)};
  \node[legende] at (3.7,-3.9) {« still » du circuit de dilution};
  \node[legende] at (3.7,-5.2) {plaque froide : filtres et atténuateurs};
  \node[legende] at (3.7,-6.5) {chambre de mélange : qubits, amplificateur paramétrique};
\end{tikzpicture}}
\caption{Réfrigérateur à dilution en coupe schématique (« lustre »). Les plaques sont de plus en plus froides vers le
bas ; les atténuateurs placés sur la ligne de commande (valeurs typiques) ramènent le bruit thermique de $300\ \mathrm{K}$ à
celui de l'étage où ils sont fixés.}
\label{fig:cryostat}
\end{figure}

\paragraph{Le goulot d'étranglement : le câblage.} Chaque qubit demande plusieurs lignes coaxiales (commande, lecture),
qui apportent de la chaleur depuis les étages chauds et du bruit thermique. Des atténuateurs ancrés à chaque étage
ramènent ce bruit : $T_{\mathrm{eff}}\approx T_{\text{étage}}+T_{\text{chaud}}\,10^{-A/10}$, donc $60\ \mathrm{dB}$ au total limitent à
$300\ \mathrm{K}\times10^{-6}=0{,}3\ \mathrm{mK}$ la contribution du bruit de l'étage ambiant \parencite{krinner2019}. Mais des machines de millions de
qubits physiques ne tiennent pas dans un cryostat avec un coaxial par qubit : les pistes de recherche sont le multiplexage des lignes,
l'électronique de commande placée à $4\ \mathrm{K}$ et le raccordement de plusieurs cryostats entre eux.

\subsubsection{Corriger les erreurs : du code à répétition au code de surface}

Classiquement, on se protège du bruit en copiant chaque bit trois fois et en prenant le vote majoritaire. Le théorème de
non-clonage interdit de copier $\alpha\ket0+\beta\ket1$, et la mesure détruirait la superposition. On \emph{intrique} plutôt
le qubit avec deux autres.

\begin{definition}{Code à répétition à trois qubits}
Le qubit logique $\alpha\ket0+\beta\ket1$ est encodé par deux CNOT dont le contrôle est $q_1$ (et les cibles $q_2$, $q_3$) :
\[
  (\alpha\ket0+\beta\ket1)\ket{00}\ \longmapsto\ \alpha\ket{000}+\beta\ket{111}.
\]
Ce n'est pas $\ket\psi^{\otimes3}$ : l'état est intriqué. Dans cette section, les kets s'écrivent $\ket{q_1q_2q_3}$ et les
fils sont étiquetés : la convention « fil du haut = dernier chiffre » ne s'applique pas. Une erreur $X_k$ (inversion de bit sur $q_k$) est détectée en mesurant les
\emph{parités} $Z_1Z_2$ et $Z_2Z_3$, de valeurs propres $\pm1$.
\end{definition}

Sur $\alpha\ket{000}+\beta\ket{111}$ les deux parités valent $+1$. Une inversion fait changer le signe des parités qui contiennent le qubit touché :

\begin{center}
\renewcommand{\arraystretch}{1.15}
\begin{tabular}{@{}lccc@{}}
\toprule
erreur & état obtenu & $(Z_1Z_2,\ Z_2Z_3)$ & correction\\
\midrule
aucune & $\alpha\ket{000}+\beta\ket{111}$ & $(+1,+1)$ & aucune\\
$X_1$ & $\alpha\ket{100}+\beta\ket{011}$ & $(-1,+1)$ & $X_1$\\
$X_2$ & $\alpha\ket{010}+\beta\ket{101}$ & $(-1,-1)$ & $X_2$\\
$X_3$ & $\alpha\ket{001}+\beta\ket{110}$ & $(+1,-1)$ & $X_3$\\
\bottomrule
\end{tabular}
\end{center}

Les deux parités sont identiques pour $\ket{000}$ et $\ket{111}$ : leur mesure localise l'erreur \emph{sans rien apprendre sur
$\alpha$ et $\beta$}, donc sans détruire la superposition. On mesure les parités avec deux qubits
auxiliaires (\emph{ancillas}) $a_1$ et $a_2$ : le CNOT inscrit dans l'ancilla le ou-exclusif des deux qubits qu'il contrôle
(figure~\ref{fig:code3}).

\begin{figure}[!ht]
\centering
\adjustbox{max width=\linewidth}{%
\begin{tikzpicture}[>=Stealth, font=\small]
  \foreach \y in {3.6, 2.4, 1.2} \draw (0,\y) -- (11.3,\y);
  \foreach \y in {0, -1.2} \draw (0,\y) -- (8.7,\y);
  \node[anchor=east] at (0,3.6) {$q_1:\ \alpha\ket0+\beta\ket1$};
  \node[anchor=east] at (0,2.4) {$q_2:\ \ket0$};
  \node[anchor=east] at (0,1.2) {$q_3:\ \ket0$};
  \node[anchor=east] at (0,0) {$a_1:\ \ket0$};
  \node[anchor=east] at (0,-1.2) {$a_2:\ \ket0$};
  % encodage : CNOT q1 -> q2, q1 -> q3
  \draw[insarouge, thick] (1.1,3.6) -- (1.1,2.4);
  \pic at (1.1,3.6) {ctl}; \pic at (1.1,2.4) {cible};
  \draw[insarouge, thick] (2.0,3.6) -- (2.0,1.2);
  \pic at (2.0,3.6) {ctl}; \pic at (2.0,1.2) {cible};
  % erreur éventuelle
  \node[draw=insagris, dashed, minimum width=9mm, minimum height=9mm] at (3.3,2.4) {$X$};
  \node[font=\footnotesize, text=insagris, anchor=south] at (3.3,3.0) {erreur ?};
  % syndrome : a1 <- q1 xor q2 ; a2 <- q2 xor q3
  \draw[insarouge, thick] (4.5,3.6) -- (4.5,0);  \pic at (4.5,3.6) {ctl}; \pic at (4.5,0) {cible};
  \draw[insarouge, thick] (5.4,2.4) -- (5.4,0);  \pic at (5.4,2.4) {ctl}; \pic at (5.4,0) {cible};
  \draw[insarouge, thick] (6.3,2.4) -- (6.3,-1.2); \pic at (6.3,2.4) {ctl}; \pic at (6.3,-1.2) {cible};
  \draw[insarouge, thick] (7.2,1.2) -- (7.2,-1.2); \pic at (7.2,1.2) {ctl}; \pic at (7.2,-1.2) {cible};
  \pic at (8.4,0) {mesure};   \pic at (8.4,-1.2) {mesure};
  \node[anchor=west, font=\footnotesize] at (9.0,-0.6) {syndrome $(s_1,s_2)$};
  \node[porte] at (10.0,3.6) {$X^{?}$}; \node[porte] at (10.0,2.4) {$X^{?}$}; \node[porte] at (10.0,1.2) {$X^{?}$};
  \draw[dashed, insagris] (8.4,0.25) -- (8.4,0.5) -- (10.0,0.5) -- (10.0,0.85);
  \foreach \x/\t in {2.65/{encodage}, 6.0/{extraction du syndrome}}
    \node[font=\footnotesize, text=insagris] at (\x,-2.1) {\t};
\end{tikzpicture}}
\caption{Code à répétition à trois qubits : encodage, erreur éventuelle (ici sur $q_2$), extraction du syndrome
$(s_1,s_2)=(q_1\oplus q_2,\,q_2\oplus q_3)$ dans les ancillas, puis correction par une porte $X$ sur le qubit désigné
par le syndrome. Les kets s'écrivent $\ket{q_1q_2q_3}$.}
\label{fig:code3}
\end{figure}

\paragraph{Gain.} La correction échoue si au moins deux qubits sont inversés. Avec une probabilité $p$ d'inversion par
qubit, la probabilité d'échec est $3p^2(1-p)+p^3=3p^2-2p^3$, inférieure à $p$ si et seulement si $(2p-1)(p-1)>0$,
c'est-à-dire $p<\frac12$. Le code est utile \emph{si le bruit est assez faible}. C'est le principe d'un
\emph{seuil}.

\paragraph{Erreurs de phase.} Une erreur $Z_1$ sur $\alpha\ket{000}+\beta\ket{111}$ donne $\alpha\ket{000}-\beta\ket{111}$ : un état
logique différent, que ce code ne voit pas. On se protège en travaillant dans la base $X$ : $H^{\otimes3}(\alpha\ket{000}+\beta\ket{111})
=\alpha\ket{+++}+\beta\ket{---}$, où une erreur $Z$ (qui échange $\ket+$ et $\ket-$) joue le rôle d'une inversion de
bit. \textcite{shor1995} a montré que l'on peut \emph{concaténer} les deux codes (neuf qubits) pour corriger n'importe quelle erreur sur un
qubit : toute erreur est une combinaison de $I,X,Y,Z$ et la mesure du syndrome l'y projette.

\paragraph{Code de surface.} Les qubits sont disposés sur un réseau plan de $d\times d$ qubits de donnée, entourés de qubits
de mesure ; les syndromes de type $X$ et de type $Z$ sont mesurés à chaque cycle. La \emph{distance} $d$ est le nombre minimal
d'erreurs qui changent l'état logique : le code corrige jusqu'à $\lfloor(d-1)/2\rfloor$ erreurs, et il faut
$2d^2-1$ qubits physiques par qubit logique. Si le taux d'erreur physique est sous le seuil ($\approx1\,\%$), l'erreur
logique \emph{décroît exponentiellement avec $d$} : à chaque passage de $d$ à $d+2$, elle est divisée par un
facteur $\Lambda>1$. Google a observé ce comportement sur des codes de distance 3, 5 puis 7, avec
$\Lambda=2{,}14\pm0{,}02$ et une erreur logique de $0{,}143\,\%$ par cycle à $d=7$ (101 qubits physiques)
\parencite{acharya2025}. En prolongeant avec ce même $\Lambda$, atteindre une erreur logique de $10^{-12}$ par cycle demande
$\Lambda^{k}\geq1{,}43\times10^{-3}/10^{-12}$, soit $k\geq\ln(1{,}43\times10^{9})/\ln(2{,}14)\approx27{,}7$, donc
$k=28$ et $d=7+2\times28=63$ : près de $8\,000$ qubits physiques par qubit logique. C'est une extrapolation (l'erreur n'est pas
exactement exponentielle), mais elle explique les millions de qubits de la section~\ref{sec:shor} : le surcoût de la
correction d'erreurs domine tout, et un meilleur $\Lambda$ (des qubits physiques moins bruités) le réduit fortement.

\subsubsection{Les qubits de chat d'Alice \& Bob}

Le code de surface doit corriger \emph{deux} types d'erreurs ($X$ et $Z$) sur un réseau plan, d'où un surcoût en $d^2$. L'idée des
\emph{qubits de chat} \parencite{mirrahimi2014} est de concevoir un qubit dont une des deux erreurs est déjà supprimée par la physique. Il ne reste
alors qu'un type d'erreur à corriger : un code \emph{unidimensionnel} (la répétition de la section précédente), de $2d-1$ qubits
au lieu de $2d^2-1$, suffit.

\paragraph{Un qubit dans un oscillateur.} On utilise un mode d'un résonateur micro-onde, oscillateur harmonique d'états de
Fock $\ket n$ ($n$ photons). Un \emph{état cohérent} $\ket\alpha=\ee^{-|\alpha|^2/2}\sum_n\frac{\alpha^n}{\sqrt{n!}}\ket n$ est vecteur propre de l'opérateur
d'annihilation : $a\ket\alpha=\alpha\ket\alpha$, avec $|\alpha|^2$ photons en moyenne. Pour $\alpha$ réel, on pose
\[
  \ket{0_L}\approx\ket{\alpha},\qquad\ket{1_L}\approx\ket{-\alpha},\qquad
  \braket{\alpha}{-\alpha}=\ee^{-|\alpha|^2-|\alpha|^2}=\ee^{-2|\alpha|^2}.
\]
Le recouvrement est $1{,}8\times10^{-2}$ pour $|\alpha|^2=2$, $3{,}4\times10^{-4}$ pour $|\alpha|^2=4$ et $1{,}1\times10^{-7}$ pour
$|\alpha|^2=8$ : les deux états sont presque orthogonaux dès que $|\alpha|^2\gtrsim4$. Les états de la base $X$ sont les
\emph{états de chat} de Schrödinger
\[
  \ket{C_\pm}=N_\pm\big(\ket\alpha\pm\ket{-\alpha}\big),\qquad N_\pm=\frac{1}{\sqrt{2\,(1\pm\ee^{-2|\alpha|^2})}},\qquad
  \ket{\pm_L}\approx\ket{C_\pm}.
\]
$\ket{C_+}$ ne contient que des nombres \emph{pairs} de photons et $\ket{C_-}$ que des nombres \emph{impairs} (les termes en
$\alpha^n$ et $(-\alpha)^n$ s'annulent ou s'ajoutent selon la parité de $n$).

\paragraph{Une perte de photon est une erreur de phase.} Le bruit principal d'un résonateur est la perte d'un photon, décrite par $a$.
Comme $a\ket{\pm\alpha}=\pm\alpha\ket{\pm\alpha}$,
\[
  a\ket{C_\pm}=\alpha N_\pm\big(\ket\alpha\mp\ket{-\alpha}\big)=\alpha\,\frac{N_\pm}{N_\mp}\,\ket{C_\mp}.
\]
Après normalisation, $\ket{C_+}\leftrightarrow\ket{C_-}$ : c'est $\ket{+_L}\leftrightarrow\ket{-_L}$, une \emph{inversion de phase} $Z$. Sur
$\ket{\pm\alpha}$, $a$ ne change rien (un simple facteur), donc la perte de photon ne provoque aucune inversion de bit. Elle survient au taux
$\kappa_1\langle n\rangle=\kappa_1|\alpha|^2$ : le taux d'inversion de phase \emph{croît linéairement} avec $|\alpha|^2$.

\paragraph{Supprimer les inversions de bit.} On ajoute une dissipation \emph{à deux photons} : le circuit éjecte des paires de photons
(taux $\kappa_2$) et en réinjecte par paires (pompe d'amplitude $\epsilon_2$). Cela équivaut à un saut d'opérateur
$\sqrt{\kappa_2}\,(a^2-\alpha^2)$ avec $|\alpha|^2$ fixé par le rapport pompe/perte. Or $(a^2-\alpha^2)\ket{\pm\alpha}=0$ car
$a^2\ket{\pm\alpha}=\alpha^2\ket{\pm\alpha}$ : l'espace $\mathrm{Vect}(\ket\alpha,\ket{-\alpha})$ est exactement le noyau de la dissipation. Elle ramène
tout état de l'oscillateur vers cet espace et le \emph{stabilise}. Retirer les photons deux par deux conserve la parité du nombre de photons, donc la
cohérence entre $\ket{C_+}$ et $\ket{C_-}$. Pour qu'un bruit transforme $\ket\alpha$ en $\ket{-\alpha}$, il doit franchir la distance $2\alpha$ dans
le plan de phase contre la dissipation, ce qui est exponentiellement improbable : le taux d'inversion de bit décroît
comme $\ee^{-2|\alpha|^2}$.

\begin{figure}[!ht]
\centering
\begin{tikzpicture}[>=Stealth, font=\small]
  \draw[->, insagris] (-4.6,0) -- (4.8,0) node[right] {$\mathrm{Re}\,\alpha$};
  \foreach \sg in {1,-1}
    \shade[inner color=insarouge!80, outer color=insarouge!5] (\sg*2.6,0) circle (0.75);
  \node[anchor=north, font=\footnotesize] at (2.6,-0.85) {$\ket{\alpha}=\ket{0_L}$};
  \node[anchor=north, font=\footnotesize] at (-2.6,-0.85) {$\ket{-\alpha}=\ket{1_L}$};
  \draw[<->, dashed, insagris, thick] (-1.8,0) -- (1.8,0);
  \node[fill=white, inner sep=1pt, text=insarouge] at (0,0) {$\times$};
  \node[anchor=north, font=\footnotesize, text=insagris, align=center] at (0,-0.2)
    {inversion de bit : $\sim\ee^{-2|\alpha|^2}$\\ (supprimée)};
  \draw[->, insarouge, thick] (-1.5,0.75) to[bend left=40] (1.5,0.75);
  \node[anchor=south, font=\footnotesize, text=insarouge, align=center] at (0,1.3)
    {photon perdu : $\ket{C_+}\leftrightarrow\ket{C_-}$\\ (inversion de phase, $\propto\kappa_1|\alpha|^2$)};
\end{tikzpicture}
\caption{Qubit de chat dans le plan de phase de l'oscillateur (axe réel) : deux états cohérents $\ket{\pm\alpha}$ représentent $\ket{0_L}$
et $\ket{1_L}$. La dissipation à deux photons les stabilise ; passer de l'un à l'autre (inversion de bit) est exponentiellement
rare, alors que la perte d'un photon échange les états de chat $\ket{C_\pm}\approx\ket{\pm_L}$, de parités opposées (inversion de phase).}
\label{fig:chat}
\end{figure}

\paragraph{Résultat : un biais de bruit énorme.} On augmente $|\alpha|^2$ pour supprimer les inversions de bit, au prix d'inversions de phase
plus fréquentes. Mais l'asymétrie est telle que le compromis est favorable. Les expériences sont de plus en plus convaincantes :
des équipes parisiennes (ENS, Inria, Alice \& Bob) ont observé la suppression exponentielle de l'inversion de bit
\parencite{lescanne2020}, puis un temps d'inversion de bit supérieur à $10\ \mathrm{s}$, pour un temps d'inversion de phase supérieur à
$490\ \mathrm{ns}$ \parencite{reglade2024}. Le rapport est de l'ordre de $2\times10^{7}$ ($10/(0{,}49\times10^{-6})$). Une \emph{annonce de l'entreprise} Alice \& Bob (septembre 2025, résultat préliminaire, non évalué par des pairs à
notre connaissance) fait état de plus d'une heure avec $\langle n\rangle=11$ \parencite{aliceandbob2025}.

\paragraph{De là à un ordinateur.} Il reste à corriger les inversions de phase, par un code à répétition de $d$ qubits de chat
\parencite{guillaud2019}, dont les portes à deux qubits (CNOT) doivent \emph{préserver le biais} : c'est le point que les expériences
doivent encore démontrer. Une estimation d'architecture évalue le calcul d'un logarithme discret sur une courbe de 256 bits à
$126\,133$ qubits de chat et $9$ heures \parencite{gouzien2023}. Alice \& Bob (start-up française) et Amazon \parencite{putterman2025} suivent cette voie ;
Alice \& Bob vise, dans sa \emph{feuille de route} d'entreprise, une machine de 100 qubits logiques en 2030. C'est un objectif commercial, pas un résultat.

% ---------------------------------------------------------------------
\subsection{Autres applications et algorithmes}\label{sec:autres}
% ---------------------------------------------------------------------

Shor et Grover sont les algorithmes les plus connus, pas les seuls. Il faut distinguer les gains \emph{prouvés} (on sait démontrer
qu'aucun algorithme classique ne fait aussi bien, ou on connaît un bien meilleur algorithme quantique), les gains
\emph{conditionnels} (valables sous hypothèses) et les méthodes \emph{heuristiques} (qui marchent peut-être, sans preuve).

\subsubsection{Panorama}

\begin{center}
\renewcommand{\arraystretch}{1.2}
\small
\begin{tabular}{@{}>{\raggedright\arraybackslash}p{3.1cm}>{\raggedright\arraybackslash}p{3.8cm}>{\raggedright\arraybackslash}p{3.0cm}>{\raggedright\arraybackslash}p{4.6cm}@{}}
\toprule
\textbf{Algorithme} & \textbf{Problème} & \textbf{Gain} & \textbf{Statut et usage}\\
\midrule
Deutsch--Jozsa, Bernstein--Vazirani, Simon (1985--1997) & fonctions « boîtes noires » & exponentiel (Simon) & prouvé ; peu d'usage direct, base de Shor\\
Shor (1994) \parencite{shor1997} & factorisation, logarithme discret & superpolynomial & prouvé ; casse RSA et courbes elliptiques\\
Grover (1996) & recherche non structurée & quadratique : $\sqrt N$ & prouvé, optimal ; clés symétriques, SAT\\
Estimation de phase \parencite{kitaev1995} & valeur propre d'une porte $U$ & précision $2^{-m}$ avec $m$ qubits & prouvé ; brique de Shor, de la chimie, de HHL\\
Estimation d'amplitude \parencite{brassard2002} & estimer une probabilité & $\varepsilon^{-1}$ au lieu de $\varepsilon^{-2}$ & prouvé ; Monte-Carlo, finance\\
Simulation hamiltonienne \parencite{lloyd1996} & évolution $\ee^{-\ii Ht}\ket\psi$ & exponentiel (systèmes quantiques) & prouvé ; chimie, matériaux : application la plus naturelle\\
HHL \parencite{harrow2009} & système linéaire $Ax=b$ & exponentiel \emph{sous conditions} & conditionnel ; en partie « déquantisé » \parencite{tang2019}\\
VQE \parencite{peruzzo2014} & énergie de l'état fondamental & aucun gain prouvé & heuristique ; machines bruitées\\
QAOA \parencite{farhi2014} & optimisation combinatoire & aucun gain prouvé & heuristique ; machines bruitées\\
\bottomrule
\end{tabular}
\end{center}

\subsubsection{Simuler la matière : l'idée de Feynman}

Pour décrire $n$ qubits, un ordinateur classique stocke $2^n$ amplitudes complexes. Avec $16$ octets par amplitude, $n=50$
demande $2^{50}\times16=2^{54}$ octets, soit $16$ pébioctets ($1{,}8\times10^{16}$ octets) ; $n=300$ demande $2^{300}\approx2\times10^{90}$
amplitudes, plus que d'atomes dans l'univers observable ($\sim10^{80}$). \textcite{feynman1982} en a tiré l'idée
inverse : un système quantique simule naturellement un autre système quantique.

La \emph{simulation hamiltonienne} calcule $\ee^{-\ii Ht}\ket{\psi}$ pour $H=A+B$, où $A$ et $B$ sont des sommes de termes simples dont on sait implémenter les exponentielles.
Si $A$ et $B$ ne commutent pas, $\ee^{-\ii(A+B)t}\neq\ee^{-\ii At}\ee^{-\ii Bt}$. La formule de Baker--Campbell--Hausdorff donne, pour un pas $\tau$,
\[
  \ee^{-\ii A\tau}\ee^{-\ii B\tau}=\exp\!\Big(-\ii(A+B)\tau-\tfrac{\tau^2}{2}[A,B]+O(\tau^3)\Big).
\]
On coupe donc $t$ en $N$ pas $\tau=t/N$ (\emph{formule de Trotter}) :
\[
  \Big\|\ee^{-\ii(A+B)t}-\big(\ee^{-\ii At/N}\ee^{-\ii Bt/N}\big)^N\Big\|\ \leq\ \frac{t^2}{2N}\,\big\|[A,B]\big\|\quad(\text{à l'ordre dominant}),
\]
l'erreur de chaque pas ($\sim\tau^2\|[A,B]\|/2$) s'additionnant sur $N$ pas. Pour $A=X$ et $B=Z$ (un qubit), $[X,Z]=-2\ii Y$ et
$\|[X,Z]\|=2$ : l'erreur vaut au plus $t^2/N$ ; pour $t=1$, $N=100$ pas donnent une erreur inférieure à $10^{-2}$ (le calcul numérique donne $7\times10^{-3}$). Les applications visées sont la chimie
(énergie d'une molécule comme le cofacteur FeMoco de la nitrogénase, qui intervient dans la fabrication des engrais \parencite{reiher2017}),
les matériaux (batteries, supraconducteurs) et la physique des particules. Les méthodes classiques (théorie de la fonctionnelle de
la densité, \emph{coupled cluster}) traitent bien les systèmes faiblement corrélés ; c'est pour les systèmes \emph{fortement
corrélés} qu'un avantage est espéré, avec des machines tolérantes aux fautes.

\subsubsection{Méthodes variationnelles et optimisation}

Sur les machines bruitées actuelles, on ne peut pas faire tourner de longs circuits ; on utilise donc des circuits courts à paramètres
$\theta$ dans une boucle avec un ordinateur classique : le circuit prépare $\ket{\psi(\theta)}$, on mesure l'énergie
$E(\theta)=\bra{\psi(\theta)}H\ket{\psi(\theta)}$, et un optimiseur classique corrige $\theta$. Cela repose sur le
\emph{principe variationnel} : en décomposant $\ket\psi=\sum_kc_k\ket{k}$ sur les états propres de $H$ (section~\ref{sec:spectral}),
$E=\sum_k|c_k|^2E_k\geq E_0\sum_k|c_k|^2=E_0$. Minimiser $E(\theta)$ approche donc l'énergie fondamentale $E_0$ par le haut. Exemple pour un qubit,
$H=X+Z$ (valeurs propres $\pm\sqrt2$) avec $\ket{\psi(\theta)}=\cos\frac\theta2\ket0+\sin\frac\theta2\ket1$ :
\[
  E(\theta)=\langle Z\rangle+\langle X\rangle=\cos\theta+\sin\theta=\sqrt2\cos\!\big(\theta-\tfrac\pi4\big),
\]
minimale pour $\theta=\frac{5\pi}{4}$, de valeur $-\sqrt2=E_0$. Le VQE \parencite{peruzzo2014} généralise cette idée à des molécules, et le QAOA
\parencite{farhi2014} l'applique à des problèmes d'optimisation en alternant deux types de portes. Aucun avantage n'est démontré pour ces méthodes :
elles sont l'objet de la recherche actuelle.

\subsubsection{Algèbre linéaire quantique : une promesse sous conditions}

L'algorithme HHL \parencite{harrow2009} prépare un état $\ket{x}\propto\sum_ix_i\ket{i}$ où $Ax=b$, en un temps polynomial en $\log N$
là où la résolution classique est polynomiale en $N$. Le gain exponentiel apparent est soumis à plusieurs conditions. Il faut
\emph{charger} $b$ dans l'état $\ket{b}$ efficacement ; $A$ doit être creuse et bien conditionnée ; et le résultat est un
\emph{état} : en lire les $N$ composantes demanderait déjà $N$ répétitions, de sorte que l'on n'obtient que des
grandeurs globales de $x$. De plus, des algorithmes classiques « inspirés du quantique » (\emph{déquantisation}) obtiennent le même
résultat sous des hypothèses d'accès analogues \parencite{tang2019}, ce qui a supprimé l'avantage exponentiel annoncé pour
certaines applications d'apprentissage. Le bilan honnête : \emph{pas d'avantage démontré en apprentissage automatique à ce jour}.

\subsubsection{Capteurs, nombres aléatoires, réseaux}

\paragraph{Métrologie quantique.} Les premières technologies quantiques commercialisées sont des capteurs (horloges atomiques,
magnétomètres, gravimètres). L'intrication y améliore la précision. Soit $N$ qubits qui subissent chacun la phase
$\ket1\mapsto\ee^{\ii\varphi}\ket1$ à mesurer. Dans l'état de GHZ $\frac{1}{\sqrt2}(\ket{0\ldots0}+\ket{1\ldots1})$ (généralisation à $N$ qubits de l'état de Bell
$\frac{1}{\sqrt2}(\ket{00}+\ket{11})$, section~\ref{sec:bell}), seuls $\ket{0\ldots0}$ et $\ket{1\ldots1}$ interviennent, et l'état devient
\[
  \frac{1}{\sqrt2}\big(\ket{0\ldots0}+\ee^{\ii N\varphi}\ket{1\ldots1}\big),\qquad
  P\big(X^{\otimes N}=+1\big)=\frac{1+\cos N\varphi}{2}=\cos^2\!\frac{N\varphi}{2}.
\]
La phase est amplifiée par $N$ : la précision atteint $\Delta\varphi\sim1/N$ (limite de Heisenberg) au lieu de $1/\sqrt N$ pour
$N$ qubits indépendants (limite quantique standard) \parencite{giovannetti2004}. En pratique, la décohérence réduit ce gain, qui
se limite souvent à de petits $N$.

\paragraph{Nombres aléatoires.} Mesurer $\ket+$ dans la base $Z$ donne $0$ ou $1$ avec la probabilité $\frac12$ : le postulat de Born
garantit un hasard \emph{intrinsèque}, contrairement à un générateur pseudo-aléatoire déterministe. Les générateurs quantiques
de nombres aléatoires (QRNG), à base de photons, sont déjà commercialisés ; ils alimentent par exemple la génération de clés cryptographiques.

\paragraph{Réseaux quantiques.} Relier des ordinateurs quantiques demande de transmettre des qubits (téléportation, section~\ref{sec:teleportation}).
Une fibre optique standard atténue de $0{,}2\ \mathrm{dB/km}$, soit une transmission $T=10^{-0{,}02L}$ ($L$ en km) :
$T=10^{-2}$ à $100\ \mathrm{km}$ et $10^{-10}$ à $500\ \mathrm{km}$. Un amplificateur est impossible (non-clonage) : on utilise des
\emph{répéteurs quantiques}, qui étendent l'intrication par \emph{échange d'intrication} (une mesure de Bell entre deux paires intriquées, comme dans la
téléportation de la section~\ref{sec:teleportation}), ou des satellites : la distribution
d'intrication sur $1\,200\ \mathrm{km}$ a été démontrée par le satellite chinois Micius \parencite{yin2017}. L'objectif est un « internet
quantique » \parencite{wehner2018} reliant des processeurs, des capteurs et des liens de QKD.

\subsubsection{Jalons et niveaux de preuve}

Le tableau place quelques jalons dans le temps et indique la \emph{nature de la source}, car une annonce d'entreprise ou une
prépublication n'a pas la valeur d'un article évalué par des pairs ou d'une norme.

\begin{center}
\renewcommand{\arraystretch}{1.2}
\small
\begin{tabular}{@{}>{\raggedright\arraybackslash}p{1.9cm}>{\raggedright\arraybackslash}p{8.8cm}>{\raggedright\arraybackslash}p{3.3cm}@{}}
\toprule
\textbf{Date} & \textbf{Jalon} & \textbf{Nature de la source}\\
\midrule
1994--1996 & algorithmes de Shor et de Grover \parencite{shor1997} & articles évalués\\
2019 & Google, $53$ qubits : échantillonnage de circuits aléatoires en $200\ \mathrm{s}$, annoncé hors de portée des machines classiques \parencite{arute2019} (contesté ensuite) & article évalué\\
août 2024 & normes NIST FIPS 203, 204, 205 \parencite{fips203} & norme\\
2024 & qubit de chat : inversion de bit $>10\ \mathrm{s}$ \parencite{reglade2024} & article évalué\\
déc. 2024 & code de surface sous le seuil, $\Lambda=2{,}14$ \parencite{acharya2025} & article évalué\\
mars 2025 & HQC choisi par le NIST \parencite{nist2025hqc} & rapport officiel\\
mai 2025 & RSA-2048 : moins d'un million de qubits \parencite{gidney2025} & prépublication\\
sept. 2025 & qubit de chat : plus d'une heure d'inversion de bit \parencite{aliceandbob2025} & annonce d'entreprise\\
mars 2026 & Shor avec $10\,000$ à $26\,000$ atomes \parencite{cain2026} & prépublication\\
août 2026 & RFC~10024 : TLS hybride X25519 + ML-KEM \parencite{rfc10024} & norme IETF\\
\bottomrule
\end{tabular}
\end{center}

Ce qui est \emph{acquis} : les gains théoriques de Shor et de Grover, la correction d'erreurs sous le seuil sur de petits codes,
et les standards post-quantiques. Ce qui \emph{ne l'est pas} : un ordinateur à plusieurs milliers de qubits \emph{logiques},
un avantage utile sur un problème pratique, et surtout la date $Z$ de l'inégalité de Mosca. C'est pourquoi la migration cryptographique se
décide sur $X+Y$, qui se connaît, et non sur une prédiction de $Z$.

\clearpage
% ---------------------------------------------------------------------
\subsection{Récapitulatif du chapitre}\label{sec:recap-appli}
% ---------------------------------------------------------------------

\begin{center}
\adjustbox{max width=\linewidth}{%
\begin{tikzpicture}[>=Stealth, font=\footnotesize,
    boite/.style={draw=insarouge, fill=insarouge!6, rounded corners=2pt, thick,
                  minimum width=4.4cm, minimum height=1.45cm, align=center, inner sep=3pt},
    lien/.style={->, very thick, insagris},
    etiq/.style={font=\scriptsize\itshape, text=insagris, inner sep=2pt}]
  \node[boite] (int) at (0,0)
    {\textbf{Interférence}\\ oracle, Deutsch, Grover : $\sqrt N$\\ {\scriptsize\S\,\ref{sec:interference}, \ref{sec:grover}}};
  \node[boite] (sho) at (6.3,0)
    {\textbf{Shor}\\ QFT, période $r$, $O(n^3)$\\ {\scriptsize\S\,\ref{sec:shor}}};
  \node[boite] (cry) at (12.6,0)
    {\textbf{Cryptographie}\\ RSA et ECC cassés ; HNDL ; PQC\\ {\scriptsize\S\,\ref{sec:crypto}}};
  \node[boite] (cryo) at (12.6,-2.3)
    {\textbf{Qubits physiques}\\ décohérence, cryostat $10\ \mathrm{mK}$\\ {\scriptsize\S\,\ref{sec:machine}}};
  \node[boite] (qec) at (6.3,-2.3)
    {\textbf{Correction d'erreurs}\\ seuil, code de surface, qubits de chat\\ {\scriptsize\S\,\ref{sec:machine}}};
  \node[boite] (aut) at (0,-2.3)
    {\textbf{Autres usages}\\ simulation, VQE, capteurs, réseaux\\ {\scriptsize\S\,\ref{sec:autres}}};
  \draw[lien] (int) -- (sho) node[etiq, midway, above=1pt] {période};
  \draw[lien] (sho) -- (cry) node[etiq, midway, above=1pt] {casse RSA};
  \draw[lien] (cry) -- (cryo) node[etiq, midway, right=1pt] {quelle machine ?};
  \draw[lien] (cryo) -- (qec) node[etiq, midway, above=1pt] {qubits fiables};
  \draw[lien] (qec) -- (aut) node[etiq, midway, above=1pt] {usages};
\end{tikzpicture}}
\end{center}

\begin{aretenir}
\begin{itemize}
  \item Un algorithme quantique organise des \emph{interférences} : l'oracle écrit $f(x)$ dans une phase ; Deutsch décide avec un appel au lieu de deux.
        Grover : $\approx\frac\pi4\sqrt N$ itérations (rotations de $2\theta$), gain quadratique et optimal ; exact en une itération pour $N=4$.
  \item Shor : la QFT transforme la période $r$ de $a^x\bmod N$ en pics en $y=jM/r$ ; $\operatorname{pgcd}(a^{r/2}\pm1,N)$ donne les facteurs ;
        coût $O(n^3)$ ($n=\log_2N$), contre un temps sous-exponentiel classique (exemple : $N=15$, $a=7$, $r=4$).
  \item RSA, Diffie--Hellman, ECDH et ECDSA tombent face à Shor ; AES-256 et SHA-256 résistent à Grover. Doubler les clés RSA ne protège pas.
  \item « Stocker maintenant, déchiffrer plus tard » : exposé dès que $X+Y>Z$ (Mosca) ; la forward secrecy n'aide pas. On migre d'abord l'échange de clés
        (hybride X25519 + ML-KEM), puis les signatures (FIPS 203--205).
  \item Qubit physique : la décohérence ($T_2$) mélange l'état ; le cryostat à $10\ \mathrm{mK}$ rend $p_1$ négligeable ($hf/k_B=0{,}24\ \mathrm{K}$ à $5\ \mathrm{GHz}$). Correction
        d'erreurs : seuil, $2d^2-1$ qubits physiques par qubit logique (code de surface).
  \item Qubit de chat : $\ket{\pm\alpha}$, perte de photon = inversion de phase, inversion de bit en $\ee^{-2|\alpha|^2}$ ; code à répétition. Une annonce d'entreprise
        n'est pas un résultat évalué.
  \item Autres usages : simulation (Trotter), VQE et QAOA (heuristiques), HHL (conditionnel), capteurs (limite de Heisenberg), réseaux (répéteurs).
\end{itemize}
\end{aretenir}

\begin{center}
\renewcommand{\arraystretch}{1.12}
\small
\begin{tabular}{@{}l l r@{}}
\toprule
\textbf{Notion} & \textbf{Formule} & \textbf{§} \\
\midrule
Rebond de phase & $U_f\ket{x}\ket{-}=(-1)^{f(x)}\ket{x}\ket{-}$ & \ref{sec:interference} \\
Grover & $P=\sin^2\!\big((2k+1)\theta\big)$, \quad $\sin\theta=\frac{1}{\sqrt N}$, \quad $k\simeq\frac\pi4\sqrt N$ & \ref{sec:grover} \\
QFT et période & $\mathrm{QFT}_M\ket x=\frac{1}{\sqrt M}\sum_y\ee^{2\ii\pi xy/M}\ket y$ ; pics en $y=jM/r$, $P=\frac1r$ & \ref{sec:shor} \\
Facteurs, coût & $\operatorname{pgcd}(a^{r/2}\pm1,N)$, \quad $O(n^3)$ & \ref{sec:shor} \\
Mosca & $X+Y>Z$ : données déjà exposées & \ref{sec:crypto} \\
Population thermique & $p_1=\frac{1}{1+\ee^{hf/k_BT}}$, \quad $\operatorname{Tr}\rho^2=\frac{1+\ee^{-2t/T_2}}{2}$ (décohérence) & \ref{sec:machine} \\
Qubit de chat & $\braket{\alpha}{-\alpha}=\ee^{-2|\alpha|^2}$, \quad $a\ket{C_\pm}=\alpha\frac{N_\pm}{N_\mp}\ket{C_\mp}$ & \ref{sec:machine} \\
\bottomrule
\end{tabular}
\end{center}

\begin{savoirfaire}
\begin{itemize}[nosep]
  \item Suivre un circuit de Deutsch ou de Grover ($N=4$) jusqu'à la mesure (sections~\ref{sec:interference}, \ref{sec:grover}).
  \item Trouver l'ordre de $a$ modulo $N$, appliquer le critère du pgcd et prévoir la loi de $y$ (section~\ref{sec:shor}, exemple $N=15$).
  \item Dire ce que casse Shor, ce que Grover affaiblit, et appliquer l'inégalité de Mosca (section~\ref{sec:crypto}).
  \item Calculer $p_1$, la pureté après décohérence et le recouvrement $\ee^{-2|\alpha|^2}$ ; lire un syndrome (section~\ref{sec:machine}).
  \item Distinguer résultat prouvé, conditionnel et annonce (section~\ref{sec:autres}).
\end{itemize}
\end{savoirfaire}
, entre chapitre5 et
%  chapitre6 (exercices). Il ne contient ni préambule ni \begin{document} : les
%  environnements (definition, resultat, remarque, aretenir, savoirfaire), les
%  macros (\ii, \ee, \dd, \ket, \bra, \braket), les symboles de circuits (ctl,
%  cible, croix, mesure, porte) et les couleurs (insarouge, insagris) sont définis
%  dans main.tex. Un chapitre = une \section.
%  Le numéro d'un chapitre est donné par l'ordre des \include, pas par le nom du
%  fichier : ce chapitre est le chapitre 6, les exercices (chapitre6.tex) le 7.
% =====================================================================

% Macros de Dirac (déjà définies dans main.tex ; ce bloc évite l'erreur
% « Undefined control sequence \ket » si main.tex est une ancienne version).
\providecommand{\ket}[1]{|#1\rangle}
\providecommand{\bra}[1]{\langle #1|}
\providecommand{\braket}[2]{\langle #1|#2\rangle}

% =====================================================================
\section{Applications et algorithmes quantiques}\label{chap:applications}
% =====================================================================

Les chapitres précédents ont construit le formalisme : qubits, portes, intrication, mesure. Ce
chapitre montre à quoi il sert. On y voit comment un circuit exploite l'interférence pour résoudre
certains problèmes plus vite qu'un ordinateur classique (sections~\ref{sec:interference}
à~\ref{sec:shor}), pourquoi cela menace une partie de la cryptographie actuelle, RSA et TLS en tête
(section~\ref{sec:crypto}), ce qu'il faut pour construire une telle machine, du cryostat aux qubits
de chat d'Alice \& Bob (section~\ref{sec:machine}), et quelles autres applications et quels autres
algorithmes existent (section~\ref{sec:autres}). Les circuits sont dessinés comme aux chapitres~4 et~5 :
le fil du haut est le dernier chiffre du ket.

\begin{remarque}[Chiffres d'actualité]
Les estimations de ressources, les calendriers de migration et les records expérimentaux cités ici
datent d'octobre 2026 et évoluent vite. Chaque chiffre est accompagné de sa source (bibliographie) ;
les annonces d'entreprises sont signalées comme telles et les prépublications aussi.
\end{remarque}

% ---------------------------------------------------------------------
\subsection{Calculer par interférence}\label{sec:interference}
% ---------------------------------------------------------------------

Un registre de $n$ qubits est décrit par $2^n$ amplitudes (chapitre~\ref{chap:operateurs}). Une porte
agit sur toutes ces amplitudes à la fois, mais la mesure n'en révèle qu'un seul résultat, tiré au
hasard avec les probabilités de Born. Dire que «~le circuit calcule toutes les valeurs en parallèle~»
est donc trompeur. Un algorithme quantique est une \emph{chorégraphie d'interférences} : les
amplitudes des mauvaises réponses se compensent (signes ou phases opposés), celles de la bonne
réponse s'additionnent, et la mesure finale donne la bonne réponse avec une probabilité proche de $1$.

\begin{definition}{Oracle}
Pour $f:\{0,1\}^n\to\{0,1\}$, l'\emph{oracle} de $f$ est la porte à $n+1$ qubits
\[
  U_f\ket{x}\ket{y}=\ket{x}\ket{y\oplus f(x)}.
\]
Elle est unitaire et $U_f^2=\mathrm{Id}$ (car $y\oplus f(x)\oplus f(x)=y$). Elle se construit à partir
d'un circuit réversible de $f$ (portes CNOT et Toffoli, section~\ref{sec:controlees}). Le coût d'un
algorithme se mesure en \emph{nombre d'appels} à $U_f$.
\end{definition}

\begin{resultat}{Rebond de phase de l'oracle}
Avec l'ancilla dans l'état $\ket{-}=\frac{1}{\sqrt2}(\ket0-\ket1)$ :
\[
  U_f\ket{x}\ket{-}=(-1)^{f(x)}\ket{x}\ket{-}.
\]
\end{resultat}

En effet $U_f\ket{x}\ket{-}=\frac{1}{\sqrt2}\big(\ket{x}\ket{f(x)}-\ket{x}\ket{1\oplus f(x)}\big)$.
Si $f(x)=0$ on retrouve $\ket{x}\ket{-}$ ; si $f(x)=1$, les deux termes sont échangés et le signe
change. La valeur $f(x)$ est écrite dans la \emph{phase} de $\ket{x}$, sans toucher à l'ancilla : c'est
le même mécanisme que le rebond de phase du CNOT (section~\ref{sec:cnot}), et la brique de tous les
algorithmes de ce chapitre.

\subsubsection{Algorithme de Deutsch}

Soit $f:\{0,1\}\to\{0,1\}$. Elle est \emph{constante} si $f(0)=f(1)$ et \emph{équilibrée} si
$f(0)\neq f(1)$. On veut savoir laquelle des deux est vraie. Classiquement, il faut évaluer $f$ deux
fois. Le circuit de la figure~\ref{fig:deutsch} répond avec \emph{un seul} appel à $U_f$.

\begin{figure}[!ht]
\centering
\begin{tikzpicture}[>=Stealth, font=\small]
  % fils : ancilla (second chiffre) en haut, requête (premier chiffre) en bas
  \draw (0,1.5) -- (8.4,1.5);
  \draw (0,0) -- (7.4,0);
  \node[anchor=east] at (0,1.5) {$\ket{1}$};
  \node[anchor=east] at (0,0) {$\ket{0}$};
  \node[porte] at (1.3,1.5) {$H$};
  \node[porte] at (1.3,0) {$H$};
  \node[draw=insarouge, fill=insarouge!8, thick, minimum width=13mm, minimum height=22mm]
    at (3.9,0.75) {$U_f$};
  \node[porte] at (6.0,0) {$H$};
  \pic at (7.2,0) {mesure};
  \node[anchor=west] at (7.7,0) {$f(0)\oplus f(1)$};
  \node[anchor=west, text=insagris, font=\footnotesize] at (8.5,1.5) {ancilla};
  \foreach \x/\lab in {0.55/0, 2.2/1, 5.15/2, 6.75/3}
    {\draw[dashed, insagris!70] (\x,2.1) -- (\x,-0.7);
     \node[anchor=north, font=\footnotesize] at (\x,-0.7) {$\Pi_\lab$};}
\end{tikzpicture}
\caption{Circuit de Deutsch. La requête $x$ est le premier chiffre du ket (fil du bas), l'ancilla
$y$ le second (fil du haut). La mesure du seul qubit de requête donne $f(0)\oplus f(1)$ : $0$ si $f$ est
constante, $1$ si elle est équilibrée.}
\label{fig:deutsch}
\end{figure}

\begin{align*}
  \ket{\Pi_0}&=\ket{0}\ket{1},\qquad
  \ket{\Pi_1}=(H\otimes H)\ket{\Pi_0}=\ket{+}\ket{-},\\[1ex]
  \ket{\Pi_2}&=U_f\ket{\Pi_1}
    =\frac{(-1)^{f(0)}\ket{0}+(-1)^{f(1)}\ket{1}}{\sqrt2}\ket{-}
    =(-1)^{f(0)}\,\frac{\ket{0}+(-1)^{f(0)\oplus f(1)}\ket{1}}{\sqrt2}\,\ket{-},\\[1ex]
  \ket{\Pi_3}&=(H\otimes I)\ket{\Pi_2}=(-1)^{f(0)}\,\ket{f(0)\oplus f(1)}\ket{-}.
\end{align*}
La deuxième ligne utilise le rebond de phase ; la troisième, $H\ket{+}=\ket0$ et $H\ket{-}=\ket1$.
Les deux amplitudes de la requête interfèrent : elles s'additionnent si $f$ est constante (état
$\ket{+}$, devenu $\ket{0}$) et se retranchent si $f$ est équilibrée (état $\ket{-}$, devenu $\ket{1}$).
Le calcul ne dit pas les valeurs $f(0)$ et $f(1)$ : il extrait seulement leur parité, et c'est ce qui
permet d'économiser un appel.

\begin{remarque}[Des oracles à $n$ qubits : Deutsch--Jozsa, Bernstein--Vazirani, Simon]
Le même schéma ($H^{\otimes n}$, oracle, $H^{\otimes n}$) généralise l'algorithme de \textcite{deutsch1985}. Pour $f$ sur $n$ bits promise
constante ou équilibrée, \textcite{deutschjozsa1992} décident avec \emph{un} appel là où un algorithme classique
déterministe en demande $2^{n-1}+1$. \textcite{bernstein1997} retrouvent en un appel un mot caché
$s$ défini par $f(x)=s\cdot x$ (classiquement : $n$ appels). \textcite{simon1997} trouve un masque caché avec $O(n)$
appels, contre $2^{\Omega(n)}$ classiquement : c'est la première séparation exponentielle, et l'idée
de l'algorithme de Shor. Ces problèmes sont artificiels (ils portent sur une boîte noire), mais
ils prouvent que l'interférence offre un gain réel.
\end{remarque}

% ---------------------------------------------------------------------
\subsection{Recherche dans une liste : l'algorithme de Grover}\label{sec:grover}
% ---------------------------------------------------------------------

Un oracle $f$ vaut $1$ en un unique élément $w$ parmi $N=2^n$, et $0$ ailleurs ; on cherche $w$.
Classiquement, il faut en moyenne $N/2$ évaluations de $f$, et $N$ dans le pire cas. \textcite{grover1996}
le trouve avec environ $\frac\pi4\sqrt N$ appels.

\begin{resultat}{Algorithme de Grover}
On part de $\ket{s}=H^{\otimes n}\ket{0\ldots0}=\frac{1}{\sqrt N}\sum_x\ket{x}$, puis on répète $k$
fois l'itération $G=D\,O_w$ avec
\[
  O_w=I-2\ket{w}\bra{w}\ \ (\text{oracle, par rebond de phase}),\qquad
  D=2\ket{s}\bra{s}-I=H^{\otimes n}\,(2\ket{0}\bra{0}-I)\,H^{\otimes n}.
\]
La mesure donne $w$ avec la probabilité $\sin^2\!\big((2k+1)\theta\big)$, où $\sin\theta=\frac{1}{\sqrt N}$.
Elle est proche de $1$ pour $k\simeq\frac\pi4\sqrt N$. Aucun algorithme quantique ne fait mieux que
$\Omega(\sqrt N)$ appels : le gain est quadratique, et optimal.
\end{resultat}

\paragraph{Pourquoi cela marche.} Posons $\ket{s'}=\frac{1}{\sqrt{N-1}}\sum_{x\neq w}\ket{x}$. Alors
$\ket{s}=\sin\theta\,\ket{w}+\cos\theta\,\ket{s'}$, et toute l'évolution reste dans le plan
engendré par $\ket{w}$ et $\ket{s'}$. L'oracle change le signe de la composante sur $\ket{w}$ : c'est
la symétrie par rapport à l'axe $\ket{s'}$. La diffusion $D=2\ket{s}\bra{s}-I$ est la symétrie par
rapport à $\ket{s}$. Deux symétries d'axes séparés de l'angle $\theta$ composent une rotation de
$2\theta$ (figure~\ref{fig:grover-geo}). Après $k$ itérations,
\[
  G^k\ket{s}=\sin\!\big((2k+1)\theta\big)\ket{w}+\cos\!\big((2k+1)\theta\big)\ket{s'}.
\]
On s'arrête quand $(2k+1)\theta\simeq\frac\pi2$. Contrairement à un algorithme classique, itérer plus
longtemps n'améliore rien : on continue de tourner et on dépasse $\ket{w}$.

\begin{figure}[!ht]
\centering
\begin{tikzpicture}[>=Stealth, font=\small]
  \draw[->, insagris] (-0.4,0) -- (4.5,0) node[right] {$\ket{s'}$};
  \draw[->, insagris] (0,-0.5) -- (0,4.5);
  \draw[insagris!50, dashed] (3.6,0) arc (0:90:3.6);
  \draw[dashed, insagris!60] (210:1.8) -- (30:3.6);
  \draw[->, very thick, insarouge] (0,0) -- (30:3.6) node[right] {$\ket{s}$};
  \draw[->, very thick, insagris] (0,0) -- (-30:3.6) node[right] {$O_w\ket{s}$};
  \draw[->, very thick, insarouge] (0,0) -- (90:3.6) node[above right] {$G\ket{s}=\ket{w}$};
  \draw[insagris] (1.1,0) arc (0:30:1.1);   \node[text=insagris] at (15:1.45) {$\theta$};
  \draw[insagris] (1.1,0) arc (0:-30:1.1);  \node[text=insagris] at (-15:1.45) {$\theta$};
  \draw[->, insarouge, thick] (30:2.2) arc (30:90:2.2);  \node[text=insarouge] at (60:2.65) {$2\theta$};
  \node[text=insagris, font=\footnotesize, anchor=north west, align=left] at (-1.7,-1.2)
    {symétrie $O_w$ : $\ket{s}\mapsto O_w\ket{s}$\\ symétrie $D$ d'axe $\ket{s}$ (pointillés)};
\end{tikzpicture}
\caption{Une itération de Grover pour $N=4$ ($\theta=\frac\pi6$). L'oracle renvoie $\ket{s}$ en
$O_w\ket{s}$ (angle $-\theta$), la diffusion le renvoie en $G\ket{s}$ (angle $3\theta=\frac\pi2$) :
une seule itération amène exactement sur $\ket{w}$.}
\label{fig:grover-geo}
\end{figure}

\subsubsection{Exemple complet : $N=4$ (deux qubits)}

Prenons $w=11$. L'oracle est la porte $Z$ contrôlée $\mathrm{CZ}=\operatorname{diag}(1,1,1,-1)$
(section~\ref{sec:controlees}). La diffusion s'écrit avec les portes du cours, car
$2\ket{00}\bra{00}-I=\operatorname{diag}(1,-1,-1,-1)=\mathrm{CZ}\,(Z\otimes Z)$ :
\[
  D=(H\otimes H)\;\mathrm{CZ}\,(Z\otimes Z)\;(H\otimes H).
\]
Le circuit de la figure~\ref{fig:grover-circ} s'écrit donc $H\otimes H$, $\mathrm{CZ}$,
$H\otimes H$, $Z\otimes Z$, $\mathrm{CZ}$, $H\otimes H$. Calcul pas à pas :
\begin{align*}
  \ket{s}&=\tfrac12\big(\ket{00}+\ket{01}+\ket{10}+\ket{11}\big),\\
  O_w\ket{s}&=\tfrac12\big(\ket{00}+\ket{01}+\ket{10}-\ket{11}\big),\\
  \braket{s}{O_w s}&=\tfrac14(1+1+1-1)=\tfrac12,\\
  D\,O_w\ket{s}&=2\braket{s}{O_w s}\ket{s}-O_w\ket{s}=\ket{s}-O_w\ket{s}
    =\tfrac12\big(0,0,0,2\big)^{\mathsf T}=\ket{11}.
\end{align*}
La mesure donne $11$ \emph{avec certitude} après un seul appel à l'oracle, là où un algorithme
classique en demande en moyenne $2{,}25$. C'est cohérent avec la formule : $\sin\theta=\frac12$,
$\theta=\frac\pi6$ et $\sin(3\theta)=1$. Pour un autre élément marqué, on change l'oracle (un
CZ entouré de portes $X$) et le résultat est le même.

\begin{figure}[!ht]
\centering
\adjustbox{max width=\linewidth}{%
\begin{tikzpicture}[>=Stealth, font=\small]
  \draw (0,1.4) -- (10.3,1.4);
  \draw (0,0) -- (10.3,0);
  \node[anchor=east] at (0,1.4) {$\ket{0}$};
  \node[anchor=east] at (0,0) {$\ket{0}$};
  \foreach \y in {0,1.4} {
    \node[porte] at (1.1,\y) {$H$};
    \node[porte] at (4.3,\y) {$H$};
    \node[porte] at (5.7,\y) {$Z$};
    \node[porte] at (8.5,\y) {$H$};
    \pic at (9.7,\y) {mesure};
  }
  \foreach \x in {2.7,7.1} {
    \draw[insarouge, thick] (\x,1.4) -- (\x,0);
    \pic at (\x,1.4) {ctl};
    \pic at (\x,0) {ctl};
  }
  \foreach \x/\lab in {1.9/{\ket{s}}, 3.5/{O_w\ket{s}}, 9.1/{\ket{11}}} {
    \draw[dashed, insagris!70] (\x,2.0) -- (\x,-0.6);
    \node[anchor=south, font=\footnotesize] at (\x,2.0) {$\lab$};
  }
  \draw[decorate, decoration={brace, amplitude=4pt, mirror}, insagris] (2.3,-0.8) -- (3.1,-0.8)
    node[midway, below=5pt, font=\footnotesize] {oracle $O_w$};
  \draw[decorate, decoration={brace, amplitude=4pt, mirror}, insagris] (3.9,-0.8) -- (9.0,-0.8)
    node[midway, below=5pt, font=\footnotesize] {diffusion $D=(H\otimes H)\,\mathrm{CZ}\,(Z\otimes Z)\,(H\otimes H)$};
\end{tikzpicture}}
\caption{Algorithme de Grover pour $N=4$ et $w=11$ : une itération suffit. Les deux points
reliés par un trait vertical forment la porte CZ.}
\label{fig:grover-circ}
\end{figure}

\paragraph{Ordres de grandeur.} Pour $N=2^{20}\approx10^6$, Grover demande
$\frac\pi4\cdot2^{10}\approx804$ itérations, contre $5\times10^5$ évaluations en moyenne. Pour une
recherche de clé AES-128 ($N=2^{128}$), il faut $\frac\pi4\,2^{64}\approx1{,}4\times10^{19}$
itérations \emph{séquentielles}, dont chacune contient un circuit AES complet ; et partager le travail
entre $m$ machines ne fait gagner qu'un facteur $\sqrt m$. L'effet est tout de même de diviser par deux le
nombre de bits de sécurité d'un chiffrement symétrique : on se protège en doublant la taille des
clés (AES-256).

% ---------------------------------------------------------------------
\subsection{Transformée de Fourier quantique et algorithme de Shor}\label{sec:shor}
% ---------------------------------------------------------------------

Le chiffrement RSA repose sur la difficulté de factoriser un entier $N$ (section~\ref{sec:crypto}).
L'algorithme de Shor, conçu en 1994 et publié en 1997 \parencite{shor1997}, factorise $N$ en un temps polynomial en $\log N$.
L'algorithme a deux étages : une réduction \emph{classique} de la factorisation à la recherche d'une
période, puis une recherche \emph{quantique} de cette période par transformée de Fourier.

\subsubsection{Réduction : de la factorisation à la période}

Soit $N$ impair et non puissance d'un nombre premier. On tire $a$ avec $1<a<N$. Si
$\operatorname{pgcd}(a,N)>1$, on tient déjà un facteur. Sinon, la suite $x\mapsto a^x\bmod N$ est
périodique : sa période $r$ est le plus petit entier $r\geq1$ tel que $a^r\equiv1\pmod N$ (l'\emph{ordre}
de $a$ modulo $N$).

\begin{resultat}{Factorisation à partir de l'ordre}
Si $r$ est pair et si $a^{r/2}\not\equiv-1\pmod N$, alors $\operatorname{pgcd}(a^{r/2}-1,N)$ et
$\operatorname{pgcd}(a^{r/2}+1,N)$ sont deux facteurs non triviaux de $N$. Pour $a$ tiré au hasard,
ces deux conditions sont réunies avec une probabilité au moins égale à $\frac12$ \parencite[chap.~5]{nielsen2010}.
\end{resultat}

En effet, posons $u=a^{r/2}$. Alors $(u-1)(u+1)=a^r-1\equiv0\pmod N$ : $N$ divise le produit. Il ne
divise pas $u-1$, sinon l'ordre de $a$ serait $r/2$, ni $u+1$ par hypothèse. Si $\operatorname{pgcd}(u-1,N)$
valait $1$, $N$ diviserait $u+1$ (lemme de Gauss) ; s'il valait $N$, $N$ diviserait $u-1$. Les deux sont exclus, donc
le pgcd est un diviseur strict de $N$ ; de même pour $u+1$. Tout se ramène donc à \emph{trouver $r$}. Classiquement, c'est
aussi difficile que factoriser. C'est la seule étape où l'ordinateur quantique intervient.

\subsubsection{La transformée de Fourier quantique}

\begin{definition}{Transformée de Fourier quantique (QFT)}
Sur l'espace de dimension $M$ de base $\ket{0},\ldots,\ket{M-1}$, avec $\omega=\ee^{2\ii\pi/M}$ :
\[
  \mathrm{QFT}_M\ket{x}=\frac{1}{\sqrt M}\sum_{y=0}^{M-1}\omega^{xy}\ket{y}.
\]
C'est une porte unitaire : ses colonnes sont orthonormées (somme géométrique de raison $\omega^{x-x'}$).
Pour $M=2^m$, elle s'implémente avec $O(m^2)$ portes élémentaires, alors que la transformée de Fourier
rapide classique demande $O(m\,2^m)$ opérations sur les $2^m$ amplitudes.
\end{definition}

Pour $M=2$ on a $\omega=-1$ et $\mathrm{QFT}_2=H$. Pour $M=4$, $\omega=\ii$ ; en écrivant $x=2x_1+x_2$ et
$y=2y_1+y_2$, on a $\ii^{xy}=(-1)^{x_1y_2+x_2y_1}\,\ii^{x_2y_2}$ (car $\ii^4=1$), d'où
\[
  \mathrm{QFT}_4\ket{x_1x_2}=\underbrace{\frac{\ket0+(-1)^{x_2}\ket1}{\sqrt2}}_{=\,H\ket{x_2}\ \text{(chiffre }y_1)}
  \otimes\ \frac{\ket0+(-1)^{x_1}\ii^{x_2}\ket1}{\sqrt2}\ \ \text{(chiffre }y_2).
\]
Le circuit $\mathrm{QFT}_4=\mathrm{SWAP}\,(I\otimes H)\,\mathrm{CS}\,(H\otimes I)$ de la
figure~\ref{fig:qft4} produit exactement cet état : $H$ sur le premier chiffre donne
$\frac{1}{\sqrt2}(\ket0+(-1)^{x_1}\ket1)$, la porte $\mathrm{CS}=\operatorname{diag}(1,1,1,\ii)$ ajoute le
facteur $\ii^{x_2}$ devant $\ket1$, $I\otimes H$ transforme $\ket{x_2}$ en $H\ket{x_2}$, et le SWAP remet les
chiffres dans l'ordre.

\begin{figure}[!ht]
\centering
\begin{tikzpicture}[>=Stealth, font=\small]
  \draw (0,1.2) -- (7.0,1.2);
  \draw (0,0) -- (7.0,0);
  \node[anchor=east] at (0,1.2) {$\ket{x_2}$};
  \node[anchor=east] at (0,0) {$\ket{x_1}$};
  \node[porte] at (1.2,0) {$H$};
  \draw[insarouge, thick] (2.7,1.2) -- (2.7,0);
  \pic at (2.7,1.2) {ctl};
  \node[porte] at (2.7,0) {$S$};
  \node[porte] at (4.2,1.2) {$H$};
  \draw[insarouge, thick] (5.7,1.2) -- (5.7,0);
  \pic at (5.7,1.2) {croix};
  \pic at (5.7,0) {croix};
  \node[anchor=west] at (7.0,1.2) {$\ket{y_2}$};
  \node[anchor=west] at (7.0,0) {$\ket{y_1}$};
\end{tikzpicture}
\caption{$\mathrm{QFT}_4$ : $H$, $S$ contrôlée par l'autre qubit, $H$, puis échange des deux fils.
Le fil du haut est le dernier chiffre.}
\label{fig:qft4}
\end{figure}

\begin{resultat}{La QFT transforme une période en pics}
Si $r$ divise $M$ et si $\ket{\psi}=\sqrt{r/M}\sum_{k=0}^{M/r-1}\ket{x_0+kr}$ (état périodique, de période
$r$), alors
\[
  \mathrm{QFT}_M\ket{\psi}=\frac{1}{\sqrt r}\sum_{j=0}^{r-1}\omega^{x_0\,jM/r}\,\ket{jM/r}.
\]
La mesure donne $y=jM/r$ avec la probabilité $\frac1r$, pour $j=0,\ldots,r-1$, quel que soit le décalage $x_0$.
\end{resultat}

Preuve : l'amplitude de $\ket{y}$ est $\frac{1}{\sqrt M}\sqrt{\frac rM}\,\omega^{x_0y}\sum_{k}\omega^{kry}$.
La somme géométrique de raison $\omega^{ry}=\ee^{2\ii\pi ry/M}$ vaut $M/r$ si $ry/M$ est entier, et $0$ sinon
(la raison est alors différente de $1$ et sa puissance $M/r$ vaut $1$). L'amplitude vaut donc
$\sqrt{r}/M\cdot M/r=1/\sqrt r$ aux multiples de $M/r$, et $0$ ailleurs. Le décalage $x_0$ n'apparaît que dans une
\emph{phase}, qui disparaît à la mesure : c'est ce qui permet de lire la période sans connaître le point de départ.

\subsubsection{L'algorithme}

\begin{resultat}{Algorithme de Shor : recherche de l'ordre de $a$ modulo $N$}
On choisit $M=2^t\geq N^2$ (soit $t\approx2n$ avec $n=\lceil\log_2N\rceil$), un registre de contrôle de $t$ qubits et un
registre de travail de $n$ qubits.
\begin{enumerate}[nosep]
  \item Préparer $\ket{0}^{\otimes t}\ket{1}$ et appliquer $H^{\otimes t}$ au registre de contrôle :
        $\ket{\varphi_1}=\frac{1}{\sqrt M}\sum_{x=0}^{M-1}\ket{x}\ket{1}$.
  \item Appliquer l'exponentiation modulaire $U_a\ket{x}\ket{z}=\ket{x}\ket{a^xz\bmod N}$ :
        $\ket{\varphi_2}=\frac{1}{\sqrt M}\sum_{x}\ket{x}\ket{a^x\bmod N}$.
  \item Appliquer $\mathrm{QFT}_M$ au registre de contrôle et le mesurer : on obtient $y$.
  \item (Classique) Développer $y/M$ en fraction continue : on y lit une fraction $j/r'$ proche de
        $y/M$ avec $r'<N$. Si $\operatorname{pgcd}(j,r)=1$, alors $r'=r$ et l'on vérifie $a^{r'}\equiv1\pmod N$.
        Sinon, $r'$ est un diviseur strict de $r$ et $a^{r'}\not\equiv1$ : on recommence.
\end{enumerate}
\end{resultat}

L'oracle $U_a$ est un produit de $t$ multiplications modulaires \emph{contrôlées} par les qubits $x_j$ de $x$ :
comme $a^x=\prod_j(a^{2^j})^{x_j}$, la $j$-ème porte multiplie le registre de travail par $a^{2^j}\bmod N$, un
nombre calculé à l'avance par un ordinateur classique. Le circuit est celui de la
figure~\ref{fig:shor-circ}.

\begin{figure}[!ht]
\centering
\begin{tikzpicture}[>=Stealth, font=\small]
  % registre de contrôle (haut) et registre de travail (bas)
  \draw (0,1.7) -- (8.3,1.7);
  \draw (0,0) -- (4.6,0);
  \draw[dashed] (4.6,0) -- (8.3,0);
  \node[anchor=east] at (0,1.7) {$\ket{0}^{\otimes t}$};
  \node[anchor=east] at (0,0) {$\ket{1}$};
  \draw (0.5,1.6) -- (0.7,1.8); \node[font=\scriptsize, anchor=south] at (0.6,1.78) {$t$};
  \draw (0.5,-0.1) -- (0.7,0.1); \node[font=\scriptsize, anchor=south] at (0.6,0.08) {$n$};
  \node[porte] at (1.5,1.7) {$H^{\otimes t}$};
  \pic at (3.3,1.7) {ctl};
  \draw[insarouge, thick] (3.3,1.7) -- (3.3,0.4);
  \node[draw=insarouge, fill=insarouge!8, thick, minimum width=17mm, minimum height=8mm] at (3.3,0) {$U_a$};
  \node[porte] at (5.4,1.7) {$\mathrm{QFT}_M$};
  \pic at (7.0,1.7) {mesure};
  \node[anchor=west] at (7.5,1.7) {$y$};
  \node[anchor=north west, text=insagris, font=\footnotesize] at (4.7,-0.3) {registre de travail : ignoré};
  \foreach \x/\lab in {2.3/1, 4.3/2, 6.3/3}
    {\draw[dashed, insagris!70] (\x,2.35) -- (\x,-0.3);
     \node[anchor=south, font=\footnotesize] at (\x,2.35) {$\ket{\varphi_\lab}$};}
\end{tikzpicture}
\caption{Circuit de Shor. Les petits traits obliques indiquent un registre de plusieurs qubits ($t$ pour le
contrôle, $n$ pour le travail). $U_a\ket{x}\ket{z}=\ket{x}\ket{a^xz\bmod N}$ est commandé par le registre de
contrôle.}
\label{fig:shor-circ}
\end{figure}

\paragraph{Pourquoi cela marche.} Après l'étape~2, mesurons (en pensée) le registre de travail : il
donne une valeur $f$ de la suite $a^x\bmod N$, et le registre de contrôle se projette sur les $x$ tels que
$a^x\equiv f$, c'est-à-dire $x\in\{x_0,x_0+r,x_0+2r,\ldots\}$ : un état périodique de période $r$. Le résultat sur la QFT
ci-dessus s'applique : la QFT concentre toute la probabilité sur les $y$ voisins de $jM/r$. La mesure du
registre de travail n'est pas nécessaire (les probabilités de $y$ sont les mêmes si on l'ignore), et le
décalage $x_0$ dépend de $f$ mais ne change pas la loi de $y$. Ainsi $y/M\approx j/r$ pour un entier $j$ \emph{aléatoire} :
on ne choisit pas $j$, mais il est premier avec $r$ dans une fraction $\Omega(1/\log\log r)$ des cas, et quelques
répétitions suffisent. Le choix $M\geq N^2$ garantit qu'une seule fraction de dénominateur inférieur à $N$ est
à distance de $\frac{1}{2M}$ de $y/M$ : le développement en fraction continue l'identifie.

\subsubsection{Exemple complet : $N=15$, $a=7$}

\paragraph{Période.} La suite $7^x\bmod15$ vaut $1,7,4,13,1,7,4,13,\ldots$ : l'ordre est $r=4$. Les
multiplicateurs de $U_a$ sont $a^{2^j}\bmod 15=7,\ 4,\ 1,\ 1$ pour $j=0,1,2,3$ (les deux derniers sont triviaux).
On prend ici $t=4$, soit $M=16$. Ce choix est trop court en général ($M\geq N^2=225$ serait exigé) mais suffit
car $r=4$ divise $M$, ce qui rend le calcul exact.

\paragraph{État avant la QFT.} $\ket{\varphi_2}=\frac14\sum_{x=0}^{15}\ket{x}\ket{7^x\bmod15}$. Si le
registre de travail est mesuré et donne $1$, le registre de contrôle est dans $\frac12(\ket0+\ket4+\ket8+\ket{12})$ ;
s'il donne $7$, dans $\frac12(\ket1+\ket5+\ket9+\ket{13})$, etc.

\paragraph{QFT.} On a $\omega=\ee^{2\ii\pi/16}$ et $\omega^4=\ii$. Pour le décalage $x_0=0$ :
\[
  \mathrm{QFT}_{16}\,\tfrac12\big(\ket0+\ket4+\ket8+\ket{12}\big)
  =\frac18\sum_{y=0}^{15}\Big(\sum_{k=0}^{3}\ii^{\,ky}\Big)\ket{y}
  =\frac12\big(\ket0+\ket4+\ket8+\ket{12}\big),
\]
car $\sum_{k=0}^{3}\ii^{\,ky}$ vaut $4$ si $y$ est multiple de $4$ et $0$ sinon. Pour les autres valeurs de
$f$, seules les phases changent. Dans tous les cas la mesure de $y$ donne $0,4,8$ ou $12$, chacun avec la probabilité $\frac14$
(figure~\ref{fig:shor-n15}).

\paragraph{Lecture classique.} On calcule $y/16$ :
\begin{center}
\renewcommand{\arraystretch}{1.2}
\begin{tabular}{@{}cccp{7.3cm}@{}}
\toprule
$y$ & $y/16$ & fraction irréductible & conclusion\\
\midrule
$0$ & $0$ & $0/1$ & inutile ($j=0$) : on recommence\\
$4$ & $\frac14$ & $1/4$ & $r=4$\ ; $7^4=2401\equiv1$ : correct\\
$8$ & $\frac12$ & $1/2$ & dénominateur $2$ : $7^2=49\equiv4\not\equiv1$, rejeté, on recommence\\
$12$ & $\frac34$ & $3/4$ & $r=4$ : correct\\
\bottomrule
\end{tabular}
\end{center}
Un essai réussit donc avec la probabilité $\frac12$ ($y=4$ ou $12$). Avec $r=4$ pair et $7^{2}=49\equiv4\not\equiv-1\pmod{15}$ :
\[
  \operatorname{pgcd}(7^2-1,15)=\operatorname{pgcd}(48,15)=3,\qquad
  \operatorname{pgcd}(7^2+1,15)=\operatorname{pgcd}(50,15)=5,\qquad 15=3\times5.
\]

\begin{figure}[!ht]
\centering
\adjustbox{max width=\linewidth}{%
\begin{tikzpicture}[>=Stealth, font=\small]
  % (a) suite 7^x mod 15, x = 0..15
  \draw[->, insagris] (-0.2,0) -- (6.7,0) node[right] {$x$};
  \draw[->, insagris] (0,-0.1) -- (0,2.6);
  \foreach \x/\v in {0/1,1/7,2/4,3/13,4/1,5/7,6/4,7/13,8/1,9/7,10/4,11/13,12/1,13/7,14/4,15/13}
    \fill[insarouge] ({\x*0.42},{\v*0.17}) circle (1.7pt);
  \foreach \v in {1,4,7,13}
    {\draw[insagris] (-0.05,{\v*0.17}) -- (0.05,{\v*0.17});
     \node[anchor=east, font=\scriptsize] at (-0.08,{\v*0.17}) {$\v$};}
  \foreach \x in {0,4,8,12}
    {\draw[insagris] ({\x*0.42},-0.05) -- ({\x*0.42},0.05);
     \node[anchor=north, font=\scriptsize] at ({\x*0.42},-0.07) {$\x$};}
  \draw[<->, insarouge] (0,2.45) -- (1.68,2.45) node[midway, above, font=\footnotesize] {période $r=4$};
  \node[font=\footnotesize, anchor=south] at (3.2,3.15) {(a) $f(x)=7^x\bmod15$};
  % flèche QFT
  \draw[->, very thick, insarouge] (7.0,0.8) -- (8.1,0.8) node[midway, above, font=\footnotesize] {$\mathrm{QFT}_{16}$};
  % (b) loi de y : P(y)=1/4 pour y = 0, 4, 8, 12
  \draw[->, insagris] (8.5,0) -- (14.9,0) node[right] {$y$};
  \draw[->, insagris] (8.6,-0.1) -- (8.6,2.6);
  \node[font=\footnotesize, anchor=south] at (11.7,3.15) {(b) loi de la mesure $y$ : $P(y)$};
  \foreach \y in {0,4,8,12}
    {\fill[insarouge] ({8.6+\y*0.4-0.1},0) rectangle ({8.6+\y*0.4+0.1},1.5);
     \node[anchor=north, font=\scriptsize] at ({8.6+\y*0.4},-0.07) {$\y$};}
  \draw[dashed, insagris] (8.55,1.5) -- (14.2,1.5);
  \node[anchor=east, font=\scriptsize] at (8.5,1.5) {$\frac14$};
  \node[text=insagris, font=\footnotesize, align=center] at (11.7,2.35) {$y/16=\frac{j}{4}$, $j=0,1,2,3$};
\end{tikzpicture}}
\caption{Shor pour $N=15$, $a=7$, $t=4$. À gauche, la suite $7^x\bmod15$ a la période $4$. À droite,
la mesure du registre de contrôle après la QFT ne donne que $y\in\{0,4,8,12\}$, avec la probabilité
$\frac14$ chacun : la période $r=4$ se lit dans l'écart $M/r=4$ entre les pics.}
\label{fig:shor-n15}
\end{figure}

\subsubsection{Ce que coûte l'algorithme}

Notons $n=\lceil\log_2N\rceil$ le nombre de bits de $N$ ($n=2048$ pour un module RSA-2048). Le coût s'exprime en
fonction de $n$, pas de $N$.
\begin{align*}
  \text{meilleur algorithme classique connu (crible algébrique, GNFS) :}&\quad
    \exp\!\Big(O\big(n^{1/3}(\log n)^{2/3}\big)\Big),\\
  \text{algorithme de Shor :}&\quad O(n^3)\ \text{portes élémentaires.}
\end{align*}
Le premier temps est sous-exponentiel en $n$ \parencite{lenstra1993}, mais supérieur à tout polynôme. Le second est un
polynôme en $n$ avec la multiplication scolaire ($O(n^2\log n\log\log n)$ avec une multiplication rapide).

\begin{remarque}[Où est le « log » ?]
On résume parfois « de $O(N^p)$ à $O(\log N)$ ». La comparaison correcte est la suivante. La division
d'essai coûte $O(\sqrt N)=O(2^{n/2})$ : exponentielle en $n$. Le crible algébrique est sous-exponentiel en $n$. Shor est
polynomial en $\log N=n\ln2$, soit $O((\log N)^3)$ et non $O(\log N)$ : rien que pour lire $N$ et multiplier
modulo $N$, il faut au moins $n$ opérations. Le « $\log$ » désigne donc la \emph{taille} du nombre.
\end{remarque}

Doubler la taille de la clé multiplie le coût de Shor par $8$, alors qu'il augmente de façon
spectaculaire celui de l'attaque classique. Le tableau compare le coût $n^3$ de Shor à la sécurité
classique que le NIST attribue aux modules RSA \parencite{nist80057}.

\begin{center}
\renewcommand{\arraystretch}{1.2}
\begin{tabular}{@{}rccc@{}}
\toprule
Module RSA $n$ (bits) & sécurité classique (NIST) & $n^3$ (portes, ordre de grandeur) & $\log_2(n^3)$\\
\midrule
$1024$ & $2^{80}$ & $1{,}1\times10^{9}$ & $30{,}0$\\
$2048$ & $2^{112}$ & $8{,}6\times10^{9}$ & $33{,}0$\\
$3072$ & $2^{128}$ & $2{,}9\times10^{10}$ & $34{,}8$\\
$7680$ & $2^{192}$ & $4{,}5\times10^{11}$ & $38{,}7$\\
$15360$ & $2^{256}$ & $3{,}6\times10^{12}$ & $41{,}7$\\
\bottomrule
\end{tabular}
\end{center}

Passer de $2048$ à $3072$ bits fait gagner un facteur $2^{16}=65\,536$ sur l'attaque classique, mais le coût de
Shor n'est multiplié que par $(3/2)^3\approx3{,}4$. Face à un ordinateur quantique, \emph{grossir
les clés RSA ne protège pas} : il faut changer d'algorithme (section~\ref{sec:crypto}). Cela dit, le nombre de portes
n'est pas le vrai coût : il faut des qubits \emph{corrigés d'erreurs}, et c'est là que se trouve le verrou technique
(section~\ref{sec:machine}). Les estimations les plus citées sont les suivantes.
\begin{itemize}[nosep]
  \item \textbf{RSA-2048} : environ $3n+0{,}002\,n\log_2n\approx6\,189$ qubits logiques et
        $0{,}3\,n^3+0{,}0005\,n^3\log_2n\approx2{,}6\times10^9$ portes de Toffoli ; avec des qubits physiques
        d'erreur $10^{-3}$ et un cycle de code de surface de $1\,\mu$s, $20$ millions de qubits physiques pendant
        $8$ heures \parencite{gidney2021}. Une optimisation de 2025 (arithmétique approchée, codes « yoked »,
        « magic state cultivation ») descend à \emph{moins d'un million} de qubits physiques, pour moins d'une
        semaine \parencite{gidney2025}.
  \item \textbf{Courbes elliptiques (ECC-256)}, utilisées dans TLS : au plus $2\,330$ qubits logiques
        \parencite{roetteler2017}. Avec des qubits de chat (section~\ref{sec:machine}), $126\,133$ qubits suffiraient
        pour $9$ heures de calcul \parencite{gouzien2023}. Une \emph{prépublication} de 2026, fondée sur des
        hypothèses optimistes pour les atomes neutres, évoque $10\,000$ à $26\,000$ atomes et quelques
        jours \parencite{cain2026}.
\end{itemize}
Ce sont des estimations, pas des mesures : les démonstrations expérimentales de $15=3\times5$ ou
$21=3\times7$ reposent sur des circuits simplifiés grâce à la connaissance préalable de la réponse
\parencite{smolin2013}. À notre connaissance, aucune n'a factorisé par l'algorithme de Shor un nombre hors de portée d'un
calcul à la main.

% ---------------------------------------------------------------------
\subsection{Cryptographie : RSA, TLS et « stocker maintenant, déchiffrer plus tard »}\label{sec:crypto}
% ---------------------------------------------------------------------

La sécurité des échanges sur Internet repose sur deux familles d'algorithmes. La cryptographie
\emph{symétrique} (AES, ChaCha20) chiffre les données avec une clé secrète partagée. La cryptographie
\emph{asymétrique} (RSA, Diffie--Hellman, courbes elliptiques) sert à \emph{établir} cette clé entre deux
personnes qui ne se sont jamais rencontrées, et à \emph{authentifier} l'interlocuteur par une signature. Les
algorithmes asymétriques s'appuient sur deux problèmes réputés difficiles : la factorisation d'un entier
(RSA) et le logarithme discret (Diffie--Hellman, courbes elliptiques). Ce sont deux problèmes de
\emph{période} dans un groupe : l'algorithme de Shor les résout tous les deux.

\subsubsection{RSA et ce que Shor en fait}

\begin{definition}{Chiffrement RSA}
Alice choisit deux nombres premiers $p\neq q$, pose $n=pq$ et $\varphi(n)=(p-1)(q-1)$, choisit $e$ premier avec
$\varphi(n)$ et calcule $d$ tel que $ed\equiv1\pmod{\varphi(n)}$. La \emph{clé publique} est $(n,e)$, la \emph{clé
privée} $d$. Pour un message $m<n$ :
\[
  c=m^e\bmod n\quad(\text{chiffrement}),\qquad m=c^d\bmod n\quad(\text{déchiffrement}).
\]
Le déchiffrement fonctionne parce que $m^{ed}\equiv m\pmod n$ (théorème d'Euler) \parencite{rivest1978}.
\end{definition}

\paragraph{Exemple jouet.} $p=3$, $q=11$, $n=33$, $\varphi(n)=20$, $e=3$ (premier avec $20$) et $d=7$ (car
$3\times7=21\equiv1\pmod{20}$). Pour $m=4$ : $c=4^3\bmod33=64-33=31$. Le déchiffrement donne
$31^7\bmod33$ ; comme $31\equiv-2$, on a $(-2)^7=-128\equiv-128+132=4=m$ : on retrouve le message.

\paragraph{Attaque.} L'attaquant connaît $(n,e)$. Il lui manque $\varphi(n)$ : le calculer équivaut à
factoriser $n$. Avec $n=33=3\times11$, il obtient $\varphi=20$ puis $d=e^{-1}\bmod20=7$ par l'algorithme d'Euclide
étendu (polynomial). Tout le secret de RSA est donc la difficulté de factoriser. Pour un module de $2048$ bits
(617 chiffres décimaux), les meilleures factorisations classiques publiques (RSA-240, $795$ bits, en 2019
\parencite{boudot2020}) sont loin du compte : on a vu à la section~\ref{sec:shor} que
l'algorithme de Shor l'attaque en $O(n^3)$ portes. Le même algorithme, appliqué au groupe d'une courbe elliptique,
résout le logarithme discret et casse Diffie--Hellman, ECDH et ECDSA.

\subsubsection{Ce que cassent Shor et Grover}

\begin{center}
\renewcommand{\arraystretch}{1.25}
\begin{tabular}{@{}>{\raggedright\arraybackslash}p{4.1cm}>{\raggedright\arraybackslash}p{2.9cm}>{\raggedright\arraybackslash}p{2.5cm}>{\raggedright\arraybackslash}p{5.2cm}@{}}
\toprule
\textbf{Primitive} & \textbf{Rôle} & \textbf{Attaque quantique} & \textbf{Conséquence}\\
\midrule
RSA, Diffie--Hellman, ECDH (X25519) & échange de clé, chiffrement & Shor & \textbf{cassés} (temps polynomial)\\
RSA, ECDSA, EdDSA & signature, certificats & Shor & \textbf{cassés}\\
AES, ChaCha20 & chiffrement symétrique & Grover & clé de $k$ bits : $2^{k/2}$ itérations ; AES-128 affaibli, \textbf{AES-256 sûr}\\
SHA-2, SHA-3 & hachage & Grover (pré-image) & SHA-256 : $2^{128}$, \textbf{sûr}\\
ML-KEM, ML-DSA, SLH-DSA & échange de clé, signature & aucune attaque efficace connue & normes NIST de 2024 (voir plus bas)\\
\bottomrule
\end{tabular}
\end{center}

Les trois premières lignes ne sont pas à égalité. Un algorithme asymétrique est \emph{détruit} par Shor : il faut
le remplacer. Un algorithme symétrique n'est que \emph{ralenti} par Grover (racine carrée, section~\ref{sec:grover}) :
il suffit de doubler la taille de la clé. « Aucune attaque connue » n'est pas une preuve de sécurité, d'où la prudence
des agences décrite plus bas.

\subsubsection{TLS 1.3 : où sont les points faibles ?}

Le protocole TLS 1.3 (RFC~8446 \parencite{rfc8446}) protège HTTPS. Il combine les trois mécanismes : un
\emph{échange de clés} éphémère (ECDHE, le plus souvent X25519 : chaque côté tire un secret $a$ ou $b$ et
publie $g^a$ ou $g^b$, d'où le secret partagé $g^{ab}$), une \emph{authentification} du serveur par signature
(certificat RSA ou ECDSA, vérifié par une chaîne d'autorités de certification, elle aussi faite de signatures),
et un \emph{chiffrement symétrique} des données (AES-GCM ou ChaCha20-Poly1305) avec des clés dérivées de $g^{ab}$
(figure~\ref{fig:tls}).

\begin{figure}[!ht]
\centering
\adjustbox{max width=\linewidth}{%
\begin{tikzpicture}[>=Stealth, font=\small,
    tag/.style={rounded corners=2pt, font=\footnotesize, inner sep=3pt, anchor=west, text width=4.6cm, align=left}]
  \node[draw=insarouge, fill=insarouge!8, thick, minimum width=20mm] (C) at (0,0) {Client};
  \node[draw=insarouge, fill=insarouge!8, thick, minimum width=20mm] (S) at (6.2,0) {Serveur};
  \draw[insagris, thick] (0,-0.4) -- (0,-6.2);
  \draw[insagris, thick] (6.2,-0.4) -- (6.2,-6.2);
  \draw[->, thick] (0,-1.0) -- (6.2,-1.0) node[midway, above, font=\footnotesize] {ClientHello + clé publique $g^a$};
  \draw[<-, thick] (0,-1.9) -- (6.2,-1.9) node[midway, above, font=\footnotesize] {ServerHello + clé publique $g^b$};
  \node[font=\footnotesize, text=insagris, anchor=east] at (-0.2,-2.35) {calcule $g^{ab}$, dérive les clés};
  \node[font=\footnotesize, text=insagris, anchor=west] at (6.4,-2.35) {calcule $g^{ab}$};
  \draw[<-, thick] (0,-3.2) -- (6.2,-3.2) node[midway, above, font=\footnotesize, align=center]
       {\{Certificat, signature (CertificateVerify),\\ Finished\}};
  \draw[->, thick] (0,-4.1) -- (6.2,-4.1) node[midway, above, font=\footnotesize] {\{Finished\}};
  \draw[<->, very thick, insarouge] (0,-5.2) -- (6.2,-5.2) node[midway, above, font=\footnotesize] {données chiffrées (AES-GCM, ChaCha20)};
  \node[tag, fill=insarouge!18] at (7.8,-1.45) {\textbf{Échange de clés} : cassé par Shor (logarithme discret)};
  \node[tag, fill=insarouge!18] at (7.8,-3.2) {\textbf{Signature et certificats} : cassés par Shor (RSA, ECDSA)};
  \node[tag, fill=insagris!20] at (7.8,-5.2) {\textbf{Chiffrement symétrique} : \mbox{Grover} seulement, AES-256 suffit};
\end{tikzpicture}}
\caption{Poignée de main TLS 1.3 simplifiée. Les accolades \{\,\} indiquent des messages déjà chiffrés avec les clés
issues de $g^{ab}$. Deux des trois briques tombent face à Shor ; la troisième résiste.}
\label{fig:tls}
\end{figure}

\begin{remarque}[La confidentialité persistante ne protège pas contre l'ordinateur quantique]
Les secrets $a$ et $b$ sont détruits après la session : si la clé privée du serveur est volée plus tard, les sessions
passées ne peuvent pas être déchiffrées (propriété dite de \emph{forward secrecy}). Mais les valeurs publiques
$g^a$ et $g^b$ ont circulé en clair. Un adversaire qui les a enregistrées peut, avec l'algorithme de Shor (logarithme
discret sur la courbe), retrouver $a$ puis $g^{ab}$, donc toutes les clés de la session. La forward secrecy protège
contre le vol d'une \emph{clé}, pas contre l'effondrement de la difficulté du \emph{problème mathématique}.
\end{remarque}

\subsubsection{« Stocker maintenant, déchiffrer plus tard »}

La menace n'attend pas l'existence de la machine. Un adversaire qui capte le trafic chiffré aujourd'hui (il lui
suffit d'écouter une liaison, ce qui est discret et peu coûteux) peut le \emph{conserver} et attendre qu'un
ordinateur quantique assez puissant existe pour le déchiffrer. C'est l'attaque \emph{Harvest Now, Decrypt Later}
(HNDL, ou « stocker maintenant, déchiffrer plus tard »). Elle vise la \emph{confidentialité de long terme} :
dossiers médicaux et génétiques, secrets industriels, secrets d'État, archives juridiques. Elle ne vise pas un
jeton de session valable dix minutes. On ne sait pas qui enregistre quoi ; il faut supposer que cela se fait.

\begin{resultat}{Inégalité de Mosca}
Soit $X$ la durée pendant laquelle les données doivent rester secrètes, $Y$ le temps nécessaire pour migrer vers une
cryptographie résistante, et $Z$ le délai avant qu'un ordinateur quantique capable de casser la cryptographie actuelle existe
(tous en années \parencite{mosca2018}). Les données sont \emph{déjà exposées} si
\[
  X+Y>Z.
\]
\end{resultat}

Les données chiffrées juste avant la fin de la migration (à la date $Y$) doivent encore être secrètes à la date $X+Y$ ;
si $Z$ tombe avant, l'adversaire les lit. Exemple (figure~\ref{fig:mosca}) : des données médicales à protéger
$X=20$ ans, une migration de $Y=7$ ans, et $Z=15$ ans : $X+Y=27>15$, avec douze années d'exposition. On ne connaît pas
$Z$ (ce chapitre ne cherche pas à le prédire), mais on connaît $X$ et $Y$ : plus ils sont grands, plus il faut
commencer tôt, et le coût d'une migration trop précoce est faible devant celui d'une divulgation irréversible.

\begin{figure}[!ht]
\centering
\begin{tikzpicture}[>=Stealth, font=\small, x=0.3cm]
  \draw[->, insagris] (0,0) -- (31,0) node[right] {années};
  \foreach \t in {0,5,...,30}
    {\draw[insagris] (\t,0.08) -- (\t,-0.08);
     \node[anchor=north, font=\scriptsize] at (\t,-0.1) {$\t$};}
  \fill[insagris!35] (0,0.5) rectangle (7,1.1);
  \fill[insarouge!18] (7,0.5) rectangle (27,1.1);
  \fill[insarouge!60] (15,0.5) rectangle (27,1.1);
  \node[font=\footnotesize] at (3.5,0.8) {$Y=7$};
  \node[font=\footnotesize] at (11,0.8) {$X=20$};
  \draw[insarouge, very thick] (15,0) -- (15,2.0);
  \node[anchor=south, font=\footnotesize, text=insarouge] at (15,2.0) {$Z=15$ : ordinateur quantique opérationnel};
  \node[anchor=west, font=\footnotesize, text=white] at (15.4,0.8) {exposé 12 ans};
  \node[anchor=north west, font=\footnotesize, text=insagris] at (0,-0.65)
    {$Y$ : migration (la cryptographie vulnérable est encore utilisée)};
  \node[anchor=north west, font=\footnotesize, text=insagris] at (0,-1.15)
    {$X$ : durée de confidentialité exigée pour les données chiffrées à la fin de la migration};
\end{tikzpicture}
\caption{Inégalité de Mosca pour $X=20$, $Y=7$, $Z=15$ ans. La zone rouge foncé est la période où des données
interceptées sont lisibles par l'adversaire alors qu'elles devraient rester secrètes.}
\label{fig:mosca}
\end{figure}

\begin{remarque}[Signatures : un risque différent]
L'attaque HNDL ne concerne que la confidentialité. Une signature ne sert qu'au moment où on la vérifie : falsifier en 2035 un
certificat TLS expiré ne change rien à une session de 2026. Les signatures à longue durée de vie (micrologiciels,
racines de confiance, documents juridiques) devront néanmoins migrer, mais plus tard : c'est l'échange de clés qui est
urgent.
\end{remarque}

\subsubsection{Première réponse : la cryptographie post-quantique}

La \emph{cryptographie post-quantique} (PQC) regroupe des algorithmes qui s'exécutent sur des ordinateurs
\emph{classiques} mais reposent sur des problèmes contre lesquels on ne connaît aucune attaque quantique efficace :
réseaux euclidiens (problème \emph{Learning With Errors}), fonctions de hachage, codes correcteurs d'erreurs. Le NIST a
publié le 13 août 2024 trois normes \parencite{fips203,fips204,fips205}, puis sélectionné un quatrième algorithme en 2025 \parencite{nist2025hqc}.

\begin{center}
\renewcommand{\arraystretch}{1.2}
\begin{tabular}{@{}llll@{}}
\toprule
\textbf{Norme} & \textbf{Nom} & \textbf{Usage} & \textbf{Problème sous-jacent}\\
\midrule
FIPS 203 & ML-KEM (ex-Kyber) & échange de clé & réseaux (module-LWE)\\
FIPS 204 & ML-DSA (ex-Dilithium) & signature & réseaux\\
FIPS 205 & SLH-DSA (ex-SPHINCS+) & signature & fonctions de hachage\\
(à venir) & HQC & échange de clé (second choix) & codes correcteurs\\
\bottomrule
\end{tabular}
\end{center}

\paragraph{Hybridation.} Dans TLS 1.3, la solution déployée est \emph{hybride} : le client envoie à la fois une clé
publique X25519 et une clé publique ML-KEM-768, et les deux secrets sont concaténés dans la dérivation des clés. L'adversaire
doit casser \emph{les deux} : X25519 (que Shor casse, mais que l'on connaît bien) et ML-KEM (que Shor ne casse pas, mais
que l'on connaît depuis moins longtemps). Le groupe \texttt{X25519MLKEM768} est normalisé par la RFC~10024 d'août 2026
\parencite{rfc10024}, et il est activé par défaut dans les navigateurs depuis fin 2024 (Chrome 131, Firefox 132). Le
prix est la taille : $32$ octets pour une clé X25519, mais $1\,184$ octets de clé publique et $1\,088$ octets de
chiffré pour ML-KEM-768.

\paragraph{Pourquoi pas ML-KEM seul ?} Les algorithmes post-quantiques sont plus jeunes et moins étudiés. En 2022, le
candidat SIKE, arrivé au quatrième tour du NIST, a été cassé par une \emph{attaque classique} : une dizaine de minutes sur un seul
cœur d'ordinateur suffisent pour son plus petit jeu de paramètres \parencite{castryck2023}. L'ANSSI recommande donc de \emph{combiner} cryptographie classique et
post-quantique pendant la transition \parencite{anssi2024}.

\paragraph{Calendrier.} Le projet de document NIST~IR~8547 (version initiale publique de novembre 2024, qui peut
avoir évolué) prévoit de déprécier RSA-2048 et les algorithmes de sécurité équivalente après 2030, et d'interdire tous les
algorithmes vulnérables après 2035 \parencite{nistir8547}. La migration est longue : il faut inventorier où la cryptographie est
utilisée (logiciels, matériel, certificats, protocoles), puis la remplacer. L'« agilité cryptographique », c'est-à-dire
la capacité de changer d'algorithme sans refondre un système, est la qualité qui coûtera le moins cher.

\subsubsection{Seconde réponse : la distribution quantique de clé}

La \emph{distribution quantique de clé} (QKD) ne cherche pas un problème difficile : elle utilise la physique. Dans le
protocole BB84 \parencite{bennett1984}, Alice envoie à Bob des qubits dont l'état est choisi au hasard parmi quatre :
$\ket0,\ket1$ (base $Z$) ou $\ket+,\ket-$ (base $X$). Bob mesure chaque qubit dans une base $Z$ ou $X$ choisie au hasard.
Ensuite, Alice et Bob annoncent \emph{publiquement leurs bases} (pas les bits) et gardent les positions où les bases
coïncident : si la base est la bonne, Bob lit exactement le bit d'Alice ; sinon son résultat est aléatoire. Exemple, sans espion :

\begin{center}
\renewcommand{\arraystretch}{1.15}
\begin{tabular}{@{}l*{8}{c}@{}}
\toprule
position & 1 & 2 & 3 & 4 & 5 & 6 & 7 & 8\\
\midrule
bit d'Alice & 0 & 1 & 1 & 0 & 1 & 0 & 0 & 1\\
base d'Alice & $Z$ & $X$ & $X$ & $Z$ & $X$ & $Z$ & $X$ & $Z$\\
base de Bob & $Z$ & $Z$ & $X$ & $X$ & $X$ & $Z$ & $Z$ & $Z$\\
\midrule
bases identiques ? & oui & non & oui & non & oui & oui & non & oui\\
clé tamisée & 0 & & 1 & & 1 & 0 & & 1\\
\bottomrule
\end{tabular}
\end{center}

La clé tamisée est $01101$ : en moyenne la moitié des qubits est conservée. Pour détecter un espion Ève, Alice et Bob
comparent publiquement un échantillon de la clé tamisée. Une attaque « intercepter puis renvoyer » est détectée parce que
\emph{mesurer perturbe} (postulat de la mesure) et que l'on ne peut pas copier un état inconnu pour mesurer la copie
(théorème de non-clonage, section~\ref{sec:teleportation}). Ève choisit la mauvaise base une fois sur deux ; dans ce cas
le qubit renvoyé est dans la mauvaise base, et Bob, s'il a la bonne base, se trompe une fois sur deux. Le taux d'erreur de la clé tamisée
(\emph{QBER}) vaut donc
\[
  \mathrm{QBER}=\underbrace{\tfrac12}_{\text{Ève se trompe de base}}\times\underbrace{\tfrac12}_{\text{Bob se trompe}}=\tfrac14=25\,\%.
\]
En pratique, quelques pour cent d'erreurs viennent du bruit ; au-delà d'un seuil de l'ordre de $11\,\%$, on abandonne la clé. La
variante E91 \parencite{ekert1991} remplace l'envoi de qubits par des paires intriquées et teste l'absence d'espion par la violation
de l'inégalité CHSH (section~\ref{sec:chsh}).

\begin{remarque}[Pourquoi les agences recommandent la PQC plutôt que la QKD]
En janvier 2024, l'ANSSI, le BSI allemand et leurs homologues néerlandais et suédois ont publié une position commune
\parencite{anssi2024qkd} : la QKD n'est pas assez mûre pour la sécurité, ne convient qu'à des cas d'usage restreints, et la priorité est donnée à
une technologie « beaucoup plus compatible avec les systèmes cryptographiques existants ». Elle demande des liaisons optiques
dédiées de portée limitée, des nœuds de confiance, une authentification classique des messages publics (qui nécessite
elle-même de la cryptographie), et ses implémentations ont des failles de matériel. La QKD protège donc un lien, la PQC un
protocole entier.
\end{remarque}

% ---------------------------------------------------------------------
\subsection{Construire un ordinateur quantique}\label{sec:machine}
% ---------------------------------------------------------------------

Les sections précédentes supposent des qubits parfaits. Les qubits réels sont bruités, et c'est le bruit, bien plus que le
nombre de portes, qui sépare les machines actuelles des estimations de la section~\ref{sec:shor}. Cette section
explique ce qu'il faut pour fabriquer un qubit, pourquoi certains tournent dans un réfrigérateur à
$10\ \mathrm{mK}$, comment on corrige les erreurs, et ce que proposent les qubits de chat d'Alice \& Bob.

\subsubsection{Que faut-il ? Les critères de DiVincenzo}

\textcite{divincenzo2000} a résumé les conditions nécessaires au calcul quantique :
\begin{enumerate}[nosep]
  \item des qubits bien définis, en nombre extensible ;
  \item une initialisation fiable dans l'état $\ket{0\ldots0}$ ;
  \item une cohérence longue devant la durée d'une porte ;
  \item un jeu universel de portes : portes à un qubit et au moins une porte à deux qubits comme le CNOT (chapitre~\ref{chap:portes}) ;
  \item une mesure individuelle des qubits.
\end{enumerate}
S'y ajoutent, pour les réseaux, deux critères de communication. Le critère~3 (cohérence) est le plus difficile à tenir quand le nombre de qubits augmente.

\subsubsection{La décohérence}

Un qubit n'est jamais seul. Supposons un qubit $S$ dans l'état $\frac{1}{\sqrt2}(\ket0+\ket1)$ qui interagit
avec son environnement $E$ de façon que celui-ci « enregistre » l'état de $S$ : $\ket0\ket{E_0}$ et $\ket1\ket{E_1}$.
L'état global est $\frac{1}{\sqrt2}(\ket0\ket{E_0}+\ket1\ket{E_1})$, intriqué. L'environnement n'est pas observé : l'état de $S$ seul
est obtenu par trace partielle (section~\ref{sec:trace}),
\[
  \rho_S=\frac12\begin{pmatrix}1&\braket{E_1}{E_0}\\\braket{E_0}{E_1}&1\end{pmatrix}.
\]
Si les deux états de l'environnement sont orthogonaux, $\rho_S=\frac12 I$ : l'état est maximalement mixte
(section~\ref{sec:pur-mixte}) et la superposition est perdue. En pratique le recouvrement décroît
comme $\braket{E_1}{E_0}=\ee^{-t/T_2}$, d'où, pour la pureté,
\[
  \operatorname{Tr}\rho_S^2=\frac{1+\ee^{-2t/T_2}}{2}\ \xrightarrow[t\to\infty]{}\ \frac12 .
\]
(à $t=T_2$, elle vaut $0{,}57$). Les termes diagonaux, eux, évoluent avec le temps de relaxation $T_1$ (perte d'énergie : $\ket1\to\ket0$). Le rapport
entre $T_2$ et la durée $t_{\mathrm{porte}}$ d'une porte fixe le nombre d'opérations avant que l'information soit perdue. Avec
$T_2=100\ \mu\mathrm{s}$ et des portes de $50\ \mathrm{ns}$, il vaut $2\,000$ : l'erreur par porte est de l'ordre de
$t_{\mathrm{porte}}/T_2=5\times10^{-4}$. Calculer sur des milliards de portes sans corriger les erreurs est
donc impossible.

\subsubsection{Les plateformes}

Plusieurs systèmes physiques servent de qubits. Chacun a ses forces et ses faiblesses ; les valeurs du tableau sont des ordres de
grandeur tirés de la littérature récente et évoluent vite.

\begin{center}
\renewcommand{\arraystretch}{1.2}
\small
\begin{tabular}{@{}>{\raggedright\arraybackslash}p{2.7cm}>{\raggedright\arraybackslash}p{2.2cm}>{\raggedright\arraybackslash}p{1.9cm}>{\raggedright\arraybackslash}p{1.9cm}>{\raggedright\arraybackslash}p{5.5cm}@{}}
\toprule
\textbf{Plateforme} & \textbf{Température} & \textbf{Porte à 2 qubits} & \textbf{Cohérence} & \textbf{Atouts / difficultés}\\
\midrule
Circuits supra\-conducteurs \parencite{kjaergaard2020} & $10$--$20\ \mathrm{mK}$ & $20$--$100\ \mathrm{ns}$ & $\sim0{,}1\ \mathrm{ms}$ & portes rapides, fabrication par lithographie ; cryogénie, câblage, voisinage limité\\
Ions piégés \parencite{bruzewicz2019} & ambiante (vide poussé) & $10\ \mu\mathrm{s}$--$1\ \mathrm{ms}$ & secondes à minutes & meilleures fidélités ($\geq99{,}9\,\%$), connexion totale ; portes lentes, passage à grande échelle\\
Atomes neutres \parencite{saffman2016} & ambiante (vide poussé) & $0{,}1$--$1\ \mu\mathrm{s}$ & $\sim1\ \mathrm{s}$ & réseaux de milliers d'atomes reconfigurables ; erreurs de porte plus élevées\\
Photons \parencite{kok2007} & ambiante (détecteurs $\sim2\ \mathrm{K}$) & $\mathrm{ns}$ & pas de décohérence en vol & communication ; pertes, portes à deux photons probabilistes\\
Spins dans le silicium \parencite{burkard2023} & $0{,}1$--$1\ \mathrm{K}$ & $10$--$100\ \mathrm{ns}$ & $\sim\mathrm{ms}$ & compatible avec la microélectronique ; peu de qubits à ce jour\\
\bottomrule
\end{tabular}
\end{center}

Aucune plateforme ne domine. Les machines actuelles comptent de quelques dizaines à quelques milliers de qubits physiques
\emph{bruités}, sans correction d'erreurs : \textcite{preskill2018} les a appelées « NISQ » (\emph{noisy intermediate-scale quantum}). Elles ne suffisent pas pour Shor.

\subsubsection{Le cryostat : pourquoi $10\ \mathrm{mK}$ ?}

Un transmon (qubit supraconducteur) a une fréquence de transition $f\approx5\ \mathrm{GHz}$ : une micro-onde. À l'équilibre thermique
à la température $T$, la loi de Boltzmann donne $p_1/p_0=\ee^{-hf/k_BT}$ avec $p_0+p_1=1$, d'où
\[
  p_1=\frac{1}{1+\ee^{\,hf/k_BT}},\qquad
  \frac{hf}{k_B}=\frac{6{,}626\times10^{-34}\times5\times10^{9}}{1{,}381\times10^{-23}}\approx0{,}24\ \mathrm{K}.
\]
Le qubit n'est « dans l'état $\ket0$ » que si $T\ll0{,}24\ \mathrm{K}$. Sinon il est dans un mélange statistique de
$\ket0$ et $\ket1$, de matrice densité $\operatorname{diag}(1-p_1,p_1)$ (section~\ref{sec:rho}) :

\begin{center}
\renewcommand{\arraystretch}{1.15}
\begin{tabular}{@{}rccc@{}}
\toprule
Température $T$ & $hf/k_BT$ & population $p_1$ de $\ket1$ & pureté $\operatorname{Tr}\rho^2$\\
\midrule
$300\ \mathrm{K}$ (ambiante) & $0{,}0008$ & $0{,}4998$ & $0{,}500$\\
$4\ \mathrm{K}$ (hélium liquide) & $0{,}060$ & $0{,}485$ & $0{,}500$\\
$1\ \mathrm{K}$ & $0{,}24$ & $0{,}44$ & $0{,}507$\\
$100\ \mathrm{mK}$ & $2{,}4$ & $0{,}083$ & $0{,}85$\\
$20\ \mathrm{mK}$ & $12$ & $6{,}2\times10^{-6}$ & $0{,}99999$\\
$10\ \mathrm{mK}$ & $24$ & $3{,}8\times10^{-11}$ & $1{,}000$\\
\bottomrule
\end{tabular}
\end{center}

À température ambiante, $p_1\approx\frac12$ : l'état est maximalement mixte (pureté minimale d'un qubit,
$\operatorname{Tr}\rho^2=\frac12$, section~\ref{sec:pur-mixte}) et le qubit ne porte aucune information. Il faut descendre vers $20\ \mathrm{mK}$ pour que
l'état initial soit pur à mieux que $10^{-5}$ près. S'y ajoute une seconde raison : les jonctions Josephson sont en aluminium, supraconducteur
seulement sous $\approx1{,}2\ \mathrm{K}$, et les photons thermiques des lignes de commande doivent être éteints. Les autres
plateformes évitent ce problème autrement : ions et atomes sont préparés dans $\ket0$ par pompage optique (un laser les force dans l'état voulu)
et refroidis par laser ; les photons ont des fréquences optiques ($\sim10^{14}\ \mathrm{Hz}$, soit $hf/k_B\approx5\times10^3\ \mathrm{K}$),
sans aucune excitation thermique.

\paragraph{Le réfrigérateur à dilution.} On atteint $10\ \mathrm{mK}$ avec un mélange d'isotopes d'hélium, $^3$He et $^4$He. Sous
$0{,}87\ \mathrm{K}$, le mélange se sépare en deux phases, l'une riche en $^3$He, l'autre diluée. Faire passer des atomes de $^3$He de
la phase riche à la phase diluée \emph{absorbe de la chaleur}, comme une évaporation : c'est le principe de la dilution. Un
refroidisseur à tube pulsé fournit le froid initial (étages à $50\ \mathrm{K}$ puis $4\ \mathrm{K}$), et le circuit de dilution
descend en étages jusqu'à la \emph{chambre de mélange} (figure~\ref{fig:cryostat}). Le prix est la \emph{puissance
froide} : de l'ordre de $10\ \mu\mathrm{W}$ à $20\ \mathrm{mK}$ pour un appareil commercial, ce qui limite la chaleur que les lignes et
l'électronique peuvent apporter à cet étage.

\begin{figure}[!ht]
\centering
\adjustbox{max width=\linewidth}{%
\begin{tikzpicture}[>=Stealth, font=\small,
    plaque/.style={fill=insagris!35, draw=insagris, thick, minimum height=3.5mm, inner sep=0pt},
    legende/.style={font=\footnotesize, anchor=west, text width=7.2cm, align=left}]
  % plaques : (hauteur, largeur, température)
  \node[plaque, minimum width=6.4cm] (p0) at (0,0)    {};
  \node[plaque, minimum width=5.6cm] (p1) at (0,-1.3) {};
  \node[plaque, minimum width=4.8cm] (p2) at (0,-2.6) {};
  \node[plaque, minimum width=4.0cm] (p3) at (0,-3.9) {};
  \node[plaque, minimum width=3.2cm] (p4) at (0,-5.2) {};
  \node[plaque, minimum width=2.4cm, fill=insarouge!30, draw=insarouge] (p5) at (0,-6.5) {};
  \foreach \a/\b in {0/-1.3, -1.3/-2.6, -2.6/-3.9, -3.9/-5.2, -5.2/-6.5}
    {\draw[insagris, thick] (-0.9,\a-0.175) -- (-0.9,\b+0.175);
     \draw[insagris, thick] (0.9,\a-0.175) -- (0.9,\b+0.175);}
  % puce
  \node[draw=insarouge, fill=insarouge!15, thick, minimum width=14mm, minimum height=5mm] (puce) at (0,-7.15) {\footnotesize puce};
  \draw[insarouge, thick] (0,-6.675) -- (0,-6.9);
  % ligne de commande (coaxial) avec atténuateurs
  \draw[insarouge, thick] (0.3,0.35) -- (0.3,-6.9);
  \foreach \y/\att in {-2.6/{$-20$ dB}, -5.2/{$-20$ dB}, -6.5/{$-20$ dB}}
    {\node[draw=insarouge, fill=white, inner sep=1.8pt] at (0.3,\y) {};
     \node[font=\scriptsize, text=insarouge, anchor=west, fill=white, inner sep=1pt] at (1.3,\y) {\att};}
  % températures (gauche)
  \foreach \y/\T in {0/{$300\ \mathrm{K}$}, -1.3/{$50\ \mathrm{K}$}, -2.6/{$4\ \mathrm{K}$}, -3.9/{$\approx1\ \mathrm{K}$}, -5.2/{$\approx100\ \mathrm{mK}$}, -6.5/{$\approx10\ \mathrm{mK}$}}
    \node[anchor=east, font=\footnotesize\bfseries] at (-3.4,\y) {\T};
  % légendes (droite)
  \node[legende] at (3.7,0)    {électronique de commande et de lecture, ambiante};
  \node[legende] at (3.7,-1.3) {1\up{er} étage du tube pulsé, ancrage thermique des câbles};
  \node[legende] at (3.7,-2.6) {2\up{e} étage du tube pulsé ; amplificateur HEMT (lecture)};
  \node[legende] at (3.7,-3.9) {« still » du circuit de dilution};
  \node[legende] at (3.7,-5.2) {plaque froide : filtres et atténuateurs};
  \node[legende] at (3.7,-6.5) {chambre de mélange : qubits, amplificateur paramétrique};
\end{tikzpicture}}
\caption{Réfrigérateur à dilution en coupe schématique (« lustre »). Les plaques sont de plus en plus froides vers le
bas ; les atténuateurs placés sur la ligne de commande (valeurs typiques) ramènent le bruit thermique de $300\ \mathrm{K}$ à
celui de l'étage où ils sont fixés.}
\label{fig:cryostat}
\end{figure}

\paragraph{Le goulot d'étranglement : le câblage.} Chaque qubit demande plusieurs lignes coaxiales (commande, lecture),
qui apportent de la chaleur depuis les étages chauds et du bruit thermique. Des atténuateurs ancrés à chaque étage
ramènent ce bruit : $T_{\mathrm{eff}}\approx T_{\text{étage}}+T_{\text{chaud}}\,10^{-A/10}$, donc $60\ \mathrm{dB}$ au total limitent à
$300\ \mathrm{K}\times10^{-6}=0{,}3\ \mathrm{mK}$ la contribution du bruit de l'étage ambiant \parencite{krinner2019}. Mais des machines de millions de
qubits physiques ne tiennent pas dans un cryostat avec un coaxial par qubit : les pistes de recherche sont le multiplexage des lignes,
l'électronique de commande placée à $4\ \mathrm{K}$ et le raccordement de plusieurs cryostats entre eux.

\subsubsection{Corriger les erreurs : du code à répétition au code de surface}

Classiquement, on se protège du bruit en copiant chaque bit trois fois et en prenant le vote majoritaire. Le théorème de
non-clonage interdit de copier $\alpha\ket0+\beta\ket1$, et la mesure détruirait la superposition. On \emph{intrique} plutôt
le qubit avec deux autres.

\begin{definition}{Code à répétition à trois qubits}
Le qubit logique $\alpha\ket0+\beta\ket1$ est encodé par deux CNOT dont le contrôle est $q_1$ (et les cibles $q_2$, $q_3$) :
\[
  (\alpha\ket0+\beta\ket1)\ket{00}\ \longmapsto\ \alpha\ket{000}+\beta\ket{111}.
\]
Ce n'est pas $\ket\psi^{\otimes3}$ : l'état est intriqué. Dans cette section, les kets s'écrivent $\ket{q_1q_2q_3}$ et les
fils sont étiquetés : la convention « fil du haut = dernier chiffre » ne s'applique pas. Une erreur $X_k$ (inversion de bit sur $q_k$) est détectée en mesurant les
\emph{parités} $Z_1Z_2$ et $Z_2Z_3$, de valeurs propres $\pm1$.
\end{definition}

Sur $\alpha\ket{000}+\beta\ket{111}$ les deux parités valent $+1$. Une inversion fait changer le signe des parités qui contiennent le qubit touché :

\begin{center}
\renewcommand{\arraystretch}{1.15}
\begin{tabular}{@{}lccc@{}}
\toprule
erreur & état obtenu & $(Z_1Z_2,\ Z_2Z_3)$ & correction\\
\midrule
aucune & $\alpha\ket{000}+\beta\ket{111}$ & $(+1,+1)$ & aucune\\
$X_1$ & $\alpha\ket{100}+\beta\ket{011}$ & $(-1,+1)$ & $X_1$\\
$X_2$ & $\alpha\ket{010}+\beta\ket{101}$ & $(-1,-1)$ & $X_2$\\
$X_3$ & $\alpha\ket{001}+\beta\ket{110}$ & $(+1,-1)$ & $X_3$\\
\bottomrule
\end{tabular}
\end{center}

Les deux parités sont identiques pour $\ket{000}$ et $\ket{111}$ : leur mesure localise l'erreur \emph{sans rien apprendre sur
$\alpha$ et $\beta$}, donc sans détruire la superposition. On mesure les parités avec deux qubits
auxiliaires (\emph{ancillas}) $a_1$ et $a_2$ : le CNOT inscrit dans l'ancilla le ou-exclusif des deux qubits qu'il contrôle
(figure~\ref{fig:code3}).

\begin{figure}[!ht]
\centering
\adjustbox{max width=\linewidth}{%
\begin{tikzpicture}[>=Stealth, font=\small]
  \foreach \y in {3.6, 2.4, 1.2} \draw (0,\y) -- (11.3,\y);
  \foreach \y in {0, -1.2} \draw (0,\y) -- (8.7,\y);
  \node[anchor=east] at (0,3.6) {$q_1:\ \alpha\ket0+\beta\ket1$};
  \node[anchor=east] at (0,2.4) {$q_2:\ \ket0$};
  \node[anchor=east] at (0,1.2) {$q_3:\ \ket0$};
  \node[anchor=east] at (0,0) {$a_1:\ \ket0$};
  \node[anchor=east] at (0,-1.2) {$a_2:\ \ket0$};
  % encodage : CNOT q1 -> q2, q1 -> q3
  \draw[insarouge, thick] (1.1,3.6) -- (1.1,2.4);
  \pic at (1.1,3.6) {ctl}; \pic at (1.1,2.4) {cible};
  \draw[insarouge, thick] (2.0,3.6) -- (2.0,1.2);
  \pic at (2.0,3.6) {ctl}; \pic at (2.0,1.2) {cible};
  % erreur éventuelle
  \node[draw=insagris, dashed, minimum width=9mm, minimum height=9mm] at (3.3,2.4) {$X$};
  \node[font=\footnotesize, text=insagris, anchor=south] at (3.3,3.0) {erreur ?};
  % syndrome : a1 <- q1 xor q2 ; a2 <- q2 xor q3
  \draw[insarouge, thick] (4.5,3.6) -- (4.5,0);  \pic at (4.5,3.6) {ctl}; \pic at (4.5,0) {cible};
  \draw[insarouge, thick] (5.4,2.4) -- (5.4,0);  \pic at (5.4,2.4) {ctl}; \pic at (5.4,0) {cible};
  \draw[insarouge, thick] (6.3,2.4) -- (6.3,-1.2); \pic at (6.3,2.4) {ctl}; \pic at (6.3,-1.2) {cible};
  \draw[insarouge, thick] (7.2,1.2) -- (7.2,-1.2); \pic at (7.2,1.2) {ctl}; \pic at (7.2,-1.2) {cible};
  \pic at (8.4,0) {mesure};   \pic at (8.4,-1.2) {mesure};
  \node[anchor=west, font=\footnotesize] at (9.0,-0.6) {syndrome $(s_1,s_2)$};
  \node[porte] at (10.0,3.6) {$X^{?}$}; \node[porte] at (10.0,2.4) {$X^{?}$}; \node[porte] at (10.0,1.2) {$X^{?}$};
  \draw[dashed, insagris] (8.4,0.25) -- (8.4,0.5) -- (10.0,0.5) -- (10.0,0.85);
  \foreach \x/\t in {2.65/{encodage}, 6.0/{extraction du syndrome}}
    \node[font=\footnotesize, text=insagris] at (\x,-2.1) {\t};
\end{tikzpicture}}
\caption{Code à répétition à trois qubits : encodage, erreur éventuelle (ici sur $q_2$), extraction du syndrome
$(s_1,s_2)=(q_1\oplus q_2,\,q_2\oplus q_3)$ dans les ancillas, puis correction par une porte $X$ sur le qubit désigné
par le syndrome. Les kets s'écrivent $\ket{q_1q_2q_3}$.}
\label{fig:code3}
\end{figure}

\paragraph{Gain.} La correction échoue si au moins deux qubits sont inversés. Avec une probabilité $p$ d'inversion par
qubit, la probabilité d'échec est $3p^2(1-p)+p^3=3p^2-2p^3$, inférieure à $p$ si et seulement si $(2p-1)(p-1)>0$,
c'est-à-dire $p<\frac12$. Le code est utile \emph{si le bruit est assez faible}. C'est le principe d'un
\emph{seuil}.

\paragraph{Erreurs de phase.} Une erreur $Z_1$ sur $\alpha\ket{000}+\beta\ket{111}$ donne $\alpha\ket{000}-\beta\ket{111}$ : un état
logique différent, que ce code ne voit pas. On se protège en travaillant dans la base $X$ : $H^{\otimes3}(\alpha\ket{000}+\beta\ket{111})
=\alpha\ket{+++}+\beta\ket{---}$, où une erreur $Z$ (qui échange $\ket+$ et $\ket-$) joue le rôle d'une inversion de
bit. \textcite{shor1995} a montré que l'on peut \emph{concaténer} les deux codes (neuf qubits) pour corriger n'importe quelle erreur sur un
qubit : toute erreur est une combinaison de $I,X,Y,Z$ et la mesure du syndrome l'y projette.

\paragraph{Code de surface.} Les qubits sont disposés sur un réseau plan de $d\times d$ qubits de donnée, entourés de qubits
de mesure ; les syndromes de type $X$ et de type $Z$ sont mesurés à chaque cycle. La \emph{distance} $d$ est le nombre minimal
d'erreurs qui changent l'état logique : le code corrige jusqu'à $\lfloor(d-1)/2\rfloor$ erreurs, et il faut
$2d^2-1$ qubits physiques par qubit logique. Si le taux d'erreur physique est sous le seuil ($\approx1\,\%$), l'erreur
logique \emph{décroît exponentiellement avec $d$} : à chaque passage de $d$ à $d+2$, elle est divisée par un
facteur $\Lambda>1$. Google a observé ce comportement sur des codes de distance 3, 5 puis 7, avec
$\Lambda=2{,}14\pm0{,}02$ et une erreur logique de $0{,}143\,\%$ par cycle à $d=7$ (101 qubits physiques)
\parencite{acharya2025}. En prolongeant avec ce même $\Lambda$, atteindre une erreur logique de $10^{-12}$ par cycle demande
$\Lambda^{k}\geq1{,}43\times10^{-3}/10^{-12}$, soit $k\geq\ln(1{,}43\times10^{9})/\ln(2{,}14)\approx27{,}7$, donc
$k=28$ et $d=7+2\times28=63$ : près de $8\,000$ qubits physiques par qubit logique. C'est une extrapolation (l'erreur n'est pas
exactement exponentielle), mais elle explique les millions de qubits de la section~\ref{sec:shor} : le surcoût de la
correction d'erreurs domine tout, et un meilleur $\Lambda$ (des qubits physiques moins bruités) le réduit fortement.

\subsubsection{Les qubits de chat d'Alice \& Bob}

Le code de surface doit corriger \emph{deux} types d'erreurs ($X$ et $Z$) sur un réseau plan, d'où un surcoût en $d^2$. L'idée des
\emph{qubits de chat} \parencite{mirrahimi2014} est de concevoir un qubit dont une des deux erreurs est déjà supprimée par la physique. Il ne reste
alors qu'un type d'erreur à corriger : un code \emph{unidimensionnel} (la répétition de la section précédente), de $2d-1$ qubits
au lieu de $2d^2-1$, suffit.

\paragraph{Un qubit dans un oscillateur.} On utilise un mode d'un résonateur micro-onde, oscillateur harmonique d'états de
Fock $\ket n$ ($n$ photons). Un \emph{état cohérent} $\ket\alpha=\ee^{-|\alpha|^2/2}\sum_n\frac{\alpha^n}{\sqrt{n!}}\ket n$ est vecteur propre de l'opérateur
d'annihilation : $a\ket\alpha=\alpha\ket\alpha$, avec $|\alpha|^2$ photons en moyenne. Pour $\alpha$ réel, on pose
\[
  \ket{0_L}\approx\ket{\alpha},\qquad\ket{1_L}\approx\ket{-\alpha},\qquad
  \braket{\alpha}{-\alpha}=\ee^{-|\alpha|^2-|\alpha|^2}=\ee^{-2|\alpha|^2}.
\]
Le recouvrement est $1{,}8\times10^{-2}$ pour $|\alpha|^2=2$, $3{,}4\times10^{-4}$ pour $|\alpha|^2=4$ et $1{,}1\times10^{-7}$ pour
$|\alpha|^2=8$ : les deux états sont presque orthogonaux dès que $|\alpha|^2\gtrsim4$. Les états de la base $X$ sont les
\emph{états de chat} de Schrödinger
\[
  \ket{C_\pm}=N_\pm\big(\ket\alpha\pm\ket{-\alpha}\big),\qquad N_\pm=\frac{1}{\sqrt{2\,(1\pm\ee^{-2|\alpha|^2})}},\qquad
  \ket{\pm_L}\approx\ket{C_\pm}.
\]
$\ket{C_+}$ ne contient que des nombres \emph{pairs} de photons et $\ket{C_-}$ que des nombres \emph{impairs} (les termes en
$\alpha^n$ et $(-\alpha)^n$ s'annulent ou s'ajoutent selon la parité de $n$).

\paragraph{Une perte de photon est une erreur de phase.} Le bruit principal d'un résonateur est la perte d'un photon, décrite par $a$.
Comme $a\ket{\pm\alpha}=\pm\alpha\ket{\pm\alpha}$,
\[
  a\ket{C_\pm}=\alpha N_\pm\big(\ket\alpha\mp\ket{-\alpha}\big)=\alpha\,\frac{N_\pm}{N_\mp}\,\ket{C_\mp}.
\]
Après normalisation, $\ket{C_+}\leftrightarrow\ket{C_-}$ : c'est $\ket{+_L}\leftrightarrow\ket{-_L}$, une \emph{inversion de phase} $Z$. Sur
$\ket{\pm\alpha}$, $a$ ne change rien (un simple facteur), donc la perte de photon ne provoque aucune inversion de bit. Elle survient au taux
$\kappa_1\langle n\rangle=\kappa_1|\alpha|^2$ : le taux d'inversion de phase \emph{croît linéairement} avec $|\alpha|^2$.

\paragraph{Supprimer les inversions de bit.} On ajoute une dissipation \emph{à deux photons} : le circuit éjecte des paires de photons
(taux $\kappa_2$) et en réinjecte par paires (pompe d'amplitude $\epsilon_2$). Cela équivaut à un saut d'opérateur
$\sqrt{\kappa_2}\,(a^2-\alpha^2)$ avec $|\alpha|^2$ fixé par le rapport pompe/perte. Or $(a^2-\alpha^2)\ket{\pm\alpha}=0$ car
$a^2\ket{\pm\alpha}=\alpha^2\ket{\pm\alpha}$ : l'espace $\mathrm{Vect}(\ket\alpha,\ket{-\alpha})$ est exactement le noyau de la dissipation. Elle ramène
tout état de l'oscillateur vers cet espace et le \emph{stabilise}. Retirer les photons deux par deux conserve la parité du nombre de photons, donc la
cohérence entre $\ket{C_+}$ et $\ket{C_-}$. Pour qu'un bruit transforme $\ket\alpha$ en $\ket{-\alpha}$, il doit franchir la distance $2\alpha$ dans
le plan de phase contre la dissipation, ce qui est exponentiellement improbable : le taux d'inversion de bit décroît
comme $\ee^{-2|\alpha|^2}$.

\begin{figure}[!ht]
\centering
\begin{tikzpicture}[>=Stealth, font=\small]
  \draw[->, insagris] (-4.6,0) -- (4.8,0) node[right] {$\mathrm{Re}\,\alpha$};
  \foreach \sg in {1,-1}
    \shade[inner color=insarouge!80, outer color=insarouge!5] (\sg*2.6,0) circle (0.75);
  \node[anchor=north, font=\footnotesize] at (2.6,-0.85) {$\ket{\alpha}=\ket{0_L}$};
  \node[anchor=north, font=\footnotesize] at (-2.6,-0.85) {$\ket{-\alpha}=\ket{1_L}$};
  \draw[<->, dashed, insagris, thick] (-1.8,0) -- (1.8,0);
  \node[fill=white, inner sep=1pt, text=insarouge] at (0,0) {$\times$};
  \node[anchor=north, font=\footnotesize, text=insagris, align=center] at (0,-0.2)
    {inversion de bit : $\sim\ee^{-2|\alpha|^2}$\\ (supprimée)};
  \draw[->, insarouge, thick] (-1.5,0.75) to[bend left=40] (1.5,0.75);
  \node[anchor=south, font=\footnotesize, text=insarouge, align=center] at (0,1.3)
    {photon perdu : $\ket{C_+}\leftrightarrow\ket{C_-}$\\ (inversion de phase, $\propto\kappa_1|\alpha|^2$)};
\end{tikzpicture}
\caption{Qubit de chat dans le plan de phase de l'oscillateur (axe réel) : deux états cohérents $\ket{\pm\alpha}$ représentent $\ket{0_L}$
et $\ket{1_L}$. La dissipation à deux photons les stabilise ; passer de l'un à l'autre (inversion de bit) est exponentiellement
rare, alors que la perte d'un photon échange les états de chat $\ket{C_\pm}\approx\ket{\pm_L}$, de parités opposées (inversion de phase).}
\label{fig:chat}
\end{figure}

\paragraph{Résultat : un biais de bruit énorme.} On augmente $|\alpha|^2$ pour supprimer les inversions de bit, au prix d'inversions de phase
plus fréquentes. Mais l'asymétrie est telle que le compromis est favorable. Les expériences sont de plus en plus convaincantes :
des équipes parisiennes (ENS, Inria, Alice \& Bob) ont observé la suppression exponentielle de l'inversion de bit
\parencite{lescanne2020}, puis un temps d'inversion de bit supérieur à $10\ \mathrm{s}$, pour un temps d'inversion de phase supérieur à
$490\ \mathrm{ns}$ \parencite{reglade2024}. Le rapport est de l'ordre de $2\times10^{7}$ ($10/(0{,}49\times10^{-6})$). Une \emph{annonce de l'entreprise} Alice \& Bob (septembre 2025, résultat préliminaire, non évalué par des pairs à
notre connaissance) fait état de plus d'une heure avec $\langle n\rangle=11$ \parencite{aliceandbob2025}.

\paragraph{De là à un ordinateur.} Il reste à corriger les inversions de phase, par un code à répétition de $d$ qubits de chat
\parencite{guillaud2019}, dont les portes à deux qubits (CNOT) doivent \emph{préserver le biais} : c'est le point que les expériences
doivent encore démontrer. Une estimation d'architecture évalue le calcul d'un logarithme discret sur une courbe de 256 bits à
$126\,133$ qubits de chat et $9$ heures \parencite{gouzien2023}. Alice \& Bob (start-up française) et Amazon \parencite{putterman2025} suivent cette voie ;
Alice \& Bob vise, dans sa \emph{feuille de route} d'entreprise, une machine de 100 qubits logiques en 2030. C'est un objectif commercial, pas un résultat.

% ---------------------------------------------------------------------
\subsection{Autres applications et algorithmes}\label{sec:autres}
% ---------------------------------------------------------------------

Shor et Grover sont les algorithmes les plus connus, pas les seuls. Il faut distinguer les gains \emph{prouvés} (on sait démontrer
qu'aucun algorithme classique ne fait aussi bien, ou on connaît un bien meilleur algorithme quantique), les gains
\emph{conditionnels} (valables sous hypothèses) et les méthodes \emph{heuristiques} (qui marchent peut-être, sans preuve).

\subsubsection{Panorama}

\begin{center}
\renewcommand{\arraystretch}{1.2}
\small
\begin{tabular}{@{}>{\raggedright\arraybackslash}p{3.1cm}>{\raggedright\arraybackslash}p{3.8cm}>{\raggedright\arraybackslash}p{3.0cm}>{\raggedright\arraybackslash}p{4.6cm}@{}}
\toprule
\textbf{Algorithme} & \textbf{Problème} & \textbf{Gain} & \textbf{Statut et usage}\\
\midrule
Deutsch--Jozsa, Bernstein--Vazirani, Simon (1985--1997) & fonctions « boîtes noires » & exponentiel (Simon) & prouvé ; peu d'usage direct, base de Shor\\
Shor (1994) \parencite{shor1997} & factorisation, logarithme discret & superpolynomial & prouvé ; casse RSA et courbes elliptiques\\
Grover (1996) & recherche non structurée & quadratique : $\sqrt N$ & prouvé, optimal ; clés symétriques, SAT\\
Estimation de phase \parencite{kitaev1995} & valeur propre d'une porte $U$ & précision $2^{-m}$ avec $m$ qubits & prouvé ; brique de Shor, de la chimie, de HHL\\
Estimation d'amplitude \parencite{brassard2002} & estimer une probabilité & $\varepsilon^{-1}$ au lieu de $\varepsilon^{-2}$ & prouvé ; Monte-Carlo, finance\\
Simulation hamiltonienne \parencite{lloyd1996} & évolution $\ee^{-\ii Ht}\ket\psi$ & exponentiel (systèmes quantiques) & prouvé ; chimie, matériaux : application la plus naturelle\\
HHL \parencite{harrow2009} & système linéaire $Ax=b$ & exponentiel \emph{sous conditions} & conditionnel ; en partie « déquantisé » \parencite{tang2019}\\
VQE \parencite{peruzzo2014} & énergie de l'état fondamental & aucun gain prouvé & heuristique ; machines bruitées\\
QAOA \parencite{farhi2014} & optimisation combinatoire & aucun gain prouvé & heuristique ; machines bruitées\\
\bottomrule
\end{tabular}
\end{center}

\subsubsection{Simuler la matière : l'idée de Feynman}

Pour décrire $n$ qubits, un ordinateur classique stocke $2^n$ amplitudes complexes. Avec $16$ octets par amplitude, $n=50$
demande $2^{50}\times16=2^{54}$ octets, soit $16$ pébioctets ($1{,}8\times10^{16}$ octets) ; $n=300$ demande $2^{300}\approx2\times10^{90}$
amplitudes, plus que d'atomes dans l'univers observable ($\sim10^{80}$). \textcite{feynman1982} en a tiré l'idée
inverse : un système quantique simule naturellement un autre système quantique.

La \emph{simulation hamiltonienne} calcule $\ee^{-\ii Ht}\ket{\psi}$ pour $H=A+B$, où $A$ et $B$ sont des sommes de termes simples dont on sait implémenter les exponentielles.
Si $A$ et $B$ ne commutent pas, $\ee^{-\ii(A+B)t}\neq\ee^{-\ii At}\ee^{-\ii Bt}$. La formule de Baker--Campbell--Hausdorff donne, pour un pas $\tau$,
\[
  \ee^{-\ii A\tau}\ee^{-\ii B\tau}=\exp\!\Big(-\ii(A+B)\tau-\tfrac{\tau^2}{2}[A,B]+O(\tau^3)\Big).
\]
On coupe donc $t$ en $N$ pas $\tau=t/N$ (\emph{formule de Trotter}) :
\[
  \Big\|\ee^{-\ii(A+B)t}-\big(\ee^{-\ii At/N}\ee^{-\ii Bt/N}\big)^N\Big\|\ \leq\ \frac{t^2}{2N}\,\big\|[A,B]\big\|\quad(\text{à l'ordre dominant}),
\]
l'erreur de chaque pas ($\sim\tau^2\|[A,B]\|/2$) s'additionnant sur $N$ pas. Pour $A=X$ et $B=Z$ (un qubit), $[X,Z]=-2\ii Y$ et
$\|[X,Z]\|=2$ : l'erreur vaut au plus $t^2/N$ ; pour $t=1$, $N=100$ pas donnent une erreur inférieure à $10^{-2}$ (le calcul numérique donne $7\times10^{-3}$). Les applications visées sont la chimie
(énergie d'une molécule comme le cofacteur FeMoco de la nitrogénase, qui intervient dans la fabrication des engrais \parencite{reiher2017}),
les matériaux (batteries, supraconducteurs) et la physique des particules. Les méthodes classiques (théorie de la fonctionnelle de
la densité, \emph{coupled cluster}) traitent bien les systèmes faiblement corrélés ; c'est pour les systèmes \emph{fortement
corrélés} qu'un avantage est espéré, avec des machines tolérantes aux fautes.

\subsubsection{Méthodes variationnelles et optimisation}

Sur les machines bruitées actuelles, on ne peut pas faire tourner de longs circuits ; on utilise donc des circuits courts à paramètres
$\theta$ dans une boucle avec un ordinateur classique : le circuit prépare $\ket{\psi(\theta)}$, on mesure l'énergie
$E(\theta)=\bra{\psi(\theta)}H\ket{\psi(\theta)}$, et un optimiseur classique corrige $\theta$. Cela repose sur le
\emph{principe variationnel} : en décomposant $\ket\psi=\sum_kc_k\ket{k}$ sur les états propres de $H$ (section~\ref{sec:spectral}),
$E=\sum_k|c_k|^2E_k\geq E_0\sum_k|c_k|^2=E_0$. Minimiser $E(\theta)$ approche donc l'énergie fondamentale $E_0$ par le haut. Exemple pour un qubit,
$H=X+Z$ (valeurs propres $\pm\sqrt2$) avec $\ket{\psi(\theta)}=\cos\frac\theta2\ket0+\sin\frac\theta2\ket1$ :
\[
  E(\theta)=\langle Z\rangle+\langle X\rangle=\cos\theta+\sin\theta=\sqrt2\cos\!\big(\theta-\tfrac\pi4\big),
\]
minimale pour $\theta=\frac{5\pi}{4}$, de valeur $-\sqrt2=E_0$. Le VQE \parencite{peruzzo2014} généralise cette idée à des molécules, et le QAOA
\parencite{farhi2014} l'applique à des problèmes d'optimisation en alternant deux types de portes. Aucun avantage n'est démontré pour ces méthodes :
elles sont l'objet de la recherche actuelle.

\subsubsection{Algèbre linéaire quantique : une promesse sous conditions}

L'algorithme HHL \parencite{harrow2009} prépare un état $\ket{x}\propto\sum_ix_i\ket{i}$ où $Ax=b$, en un temps polynomial en $\log N$
là où la résolution classique est polynomiale en $N$. Le gain exponentiel apparent est soumis à plusieurs conditions. Il faut
\emph{charger} $b$ dans l'état $\ket{b}$ efficacement ; $A$ doit être creuse et bien conditionnée ; et le résultat est un
\emph{état} : en lire les $N$ composantes demanderait déjà $N$ répétitions, de sorte que l'on n'obtient que des
grandeurs globales de $x$. De plus, des algorithmes classiques « inspirés du quantique » (\emph{déquantisation}) obtiennent le même
résultat sous des hypothèses d'accès analogues \parencite{tang2019}, ce qui a supprimé l'avantage exponentiel annoncé pour
certaines applications d'apprentissage. Le bilan honnête : \emph{pas d'avantage démontré en apprentissage automatique à ce jour}.

\subsubsection{Capteurs, nombres aléatoires, réseaux}

\paragraph{Métrologie quantique.} Les premières technologies quantiques commercialisées sont des capteurs (horloges atomiques,
magnétomètres, gravimètres). L'intrication y améliore la précision. Soit $N$ qubits qui subissent chacun la phase
$\ket1\mapsto\ee^{\ii\varphi}\ket1$ à mesurer. Dans l'état de GHZ $\frac{1}{\sqrt2}(\ket{0\ldots0}+\ket{1\ldots1})$ (généralisation à $N$ qubits de l'état de Bell
$\frac{1}{\sqrt2}(\ket{00}+\ket{11})$, section~\ref{sec:bell}), seuls $\ket{0\ldots0}$ et $\ket{1\ldots1}$ interviennent, et l'état devient
\[
  \frac{1}{\sqrt2}\big(\ket{0\ldots0}+\ee^{\ii N\varphi}\ket{1\ldots1}\big),\qquad
  P\big(X^{\otimes N}=+1\big)=\frac{1+\cos N\varphi}{2}=\cos^2\!\frac{N\varphi}{2}.
\]
La phase est amplifiée par $N$ : la précision atteint $\Delta\varphi\sim1/N$ (limite de Heisenberg) au lieu de $1/\sqrt N$ pour
$N$ qubits indépendants (limite quantique standard) \parencite{giovannetti2004}. En pratique, la décohérence réduit ce gain, qui
se limite souvent à de petits $N$.

\paragraph{Nombres aléatoires.} Mesurer $\ket+$ dans la base $Z$ donne $0$ ou $1$ avec la probabilité $\frac12$ : le postulat de Born
garantit un hasard \emph{intrinsèque}, contrairement à un générateur pseudo-aléatoire déterministe. Les générateurs quantiques
de nombres aléatoires (QRNG), à base de photons, sont déjà commercialisés ; ils alimentent par exemple la génération de clés cryptographiques.

\paragraph{Réseaux quantiques.} Relier des ordinateurs quantiques demande de transmettre des qubits (téléportation, section~\ref{sec:teleportation}).
Une fibre optique standard atténue de $0{,}2\ \mathrm{dB/km}$, soit une transmission $T=10^{-0{,}02L}$ ($L$ en km) :
$T=10^{-2}$ à $100\ \mathrm{km}$ et $10^{-10}$ à $500\ \mathrm{km}$. Un amplificateur est impossible (non-clonage) : on utilise des
\emph{répéteurs quantiques}, qui étendent l'intrication par \emph{échange d'intrication} (une mesure de Bell entre deux paires intriquées, comme dans la
téléportation de la section~\ref{sec:teleportation}), ou des satellites : la distribution
d'intrication sur $1\,200\ \mathrm{km}$ a été démontrée par le satellite chinois Micius \parencite{yin2017}. L'objectif est un « internet
quantique » \parencite{wehner2018} reliant des processeurs, des capteurs et des liens de QKD.

\subsubsection{Jalons et niveaux de preuve}

Le tableau place quelques jalons dans le temps et indique la \emph{nature de la source}, car une annonce d'entreprise ou une
prépublication n'a pas la valeur d'un article évalué par des pairs ou d'une norme.

\begin{center}
\renewcommand{\arraystretch}{1.2}
\small
\begin{tabular}{@{}>{\raggedright\arraybackslash}p{1.9cm}>{\raggedright\arraybackslash}p{8.8cm}>{\raggedright\arraybackslash}p{3.3cm}@{}}
\toprule
\textbf{Date} & \textbf{Jalon} & \textbf{Nature de la source}\\
\midrule
1994--1996 & algorithmes de Shor et de Grover \parencite{shor1997} & articles évalués\\
2019 & Google, $53$ qubits : échantillonnage de circuits aléatoires en $200\ \mathrm{s}$, annoncé hors de portée des machines classiques \parencite{arute2019} (contesté ensuite) & article évalué\\
août 2024 & normes NIST FIPS 203, 204, 205 \parencite{fips203} & norme\\
2024 & qubit de chat : inversion de bit $>10\ \mathrm{s}$ \parencite{reglade2024} & article évalué\\
déc. 2024 & code de surface sous le seuil, $\Lambda=2{,}14$ \parencite{acharya2025} & article évalué\\
mars 2025 & HQC choisi par le NIST \parencite{nist2025hqc} & rapport officiel\\
mai 2025 & RSA-2048 : moins d'un million de qubits \parencite{gidney2025} & prépublication\\
sept. 2025 & qubit de chat : plus d'une heure d'inversion de bit \parencite{aliceandbob2025} & annonce d'entreprise\\
mars 2026 & Shor avec $10\,000$ à $26\,000$ atomes \parencite{cain2026} & prépublication\\
août 2026 & RFC~10024 : TLS hybride X25519 + ML-KEM \parencite{rfc10024} & norme IETF\\
\bottomrule
\end{tabular}
\end{center}

Ce qui est \emph{acquis} : les gains théoriques de Shor et de Grover, la correction d'erreurs sous le seuil sur de petits codes,
et les standards post-quantiques. Ce qui \emph{ne l'est pas} : un ordinateur à plusieurs milliers de qubits \emph{logiques},
un avantage utile sur un problème pratique, et surtout la date $Z$ de l'inégalité de Mosca. C'est pourquoi la migration cryptographique se
décide sur $X+Y$, qui se connaît, et non sur une prédiction de $Z$.

\clearpage
% ---------------------------------------------------------------------
\subsection{Récapitulatif du chapitre}\label{sec:recap-appli}
% ---------------------------------------------------------------------

\begin{center}
\adjustbox{max width=\linewidth}{%
\begin{tikzpicture}[>=Stealth, font=\footnotesize,
    boite/.style={draw=insarouge, fill=insarouge!6, rounded corners=2pt, thick,
                  minimum width=4.4cm, minimum height=1.45cm, align=center, inner sep=3pt},
    lien/.style={->, very thick, insagris},
    etiq/.style={font=\scriptsize\itshape, text=insagris, inner sep=2pt}]
  \node[boite] (int) at (0,0)
    {\textbf{Interférence}\\ oracle, Deutsch, Grover : $\sqrt N$\\ {\scriptsize\S\,\ref{sec:interference}, \ref{sec:grover}}};
  \node[boite] (sho) at (6.3,0)
    {\textbf{Shor}\\ QFT, période $r$, $O(n^3)$\\ {\scriptsize\S\,\ref{sec:shor}}};
  \node[boite] (cry) at (12.6,0)
    {\textbf{Cryptographie}\\ RSA et ECC cassés ; HNDL ; PQC\\ {\scriptsize\S\,\ref{sec:crypto}}};
  \node[boite] (cryo) at (12.6,-2.3)
    {\textbf{Qubits physiques}\\ décohérence, cryostat $10\ \mathrm{mK}$\\ {\scriptsize\S\,\ref{sec:machine}}};
  \node[boite] (qec) at (6.3,-2.3)
    {\textbf{Correction d'erreurs}\\ seuil, code de surface, qubits de chat\\ {\scriptsize\S\,\ref{sec:machine}}};
  \node[boite] (aut) at (0,-2.3)
    {\textbf{Autres usages}\\ simulation, VQE, capteurs, réseaux\\ {\scriptsize\S\,\ref{sec:autres}}};
  \draw[lien] (int) -- (sho) node[etiq, midway, above=1pt] {période};
  \draw[lien] (sho) -- (cry) node[etiq, midway, above=1pt] {casse RSA};
  \draw[lien] (cry) -- (cryo) node[etiq, midway, right=1pt] {quelle machine ?};
  \draw[lien] (cryo) -- (qec) node[etiq, midway, above=1pt] {qubits fiables};
  \draw[lien] (qec) -- (aut) node[etiq, midway, above=1pt] {usages};
\end{tikzpicture}}
\end{center}

\begin{aretenir}
\begin{itemize}
  \item Un algorithme quantique organise des \emph{interférences} : l'oracle écrit $f(x)$ dans une phase ; Deutsch décide avec un appel au lieu de deux.
        Grover : $\approx\frac\pi4\sqrt N$ itérations (rotations de $2\theta$), gain quadratique et optimal ; exact en une itération pour $N=4$.
  \item Shor : la QFT transforme la période $r$ de $a^x\bmod N$ en pics en $y=jM/r$ ; $\operatorname{pgcd}(a^{r/2}\pm1,N)$ donne les facteurs ;
        coût $O(n^3)$ ($n=\log_2N$), contre un temps sous-exponentiel classique (exemple : $N=15$, $a=7$, $r=4$).
  \item RSA, Diffie--Hellman, ECDH et ECDSA tombent face à Shor ; AES-256 et SHA-256 résistent à Grover. Doubler les clés RSA ne protège pas.
  \item « Stocker maintenant, déchiffrer plus tard » : exposé dès que $X+Y>Z$ (Mosca) ; la forward secrecy n'aide pas. On migre d'abord l'échange de clés
        (hybride X25519 + ML-KEM), puis les signatures (FIPS 203--205).
  \item Qubit physique : la décohérence ($T_2$) mélange l'état ; le cryostat à $10\ \mathrm{mK}$ rend $p_1$ négligeable ($hf/k_B=0{,}24\ \mathrm{K}$ à $5\ \mathrm{GHz}$). Correction
        d'erreurs : seuil, $2d^2-1$ qubits physiques par qubit logique (code de surface).
  \item Qubit de chat : $\ket{\pm\alpha}$, perte de photon = inversion de phase, inversion de bit en $\ee^{-2|\alpha|^2}$ ; code à répétition. Une annonce d'entreprise
        n'est pas un résultat évalué.
  \item Autres usages : simulation (Trotter), VQE et QAOA (heuristiques), HHL (conditionnel), capteurs (limite de Heisenberg), réseaux (répéteurs).
\end{itemize}
\end{aretenir}

\begin{center}
\renewcommand{\arraystretch}{1.12}
\small
\begin{tabular}{@{}l l r@{}}
\toprule
\textbf{Notion} & \textbf{Formule} & \textbf{§} \\
\midrule
Rebond de phase & $U_f\ket{x}\ket{-}=(-1)^{f(x)}\ket{x}\ket{-}$ & \ref{sec:interference} \\
Grover & $P=\sin^2\!\big((2k+1)\theta\big)$, \quad $\sin\theta=\frac{1}{\sqrt N}$, \quad $k\simeq\frac\pi4\sqrt N$ & \ref{sec:grover} \\
QFT et période & $\mathrm{QFT}_M\ket x=\frac{1}{\sqrt M}\sum_y\ee^{2\ii\pi xy/M}\ket y$ ; pics en $y=jM/r$, $P=\frac1r$ & \ref{sec:shor} \\
Facteurs, coût & $\operatorname{pgcd}(a^{r/2}\pm1,N)$, \quad $O(n^3)$ & \ref{sec:shor} \\
Mosca & $X+Y>Z$ : données déjà exposées & \ref{sec:crypto} \\
Population thermique & $p_1=\frac{1}{1+\ee^{hf/k_BT}}$, \quad $\operatorname{Tr}\rho^2=\frac{1+\ee^{-2t/T_2}}{2}$ (décohérence) & \ref{sec:machine} \\
Qubit de chat & $\braket{\alpha}{-\alpha}=\ee^{-2|\alpha|^2}$, \quad $a\ket{C_\pm}=\alpha\frac{N_\pm}{N_\mp}\ket{C_\mp}$ & \ref{sec:machine} \\
\bottomrule
\end{tabular}
\end{center}

\begin{savoirfaire}
\begin{itemize}[nosep]
  \item Suivre un circuit de Deutsch ou de Grover ($N=4$) jusqu'à la mesure (sections~\ref{sec:interference}, \ref{sec:grover}).
  \item Trouver l'ordre de $a$ modulo $N$, appliquer le critère du pgcd et prévoir la loi de $y$ (section~\ref{sec:shor}, exemple $N=15$).
  \item Dire ce que casse Shor, ce que Grover affaiblit, et appliquer l'inégalité de Mosca (section~\ref{sec:crypto}).
  \item Calculer $p_1$, la pureté après décohérence et le recouvrement $\ee^{-2|\alpha|^2}$ ; lire un syndrome (section~\ref{sec:machine}).
  \item Distinguer résultat prouvé, conditionnel et annonce (section~\ref{sec:autres}).
\end{itemize}
\end{savoirfaire}
, entre chapitre5 et
%  chapitre6 (exercices). Il ne contient ni préambule ni \begin{document} : les
%  environnements (definition, resultat, remarque, aretenir, savoirfaire), les
%  macros (\ii, \ee, \dd, \ket, \bra, \braket), les symboles de circuits (ctl,
%  cible, croix, mesure, porte) et les couleurs (insarouge, insagris) sont définis
%  dans main.tex. Un chapitre = une \section.
%  Le numéro d'un chapitre est donné par l'ordre des \include, pas par le nom du
%  fichier : ce chapitre est le chapitre 6, les exercices (chapitre6.tex) le 7.
% =====================================================================

% Macros de Dirac (déjà définies dans main.tex ; ce bloc évite l'erreur
% « Undefined control sequence \ket » si main.tex est une ancienne version).
\providecommand{\ket}[1]{|#1\rangle}
\providecommand{\bra}[1]{\langle #1|}
\providecommand{\braket}[2]{\langle #1|#2\rangle}

% =====================================================================
\section{Applications et algorithmes quantiques}\label{chap:applications}
% =====================================================================

Les chapitres précédents ont construit le formalisme : qubits, portes, intrication, mesure. Ce
chapitre montre à quoi il sert. On y voit comment un circuit exploite l'interférence pour résoudre
certains problèmes plus vite qu'un ordinateur classique (sections~\ref{sec:interference}
à~\ref{sec:shor}), pourquoi cela menace une partie de la cryptographie actuelle, RSA et TLS en tête
(section~\ref{sec:crypto}), ce qu'il faut pour construire une telle machine, du cryostat aux qubits
de chat d'Alice \& Bob (section~\ref{sec:machine}), et quelles autres applications et quels autres
algorithmes existent (section~\ref{sec:autres}). Les circuits sont dessinés comme aux chapitres~4 et~5 :
le fil du haut est le dernier chiffre du ket.

\begin{remarque}[Chiffres d'actualité]
Les estimations de ressources, les calendriers de migration et les records expérimentaux cités ici
datent d'octobre 2026 et évoluent vite. Chaque chiffre est accompagné de sa source (bibliographie) ;
les annonces d'entreprises sont signalées comme telles et les prépublications aussi.
\end{remarque}

% ---------------------------------------------------------------------
\subsection{Calculer par interférence}\label{sec:interference}
% ---------------------------------------------------------------------

Un registre de $n$ qubits est décrit par $2^n$ amplitudes (chapitre~\ref{chap:operateurs}). Une porte
agit sur toutes ces amplitudes à la fois, mais la mesure n'en révèle qu'un seul résultat, tiré au
hasard avec les probabilités de Born. Dire que «~le circuit calcule toutes les valeurs en parallèle~»
est donc trompeur. Un algorithme quantique est une \emph{chorégraphie d'interférences} : les
amplitudes des mauvaises réponses se compensent (signes ou phases opposés), celles de la bonne
réponse s'additionnent, et la mesure finale donne la bonne réponse avec une probabilité proche de $1$.

\begin{definition}{Oracle}
Pour $f:\{0,1\}^n\to\{0,1\}$, l'\emph{oracle} de $f$ est la porte à $n+1$ qubits
\[
  U_f\ket{x}\ket{y}=\ket{x}\ket{y\oplus f(x)}.
\]
Elle est unitaire et $U_f^2=\mathrm{Id}$ (car $y\oplus f(x)\oplus f(x)=y$). Elle se construit à partir
d'un circuit réversible de $f$ (portes CNOT et Toffoli, section~\ref{sec:controlees}). Le coût d'un
algorithme se mesure en \emph{nombre d'appels} à $U_f$.
\end{definition}

\begin{resultat}{Rebond de phase de l'oracle}
Avec l'ancilla dans l'état $\ket{-}=\frac{1}{\sqrt2}(\ket0-\ket1)$ :
\[
  U_f\ket{x}\ket{-}=(-1)^{f(x)}\ket{x}\ket{-}.
\]
\end{resultat}

En effet $U_f\ket{x}\ket{-}=\frac{1}{\sqrt2}\big(\ket{x}\ket{f(x)}-\ket{x}\ket{1\oplus f(x)}\big)$.
Si $f(x)=0$ on retrouve $\ket{x}\ket{-}$ ; si $f(x)=1$, les deux termes sont échangés et le signe
change. La valeur $f(x)$ est écrite dans la \emph{phase} de $\ket{x}$, sans toucher à l'ancilla : c'est
le même mécanisme que le rebond de phase du CNOT (section~\ref{sec:cnot}), et la brique de tous les
algorithmes de ce chapitre.

\subsubsection{Algorithme de Deutsch}

Soit $f:\{0,1\}\to\{0,1\}$. Elle est \emph{constante} si $f(0)=f(1)$ et \emph{équilibrée} si
$f(0)\neq f(1)$. On veut savoir laquelle des deux est vraie. Classiquement, il faut évaluer $f$ deux
fois. Le circuit de la figure~\ref{fig:deutsch} répond avec \emph{un seul} appel à $U_f$.

\begin{figure}[!ht]
\centering
\begin{tikzpicture}[>=Stealth, font=\small]
  % fils : ancilla (second chiffre) en haut, requête (premier chiffre) en bas
  \draw (0,1.5) -- (8.4,1.5);
  \draw (0,0) -- (7.4,0);
  \node[anchor=east] at (0,1.5) {$\ket{1}$};
  \node[anchor=east] at (0,0) {$\ket{0}$};
  \node[porte] at (1.3,1.5) {$H$};
  \node[porte] at (1.3,0) {$H$};
  \node[draw=insarouge, fill=insarouge!8, thick, minimum width=13mm, minimum height=22mm]
    at (3.9,0.75) {$U_f$};
  \node[porte] at (6.0,0) {$H$};
  \pic at (7.2,0) {mesure};
  \node[anchor=west] at (7.7,0) {$f(0)\oplus f(1)$};
  \node[anchor=west, text=insagris, font=\footnotesize] at (8.5,1.5) {ancilla};
  \foreach \x/\lab in {0.55/0, 2.2/1, 5.15/2, 6.75/3}
    {\draw[dashed, insagris!70] (\x,2.1) -- (\x,-0.7);
     \node[anchor=north, font=\footnotesize] at (\x,-0.7) {$\Pi_\lab$};}
\end{tikzpicture}
\caption{Circuit de Deutsch. La requête $x$ est le premier chiffre du ket (fil du bas), l'ancilla
$y$ le second (fil du haut). La mesure du seul qubit de requête donne $f(0)\oplus f(1)$ : $0$ si $f$ est
constante, $1$ si elle est équilibrée.}
\label{fig:deutsch}
\end{figure}

\begin{align*}
  \ket{\Pi_0}&=\ket{0}\ket{1},\qquad
  \ket{\Pi_1}=(H\otimes H)\ket{\Pi_0}=\ket{+}\ket{-},\\[1ex]
  \ket{\Pi_2}&=U_f\ket{\Pi_1}
    =\frac{(-1)^{f(0)}\ket{0}+(-1)^{f(1)}\ket{1}}{\sqrt2}\ket{-}
    =(-1)^{f(0)}\,\frac{\ket{0}+(-1)^{f(0)\oplus f(1)}\ket{1}}{\sqrt2}\,\ket{-},\\[1ex]
  \ket{\Pi_3}&=(H\otimes I)\ket{\Pi_2}=(-1)^{f(0)}\,\ket{f(0)\oplus f(1)}\ket{-}.
\end{align*}
La deuxième ligne utilise le rebond de phase ; la troisième, $H\ket{+}=\ket0$ et $H\ket{-}=\ket1$.
Les deux amplitudes de la requête interfèrent : elles s'additionnent si $f$ est constante (état
$\ket{+}$, devenu $\ket{0}$) et se retranchent si $f$ est équilibrée (état $\ket{-}$, devenu $\ket{1}$).
Le calcul ne dit pas les valeurs $f(0)$ et $f(1)$ : il extrait seulement leur parité, et c'est ce qui
permet d'économiser un appel.

\begin{remarque}[Des oracles à $n$ qubits : Deutsch--Jozsa, Bernstein--Vazirani, Simon]
Le même schéma ($H^{\otimes n}$, oracle, $H^{\otimes n}$) généralise l'algorithme de \textcite{deutsch1985}. Pour $f$ sur $n$ bits promise
constante ou équilibrée, \textcite{deutschjozsa1992} décident avec \emph{un} appel là où un algorithme classique
déterministe en demande $2^{n-1}+1$. \textcite{bernstein1997} retrouvent en un appel un mot caché
$s$ défini par $f(x)=s\cdot x$ (classiquement : $n$ appels). \textcite{simon1997} trouve un masque caché avec $O(n)$
appels, contre $2^{\Omega(n)}$ classiquement : c'est la première séparation exponentielle, et l'idée
de l'algorithme de Shor. Ces problèmes sont artificiels (ils portent sur une boîte noire), mais
ils prouvent que l'interférence offre un gain réel.
\end{remarque}

% ---------------------------------------------------------------------
\subsection{Recherche dans une liste : l'algorithme de Grover}\label{sec:grover}
% ---------------------------------------------------------------------

Un oracle $f$ vaut $1$ en un unique élément $w$ parmi $N=2^n$, et $0$ ailleurs ; on cherche $w$.
Classiquement, il faut en moyenne $N/2$ évaluations de $f$, et $N$ dans le pire cas. \textcite{grover1996}
le trouve avec environ $\frac\pi4\sqrt N$ appels.

\begin{resultat}{Algorithme de Grover}
On part de $\ket{s}=H^{\otimes n}\ket{0\ldots0}=\frac{1}{\sqrt N}\sum_x\ket{x}$, puis on répète $k$
fois l'itération $G=D\,O_w$ avec
\[
  O_w=I-2\ket{w}\bra{w}\ \ (\text{oracle, par rebond de phase}),\qquad
  D=2\ket{s}\bra{s}-I=H^{\otimes n}\,(2\ket{0}\bra{0}-I)\,H^{\otimes n}.
\]
La mesure donne $w$ avec la probabilité $\sin^2\!\big((2k+1)\theta\big)$, où $\sin\theta=\frac{1}{\sqrt N}$.
Elle est proche de $1$ pour $k\simeq\frac\pi4\sqrt N$. Aucun algorithme quantique ne fait mieux que
$\Omega(\sqrt N)$ appels : le gain est quadratique, et optimal.
\end{resultat}

\paragraph{Pourquoi cela marche.} Posons $\ket{s'}=\frac{1}{\sqrt{N-1}}\sum_{x\neq w}\ket{x}$. Alors
$\ket{s}=\sin\theta\,\ket{w}+\cos\theta\,\ket{s'}$, et toute l'évolution reste dans le plan
engendré par $\ket{w}$ et $\ket{s'}$. L'oracle change le signe de la composante sur $\ket{w}$ : c'est
la symétrie par rapport à l'axe $\ket{s'}$. La diffusion $D=2\ket{s}\bra{s}-I$ est la symétrie par
rapport à $\ket{s}$. Deux symétries d'axes séparés de l'angle $\theta$ composent une rotation de
$2\theta$ (figure~\ref{fig:grover-geo}). Après $k$ itérations,
\[
  G^k\ket{s}=\sin\!\big((2k+1)\theta\big)\ket{w}+\cos\!\big((2k+1)\theta\big)\ket{s'}.
\]
On s'arrête quand $(2k+1)\theta\simeq\frac\pi2$. Contrairement à un algorithme classique, itérer plus
longtemps n'améliore rien : on continue de tourner et on dépasse $\ket{w}$.

\begin{figure}[!ht]
\centering
\begin{tikzpicture}[>=Stealth, font=\small]
  \draw[->, insagris] (-0.4,0) -- (4.5,0) node[right] {$\ket{s'}$};
  \draw[->, insagris] (0,-0.5) -- (0,4.5);
  \draw[insagris!50, dashed] (3.6,0) arc (0:90:3.6);
  \draw[dashed, insagris!60] (210:1.8) -- (30:3.6);
  \draw[->, very thick, insarouge] (0,0) -- (30:3.6) node[right] {$\ket{s}$};
  \draw[->, very thick, insagris] (0,0) -- (-30:3.6) node[right] {$O_w\ket{s}$};
  \draw[->, very thick, insarouge] (0,0) -- (90:3.6) node[above right] {$G\ket{s}=\ket{w}$};
  \draw[insagris] (1.1,0) arc (0:30:1.1);   \node[text=insagris] at (15:1.45) {$\theta$};
  \draw[insagris] (1.1,0) arc (0:-30:1.1);  \node[text=insagris] at (-15:1.45) {$\theta$};
  \draw[->, insarouge, thick] (30:2.2) arc (30:90:2.2);  \node[text=insarouge] at (60:2.65) {$2\theta$};
  \node[text=insagris, font=\footnotesize, anchor=north west, align=left] at (-1.7,-1.2)
    {symétrie $O_w$ : $\ket{s}\mapsto O_w\ket{s}$\\ symétrie $D$ d'axe $\ket{s}$ (pointillés)};
\end{tikzpicture}
\caption{Une itération de Grover pour $N=4$ ($\theta=\frac\pi6$). L'oracle renvoie $\ket{s}$ en
$O_w\ket{s}$ (angle $-\theta$), la diffusion le renvoie en $G\ket{s}$ (angle $3\theta=\frac\pi2$) :
une seule itération amène exactement sur $\ket{w}$.}
\label{fig:grover-geo}
\end{figure}

\subsubsection{Exemple complet : $N=4$ (deux qubits)}

Prenons $w=11$. L'oracle est la porte $Z$ contrôlée $\mathrm{CZ}=\operatorname{diag}(1,1,1,-1)$
(section~\ref{sec:controlees}). La diffusion s'écrit avec les portes du cours, car
$2\ket{00}\bra{00}-I=\operatorname{diag}(1,-1,-1,-1)=\mathrm{CZ}\,(Z\otimes Z)$ :
\[
  D=(H\otimes H)\;\mathrm{CZ}\,(Z\otimes Z)\;(H\otimes H).
\]
Le circuit de la figure~\ref{fig:grover-circ} s'écrit donc $H\otimes H$, $\mathrm{CZ}$,
$H\otimes H$, $Z\otimes Z$, $\mathrm{CZ}$, $H\otimes H$. Calcul pas à pas :
\begin{align*}
  \ket{s}&=\tfrac12\big(\ket{00}+\ket{01}+\ket{10}+\ket{11}\big),\\
  O_w\ket{s}&=\tfrac12\big(\ket{00}+\ket{01}+\ket{10}-\ket{11}\big),\\
  \braket{s}{O_w s}&=\tfrac14(1+1+1-1)=\tfrac12,\\
  D\,O_w\ket{s}&=2\braket{s}{O_w s}\ket{s}-O_w\ket{s}=\ket{s}-O_w\ket{s}
    =\tfrac12\big(0,0,0,2\big)^{\mathsf T}=\ket{11}.
\end{align*}
La mesure donne $11$ \emph{avec certitude} après un seul appel à l'oracle, là où un algorithme
classique en demande en moyenne $2{,}25$. C'est cohérent avec la formule : $\sin\theta=\frac12$,
$\theta=\frac\pi6$ et $\sin(3\theta)=1$. Pour un autre élément marqué, on change l'oracle (un
CZ entouré de portes $X$) et le résultat est le même.

\begin{figure}[!ht]
\centering
\adjustbox{max width=\linewidth}{%
\begin{tikzpicture}[>=Stealth, font=\small]
  \draw (0,1.4) -- (10.3,1.4);
  \draw (0,0) -- (10.3,0);
  \node[anchor=east] at (0,1.4) {$\ket{0}$};
  \node[anchor=east] at (0,0) {$\ket{0}$};
  \foreach \y in {0,1.4} {
    \node[porte] at (1.1,\y) {$H$};
    \node[porte] at (4.3,\y) {$H$};
    \node[porte] at (5.7,\y) {$Z$};
    \node[porte] at (8.5,\y) {$H$};
    \pic at (9.7,\y) {mesure};
  }
  \foreach \x in {2.7,7.1} {
    \draw[insarouge, thick] (\x,1.4) -- (\x,0);
    \pic at (\x,1.4) {ctl};
    \pic at (\x,0) {ctl};
  }
  \foreach \x/\lab in {1.9/{\ket{s}}, 3.5/{O_w\ket{s}}, 9.1/{\ket{11}}} {
    \draw[dashed, insagris!70] (\x,2.0) -- (\x,-0.6);
    \node[anchor=south, font=\footnotesize] at (\x,2.0) {$\lab$};
  }
  \draw[decorate, decoration={brace, amplitude=4pt, mirror}, insagris] (2.3,-0.8) -- (3.1,-0.8)
    node[midway, below=5pt, font=\footnotesize] {oracle $O_w$};
  \draw[decorate, decoration={brace, amplitude=4pt, mirror}, insagris] (3.9,-0.8) -- (9.0,-0.8)
    node[midway, below=5pt, font=\footnotesize] {diffusion $D=(H\otimes H)\,\mathrm{CZ}\,(Z\otimes Z)\,(H\otimes H)$};
\end{tikzpicture}}
\caption{Algorithme de Grover pour $N=4$ et $w=11$ : une itération suffit. Les deux points
reliés par un trait vertical forment la porte CZ.}
\label{fig:grover-circ}
\end{figure}

\paragraph{Ordres de grandeur.} Pour $N=2^{20}\approx10^6$, Grover demande
$\frac\pi4\cdot2^{10}\approx804$ itérations, contre $5\times10^5$ évaluations en moyenne. Pour une
recherche de clé AES-128 ($N=2^{128}$), il faut $\frac\pi4\,2^{64}\approx1{,}4\times10^{19}$
itérations \emph{séquentielles}, dont chacune contient un circuit AES complet ; et partager le travail
entre $m$ machines ne fait gagner qu'un facteur $\sqrt m$. L'effet est tout de même de diviser par deux le
nombre de bits de sécurité d'un chiffrement symétrique : on se protège en doublant la taille des
clés (AES-256).

% ---------------------------------------------------------------------
\subsection{Transformée de Fourier quantique et algorithme de Shor}\label{sec:shor}
% ---------------------------------------------------------------------

Le chiffrement RSA repose sur la difficulté de factoriser un entier $N$ (section~\ref{sec:crypto}).
L'algorithme de Shor, conçu en 1994 et publié en 1997 \parencite{shor1997}, factorise $N$ en un temps polynomial en $\log N$.
L'algorithme a deux étages : une réduction \emph{classique} de la factorisation à la recherche d'une
période, puis une recherche \emph{quantique} de cette période par transformée de Fourier.

\subsubsection{Réduction : de la factorisation à la période}

Soit $N$ impair et non puissance d'un nombre premier. On tire $a$ avec $1<a<N$. Si
$\operatorname{pgcd}(a,N)>1$, on tient déjà un facteur. Sinon, la suite $x\mapsto a^x\bmod N$ est
périodique : sa période $r$ est le plus petit entier $r\geq1$ tel que $a^r\equiv1\pmod N$ (l'\emph{ordre}
de $a$ modulo $N$).

\begin{resultat}{Factorisation à partir de l'ordre}
Si $r$ est pair et si $a^{r/2}\not\equiv-1\pmod N$, alors $\operatorname{pgcd}(a^{r/2}-1,N)$ et
$\operatorname{pgcd}(a^{r/2}+1,N)$ sont deux facteurs non triviaux de $N$. Pour $a$ tiré au hasard,
ces deux conditions sont réunies avec une probabilité au moins égale à $\frac12$ \parencite[chap.~5]{nielsen2010}.
\end{resultat}

En effet, posons $u=a^{r/2}$. Alors $(u-1)(u+1)=a^r-1\equiv0\pmod N$ : $N$ divise le produit. Il ne
divise pas $u-1$, sinon l'ordre de $a$ serait $r/2$, ni $u+1$ par hypothèse. Si $\operatorname{pgcd}(u-1,N)$
valait $1$, $N$ diviserait $u+1$ (lemme de Gauss) ; s'il valait $N$, $N$ diviserait $u-1$. Les deux sont exclus, donc
le pgcd est un diviseur strict de $N$ ; de même pour $u+1$. Tout se ramène donc à \emph{trouver $r$}. Classiquement, c'est
aussi difficile que factoriser. C'est la seule étape où l'ordinateur quantique intervient.

\subsubsection{La transformée de Fourier quantique}

\begin{definition}{Transformée de Fourier quantique (QFT)}
Sur l'espace de dimension $M$ de base $\ket{0},\ldots,\ket{M-1}$, avec $\omega=\ee^{2\ii\pi/M}$ :
\[
  \mathrm{QFT}_M\ket{x}=\frac{1}{\sqrt M}\sum_{y=0}^{M-1}\omega^{xy}\ket{y}.
\]
C'est une porte unitaire : ses colonnes sont orthonormées (somme géométrique de raison $\omega^{x-x'}$).
Pour $M=2^m$, elle s'implémente avec $O(m^2)$ portes élémentaires, alors que la transformée de Fourier
rapide classique demande $O(m\,2^m)$ opérations sur les $2^m$ amplitudes.
\end{definition}

Pour $M=2$ on a $\omega=-1$ et $\mathrm{QFT}_2=H$. Pour $M=4$, $\omega=\ii$ ; en écrivant $x=2x_1+x_2$ et
$y=2y_1+y_2$, on a $\ii^{xy}=(-1)^{x_1y_2+x_2y_1}\,\ii^{x_2y_2}$ (car $\ii^4=1$), d'où
\[
  \mathrm{QFT}_4\ket{x_1x_2}=\underbrace{\frac{\ket0+(-1)^{x_2}\ket1}{\sqrt2}}_{=\,H\ket{x_2}\ \text{(chiffre }y_1)}
  \otimes\ \frac{\ket0+(-1)^{x_1}\ii^{x_2}\ket1}{\sqrt2}\ \ \text{(chiffre }y_2).
\]
Le circuit $\mathrm{QFT}_4=\mathrm{SWAP}\,(I\otimes H)\,\mathrm{CS}\,(H\otimes I)$ de la
figure~\ref{fig:qft4} produit exactement cet état : $H$ sur le premier chiffre donne
$\frac{1}{\sqrt2}(\ket0+(-1)^{x_1}\ket1)$, la porte $\mathrm{CS}=\operatorname{diag}(1,1,1,\ii)$ ajoute le
facteur $\ii^{x_2}$ devant $\ket1$, $I\otimes H$ transforme $\ket{x_2}$ en $H\ket{x_2}$, et le SWAP remet les
chiffres dans l'ordre.

\begin{figure}[!ht]
\centering
\begin{tikzpicture}[>=Stealth, font=\small]
  \draw (0,1.2) -- (7.0,1.2);
  \draw (0,0) -- (7.0,0);
  \node[anchor=east] at (0,1.2) {$\ket{x_2}$};
  \node[anchor=east] at (0,0) {$\ket{x_1}$};
  \node[porte] at (1.2,0) {$H$};
  \draw[insarouge, thick] (2.7,1.2) -- (2.7,0);
  \pic at (2.7,1.2) {ctl};
  \node[porte] at (2.7,0) {$S$};
  \node[porte] at (4.2,1.2) {$H$};
  \draw[insarouge, thick] (5.7,1.2) -- (5.7,0);
  \pic at (5.7,1.2) {croix};
  \pic at (5.7,0) {croix};
  \node[anchor=west] at (7.0,1.2) {$\ket{y_2}$};
  \node[anchor=west] at (7.0,0) {$\ket{y_1}$};
\end{tikzpicture}
\caption{$\mathrm{QFT}_4$ : $H$, $S$ contrôlée par l'autre qubit, $H$, puis échange des deux fils.
Le fil du haut est le dernier chiffre.}
\label{fig:qft4}
\end{figure}

\begin{resultat}{La QFT transforme une période en pics}
Si $r$ divise $M$ et si $\ket{\psi}=\sqrt{r/M}\sum_{k=0}^{M/r-1}\ket{x_0+kr}$ (état périodique, de période
$r$), alors
\[
  \mathrm{QFT}_M\ket{\psi}=\frac{1}{\sqrt r}\sum_{j=0}^{r-1}\omega^{x_0\,jM/r}\,\ket{jM/r}.
\]
La mesure donne $y=jM/r$ avec la probabilité $\frac1r$, pour $j=0,\ldots,r-1$, quel que soit le décalage $x_0$.
\end{resultat}

Preuve : l'amplitude de $\ket{y}$ est $\frac{1}{\sqrt M}\sqrt{\frac rM}\,\omega^{x_0y}\sum_{k}\omega^{kry}$.
La somme géométrique de raison $\omega^{ry}=\ee^{2\ii\pi ry/M}$ vaut $M/r$ si $ry/M$ est entier, et $0$ sinon
(la raison est alors différente de $1$ et sa puissance $M/r$ vaut $1$). L'amplitude vaut donc
$\sqrt{r}/M\cdot M/r=1/\sqrt r$ aux multiples de $M/r$, et $0$ ailleurs. Le décalage $x_0$ n'apparaît que dans une
\emph{phase}, qui disparaît à la mesure : c'est ce qui permet de lire la période sans connaître le point de départ.

\subsubsection{L'algorithme}

\begin{resultat}{Algorithme de Shor : recherche de l'ordre de $a$ modulo $N$}
On choisit $M=2^t\geq N^2$ (soit $t\approx2n$ avec $n=\lceil\log_2N\rceil$), un registre de contrôle de $t$ qubits et un
registre de travail de $n$ qubits.
\begin{enumerate}[nosep]
  \item Préparer $\ket{0}^{\otimes t}\ket{1}$ et appliquer $H^{\otimes t}$ au registre de contrôle :
        $\ket{\varphi_1}=\frac{1}{\sqrt M}\sum_{x=0}^{M-1}\ket{x}\ket{1}$.
  \item Appliquer l'exponentiation modulaire $U_a\ket{x}\ket{z}=\ket{x}\ket{a^xz\bmod N}$ :
        $\ket{\varphi_2}=\frac{1}{\sqrt M}\sum_{x}\ket{x}\ket{a^x\bmod N}$.
  \item Appliquer $\mathrm{QFT}_M$ au registre de contrôle et le mesurer : on obtient $y$.
  \item (Classique) Développer $y/M$ en fraction continue : on y lit une fraction $j/r'$ proche de
        $y/M$ avec $r'<N$. Si $\operatorname{pgcd}(j,r)=1$, alors $r'=r$ et l'on vérifie $a^{r'}\equiv1\pmod N$.
        Sinon, $r'$ est un diviseur strict de $r$ et $a^{r'}\not\equiv1$ : on recommence.
\end{enumerate}
\end{resultat}

L'oracle $U_a$ est un produit de $t$ multiplications modulaires \emph{contrôlées} par les qubits $x_j$ de $x$ :
comme $a^x=\prod_j(a^{2^j})^{x_j}$, la $j$-ème porte multiplie le registre de travail par $a^{2^j}\bmod N$, un
nombre calculé à l'avance par un ordinateur classique. Le circuit est celui de la
figure~\ref{fig:shor-circ}.

\begin{figure}[!ht]
\centering
\begin{tikzpicture}[>=Stealth, font=\small]
  % registre de contrôle (haut) et registre de travail (bas)
  \draw (0,1.7) -- (8.3,1.7);
  \draw (0,0) -- (4.6,0);
  \draw[dashed] (4.6,0) -- (8.3,0);
  \node[anchor=east] at (0,1.7) {$\ket{0}^{\otimes t}$};
  \node[anchor=east] at (0,0) {$\ket{1}$};
  \draw (0.5,1.6) -- (0.7,1.8); \node[font=\scriptsize, anchor=south] at (0.6,1.78) {$t$};
  \draw (0.5,-0.1) -- (0.7,0.1); \node[font=\scriptsize, anchor=south] at (0.6,0.08) {$n$};
  \node[porte] at (1.5,1.7) {$H^{\otimes t}$};
  \pic at (3.3,1.7) {ctl};
  \draw[insarouge, thick] (3.3,1.7) -- (3.3,0.4);
  \node[draw=insarouge, fill=insarouge!8, thick, minimum width=17mm, minimum height=8mm] at (3.3,0) {$U_a$};
  \node[porte] at (5.4,1.7) {$\mathrm{QFT}_M$};
  \pic at (7.0,1.7) {mesure};
  \node[anchor=west] at (7.5,1.7) {$y$};
  \node[anchor=north west, text=insagris, font=\footnotesize] at (4.7,-0.3) {registre de travail : ignoré};
  \foreach \x/\lab in {2.3/1, 4.3/2, 6.3/3}
    {\draw[dashed, insagris!70] (\x,2.35) -- (\x,-0.3);
     \node[anchor=south, font=\footnotesize] at (\x,2.35) {$\ket{\varphi_\lab}$};}
\end{tikzpicture}
\caption{Circuit de Shor. Les petits traits obliques indiquent un registre de plusieurs qubits ($t$ pour le
contrôle, $n$ pour le travail). $U_a\ket{x}\ket{z}=\ket{x}\ket{a^xz\bmod N}$ est commandé par le registre de
contrôle.}
\label{fig:shor-circ}
\end{figure}

\paragraph{Pourquoi cela marche.} Après l'étape~2, mesurons (en pensée) le registre de travail : il
donne une valeur $f$ de la suite $a^x\bmod N$, et le registre de contrôle se projette sur les $x$ tels que
$a^x\equiv f$, c'est-à-dire $x\in\{x_0,x_0+r,x_0+2r,\ldots\}$ : un état périodique de période $r$. Le résultat sur la QFT
ci-dessus s'applique : la QFT concentre toute la probabilité sur les $y$ voisins de $jM/r$. La mesure du
registre de travail n'est pas nécessaire (les probabilités de $y$ sont les mêmes si on l'ignore), et le
décalage $x_0$ dépend de $f$ mais ne change pas la loi de $y$. Ainsi $y/M\approx j/r$ pour un entier $j$ \emph{aléatoire} :
on ne choisit pas $j$, mais il est premier avec $r$ dans une fraction $\Omega(1/\log\log r)$ des cas, et quelques
répétitions suffisent. Le choix $M\geq N^2$ garantit qu'une seule fraction de dénominateur inférieur à $N$ est
à distance de $\frac{1}{2M}$ de $y/M$ : le développement en fraction continue l'identifie.

\subsubsection{Exemple complet : $N=15$, $a=7$}

\paragraph{Période.} La suite $7^x\bmod15$ vaut $1,7,4,13,1,7,4,13,\ldots$ : l'ordre est $r=4$. Les
multiplicateurs de $U_a$ sont $a^{2^j}\bmod 15=7,\ 4,\ 1,\ 1$ pour $j=0,1,2,3$ (les deux derniers sont triviaux).
On prend ici $t=4$, soit $M=16$. Ce choix est trop court en général ($M\geq N^2=225$ serait exigé) mais suffit
car $r=4$ divise $M$, ce qui rend le calcul exact.

\paragraph{État avant la QFT.} $\ket{\varphi_2}=\frac14\sum_{x=0}^{15}\ket{x}\ket{7^x\bmod15}$. Si le
registre de travail est mesuré et donne $1$, le registre de contrôle est dans $\frac12(\ket0+\ket4+\ket8+\ket{12})$ ;
s'il donne $7$, dans $\frac12(\ket1+\ket5+\ket9+\ket{13})$, etc.

\paragraph{QFT.} On a $\omega=\ee^{2\ii\pi/16}$ et $\omega^4=\ii$. Pour le décalage $x_0=0$ :
\[
  \mathrm{QFT}_{16}\,\tfrac12\big(\ket0+\ket4+\ket8+\ket{12}\big)
  =\frac18\sum_{y=0}^{15}\Big(\sum_{k=0}^{3}\ii^{\,ky}\Big)\ket{y}
  =\frac12\big(\ket0+\ket4+\ket8+\ket{12}\big),
\]
car $\sum_{k=0}^{3}\ii^{\,ky}$ vaut $4$ si $y$ est multiple de $4$ et $0$ sinon. Pour les autres valeurs de
$f$, seules les phases changent. Dans tous les cas la mesure de $y$ donne $0,4,8$ ou $12$, chacun avec la probabilité $\frac14$
(figure~\ref{fig:shor-n15}).

\paragraph{Lecture classique.} On calcule $y/16$ :
\begin{center}
\renewcommand{\arraystretch}{1.2}
\begin{tabular}{@{}cccp{7.3cm}@{}}
\toprule
$y$ & $y/16$ & fraction irréductible & conclusion\\
\midrule
$0$ & $0$ & $0/1$ & inutile ($j=0$) : on recommence\\
$4$ & $\frac14$ & $1/4$ & $r=4$\ ; $7^4=2401\equiv1$ : correct\\
$8$ & $\frac12$ & $1/2$ & dénominateur $2$ : $7^2=49\equiv4\not\equiv1$, rejeté, on recommence\\
$12$ & $\frac34$ & $3/4$ & $r=4$ : correct\\
\bottomrule
\end{tabular}
\end{center}
Un essai réussit donc avec la probabilité $\frac12$ ($y=4$ ou $12$). Avec $r=4$ pair et $7^{2}=49\equiv4\not\equiv-1\pmod{15}$ :
\[
  \operatorname{pgcd}(7^2-1,15)=\operatorname{pgcd}(48,15)=3,\qquad
  \operatorname{pgcd}(7^2+1,15)=\operatorname{pgcd}(50,15)=5,\qquad 15=3\times5.
\]

\begin{figure}[!ht]
\centering
\adjustbox{max width=\linewidth}{%
\begin{tikzpicture}[>=Stealth, font=\small]
  % (a) suite 7^x mod 15, x = 0..15
  \draw[->, insagris] (-0.2,0) -- (6.7,0) node[right] {$x$};
  \draw[->, insagris] (0,-0.1) -- (0,2.6);
  \foreach \x/\v in {0/1,1/7,2/4,3/13,4/1,5/7,6/4,7/13,8/1,9/7,10/4,11/13,12/1,13/7,14/4,15/13}
    \fill[insarouge] ({\x*0.42},{\v*0.17}) circle (1.7pt);
  \foreach \v in {1,4,7,13}
    {\draw[insagris] (-0.05,{\v*0.17}) -- (0.05,{\v*0.17});
     \node[anchor=east, font=\scriptsize] at (-0.08,{\v*0.17}) {$\v$};}
  \foreach \x in {0,4,8,12}
    {\draw[insagris] ({\x*0.42},-0.05) -- ({\x*0.42},0.05);
     \node[anchor=north, font=\scriptsize] at ({\x*0.42},-0.07) {$\x$};}
  \draw[<->, insarouge] (0,2.45) -- (1.68,2.45) node[midway, above, font=\footnotesize] {période $r=4$};
  \node[font=\footnotesize, anchor=south] at (3.2,3.15) {(a) $f(x)=7^x\bmod15$};
  % flèche QFT
  \draw[->, very thick, insarouge] (7.0,0.8) -- (8.1,0.8) node[midway, above, font=\footnotesize] {$\mathrm{QFT}_{16}$};
  % (b) loi de y : P(y)=1/4 pour y = 0, 4, 8, 12
  \draw[->, insagris] (8.5,0) -- (14.9,0) node[right] {$y$};
  \draw[->, insagris] (8.6,-0.1) -- (8.6,2.6);
  \node[font=\footnotesize, anchor=south] at (11.7,3.15) {(b) loi de la mesure $y$ : $P(y)$};
  \foreach \y in {0,4,8,12}
    {\fill[insarouge] ({8.6+\y*0.4-0.1},0) rectangle ({8.6+\y*0.4+0.1},1.5);
     \node[anchor=north, font=\scriptsize] at ({8.6+\y*0.4},-0.07) {$\y$};}
  \draw[dashed, insagris] (8.55,1.5) -- (14.2,1.5);
  \node[anchor=east, font=\scriptsize] at (8.5,1.5) {$\frac14$};
  \node[text=insagris, font=\footnotesize, align=center] at (11.7,2.35) {$y/16=\frac{j}{4}$, $j=0,1,2,3$};
\end{tikzpicture}}
\caption{Shor pour $N=15$, $a=7$, $t=4$. À gauche, la suite $7^x\bmod15$ a la période $4$. À droite,
la mesure du registre de contrôle après la QFT ne donne que $y\in\{0,4,8,12\}$, avec la probabilité
$\frac14$ chacun : la période $r=4$ se lit dans l'écart $M/r=4$ entre les pics.}
\label{fig:shor-n15}
\end{figure}

\subsubsection{Ce que coûte l'algorithme}

Notons $n=\lceil\log_2N\rceil$ le nombre de bits de $N$ ($n=2048$ pour un module RSA-2048). Le coût s'exprime en
fonction de $n$, pas de $N$.
\begin{align*}
  \text{meilleur algorithme classique connu (crible algébrique, GNFS) :}&\quad
    \exp\!\Big(O\big(n^{1/3}(\log n)^{2/3}\big)\Big),\\
  \text{algorithme de Shor :}&\quad O(n^3)\ \text{portes élémentaires.}
\end{align*}
Le premier temps est sous-exponentiel en $n$ \parencite{lenstra1993}, mais supérieur à tout polynôme. Le second est un
polynôme en $n$ avec la multiplication scolaire ($O(n^2\log n\log\log n)$ avec une multiplication rapide).

\begin{remarque}[Où est le « log » ?]
On résume parfois « de $O(N^p)$ à $O(\log N)$ ». La comparaison correcte est la suivante. La division
d'essai coûte $O(\sqrt N)=O(2^{n/2})$ : exponentielle en $n$. Le crible algébrique est sous-exponentiel en $n$. Shor est
polynomial en $\log N=n\ln2$, soit $O((\log N)^3)$ et non $O(\log N)$ : rien que pour lire $N$ et multiplier
modulo $N$, il faut au moins $n$ opérations. Le « $\log$ » désigne donc la \emph{taille} du nombre.
\end{remarque}

Doubler la taille de la clé multiplie le coût de Shor par $8$, alors qu'il augmente de façon
spectaculaire celui de l'attaque classique. Le tableau compare le coût $n^3$ de Shor à la sécurité
classique que le NIST attribue aux modules RSA \parencite{nist80057}.

\begin{center}
\renewcommand{\arraystretch}{1.2}
\begin{tabular}{@{}rccc@{}}
\toprule
Module RSA $n$ (bits) & sécurité classique (NIST) & $n^3$ (portes, ordre de grandeur) & $\log_2(n^3)$\\
\midrule
$1024$ & $2^{80}$ & $1{,}1\times10^{9}$ & $30{,}0$\\
$2048$ & $2^{112}$ & $8{,}6\times10^{9}$ & $33{,}0$\\
$3072$ & $2^{128}$ & $2{,}9\times10^{10}$ & $34{,}8$\\
$7680$ & $2^{192}$ & $4{,}5\times10^{11}$ & $38{,}7$\\
$15360$ & $2^{256}$ & $3{,}6\times10^{12}$ & $41{,}7$\\
\bottomrule
\end{tabular}
\end{center}

Passer de $2048$ à $3072$ bits fait gagner un facteur $2^{16}=65\,536$ sur l'attaque classique, mais le coût de
Shor n'est multiplié que par $(3/2)^3\approx3{,}4$. Face à un ordinateur quantique, \emph{grossir
les clés RSA ne protège pas} : il faut changer d'algorithme (section~\ref{sec:crypto}). Cela dit, le nombre de portes
n'est pas le vrai coût : il faut des qubits \emph{corrigés d'erreurs}, et c'est là que se trouve le verrou technique
(section~\ref{sec:machine}). Les estimations les plus citées sont les suivantes.
\begin{itemize}[nosep]
  \item \textbf{RSA-2048} : environ $3n+0{,}002\,n\log_2n\approx6\,189$ qubits logiques et
        $0{,}3\,n^3+0{,}0005\,n^3\log_2n\approx2{,}6\times10^9$ portes de Toffoli ; avec des qubits physiques
        d'erreur $10^{-3}$ et un cycle de code de surface de $1\,\mu$s, $20$ millions de qubits physiques pendant
        $8$ heures \parencite{gidney2021}. Une optimisation de 2025 (arithmétique approchée, codes « yoked »,
        « magic state cultivation ») descend à \emph{moins d'un million} de qubits physiques, pour moins d'une
        semaine \parencite{gidney2025}.
  \item \textbf{Courbes elliptiques (ECC-256)}, utilisées dans TLS : au plus $2\,330$ qubits logiques
        \parencite{roetteler2017}. Avec des qubits de chat (section~\ref{sec:machine}), $126\,133$ qubits suffiraient
        pour $9$ heures de calcul \parencite{gouzien2023}. Une \emph{prépublication} de 2026, fondée sur des
        hypothèses optimistes pour les atomes neutres, évoque $10\,000$ à $26\,000$ atomes et quelques
        jours \parencite{cain2026}.
\end{itemize}
Ce sont des estimations, pas des mesures : les démonstrations expérimentales de $15=3\times5$ ou
$21=3\times7$ reposent sur des circuits simplifiés grâce à la connaissance préalable de la réponse
\parencite{smolin2013}. À notre connaissance, aucune n'a factorisé par l'algorithme de Shor un nombre hors de portée d'un
calcul à la main.

% ---------------------------------------------------------------------
\subsection{Cryptographie : RSA, TLS et « stocker maintenant, déchiffrer plus tard »}\label{sec:crypto}
% ---------------------------------------------------------------------

La sécurité des échanges sur Internet repose sur deux familles d'algorithmes. La cryptographie
\emph{symétrique} (AES, ChaCha20) chiffre les données avec une clé secrète partagée. La cryptographie
\emph{asymétrique} (RSA, Diffie--Hellman, courbes elliptiques) sert à \emph{établir} cette clé entre deux
personnes qui ne se sont jamais rencontrées, et à \emph{authentifier} l'interlocuteur par une signature. Les
algorithmes asymétriques s'appuient sur deux problèmes réputés difficiles : la factorisation d'un entier
(RSA) et le logarithme discret (Diffie--Hellman, courbes elliptiques). Ce sont deux problèmes de
\emph{période} dans un groupe : l'algorithme de Shor les résout tous les deux.

\subsubsection{RSA et ce que Shor en fait}

\begin{definition}{Chiffrement RSA}
Alice choisit deux nombres premiers $p\neq q$, pose $n=pq$ et $\varphi(n)=(p-1)(q-1)$, choisit $e$ premier avec
$\varphi(n)$ et calcule $d$ tel que $ed\equiv1\pmod{\varphi(n)}$. La \emph{clé publique} est $(n,e)$, la \emph{clé
privée} $d$. Pour un message $m<n$ :
\[
  c=m^e\bmod n\quad(\text{chiffrement}),\qquad m=c^d\bmod n\quad(\text{déchiffrement}).
\]
Le déchiffrement fonctionne parce que $m^{ed}\equiv m\pmod n$ (théorème d'Euler) \parencite{rivest1978}.
\end{definition}

\paragraph{Exemple jouet.} $p=3$, $q=11$, $n=33$, $\varphi(n)=20$, $e=3$ (premier avec $20$) et $d=7$ (car
$3\times7=21\equiv1\pmod{20}$). Pour $m=4$ : $c=4^3\bmod33=64-33=31$. Le déchiffrement donne
$31^7\bmod33$ ; comme $31\equiv-2$, on a $(-2)^7=-128\equiv-128+132=4=m$ : on retrouve le message.

\paragraph{Attaque.} L'attaquant connaît $(n,e)$. Il lui manque $\varphi(n)$ : le calculer équivaut à
factoriser $n$. Avec $n=33=3\times11$, il obtient $\varphi=20$ puis $d=e^{-1}\bmod20=7$ par l'algorithme d'Euclide
étendu (polynomial). Tout le secret de RSA est donc la difficulté de factoriser. Pour un module de $2048$ bits
(617 chiffres décimaux), les meilleures factorisations classiques publiques (RSA-240, $795$ bits, en 2019
\parencite{boudot2020}) sont loin du compte : on a vu à la section~\ref{sec:shor} que
l'algorithme de Shor l'attaque en $O(n^3)$ portes. Le même algorithme, appliqué au groupe d'une courbe elliptique,
résout le logarithme discret et casse Diffie--Hellman, ECDH et ECDSA.

\subsubsection{Ce que cassent Shor et Grover}

\begin{center}
\renewcommand{\arraystretch}{1.25}
\begin{tabular}{@{}>{\raggedright\arraybackslash}p{4.1cm}>{\raggedright\arraybackslash}p{2.9cm}>{\raggedright\arraybackslash}p{2.5cm}>{\raggedright\arraybackslash}p{5.2cm}@{}}
\toprule
\textbf{Primitive} & \textbf{Rôle} & \textbf{Attaque quantique} & \textbf{Conséquence}\\
\midrule
RSA, Diffie--Hellman, ECDH (X25519) & échange de clé, chiffrement & Shor & \textbf{cassés} (temps polynomial)\\
RSA, ECDSA, EdDSA & signature, certificats & Shor & \textbf{cassés}\\
AES, ChaCha20 & chiffrement symétrique & Grover & clé de $k$ bits : $2^{k/2}$ itérations ; AES-128 affaibli, \textbf{AES-256 sûr}\\
SHA-2, SHA-3 & hachage & Grover (pré-image) & SHA-256 : $2^{128}$, \textbf{sûr}\\
ML-KEM, ML-DSA, SLH-DSA & échange de clé, signature & aucune attaque efficace connue & normes NIST de 2024 (voir plus bas)\\
\bottomrule
\end{tabular}
\end{center}

Les trois premières lignes ne sont pas à égalité. Un algorithme asymétrique est \emph{détruit} par Shor : il faut
le remplacer. Un algorithme symétrique n'est que \emph{ralenti} par Grover (racine carrée, section~\ref{sec:grover}) :
il suffit de doubler la taille de la clé. « Aucune attaque connue » n'est pas une preuve de sécurité, d'où la prudence
des agences décrite plus bas.

\subsubsection{TLS 1.3 : où sont les points faibles ?}

Le protocole TLS 1.3 (RFC~8446 \parencite{rfc8446}) protège HTTPS. Il combine les trois mécanismes : un
\emph{échange de clés} éphémère (ECDHE, le plus souvent X25519 : chaque côté tire un secret $a$ ou $b$ et
publie $g^a$ ou $g^b$, d'où le secret partagé $g^{ab}$), une \emph{authentification} du serveur par signature
(certificat RSA ou ECDSA, vérifié par une chaîne d'autorités de certification, elle aussi faite de signatures),
et un \emph{chiffrement symétrique} des données (AES-GCM ou ChaCha20-Poly1305) avec des clés dérivées de $g^{ab}$
(figure~\ref{fig:tls}).

\begin{figure}[!ht]
\centering
\adjustbox{max width=\linewidth}{%
\begin{tikzpicture}[>=Stealth, font=\small,
    tag/.style={rounded corners=2pt, font=\footnotesize, inner sep=3pt, anchor=west, text width=4.6cm, align=left}]
  \node[draw=insarouge, fill=insarouge!8, thick, minimum width=20mm] (C) at (0,0) {Client};
  \node[draw=insarouge, fill=insarouge!8, thick, minimum width=20mm] (S) at (6.2,0) {Serveur};
  \draw[insagris, thick] (0,-0.4) -- (0,-6.2);
  \draw[insagris, thick] (6.2,-0.4) -- (6.2,-6.2);
  \draw[->, thick] (0,-1.0) -- (6.2,-1.0) node[midway, above, font=\footnotesize] {ClientHello + clé publique $g^a$};
  \draw[<-, thick] (0,-1.9) -- (6.2,-1.9) node[midway, above, font=\footnotesize] {ServerHello + clé publique $g^b$};
  \node[font=\footnotesize, text=insagris, anchor=east] at (-0.2,-2.35) {calcule $g^{ab}$, dérive les clés};
  \node[font=\footnotesize, text=insagris, anchor=west] at (6.4,-2.35) {calcule $g^{ab}$};
  \draw[<-, thick] (0,-3.2) -- (6.2,-3.2) node[midway, above, font=\footnotesize, align=center]
       {\{Certificat, signature (CertificateVerify),\\ Finished\}};
  \draw[->, thick] (0,-4.1) -- (6.2,-4.1) node[midway, above, font=\footnotesize] {\{Finished\}};
  \draw[<->, very thick, insarouge] (0,-5.2) -- (6.2,-5.2) node[midway, above, font=\footnotesize] {données chiffrées (AES-GCM, ChaCha20)};
  \node[tag, fill=insarouge!18] at (7.8,-1.45) {\textbf{Échange de clés} : cassé par Shor (logarithme discret)};
  \node[tag, fill=insarouge!18] at (7.8,-3.2) {\textbf{Signature et certificats} : cassés par Shor (RSA, ECDSA)};
  \node[tag, fill=insagris!20] at (7.8,-5.2) {\textbf{Chiffrement symétrique} : \mbox{Grover} seulement, AES-256 suffit};
\end{tikzpicture}}
\caption{Poignée de main TLS 1.3 simplifiée. Les accolades \{\,\} indiquent des messages déjà chiffrés avec les clés
issues de $g^{ab}$. Deux des trois briques tombent face à Shor ; la troisième résiste.}
\label{fig:tls}
\end{figure}

\begin{remarque}[La confidentialité persistante ne protège pas contre l'ordinateur quantique]
Les secrets $a$ et $b$ sont détruits après la session : si la clé privée du serveur est volée plus tard, les sessions
passées ne peuvent pas être déchiffrées (propriété dite de \emph{forward secrecy}). Mais les valeurs publiques
$g^a$ et $g^b$ ont circulé en clair. Un adversaire qui les a enregistrées peut, avec l'algorithme de Shor (logarithme
discret sur la courbe), retrouver $a$ puis $g^{ab}$, donc toutes les clés de la session. La forward secrecy protège
contre le vol d'une \emph{clé}, pas contre l'effondrement de la difficulté du \emph{problème mathématique}.
\end{remarque}

\subsubsection{« Stocker maintenant, déchiffrer plus tard »}

La menace n'attend pas l'existence de la machine. Un adversaire qui capte le trafic chiffré aujourd'hui (il lui
suffit d'écouter une liaison, ce qui est discret et peu coûteux) peut le \emph{conserver} et attendre qu'un
ordinateur quantique assez puissant existe pour le déchiffrer. C'est l'attaque \emph{Harvest Now, Decrypt Later}
(HNDL, ou « stocker maintenant, déchiffrer plus tard »). Elle vise la \emph{confidentialité de long terme} :
dossiers médicaux et génétiques, secrets industriels, secrets d'État, archives juridiques. Elle ne vise pas un
jeton de session valable dix minutes. On ne sait pas qui enregistre quoi ; il faut supposer que cela se fait.

\begin{resultat}{Inégalité de Mosca}
Soit $X$ la durée pendant laquelle les données doivent rester secrètes, $Y$ le temps nécessaire pour migrer vers une
cryptographie résistante, et $Z$ le délai avant qu'un ordinateur quantique capable de casser la cryptographie actuelle existe
(tous en années \parencite{mosca2018}). Les données sont \emph{déjà exposées} si
\[
  X+Y>Z.
\]
\end{resultat}

Les données chiffrées juste avant la fin de la migration (à la date $Y$) doivent encore être secrètes à la date $X+Y$ ;
si $Z$ tombe avant, l'adversaire les lit. Exemple (figure~\ref{fig:mosca}) : des données médicales à protéger
$X=20$ ans, une migration de $Y=7$ ans, et $Z=15$ ans : $X+Y=27>15$, avec douze années d'exposition. On ne connaît pas
$Z$ (ce chapitre ne cherche pas à le prédire), mais on connaît $X$ et $Y$ : plus ils sont grands, plus il faut
commencer tôt, et le coût d'une migration trop précoce est faible devant celui d'une divulgation irréversible.

\begin{figure}[!ht]
\centering
\begin{tikzpicture}[>=Stealth, font=\small, x=0.3cm]
  \draw[->, insagris] (0,0) -- (31,0) node[right] {années};
  \foreach \t in {0,5,...,30}
    {\draw[insagris] (\t,0.08) -- (\t,-0.08);
     \node[anchor=north, font=\scriptsize] at (\t,-0.1) {$\t$};}
  \fill[insagris!35] (0,0.5) rectangle (7,1.1);
  \fill[insarouge!18] (7,0.5) rectangle (27,1.1);
  \fill[insarouge!60] (15,0.5) rectangle (27,1.1);
  \node[font=\footnotesize] at (3.5,0.8) {$Y=7$};
  \node[font=\footnotesize] at (11,0.8) {$X=20$};
  \draw[insarouge, very thick] (15,0) -- (15,2.0);
  \node[anchor=south, font=\footnotesize, text=insarouge] at (15,2.0) {$Z=15$ : ordinateur quantique opérationnel};
  \node[anchor=west, font=\footnotesize, text=white] at (15.4,0.8) {exposé 12 ans};
  \node[anchor=north west, font=\footnotesize, text=insagris] at (0,-0.65)
    {$Y$ : migration (la cryptographie vulnérable est encore utilisée)};
  \node[anchor=north west, font=\footnotesize, text=insagris] at (0,-1.15)
    {$X$ : durée de confidentialité exigée pour les données chiffrées à la fin de la migration};
\end{tikzpicture}
\caption{Inégalité de Mosca pour $X=20$, $Y=7$, $Z=15$ ans. La zone rouge foncé est la période où des données
interceptées sont lisibles par l'adversaire alors qu'elles devraient rester secrètes.}
\label{fig:mosca}
\end{figure}

\begin{remarque}[Signatures : un risque différent]
L'attaque HNDL ne concerne que la confidentialité. Une signature ne sert qu'au moment où on la vérifie : falsifier en 2035 un
certificat TLS expiré ne change rien à une session de 2026. Les signatures à longue durée de vie (micrologiciels,
racines de confiance, documents juridiques) devront néanmoins migrer, mais plus tard : c'est l'échange de clés qui est
urgent.
\end{remarque}

\subsubsection{Première réponse : la cryptographie post-quantique}

La \emph{cryptographie post-quantique} (PQC) regroupe des algorithmes qui s'exécutent sur des ordinateurs
\emph{classiques} mais reposent sur des problèmes contre lesquels on ne connaît aucune attaque quantique efficace :
réseaux euclidiens (problème \emph{Learning With Errors}), fonctions de hachage, codes correcteurs d'erreurs. Le NIST a
publié le 13 août 2024 trois normes \parencite{fips203,fips204,fips205}, puis sélectionné un quatrième algorithme en 2025 \parencite{nist2025hqc}.

\begin{center}
\renewcommand{\arraystretch}{1.2}
\begin{tabular}{@{}llll@{}}
\toprule
\textbf{Norme} & \textbf{Nom} & \textbf{Usage} & \textbf{Problème sous-jacent}\\
\midrule
FIPS 203 & ML-KEM (ex-Kyber) & échange de clé & réseaux (module-LWE)\\
FIPS 204 & ML-DSA (ex-Dilithium) & signature & réseaux\\
FIPS 205 & SLH-DSA (ex-SPHINCS+) & signature & fonctions de hachage\\
(à venir) & HQC & échange de clé (second choix) & codes correcteurs\\
\bottomrule
\end{tabular}
\end{center}

\paragraph{Hybridation.} Dans TLS 1.3, la solution déployée est \emph{hybride} : le client envoie à la fois une clé
publique X25519 et une clé publique ML-KEM-768, et les deux secrets sont concaténés dans la dérivation des clés. L'adversaire
doit casser \emph{les deux} : X25519 (que Shor casse, mais que l'on connaît bien) et ML-KEM (que Shor ne casse pas, mais
que l'on connaît depuis moins longtemps). Le groupe \texttt{X25519MLKEM768} est normalisé par la RFC~10024 d'août 2026
\parencite{rfc10024}, et il est activé par défaut dans les navigateurs depuis fin 2024 (Chrome 131, Firefox 132). Le
prix est la taille : $32$ octets pour une clé X25519, mais $1\,184$ octets de clé publique et $1\,088$ octets de
chiffré pour ML-KEM-768.

\paragraph{Pourquoi pas ML-KEM seul ?} Les algorithmes post-quantiques sont plus jeunes et moins étudiés. En 2022, le
candidat SIKE, arrivé au quatrième tour du NIST, a été cassé par une \emph{attaque classique} : une dizaine de minutes sur un seul
cœur d'ordinateur suffisent pour son plus petit jeu de paramètres \parencite{castryck2023}. L'ANSSI recommande donc de \emph{combiner} cryptographie classique et
post-quantique pendant la transition \parencite{anssi2024}.

\paragraph{Calendrier.} Le projet de document NIST~IR~8547 (version initiale publique de novembre 2024, qui peut
avoir évolué) prévoit de déprécier RSA-2048 et les algorithmes de sécurité équivalente après 2030, et d'interdire tous les
algorithmes vulnérables après 2035 \parencite{nistir8547}. La migration est longue : il faut inventorier où la cryptographie est
utilisée (logiciels, matériel, certificats, protocoles), puis la remplacer. L'« agilité cryptographique », c'est-à-dire
la capacité de changer d'algorithme sans refondre un système, est la qualité qui coûtera le moins cher.

\subsubsection{Seconde réponse : la distribution quantique de clé}

La \emph{distribution quantique de clé} (QKD) ne cherche pas un problème difficile : elle utilise la physique. Dans le
protocole BB84 \parencite{bennett1984}, Alice envoie à Bob des qubits dont l'état est choisi au hasard parmi quatre :
$\ket0,\ket1$ (base $Z$) ou $\ket+,\ket-$ (base $X$). Bob mesure chaque qubit dans une base $Z$ ou $X$ choisie au hasard.
Ensuite, Alice et Bob annoncent \emph{publiquement leurs bases} (pas les bits) et gardent les positions où les bases
coïncident : si la base est la bonne, Bob lit exactement le bit d'Alice ; sinon son résultat est aléatoire. Exemple, sans espion :

\begin{center}
\renewcommand{\arraystretch}{1.15}
\begin{tabular}{@{}l*{8}{c}@{}}
\toprule
position & 1 & 2 & 3 & 4 & 5 & 6 & 7 & 8\\
\midrule
bit d'Alice & 0 & 1 & 1 & 0 & 1 & 0 & 0 & 1\\
base d'Alice & $Z$ & $X$ & $X$ & $Z$ & $X$ & $Z$ & $X$ & $Z$\\
base de Bob & $Z$ & $Z$ & $X$ & $X$ & $X$ & $Z$ & $Z$ & $Z$\\
\midrule
bases identiques ? & oui & non & oui & non & oui & oui & non & oui\\
clé tamisée & 0 & & 1 & & 1 & 0 & & 1\\
\bottomrule
\end{tabular}
\end{center}

La clé tamisée est $01101$ : en moyenne la moitié des qubits est conservée. Pour détecter un espion Ève, Alice et Bob
comparent publiquement un échantillon de la clé tamisée. Une attaque « intercepter puis renvoyer » est détectée parce que
\emph{mesurer perturbe} (postulat de la mesure) et que l'on ne peut pas copier un état inconnu pour mesurer la copie
(théorème de non-clonage, section~\ref{sec:teleportation}). Ève choisit la mauvaise base une fois sur deux ; dans ce cas
le qubit renvoyé est dans la mauvaise base, et Bob, s'il a la bonne base, se trompe une fois sur deux. Le taux d'erreur de la clé tamisée
(\emph{QBER}) vaut donc
\[
  \mathrm{QBER}=\underbrace{\tfrac12}_{\text{Ève se trompe de base}}\times\underbrace{\tfrac12}_{\text{Bob se trompe}}=\tfrac14=25\,\%.
\]
En pratique, quelques pour cent d'erreurs viennent du bruit ; au-delà d'un seuil de l'ordre de $11\,\%$, on abandonne la clé. La
variante E91 \parencite{ekert1991} remplace l'envoi de qubits par des paires intriquées et teste l'absence d'espion par la violation
de l'inégalité CHSH (section~\ref{sec:chsh}).

\begin{remarque}[Pourquoi les agences recommandent la PQC plutôt que la QKD]
En janvier 2024, l'ANSSI, le BSI allemand et leurs homologues néerlandais et suédois ont publié une position commune
\parencite{anssi2024qkd} : la QKD n'est pas assez mûre pour la sécurité, ne convient qu'à des cas d'usage restreints, et la priorité est donnée à
une technologie « beaucoup plus compatible avec les systèmes cryptographiques existants ». Elle demande des liaisons optiques
dédiées de portée limitée, des nœuds de confiance, une authentification classique des messages publics (qui nécessite
elle-même de la cryptographie), et ses implémentations ont des failles de matériel. La QKD protège donc un lien, la PQC un
protocole entier.
\end{remarque}

% ---------------------------------------------------------------------
\subsection{Construire un ordinateur quantique}\label{sec:machine}
% ---------------------------------------------------------------------

Les sections précédentes supposent des qubits parfaits. Les qubits réels sont bruités, et c'est le bruit, bien plus que le
nombre de portes, qui sépare les machines actuelles des estimations de la section~\ref{sec:shor}. Cette section
explique ce qu'il faut pour fabriquer un qubit, pourquoi certains tournent dans un réfrigérateur à
$10\ \mathrm{mK}$, comment on corrige les erreurs, et ce que proposent les qubits de chat d'Alice \& Bob.

\subsubsection{Que faut-il ? Les critères de DiVincenzo}

\textcite{divincenzo2000} a résumé les conditions nécessaires au calcul quantique :
\begin{enumerate}[nosep]
  \item des qubits bien définis, en nombre extensible ;
  \item une initialisation fiable dans l'état $\ket{0\ldots0}$ ;
  \item une cohérence longue devant la durée d'une porte ;
  \item un jeu universel de portes : portes à un qubit et au moins une porte à deux qubits comme le CNOT (chapitre~\ref{chap:portes}) ;
  \item une mesure individuelle des qubits.
\end{enumerate}
S'y ajoutent, pour les réseaux, deux critères de communication. Le critère~3 (cohérence) est le plus difficile à tenir quand le nombre de qubits augmente.

\subsubsection{La décohérence}

Un qubit n'est jamais seul. Supposons un qubit $S$ dans l'état $\frac{1}{\sqrt2}(\ket0+\ket1)$ qui interagit
avec son environnement $E$ de façon que celui-ci « enregistre » l'état de $S$ : $\ket0\ket{E_0}$ et $\ket1\ket{E_1}$.
L'état global est $\frac{1}{\sqrt2}(\ket0\ket{E_0}+\ket1\ket{E_1})$, intriqué. L'environnement n'est pas observé : l'état de $S$ seul
est obtenu par trace partielle (section~\ref{sec:trace}),
\[
  \rho_S=\frac12\begin{pmatrix}1&\braket{E_1}{E_0}\\\braket{E_0}{E_1}&1\end{pmatrix}.
\]
Si les deux états de l'environnement sont orthogonaux, $\rho_S=\frac12 I$ : l'état est maximalement mixte
(section~\ref{sec:pur-mixte}) et la superposition est perdue. En pratique le recouvrement décroît
comme $\braket{E_1}{E_0}=\ee^{-t/T_2}$, d'où, pour la pureté,
\[
  \operatorname{Tr}\rho_S^2=\frac{1+\ee^{-2t/T_2}}{2}\ \xrightarrow[t\to\infty]{}\ \frac12 .
\]
(à $t=T_2$, elle vaut $0{,}57$). Les termes diagonaux, eux, évoluent avec le temps de relaxation $T_1$ (perte d'énergie : $\ket1\to\ket0$). Le rapport
entre $T_2$ et la durée $t_{\mathrm{porte}}$ d'une porte fixe le nombre d'opérations avant que l'information soit perdue. Avec
$T_2=100\ \mu\mathrm{s}$ et des portes de $50\ \mathrm{ns}$, il vaut $2\,000$ : l'erreur par porte est de l'ordre de
$t_{\mathrm{porte}}/T_2=5\times10^{-4}$. Calculer sur des milliards de portes sans corriger les erreurs est
donc impossible.

\subsubsection{Les plateformes}

Plusieurs systèmes physiques servent de qubits. Chacun a ses forces et ses faiblesses ; les valeurs du tableau sont des ordres de
grandeur tirés de la littérature récente et évoluent vite.

\begin{center}
\renewcommand{\arraystretch}{1.2}
\small
\begin{tabular}{@{}>{\raggedright\arraybackslash}p{2.7cm}>{\raggedright\arraybackslash}p{2.2cm}>{\raggedright\arraybackslash}p{1.9cm}>{\raggedright\arraybackslash}p{1.9cm}>{\raggedright\arraybackslash}p{5.5cm}@{}}
\toprule
\textbf{Plateforme} & \textbf{Température} & \textbf{Porte à 2 qubits} & \textbf{Cohérence} & \textbf{Atouts / difficultés}\\
\midrule
Circuits supra\-conducteurs \parencite{kjaergaard2020} & $10$--$20\ \mathrm{mK}$ & $20$--$100\ \mathrm{ns}$ & $\sim0{,}1\ \mathrm{ms}$ & portes rapides, fabrication par lithographie ; cryogénie, câblage, voisinage limité\\
Ions piégés \parencite{bruzewicz2019} & ambiante (vide poussé) & $10\ \mu\mathrm{s}$--$1\ \mathrm{ms}$ & secondes à minutes & meilleures fidélités ($\geq99{,}9\,\%$), connexion totale ; portes lentes, passage à grande échelle\\
Atomes neutres \parencite{saffman2016} & ambiante (vide poussé) & $0{,}1$--$1\ \mu\mathrm{s}$ & $\sim1\ \mathrm{s}$ & réseaux de milliers d'atomes reconfigurables ; erreurs de porte plus élevées\\
Photons \parencite{kok2007} & ambiante (détecteurs $\sim2\ \mathrm{K}$) & $\mathrm{ns}$ & pas de décohérence en vol & communication ; pertes, portes à deux photons probabilistes\\
Spins dans le silicium \parencite{burkard2023} & $0{,}1$--$1\ \mathrm{K}$ & $10$--$100\ \mathrm{ns}$ & $\sim\mathrm{ms}$ & compatible avec la microélectronique ; peu de qubits à ce jour\\
\bottomrule
\end{tabular}
\end{center}

Aucune plateforme ne domine. Les machines actuelles comptent de quelques dizaines à quelques milliers de qubits physiques
\emph{bruités}, sans correction d'erreurs : \textcite{preskill2018} les a appelées « NISQ » (\emph{noisy intermediate-scale quantum}). Elles ne suffisent pas pour Shor.

\subsubsection{Le cryostat : pourquoi $10\ \mathrm{mK}$ ?}

Un transmon (qubit supraconducteur) a une fréquence de transition $f\approx5\ \mathrm{GHz}$ : une micro-onde. À l'équilibre thermique
à la température $T$, la loi de Boltzmann donne $p_1/p_0=\ee^{-hf/k_BT}$ avec $p_0+p_1=1$, d'où
\[
  p_1=\frac{1}{1+\ee^{\,hf/k_BT}},\qquad
  \frac{hf}{k_B}=\frac{6{,}626\times10^{-34}\times5\times10^{9}}{1{,}381\times10^{-23}}\approx0{,}24\ \mathrm{K}.
\]
Le qubit n'est « dans l'état $\ket0$ » que si $T\ll0{,}24\ \mathrm{K}$. Sinon il est dans un mélange statistique de
$\ket0$ et $\ket1$, de matrice densité $\operatorname{diag}(1-p_1,p_1)$ (section~\ref{sec:rho}) :

\begin{center}
\renewcommand{\arraystretch}{1.15}
\begin{tabular}{@{}rccc@{}}
\toprule
Température $T$ & $hf/k_BT$ & population $p_1$ de $\ket1$ & pureté $\operatorname{Tr}\rho^2$\\
\midrule
$300\ \mathrm{K}$ (ambiante) & $0{,}0008$ & $0{,}4998$ & $0{,}500$\\
$4\ \mathrm{K}$ (hélium liquide) & $0{,}060$ & $0{,}485$ & $0{,}500$\\
$1\ \mathrm{K}$ & $0{,}24$ & $0{,}44$ & $0{,}507$\\
$100\ \mathrm{mK}$ & $2{,}4$ & $0{,}083$ & $0{,}85$\\
$20\ \mathrm{mK}$ & $12$ & $6{,}2\times10^{-6}$ & $0{,}99999$\\
$10\ \mathrm{mK}$ & $24$ & $3{,}8\times10^{-11}$ & $1{,}000$\\
\bottomrule
\end{tabular}
\end{center}

À température ambiante, $p_1\approx\frac12$ : l'état est maximalement mixte (pureté minimale d'un qubit,
$\operatorname{Tr}\rho^2=\frac12$, section~\ref{sec:pur-mixte}) et le qubit ne porte aucune information. Il faut descendre vers $20\ \mathrm{mK}$ pour que
l'état initial soit pur à mieux que $10^{-5}$ près. S'y ajoute une seconde raison : les jonctions Josephson sont en aluminium, supraconducteur
seulement sous $\approx1{,}2\ \mathrm{K}$, et les photons thermiques des lignes de commande doivent être éteints. Les autres
plateformes évitent ce problème autrement : ions et atomes sont préparés dans $\ket0$ par pompage optique (un laser les force dans l'état voulu)
et refroidis par laser ; les photons ont des fréquences optiques ($\sim10^{14}\ \mathrm{Hz}$, soit $hf/k_B\approx5\times10^3\ \mathrm{K}$),
sans aucune excitation thermique.

\paragraph{Le réfrigérateur à dilution.} On atteint $10\ \mathrm{mK}$ avec un mélange d'isotopes d'hélium, $^3$He et $^4$He. Sous
$0{,}87\ \mathrm{K}$, le mélange se sépare en deux phases, l'une riche en $^3$He, l'autre diluée. Faire passer des atomes de $^3$He de
la phase riche à la phase diluée \emph{absorbe de la chaleur}, comme une évaporation : c'est le principe de la dilution. Un
refroidisseur à tube pulsé fournit le froid initial (étages à $50\ \mathrm{K}$ puis $4\ \mathrm{K}$), et le circuit de dilution
descend en étages jusqu'à la \emph{chambre de mélange} (figure~\ref{fig:cryostat}). Le prix est la \emph{puissance
froide} : de l'ordre de $10\ \mu\mathrm{W}$ à $20\ \mathrm{mK}$ pour un appareil commercial, ce qui limite la chaleur que les lignes et
l'électronique peuvent apporter à cet étage.

\begin{figure}[!ht]
\centering
\adjustbox{max width=\linewidth}{%
\begin{tikzpicture}[>=Stealth, font=\small,
    plaque/.style={fill=insagris!35, draw=insagris, thick, minimum height=3.5mm, inner sep=0pt},
    legende/.style={font=\footnotesize, anchor=west, text width=7.2cm, align=left}]
  % plaques : (hauteur, largeur, température)
  \node[plaque, minimum width=6.4cm] (p0) at (0,0)    {};
  \node[plaque, minimum width=5.6cm] (p1) at (0,-1.3) {};
  \node[plaque, minimum width=4.8cm] (p2) at (0,-2.6) {};
  \node[plaque, minimum width=4.0cm] (p3) at (0,-3.9) {};
  \node[plaque, minimum width=3.2cm] (p4) at (0,-5.2) {};
  \node[plaque, minimum width=2.4cm, fill=insarouge!30, draw=insarouge] (p5) at (0,-6.5) {};
  \foreach \a/\b in {0/-1.3, -1.3/-2.6, -2.6/-3.9, -3.9/-5.2, -5.2/-6.5}
    {\draw[insagris, thick] (-0.9,\a-0.175) -- (-0.9,\b+0.175);
     \draw[insagris, thick] (0.9,\a-0.175) -- (0.9,\b+0.175);}
  % puce
  \node[draw=insarouge, fill=insarouge!15, thick, minimum width=14mm, minimum height=5mm] (puce) at (0,-7.15) {\footnotesize puce};
  \draw[insarouge, thick] (0,-6.675) -- (0,-6.9);
  % ligne de commande (coaxial) avec atténuateurs
  \draw[insarouge, thick] (0.3,0.35) -- (0.3,-6.9);
  \foreach \y/\att in {-2.6/{$-20$ dB}, -5.2/{$-20$ dB}, -6.5/{$-20$ dB}}
    {\node[draw=insarouge, fill=white, inner sep=1.8pt] at (0.3,\y) {};
     \node[font=\scriptsize, text=insarouge, anchor=west, fill=white, inner sep=1pt] at (1.3,\y) {\att};}
  % températures (gauche)
  \foreach \y/\T in {0/{$300\ \mathrm{K}$}, -1.3/{$50\ \mathrm{K}$}, -2.6/{$4\ \mathrm{K}$}, -3.9/{$\approx1\ \mathrm{K}$}, -5.2/{$\approx100\ \mathrm{mK}$}, -6.5/{$\approx10\ \mathrm{mK}$}}
    \node[anchor=east, font=\footnotesize\bfseries] at (-3.4,\y) {\T};
  % légendes (droite)
  \node[legende] at (3.7,0)    {électronique de commande et de lecture, ambiante};
  \node[legende] at (3.7,-1.3) {1\up{er} étage du tube pulsé, ancrage thermique des câbles};
  \node[legende] at (3.7,-2.6) {2\up{e} étage du tube pulsé ; amplificateur HEMT (lecture)};
  \node[legende] at (3.7,-3.9) {« still » du circuit de dilution};
  \node[legende] at (3.7,-5.2) {plaque froide : filtres et atténuateurs};
  \node[legende] at (3.7,-6.5) {chambre de mélange : qubits, amplificateur paramétrique};
\end{tikzpicture}}
\caption{Réfrigérateur à dilution en coupe schématique (« lustre »). Les plaques sont de plus en plus froides vers le
bas ; les atténuateurs placés sur la ligne de commande (valeurs typiques) ramènent le bruit thermique de $300\ \mathrm{K}$ à
celui de l'étage où ils sont fixés.}
\label{fig:cryostat}
\end{figure}

\paragraph{Le goulot d'étranglement : le câblage.} Chaque qubit demande plusieurs lignes coaxiales (commande, lecture),
qui apportent de la chaleur depuis les étages chauds et du bruit thermique. Des atténuateurs ancrés à chaque étage
ramènent ce bruit : $T_{\mathrm{eff}}\approx T_{\text{étage}}+T_{\text{chaud}}\,10^{-A/10}$, donc $60\ \mathrm{dB}$ au total limitent à
$300\ \mathrm{K}\times10^{-6}=0{,}3\ \mathrm{mK}$ la contribution du bruit de l'étage ambiant \parencite{krinner2019}. Mais des machines de millions de
qubits physiques ne tiennent pas dans un cryostat avec un coaxial par qubit : les pistes de recherche sont le multiplexage des lignes,
l'électronique de commande placée à $4\ \mathrm{K}$ et le raccordement de plusieurs cryostats entre eux.

\subsubsection{Corriger les erreurs : du code à répétition au code de surface}

Classiquement, on se protège du bruit en copiant chaque bit trois fois et en prenant le vote majoritaire. Le théorème de
non-clonage interdit de copier $\alpha\ket0+\beta\ket1$, et la mesure détruirait la superposition. On \emph{intrique} plutôt
le qubit avec deux autres.

\begin{definition}{Code à répétition à trois qubits}
Le qubit logique $\alpha\ket0+\beta\ket1$ est encodé par deux CNOT dont le contrôle est $q_1$ (et les cibles $q_2$, $q_3$) :
\[
  (\alpha\ket0+\beta\ket1)\ket{00}\ \longmapsto\ \alpha\ket{000}+\beta\ket{111}.
\]
Ce n'est pas $\ket\psi^{\otimes3}$ : l'état est intriqué. Dans cette section, les kets s'écrivent $\ket{q_1q_2q_3}$ et les
fils sont étiquetés : la convention « fil du haut = dernier chiffre » ne s'applique pas. Une erreur $X_k$ (inversion de bit sur $q_k$) est détectée en mesurant les
\emph{parités} $Z_1Z_2$ et $Z_2Z_3$, de valeurs propres $\pm1$.
\end{definition}

Sur $\alpha\ket{000}+\beta\ket{111}$ les deux parités valent $+1$. Une inversion fait changer le signe des parités qui contiennent le qubit touché :

\begin{center}
\renewcommand{\arraystretch}{1.15}
\begin{tabular}{@{}lccc@{}}
\toprule
erreur & état obtenu & $(Z_1Z_2,\ Z_2Z_3)$ & correction\\
\midrule
aucune & $\alpha\ket{000}+\beta\ket{111}$ & $(+1,+1)$ & aucune\\
$X_1$ & $\alpha\ket{100}+\beta\ket{011}$ & $(-1,+1)$ & $X_1$\\
$X_2$ & $\alpha\ket{010}+\beta\ket{101}$ & $(-1,-1)$ & $X_2$\\
$X_3$ & $\alpha\ket{001}+\beta\ket{110}$ & $(+1,-1)$ & $X_3$\\
\bottomrule
\end{tabular}
\end{center}

Les deux parités sont identiques pour $\ket{000}$ et $\ket{111}$ : leur mesure localise l'erreur \emph{sans rien apprendre sur
$\alpha$ et $\beta$}, donc sans détruire la superposition. On mesure les parités avec deux qubits
auxiliaires (\emph{ancillas}) $a_1$ et $a_2$ : le CNOT inscrit dans l'ancilla le ou-exclusif des deux qubits qu'il contrôle
(figure~\ref{fig:code3}).

\begin{figure}[!ht]
\centering
\adjustbox{max width=\linewidth}{%
\begin{tikzpicture}[>=Stealth, font=\small]
  \foreach \y in {3.6, 2.4, 1.2} \draw (0,\y) -- (11.3,\y);
  \foreach \y in {0, -1.2} \draw (0,\y) -- (8.7,\y);
  \node[anchor=east] at (0,3.6) {$q_1:\ \alpha\ket0+\beta\ket1$};
  \node[anchor=east] at (0,2.4) {$q_2:\ \ket0$};
  \node[anchor=east] at (0,1.2) {$q_3:\ \ket0$};
  \node[anchor=east] at (0,0) {$a_1:\ \ket0$};
  \node[anchor=east] at (0,-1.2) {$a_2:\ \ket0$};
  % encodage : CNOT q1 -> q2, q1 -> q3
  \draw[insarouge, thick] (1.1,3.6) -- (1.1,2.4);
  \pic at (1.1,3.6) {ctl}; \pic at (1.1,2.4) {cible};
  \draw[insarouge, thick] (2.0,3.6) -- (2.0,1.2);
  \pic at (2.0,3.6) {ctl}; \pic at (2.0,1.2) {cible};
  % erreur éventuelle
  \node[draw=insagris, dashed, minimum width=9mm, minimum height=9mm] at (3.3,2.4) {$X$};
  \node[font=\footnotesize, text=insagris, anchor=south] at (3.3,3.0) {erreur ?};
  % syndrome : a1 <- q1 xor q2 ; a2 <- q2 xor q3
  \draw[insarouge, thick] (4.5,3.6) -- (4.5,0);  \pic at (4.5,3.6) {ctl}; \pic at (4.5,0) {cible};
  \draw[insarouge, thick] (5.4,2.4) -- (5.4,0);  \pic at (5.4,2.4) {ctl}; \pic at (5.4,0) {cible};
  \draw[insarouge, thick] (6.3,2.4) -- (6.3,-1.2); \pic at (6.3,2.4) {ctl}; \pic at (6.3,-1.2) {cible};
  \draw[insarouge, thick] (7.2,1.2) -- (7.2,-1.2); \pic at (7.2,1.2) {ctl}; \pic at (7.2,-1.2) {cible};
  \pic at (8.4,0) {mesure};   \pic at (8.4,-1.2) {mesure};
  \node[anchor=west, font=\footnotesize] at (9.0,-0.6) {syndrome $(s_1,s_2)$};
  \node[porte] at (10.0,3.6) {$X^{?}$}; \node[porte] at (10.0,2.4) {$X^{?}$}; \node[porte] at (10.0,1.2) {$X^{?}$};
  \draw[dashed, insagris] (8.4,0.25) -- (8.4,0.5) -- (10.0,0.5) -- (10.0,0.85);
  \foreach \x/\t in {2.65/{encodage}, 6.0/{extraction du syndrome}}
    \node[font=\footnotesize, text=insagris] at (\x,-2.1) {\t};
\end{tikzpicture}}
\caption{Code à répétition à trois qubits : encodage, erreur éventuelle (ici sur $q_2$), extraction du syndrome
$(s_1,s_2)=(q_1\oplus q_2,\,q_2\oplus q_3)$ dans les ancillas, puis correction par une porte $X$ sur le qubit désigné
par le syndrome. Les kets s'écrivent $\ket{q_1q_2q_3}$.}
\label{fig:code3}
\end{figure}

\paragraph{Gain.} La correction échoue si au moins deux qubits sont inversés. Avec une probabilité $p$ d'inversion par
qubit, la probabilité d'échec est $3p^2(1-p)+p^3=3p^2-2p^3$, inférieure à $p$ si et seulement si $(2p-1)(p-1)>0$,
c'est-à-dire $p<\frac12$. Le code est utile \emph{si le bruit est assez faible}. C'est le principe d'un
\emph{seuil}.

\paragraph{Erreurs de phase.} Une erreur $Z_1$ sur $\alpha\ket{000}+\beta\ket{111}$ donne $\alpha\ket{000}-\beta\ket{111}$ : un état
logique différent, que ce code ne voit pas. On se protège en travaillant dans la base $X$ : $H^{\otimes3}(\alpha\ket{000}+\beta\ket{111})
=\alpha\ket{+++}+\beta\ket{---}$, où une erreur $Z$ (qui échange $\ket+$ et $\ket-$) joue le rôle d'une inversion de
bit. \textcite{shor1995} a montré que l'on peut \emph{concaténer} les deux codes (neuf qubits) pour corriger n'importe quelle erreur sur un
qubit : toute erreur est une combinaison de $I,X,Y,Z$ et la mesure du syndrome l'y projette.

\paragraph{Code de surface.} Les qubits sont disposés sur un réseau plan de $d\times d$ qubits de donnée, entourés de qubits
de mesure ; les syndromes de type $X$ et de type $Z$ sont mesurés à chaque cycle. La \emph{distance} $d$ est le nombre minimal
d'erreurs qui changent l'état logique : le code corrige jusqu'à $\lfloor(d-1)/2\rfloor$ erreurs, et il faut
$2d^2-1$ qubits physiques par qubit logique. Si le taux d'erreur physique est sous le seuil ($\approx1\,\%$), l'erreur
logique \emph{décroît exponentiellement avec $d$} : à chaque passage de $d$ à $d+2$, elle est divisée par un
facteur $\Lambda>1$. Google a observé ce comportement sur des codes de distance 3, 5 puis 7, avec
$\Lambda=2{,}14\pm0{,}02$ et une erreur logique de $0{,}143\,\%$ par cycle à $d=7$ (101 qubits physiques)
\parencite{acharya2025}. En prolongeant avec ce même $\Lambda$, atteindre une erreur logique de $10^{-12}$ par cycle demande
$\Lambda^{k}\geq1{,}43\times10^{-3}/10^{-12}$, soit $k\geq\ln(1{,}43\times10^{9})/\ln(2{,}14)\approx27{,}7$, donc
$k=28$ et $d=7+2\times28=63$ : près de $8\,000$ qubits physiques par qubit logique. C'est une extrapolation (l'erreur n'est pas
exactement exponentielle), mais elle explique les millions de qubits de la section~\ref{sec:shor} : le surcoût de la
correction d'erreurs domine tout, et un meilleur $\Lambda$ (des qubits physiques moins bruités) le réduit fortement.

\subsubsection{Les qubits de chat d'Alice \& Bob}

Le code de surface doit corriger \emph{deux} types d'erreurs ($X$ et $Z$) sur un réseau plan, d'où un surcoût en $d^2$. L'idée des
\emph{qubits de chat} \parencite{mirrahimi2014} est de concevoir un qubit dont une des deux erreurs est déjà supprimée par la physique. Il ne reste
alors qu'un type d'erreur à corriger : un code \emph{unidimensionnel} (la répétition de la section précédente), de $2d-1$ qubits
au lieu de $2d^2-1$, suffit.

\paragraph{Un qubit dans un oscillateur.} On utilise un mode d'un résonateur micro-onde, oscillateur harmonique d'états de
Fock $\ket n$ ($n$ photons). Un \emph{état cohérent} $\ket\alpha=\ee^{-|\alpha|^2/2}\sum_n\frac{\alpha^n}{\sqrt{n!}}\ket n$ est vecteur propre de l'opérateur
d'annihilation : $a\ket\alpha=\alpha\ket\alpha$, avec $|\alpha|^2$ photons en moyenne. Pour $\alpha$ réel, on pose
\[
  \ket{0_L}\approx\ket{\alpha},\qquad\ket{1_L}\approx\ket{-\alpha},\qquad
  \braket{\alpha}{-\alpha}=\ee^{-|\alpha|^2-|\alpha|^2}=\ee^{-2|\alpha|^2}.
\]
Le recouvrement est $1{,}8\times10^{-2}$ pour $|\alpha|^2=2$, $3{,}4\times10^{-4}$ pour $|\alpha|^2=4$ et $1{,}1\times10^{-7}$ pour
$|\alpha|^2=8$ : les deux états sont presque orthogonaux dès que $|\alpha|^2\gtrsim4$. Les états de la base $X$ sont les
\emph{états de chat} de Schrödinger
\[
  \ket{C_\pm}=N_\pm\big(\ket\alpha\pm\ket{-\alpha}\big),\qquad N_\pm=\frac{1}{\sqrt{2\,(1\pm\ee^{-2|\alpha|^2})}},\qquad
  \ket{\pm_L}\approx\ket{C_\pm}.
\]
$\ket{C_+}$ ne contient que des nombres \emph{pairs} de photons et $\ket{C_-}$ que des nombres \emph{impairs} (les termes en
$\alpha^n$ et $(-\alpha)^n$ s'annulent ou s'ajoutent selon la parité de $n$).

\paragraph{Une perte de photon est une erreur de phase.} Le bruit principal d'un résonateur est la perte d'un photon, décrite par $a$.
Comme $a\ket{\pm\alpha}=\pm\alpha\ket{\pm\alpha}$,
\[
  a\ket{C_\pm}=\alpha N_\pm\big(\ket\alpha\mp\ket{-\alpha}\big)=\alpha\,\frac{N_\pm}{N_\mp}\,\ket{C_\mp}.
\]
Après normalisation, $\ket{C_+}\leftrightarrow\ket{C_-}$ : c'est $\ket{+_L}\leftrightarrow\ket{-_L}$, une \emph{inversion de phase} $Z$. Sur
$\ket{\pm\alpha}$, $a$ ne change rien (un simple facteur), donc la perte de photon ne provoque aucune inversion de bit. Elle survient au taux
$\kappa_1\langle n\rangle=\kappa_1|\alpha|^2$ : le taux d'inversion de phase \emph{croît linéairement} avec $|\alpha|^2$.

\paragraph{Supprimer les inversions de bit.} On ajoute une dissipation \emph{à deux photons} : le circuit éjecte des paires de photons
(taux $\kappa_2$) et en réinjecte par paires (pompe d'amplitude $\epsilon_2$). Cela équivaut à un saut d'opérateur
$\sqrt{\kappa_2}\,(a^2-\alpha^2)$ avec $|\alpha|^2$ fixé par le rapport pompe/perte. Or $(a^2-\alpha^2)\ket{\pm\alpha}=0$ car
$a^2\ket{\pm\alpha}=\alpha^2\ket{\pm\alpha}$ : l'espace $\mathrm{Vect}(\ket\alpha,\ket{-\alpha})$ est exactement le noyau de la dissipation. Elle ramène
tout état de l'oscillateur vers cet espace et le \emph{stabilise}. Retirer les photons deux par deux conserve la parité du nombre de photons, donc la
cohérence entre $\ket{C_+}$ et $\ket{C_-}$. Pour qu'un bruit transforme $\ket\alpha$ en $\ket{-\alpha}$, il doit franchir la distance $2\alpha$ dans
le plan de phase contre la dissipation, ce qui est exponentiellement improbable : le taux d'inversion de bit décroît
comme $\ee^{-2|\alpha|^2}$.

\begin{figure}[!ht]
\centering
\begin{tikzpicture}[>=Stealth, font=\small]
  \draw[->, insagris] (-4.6,0) -- (4.8,0) node[right] {$\mathrm{Re}\,\alpha$};
  \foreach \sg in {1,-1}
    \shade[inner color=insarouge!80, outer color=insarouge!5] (\sg*2.6,0) circle (0.75);
  \node[anchor=north, font=\footnotesize] at (2.6,-0.85) {$\ket{\alpha}=\ket{0_L}$};
  \node[anchor=north, font=\footnotesize] at (-2.6,-0.85) {$\ket{-\alpha}=\ket{1_L}$};
  \draw[<->, dashed, insagris, thick] (-1.8,0) -- (1.8,0);
  \node[fill=white, inner sep=1pt, text=insarouge] at (0,0) {$\times$};
  \node[anchor=north, font=\footnotesize, text=insagris, align=center] at (0,-0.2)
    {inversion de bit : $\sim\ee^{-2|\alpha|^2}$\\ (supprimée)};
  \draw[->, insarouge, thick] (-1.5,0.75) to[bend left=40] (1.5,0.75);
  \node[anchor=south, font=\footnotesize, text=insarouge, align=center] at (0,1.3)
    {photon perdu : $\ket{C_+}\leftrightarrow\ket{C_-}$\\ (inversion de phase, $\propto\kappa_1|\alpha|^2$)};
\end{tikzpicture}
\caption{Qubit de chat dans le plan de phase de l'oscillateur (axe réel) : deux états cohérents $\ket{\pm\alpha}$ représentent $\ket{0_L}$
et $\ket{1_L}$. La dissipation à deux photons les stabilise ; passer de l'un à l'autre (inversion de bit) est exponentiellement
rare, alors que la perte d'un photon échange les états de chat $\ket{C_\pm}\approx\ket{\pm_L}$, de parités opposées (inversion de phase).}
\label{fig:chat}
\end{figure}

\paragraph{Résultat : un biais de bruit énorme.} On augmente $|\alpha|^2$ pour supprimer les inversions de bit, au prix d'inversions de phase
plus fréquentes. Mais l'asymétrie est telle que le compromis est favorable. Les expériences sont de plus en plus convaincantes :
des équipes parisiennes (ENS, Inria, Alice \& Bob) ont observé la suppression exponentielle de l'inversion de bit
\parencite{lescanne2020}, puis un temps d'inversion de bit supérieur à $10\ \mathrm{s}$, pour un temps d'inversion de phase supérieur à
$490\ \mathrm{ns}$ \parencite{reglade2024}. Le rapport est de l'ordre de $2\times10^{7}$ ($10/(0{,}49\times10^{-6})$). Une \emph{annonce de l'entreprise} Alice \& Bob (septembre 2025, résultat préliminaire, non évalué par des pairs à
notre connaissance) fait état de plus d'une heure avec $\langle n\rangle=11$ \parencite{aliceandbob2025}.

\paragraph{De là à un ordinateur.} Il reste à corriger les inversions de phase, par un code à répétition de $d$ qubits de chat
\parencite{guillaud2019}, dont les portes à deux qubits (CNOT) doivent \emph{préserver le biais} : c'est le point que les expériences
doivent encore démontrer. Une estimation d'architecture évalue le calcul d'un logarithme discret sur une courbe de 256 bits à
$126\,133$ qubits de chat et $9$ heures \parencite{gouzien2023}. Alice \& Bob (start-up française) et Amazon \parencite{putterman2025} suivent cette voie ;
Alice \& Bob vise, dans sa \emph{feuille de route} d'entreprise, une machine de 100 qubits logiques en 2030. C'est un objectif commercial, pas un résultat.

% ---------------------------------------------------------------------
\subsection{Autres applications et algorithmes}\label{sec:autres}
% ---------------------------------------------------------------------

Shor et Grover sont les algorithmes les plus connus, pas les seuls. Il faut distinguer les gains \emph{prouvés} (on sait démontrer
qu'aucun algorithme classique ne fait aussi bien, ou on connaît un bien meilleur algorithme quantique), les gains
\emph{conditionnels} (valables sous hypothèses) et les méthodes \emph{heuristiques} (qui marchent peut-être, sans preuve).

\subsubsection{Panorama}

\begin{center}
\renewcommand{\arraystretch}{1.2}
\small
\begin{tabular}{@{}>{\raggedright\arraybackslash}p{3.1cm}>{\raggedright\arraybackslash}p{3.8cm}>{\raggedright\arraybackslash}p{3.0cm}>{\raggedright\arraybackslash}p{4.6cm}@{}}
\toprule
\textbf{Algorithme} & \textbf{Problème} & \textbf{Gain} & \textbf{Statut et usage}\\
\midrule
Deutsch--Jozsa, Bernstein--Vazirani, Simon (1985--1997) & fonctions « boîtes noires » & exponentiel (Simon) & prouvé ; peu d'usage direct, base de Shor\\
Shor (1994) \parencite{shor1997} & factorisation, logarithme discret & superpolynomial & prouvé ; casse RSA et courbes elliptiques\\
Grover (1996) & recherche non structurée & quadratique : $\sqrt N$ & prouvé, optimal ; clés symétriques, SAT\\
Estimation de phase \parencite{kitaev1995} & valeur propre d'une porte $U$ & précision $2^{-m}$ avec $m$ qubits & prouvé ; brique de Shor, de la chimie, de HHL\\
Estimation d'amplitude \parencite{brassard2002} & estimer une probabilité & $\varepsilon^{-1}$ au lieu de $\varepsilon^{-2}$ & prouvé ; Monte-Carlo, finance\\
Simulation hamiltonienne \parencite{lloyd1996} & évolution $\ee^{-\ii Ht}\ket\psi$ & exponentiel (systèmes quantiques) & prouvé ; chimie, matériaux : application la plus naturelle\\
HHL \parencite{harrow2009} & système linéaire $Ax=b$ & exponentiel \emph{sous conditions} & conditionnel ; en partie « déquantisé » \parencite{tang2019}\\
VQE \parencite{peruzzo2014} & énergie de l'état fondamental & aucun gain prouvé & heuristique ; machines bruitées\\
QAOA \parencite{farhi2014} & optimisation combinatoire & aucun gain prouvé & heuristique ; machines bruitées\\
\bottomrule
\end{tabular}
\end{center}

\subsubsection{Simuler la matière : l'idée de Feynman}

Pour décrire $n$ qubits, un ordinateur classique stocke $2^n$ amplitudes complexes. Avec $16$ octets par amplitude, $n=50$
demande $2^{50}\times16=2^{54}$ octets, soit $16$ pébioctets ($1{,}8\times10^{16}$ octets) ; $n=300$ demande $2^{300}\approx2\times10^{90}$
amplitudes, plus que d'atomes dans l'univers observable ($\sim10^{80}$). \textcite{feynman1982} en a tiré l'idée
inverse : un système quantique simule naturellement un autre système quantique.

La \emph{simulation hamiltonienne} calcule $\ee^{-\ii Ht}\ket{\psi}$ pour $H=A+B$, où $A$ et $B$ sont des sommes de termes simples dont on sait implémenter les exponentielles.
Si $A$ et $B$ ne commutent pas, $\ee^{-\ii(A+B)t}\neq\ee^{-\ii At}\ee^{-\ii Bt}$. La formule de Baker--Campbell--Hausdorff donne, pour un pas $\tau$,
\[
  \ee^{-\ii A\tau}\ee^{-\ii B\tau}=\exp\!\Big(-\ii(A+B)\tau-\tfrac{\tau^2}{2}[A,B]+O(\tau^3)\Big).
\]
On coupe donc $t$ en $N$ pas $\tau=t/N$ (\emph{formule de Trotter}) :
\[
  \Big\|\ee^{-\ii(A+B)t}-\big(\ee^{-\ii At/N}\ee^{-\ii Bt/N}\big)^N\Big\|\ \leq\ \frac{t^2}{2N}\,\big\|[A,B]\big\|\quad(\text{à l'ordre dominant}),
\]
l'erreur de chaque pas ($\sim\tau^2\|[A,B]\|/2$) s'additionnant sur $N$ pas. Pour $A=X$ et $B=Z$ (un qubit), $[X,Z]=-2\ii Y$ et
$\|[X,Z]\|=2$ : l'erreur vaut au plus $t^2/N$ ; pour $t=1$, $N=100$ pas donnent une erreur inférieure à $10^{-2}$ (le calcul numérique donne $7\times10^{-3}$). Les applications visées sont la chimie
(énergie d'une molécule comme le cofacteur FeMoco de la nitrogénase, qui intervient dans la fabrication des engrais \parencite{reiher2017}),
les matériaux (batteries, supraconducteurs) et la physique des particules. Les méthodes classiques (théorie de la fonctionnelle de
la densité, \emph{coupled cluster}) traitent bien les systèmes faiblement corrélés ; c'est pour les systèmes \emph{fortement
corrélés} qu'un avantage est espéré, avec des machines tolérantes aux fautes.

\subsubsection{Méthodes variationnelles et optimisation}

Sur les machines bruitées actuelles, on ne peut pas faire tourner de longs circuits ; on utilise donc des circuits courts à paramètres
$\theta$ dans une boucle avec un ordinateur classique : le circuit prépare $\ket{\psi(\theta)}$, on mesure l'énergie
$E(\theta)=\bra{\psi(\theta)}H\ket{\psi(\theta)}$, et un optimiseur classique corrige $\theta$. Cela repose sur le
\emph{principe variationnel} : en décomposant $\ket\psi=\sum_kc_k\ket{k}$ sur les états propres de $H$ (section~\ref{sec:spectral}),
$E=\sum_k|c_k|^2E_k\geq E_0\sum_k|c_k|^2=E_0$. Minimiser $E(\theta)$ approche donc l'énergie fondamentale $E_0$ par le haut. Exemple pour un qubit,
$H=X+Z$ (valeurs propres $\pm\sqrt2$) avec $\ket{\psi(\theta)}=\cos\frac\theta2\ket0+\sin\frac\theta2\ket1$ :
\[
  E(\theta)=\langle Z\rangle+\langle X\rangle=\cos\theta+\sin\theta=\sqrt2\cos\!\big(\theta-\tfrac\pi4\big),
\]
minimale pour $\theta=\frac{5\pi}{4}$, de valeur $-\sqrt2=E_0$. Le VQE \parencite{peruzzo2014} généralise cette idée à des molécules, et le QAOA
\parencite{farhi2014} l'applique à des problèmes d'optimisation en alternant deux types de portes. Aucun avantage n'est démontré pour ces méthodes :
elles sont l'objet de la recherche actuelle.

\subsubsection{Algèbre linéaire quantique : une promesse sous conditions}

L'algorithme HHL \parencite{harrow2009} prépare un état $\ket{x}\propto\sum_ix_i\ket{i}$ où $Ax=b$, en un temps polynomial en $\log N$
là où la résolution classique est polynomiale en $N$. Le gain exponentiel apparent est soumis à plusieurs conditions. Il faut
\emph{charger} $b$ dans l'état $\ket{b}$ efficacement ; $A$ doit être creuse et bien conditionnée ; et le résultat est un
\emph{état} : en lire les $N$ composantes demanderait déjà $N$ répétitions, de sorte que l'on n'obtient que des
grandeurs globales de $x$. De plus, des algorithmes classiques « inspirés du quantique » (\emph{déquantisation}) obtiennent le même
résultat sous des hypothèses d'accès analogues \parencite{tang2019}, ce qui a supprimé l'avantage exponentiel annoncé pour
certaines applications d'apprentissage. Le bilan honnête : \emph{pas d'avantage démontré en apprentissage automatique à ce jour}.

\subsubsection{Capteurs, nombres aléatoires, réseaux}

\paragraph{Métrologie quantique.} Les premières technologies quantiques commercialisées sont des capteurs (horloges atomiques,
magnétomètres, gravimètres). L'intrication y améliore la précision. Soit $N$ qubits qui subissent chacun la phase
$\ket1\mapsto\ee^{\ii\varphi}\ket1$ à mesurer. Dans l'état de GHZ $\frac{1}{\sqrt2}(\ket{0\ldots0}+\ket{1\ldots1})$ (généralisation à $N$ qubits de l'état de Bell
$\frac{1}{\sqrt2}(\ket{00}+\ket{11})$, section~\ref{sec:bell}), seuls $\ket{0\ldots0}$ et $\ket{1\ldots1}$ interviennent, et l'état devient
\[
  \frac{1}{\sqrt2}\big(\ket{0\ldots0}+\ee^{\ii N\varphi}\ket{1\ldots1}\big),\qquad
  P\big(X^{\otimes N}=+1\big)=\frac{1+\cos N\varphi}{2}=\cos^2\!\frac{N\varphi}{2}.
\]
La phase est amplifiée par $N$ : la précision atteint $\Delta\varphi\sim1/N$ (limite de Heisenberg) au lieu de $1/\sqrt N$ pour
$N$ qubits indépendants (limite quantique standard) \parencite{giovannetti2004}. En pratique, la décohérence réduit ce gain, qui
se limite souvent à de petits $N$.

\paragraph{Nombres aléatoires.} Mesurer $\ket+$ dans la base $Z$ donne $0$ ou $1$ avec la probabilité $\frac12$ : le postulat de Born
garantit un hasard \emph{intrinsèque}, contrairement à un générateur pseudo-aléatoire déterministe. Les générateurs quantiques
de nombres aléatoires (QRNG), à base de photons, sont déjà commercialisés ; ils alimentent par exemple la génération de clés cryptographiques.

\paragraph{Réseaux quantiques.} Relier des ordinateurs quantiques demande de transmettre des qubits (téléportation, section~\ref{sec:teleportation}).
Une fibre optique standard atténue de $0{,}2\ \mathrm{dB/km}$, soit une transmission $T=10^{-0{,}02L}$ ($L$ en km) :
$T=10^{-2}$ à $100\ \mathrm{km}$ et $10^{-10}$ à $500\ \mathrm{km}$. Un amplificateur est impossible (non-clonage) : on utilise des
\emph{répéteurs quantiques}, qui étendent l'intrication par \emph{échange d'intrication} (une mesure de Bell entre deux paires intriquées, comme dans la
téléportation de la section~\ref{sec:teleportation}), ou des satellites : la distribution
d'intrication sur $1\,200\ \mathrm{km}$ a été démontrée par le satellite chinois Micius \parencite{yin2017}. L'objectif est un « internet
quantique » \parencite{wehner2018} reliant des processeurs, des capteurs et des liens de QKD.

\subsubsection{Jalons et niveaux de preuve}

Le tableau place quelques jalons dans le temps et indique la \emph{nature de la source}, car une annonce d'entreprise ou une
prépublication n'a pas la valeur d'un article évalué par des pairs ou d'une norme.

\begin{center}
\renewcommand{\arraystretch}{1.2}
\small
\begin{tabular}{@{}>{\raggedright\arraybackslash}p{1.9cm}>{\raggedright\arraybackslash}p{8.8cm}>{\raggedright\arraybackslash}p{3.3cm}@{}}
\toprule
\textbf{Date} & \textbf{Jalon} & \textbf{Nature de la source}\\
\midrule
1994--1996 & algorithmes de Shor et de Grover \parencite{shor1997} & articles évalués\\
2019 & Google, $53$ qubits : échantillonnage de circuits aléatoires en $200\ \mathrm{s}$, annoncé hors de portée des machines classiques \parencite{arute2019} (contesté ensuite) & article évalué\\
août 2024 & normes NIST FIPS 203, 204, 205 \parencite{fips203} & norme\\
2024 & qubit de chat : inversion de bit $>10\ \mathrm{s}$ \parencite{reglade2024} & article évalué\\
déc. 2024 & code de surface sous le seuil, $\Lambda=2{,}14$ \parencite{acharya2025} & article évalué\\
mars 2025 & HQC choisi par le NIST \parencite{nist2025hqc} & rapport officiel\\
mai 2025 & RSA-2048 : moins d'un million de qubits \parencite{gidney2025} & prépublication\\
sept. 2025 & qubit de chat : plus d'une heure d'inversion de bit \parencite{aliceandbob2025} & annonce d'entreprise\\
mars 2026 & Shor avec $10\,000$ à $26\,000$ atomes \parencite{cain2026} & prépublication\\
août 2026 & RFC~10024 : TLS hybride X25519 + ML-KEM \parencite{rfc10024} & norme IETF\\
\bottomrule
\end{tabular}
\end{center}

Ce qui est \emph{acquis} : les gains théoriques de Shor et de Grover, la correction d'erreurs sous le seuil sur de petits codes,
et les standards post-quantiques. Ce qui \emph{ne l'est pas} : un ordinateur à plusieurs milliers de qubits \emph{logiques},
un avantage utile sur un problème pratique, et surtout la date $Z$ de l'inégalité de Mosca. C'est pourquoi la migration cryptographique se
décide sur $X+Y$, qui se connaît, et non sur une prédiction de $Z$.

\clearpage
% ---------------------------------------------------------------------
\subsection{Récapitulatif du chapitre}\label{sec:recap-appli}
% ---------------------------------------------------------------------

\begin{center}
\adjustbox{max width=\linewidth}{%
\begin{tikzpicture}[>=Stealth, font=\footnotesize,
    boite/.style={draw=insarouge, fill=insarouge!6, rounded corners=2pt, thick,
                  minimum width=4.4cm, minimum height=1.45cm, align=center, inner sep=3pt},
    lien/.style={->, very thick, insagris},
    etiq/.style={font=\scriptsize\itshape, text=insagris, inner sep=2pt}]
  \node[boite] (int) at (0,0)
    {\textbf{Interférence}\\ oracle, Deutsch, Grover : $\sqrt N$\\ {\scriptsize\S\,\ref{sec:interference}, \ref{sec:grover}}};
  \node[boite] (sho) at (6.3,0)
    {\textbf{Shor}\\ QFT, période $r$, $O(n^3)$\\ {\scriptsize\S\,\ref{sec:shor}}};
  \node[boite] (cry) at (12.6,0)
    {\textbf{Cryptographie}\\ RSA et ECC cassés ; HNDL ; PQC\\ {\scriptsize\S\,\ref{sec:crypto}}};
  \node[boite] (cryo) at (12.6,-2.3)
    {\textbf{Qubits physiques}\\ décohérence, cryostat $10\ \mathrm{mK}$\\ {\scriptsize\S\,\ref{sec:machine}}};
  \node[boite] (qec) at (6.3,-2.3)
    {\textbf{Correction d'erreurs}\\ seuil, code de surface, qubits de chat\\ {\scriptsize\S\,\ref{sec:machine}}};
  \node[boite] (aut) at (0,-2.3)
    {\textbf{Autres usages}\\ simulation, VQE, capteurs, réseaux\\ {\scriptsize\S\,\ref{sec:autres}}};
  \draw[lien] (int) -- (sho) node[etiq, midway, above=1pt] {période};
  \draw[lien] (sho) -- (cry) node[etiq, midway, above=1pt] {casse RSA};
  \draw[lien] (cry) -- (cryo) node[etiq, midway, right=1pt] {quelle machine ?};
  \draw[lien] (cryo) -- (qec) node[etiq, midway, above=1pt] {qubits fiables};
  \draw[lien] (qec) -- (aut) node[etiq, midway, above=1pt] {usages};
\end{tikzpicture}}
\end{center}

\begin{aretenir}
\begin{itemize}
  \item Un algorithme quantique organise des \emph{interférences} : l'oracle écrit $f(x)$ dans une phase ; Deutsch décide avec un appel au lieu de deux.
        Grover : $\approx\frac\pi4\sqrt N$ itérations (rotations de $2\theta$), gain quadratique et optimal ; exact en une itération pour $N=4$.
  \item Shor : la QFT transforme la période $r$ de $a^x\bmod N$ en pics en $y=jM/r$ ; $\operatorname{pgcd}(a^{r/2}\pm1,N)$ donne les facteurs ;
        coût $O(n^3)$ ($n=\log_2N$), contre un temps sous-exponentiel classique (exemple : $N=15$, $a=7$, $r=4$).
  \item RSA, Diffie--Hellman, ECDH et ECDSA tombent face à Shor ; AES-256 et SHA-256 résistent à Grover. Doubler les clés RSA ne protège pas.
  \item « Stocker maintenant, déchiffrer plus tard » : exposé dès que $X+Y>Z$ (Mosca) ; la forward secrecy n'aide pas. On migre d'abord l'échange de clés
        (hybride X25519 + ML-KEM), puis les signatures (FIPS 203--205).
  \item Qubit physique : la décohérence ($T_2$) mélange l'état ; le cryostat à $10\ \mathrm{mK}$ rend $p_1$ négligeable ($hf/k_B=0{,}24\ \mathrm{K}$ à $5\ \mathrm{GHz}$). Correction
        d'erreurs : seuil, $2d^2-1$ qubits physiques par qubit logique (code de surface).
  \item Qubit de chat : $\ket{\pm\alpha}$, perte de photon = inversion de phase, inversion de bit en $\ee^{-2|\alpha|^2}$ ; code à répétition. Une annonce d'entreprise
        n'est pas un résultat évalué.
  \item Autres usages : simulation (Trotter), VQE et QAOA (heuristiques), HHL (conditionnel), capteurs (limite de Heisenberg), réseaux (répéteurs).
\end{itemize}
\end{aretenir}

\begin{center}
\renewcommand{\arraystretch}{1.12}
\small
\begin{tabular}{@{}l l r@{}}
\toprule
\textbf{Notion} & \textbf{Formule} & \textbf{§} \\
\midrule
Rebond de phase & $U_f\ket{x}\ket{-}=(-1)^{f(x)}\ket{x}\ket{-}$ & \ref{sec:interference} \\
Grover & $P=\sin^2\!\big((2k+1)\theta\big)$, \quad $\sin\theta=\frac{1}{\sqrt N}$, \quad $k\simeq\frac\pi4\sqrt N$ & \ref{sec:grover} \\
QFT et période & $\mathrm{QFT}_M\ket x=\frac{1}{\sqrt M}\sum_y\ee^{2\ii\pi xy/M}\ket y$ ; pics en $y=jM/r$, $P=\frac1r$ & \ref{sec:shor} \\
Facteurs, coût & $\operatorname{pgcd}(a^{r/2}\pm1,N)$, \quad $O(n^3)$ & \ref{sec:shor} \\
Mosca & $X+Y>Z$ : données déjà exposées & \ref{sec:crypto} \\
Population thermique & $p_1=\frac{1}{1+\ee^{hf/k_BT}}$, \quad $\operatorname{Tr}\rho^2=\frac{1+\ee^{-2t/T_2}}{2}$ (décohérence) & \ref{sec:machine} \\
Qubit de chat & $\braket{\alpha}{-\alpha}=\ee^{-2|\alpha|^2}$, \quad $a\ket{C_\pm}=\alpha\frac{N_\pm}{N_\mp}\ket{C_\mp}$ & \ref{sec:machine} \\
\bottomrule
\end{tabular}
\end{center}

\begin{savoirfaire}
\begin{itemize}[nosep]
  \item Suivre un circuit de Deutsch ou de Grover ($N=4$) jusqu'à la mesure (sections~\ref{sec:interference}, \ref{sec:grover}).
  \item Trouver l'ordre de $a$ modulo $N$, appliquer le critère du pgcd et prévoir la loi de $y$ (section~\ref{sec:shor}, exemple $N=15$).
  \item Dire ce que casse Shor, ce que Grover affaiblit, et appliquer l'inégalité de Mosca (section~\ref{sec:crypto}).
  \item Calculer $p_1$, la pureté après décohérence et le recouvrement $\ee^{-2|\alpha|^2}$ ; lire un syndrome (section~\ref{sec:machine}).
  \item Distinguer résultat prouvé, conditionnel et annonce (section~\ref{sec:autres}).
\end{itemize}
\end{savoirfaire}
